\documentclass[10pt,a4paper,oneside,open=any]{scrbook}
\usepackage{style/germanbook}
\DeclareTOCStyleEntry[numwidth=3em]{tocline}{section}
\hypersetup{
  pdftitle={German Verb Reference},
  pdfauthor={Jason Laurie}
}

\title{German Verb Reference}
\subtitle{Conjugation, Government, and Grammar Tables}
\author{Jason Laurie}
\date{September 2026}

\begin{document}

\frontmatter
\maketitle
\tableofcontents

\mainmatter
\part{Grammar Reference}
\chapter{Cases and Declension}

\section{The four cases}

German case marks the job performed by a noun phrase. The tables below are a
compact reference; the examples show the same masculine noun changing case.

\begin{center}
\begin{tabularx}{\textwidth}{@{}llllX@{}}
\toprule
Case & Typical function & Question & Masculine example & Example sentence \\
\midrule
\rowcolor{caseNominativeBg}
Nominativ & subject & wer?/was? & der Mann & \textit{Der Mann wartet.} \\
\rowcolor{caseAccusativeBg}
Akkusativ & direct object & wen?/was? & den Mann & \textit{Ich sehe den Mann.} \\
\rowcolor{caseDativeBg}
Dativ & indirect object & wem? & dem Mann & \textit{Ich helfe dem Mann.} \\
\rowcolor{caseGenitiveBg}
Genitiv & possession/association & wessen? & des Mannes & \textit{Das Auto des Mannes.} \\
\bottomrule
\end{tabularx}
\end{center}

\section{Definite and indefinite articles}

\subsection{Definite article}

\begin{center}
\begin{tabular}{@{}lMNFP@{}}
\toprule
 & Maskulin & Neutrum & Feminin & Plural \\
\midrule
\rowcolor{caseNominativeBg}
Nominativ & der & das & die & die \\
\rowcolor{caseAccusativeBg}
Akkusativ & den & das & die & die \\
\rowcolor{caseDativeBg}
Dativ & dem & dem & der & den \\
\rowcolor{caseGenitiveBg}
Genitiv & des & des & der & der \\
\bottomrule
\end{tabular}
\end{center}

\subsection{Indefinite article and \textit{ein}-words}

The same endings are used by \textit{kein} and possessive determiners such as
\textit{mein}, \textit{dein}, \textit{sein}, \textit{ihr}, \textit{unser}, and
\textit{euer}. A dash marks the absence of an indefinite plural article.

\begin{center}
\begin{tabular}{@{}lMNFP@{}}
\toprule
 & Maskulin & Neutrum & Feminin & Plural \\
\midrule
\rowcolor{caseNominativeBg}
Nominativ & ein & ein & eine & -- \\
\rowcolor{caseAccusativeBg}
Akkusativ & einen & ein & eine & -- \\
\rowcolor{caseDativeBg}
Dativ & einem & einem & einer & -- \\
\rowcolor{caseGenitiveBg}
Genitiv & eines & eines & einer & -- \\
\bottomrule
\end{tabular}
\end{center}

\section{Adjective endings}

The determiner and adjective share the job of signalling gender, number, and
case. After a definite article the adjective normally has a weak ending. After
an \textit{ein}-word it supplies information missing from the determiner. With
no determiner, the adjective carries the strong ending itself.

\subsection{Weak endings after the definite article}

\begin{center}
\begin{tabular}{@{}lMNFP@{}}
\toprule
 & Maskulin & Neutrum & Feminin & Plural \\
\midrule
\rowcolor{caseNominativeBg}
Nominativ & -e & -e & -e & -en \\
\rowcolor{caseAccusativeBg}
Akkusativ & -en & -e & -e & -en \\
\rowcolor{caseDativeBg}
Dativ & -en & -en & -en & -en \\
\rowcolor{caseGenitiveBg}
Genitiv & -en & -en & -en & -en \\
\bottomrule
\end{tabular}
\end{center}

\subsection{Mixed endings after an \textit{ein}-word}

\begin{center}
\begin{tabular}{@{}lMNFP@{}}
\toprule
 & Maskulin & Neutrum & Feminin & Plural with possessive/\textit{kein} \\
\midrule
\rowcolor{caseNominativeBg}
Nominativ & -er & -es & -e & -en \\
\rowcolor{caseAccusativeBg}
Akkusativ & -en & -es & -e & -en \\
\rowcolor{caseDativeBg}
Dativ & -en & -en & -en & -en \\
\rowcolor{caseGenitiveBg}
Genitiv & -en & -en & -en & -en \\
\bottomrule
\end{tabular}
\end{center}

\subsection{Strong endings without a determiner}

\begin{center}
\begin{tabular}{@{}lMNFP@{}}
\toprule
 & Maskulin & Neutrum & Feminin & Plural \\
\midrule
\rowcolor{caseNominativeBg}
Nominativ & -er & -es & -e & -e \\
\rowcolor{caseAccusativeBg}
Akkusativ & -en & -es & -e & -e \\
\rowcolor{caseDativeBg}
Dativ & -em & -em & -er & -en \\
\rowcolor{caseGenitiveBg}
Genitiv & -en & -en & -er & -er \\
\bottomrule
\end{tabular}
\end{center}

\subsection{Model noun phrases}

\begin{center}
\begin{tabularx}{\textwidth}{@{}lXXX@{}}
\toprule
 & Definite article & \textit{ein}-word & No article \\
\midrule
\rowcolor{caseNominativeBg}
Nominativ & der gute Wein & ein guter Wein & guter Wein \\
\rowcolor{caseAccusativeBg}
Akkusativ & den guten Wein & einen guten Wein & guten Wein \\
\rowcolor{caseDativeBg}
Dativ & mit dem guten Wein & mit einem guten Wein & mit gutem Wein \\
\rowcolor{caseGenitiveBg}
Genitiv & wegen des guten Weines & wegen eines guten Weines & wegen guten Weines \\
\bottomrule
\end{tabularx}
\end{center}

\section{Personal and reflexive pronouns}

\begin{center}
\begin{tabular}{@{}lcccccc@{}}
\toprule
 & ich & du & er/sie/es & wir & ihr & sie/Sie \\
\midrule
\rowcolor{caseAccusativeBg}
Akkusativ & mich & dich & ihn/sie/es & uns & euch & sie/Sie \\
\rowcolor{caseDativeBg}
Dativ & mir & dir & ihm/ihr/ihm & uns & euch & ihnen/Ihnen \\
\rowcolor{caseAccusativeBg}
Reflexiv Akk. & mich & dich & sich & uns & euch & sich \\
\rowcolor{caseDativeBg}
Reflexiv Dat. & mir & dir & sich & uns & euch & sich \\
\bottomrule
\end{tabular}
\end{center}

\section{Important noun endings}

\begin{itemize}[leftmargin=*]
  \item Dative plural nouns normally add \textbf{-n} when the plural does not
        already end in \textit{-n} or \textit{-s}: \textit{mit den Kindern}.
  \item Masculine and neuter nouns normally add \textbf{-(e)s} in the genitive:
        \textit{des Mannes}, \textit{des Kindes}.
  \item Weak masculine nouns add \textbf{-(e)n} outside the nominative singular:
        \textit{der Student}, but \textit{den/dem/des Studenten}.
\end{itemize}

\chapter{Greetings, Farewells, Letters, and Emails}

The appropriate greeting depends on the relationship, situation, time of day,
and region. In unfamiliar, official, or professional situations, begin with
\textit{Sie} unless the other person or the organisation clearly uses
\textit{du}. It is normally safe to match the level of formality used by the
other person.

\section{Everyday spoken greetings}

\begin{center}
\small
\rowcolors{2}{referenceBlueBg}{white}
\begin{tabularx}{\textwidth}{@{}>{\bfseries}l l l X@{}}
\toprule
German & English & Register & Typical use \\
\midrule
Guten Morgen! & Good morning! & neutral/formal & morning \\
Guten Tag! & Good day/hello! & neutral/formal & daytime \\
Guten Abend! & Good evening! & neutral/formal & evening \\
Hallo! & Hello! & neutral/informal & very widely used \\
Hi! & Hi! & informal & friends and familiar contacts \\
Grüß dich! & Hello! & informal & one person, especially regionally \\
Grüß euch! & Hello! & informal & several people, especially regionally \\
Herzlich willkommen! & A warm welcome! & all registers & welcoming a visitor \\
\bottomrule
\end{tabularx}
\end{center}

German normally uses \textit{Guten Tag}, not a literal translation of
\enquote{good afternoon}. \textit{Gute Nacht} is normally said when leaving late or
going to bed, rather than when first greeting somebody in the evening.

\subsection{Regional greetings}

\begin{center}
\small
\rowcolors{2}{referenceGreenBg}{white}
\begin{tabularx}{\textwidth}{@{}>{\bfseries}l l X@{}}
\toprule
Expression & Region/register & Note \\
\midrule
Moin! & northern Germany, informal/neutral & commonly used throughout the day \\
Grüß Gott! & southern Germany and Austria, polite & literally \enquote{greet God} \\
Servus! & southern Germany and Austria, informal & used for hello and goodbye \\
Mahlzeit! & workplace, informal & greeting around lunchtime \\
\bottomrule
\end{tabularx}
\end{center}

Regional expressions are useful to recognise, but the neutral forms
\textit{Guten Tag} and \textit{Hallo} work more widely.

\section{Introductions and small talk}

\begin{center}
\small
\rowcolors{2}{referenceYellowBg}{white}
\begin{tabularx}{\textwidth}{@{}>{\bfseries}X X@{}}
\toprule
German & English \\
\midrule
Wie geht es dir? & How are you? (informal) \\
Wie geht es Ihnen? & How are you? (formal) \\
Danke, gut. Und dir/Ihnen? & Fine, thank you. And you? \\
Darf ich mich vorstellen? & May I introduce myself? \\
Ich heiße \ldots{} / Mein Name ist \ldots{} & My name is \ldots{} \\
Freut mich, dich/Sie kennenzulernen. & Pleased to meet you. \\
Schön, dich/Sie wiederzusehen. & Nice to see you again. \\
Wie war Ihre/deine Reise? & How was your journey? \\
\bottomrule
\end{tabularx}
\end{center}

Use \textit{Herr} or \textit{Frau} with the surname:
\textit{Guten Tag, Frau Wagner}. \textit{Fräulein} is outdated and should not
be used as the normal form of address. Retain an academic title when it is
given, for example \textit{Sehr geehrte Frau Dr.~Weber}.

\section{Saying goodbye}

\begin{center}
\small
\rowcolors{2}{referenceRedBg}{white}
\begin{tabularx}{\textwidth}{@{}>{\bfseries}l l X@{}}
\toprule
German & English & Register/use \\
\midrule
Auf Wiedersehen! & Goodbye! & neutral/formal, face to face \\
Tschüss! & Bye! & neutral/informal \\
Bis bald! & See you soon! & neutral/informal \\
Bis später! & See you later! & neutral/informal \\
Bis morgen/Montag! & See you tomorrow/Monday! & neutral/informal \\
Bis dann! & See you then! & informal \\
Mach's gut! / Macht's gut! & Take care! & informal singular/plural \\
Einen schönen Tag noch! & Have a nice day! & polite/neutral \\
Gute Nacht! & Good night! & late evening/bedtime \\
Auf Wiederhören! & Goodbye! & telephone calls \\
\bottomrule
\end{tabularx}
\end{center}

A natural response is \textit{Danke, gleichfalls!} or
\textit{Danke, Ihnen/dir auch!} To suggest using first names and
\textit{du}, ask \textit{Wollen wir uns duzen?} Until that is established,
continue using \textit{Sie} in a formal relationship.

\section{Useful telephone language}

\begin{center}
\small
\rowcolors{2}{referenceBlueBg}{white}
\begin{tabularx}{\textwidth}{@{}>{\bfseries}X X@{}}
\toprule
German & English \\
\midrule
Guten Tag, hier spricht Alex Morgan. & Hello, Alex Morgan speaking. \\
Könnte ich bitte Frau Klein sprechen? & Could I speak to Ms Klein, please? \\
Einen Moment bitte. & One moment, please. \\
Sie ist leider gerade nicht erreichbar. & Unfortunately, she is unavailable. \\
Kann ich etwas ausrichten? & Can I take a message? \\
Ich rufe später noch einmal an. & I will call again later. \\
Vielen Dank für Ihren Anruf. & Thank you for your call. \\
Auf Wiederhören! & Goodbye! \\
\bottomrule
\end{tabularx}
\end{center}

\section{Choosing a written register}

\begin{center}
\small
\rowcolors{2}{referenceGreenBg}{white}
\begin{tabularx}{\textwidth}{@{}>{\bfseries}l X X@{}}
\toprule
Register & Opening & Closing \\
\midrule
Formal & Sehr geehrte Frau Müller, & Mit freundlichen Grüßen \\
Formal, unknown recipient(s) & Sehr geehrte Damen und Herren, & Mit freundlichen Grüßen \\
Established professional contact & Liebe Frau Müller, & Beste/Freundliche Grüße \\
Neutral/informal & Hallo Anna, & Viele Grüße \\
Personal & Liebe Anna, / Lieber Paul, & Herzliche/Liebe Grüße \\
\bottomrule
\end{tabularx}
\end{center}

\textit{Hallo} is common in ordinary emails, but
\textit{Sehr geehrte Frau + surname} or
\textit{Sehr geehrter Herr + surname} remains the safest opening for a
first formal contact, an application, or communication with an authority.

\section{Structure of a formal letter or email}

\begin{enumerate}[leftmargin=*]
  \item Give the email a short, informative subject. A modern subject line
        does not need the word \textit{Betreff} and normally has no final full
        stop.
  \item Use the person's name when it is known. Put a comma after the
        salutation.
  \item Begin the following line with a lowercase letter unless its first word
        must otherwise be capitalised: \textit{Sehr geehrte Frau Klein,
        vielen Dank \ldots{}}
  \item State the reason for writing early, then divide longer text into clear
        paragraphs.
  \item End with an appropriate closing sentence and a separate complimentary
        close.
  \item Do not put a comma or full stop after a standalone closing such as
        \textit{Mit freundlichen Grüßen}.
  \item Add your full name. A professional email may also include a signature
        block with contact details.
\end{enumerate}

In formal correspondence, the pronouns \textit{Sie, Ihnen}, and
\textit{Ihr/Ihre} are always capitalised. Informal \textit{du, dir, dich,
ihr}, and the corresponding possessives are normally lowercase; in personal
letters, emails, and messages they may alternatively be capitalised. Whichever
option is chosen should be used consistently.

For a printed formal letter, also include the sender and recipient addresses,
place and date, subject, and a handwritten signature. List enclosures beneath
the signature if necessary. Exact positioning may follow an employer's
template or the DIN~5008 business-letter standard.

\section{Useful written phrases}

\begin{center}
\small
\rowcolors{2}{referenceYellowBg}{white}
\begin{tabularx}{\textwidth}{@{}>{\bfseries}l X X@{}}
\toprule
Purpose & German & English \\
\midrule
Opening & Vielen Dank für Ihre Nachricht. & Thank you for your message. \\
Reason & Ich schreibe Ihnen, weil \ldots{} & I am writing because \ldots{} \\
Enquiry & Ich möchte mich nach \ldots{} erkundigen. & I would like to ask about \ldots{} \\
Request & Könnten Sie mir bitte mitteilen, ob \ldots{}? & Could you please tell me whether \ldots{}? \\
Attachment & Im Anhang finden Sie \ldots{} & Please find \ldots{} attached. \\
Apology & Bitte entschuldigen Sie die verspätete Antwort. & Please excuse the late reply. \\
Problem & Leider kann ich den Termin nicht wahrnehmen. & Unfortunately, I cannot attend the appointment. \\
Thanks & Vielen Dank im Voraus. & Thank you in advance. \\
Closing & Ich freue mich auf Ihre Antwort. & I look forward to your reply. \\
Offer & Für Rückfragen stehe ich gern zur Verfügung. & I am happy to answer any questions. \\
\bottomrule
\end{tabularx}
\end{center}

\section{Model emails}

\subsection{Formal email}

\begin{tcolorbox}[
  colback=referenceBlueBg,
  colframe=genderMasculine,
  title={Anfrage zum B1-Abendkurs},
  fonttitle=\bfseries
]
Sehr geehrte Frau Schneider,

\medskip
ich interessiere mich für Ihren Deutschkurs auf Niveau B1. Könnten Sie mir
bitte mitteilen, wann der nächste Abendkurs beginnt und wie viel er kostet?

\medskip
Vielen Dank im Voraus. Ich freue mich auf Ihre Antwort.

\medskip
Mit freundlichen Grüßen

Alex Morgan
\end{tcolorbox}

\subsection{Informal email}

\begin{tcolorbox}[
  colback=referenceRedBg,
  colframe=genderFeminine,
  title={Treffen am Samstag},
  fonttitle=\bfseries
]
Hallo Lena,

\medskip
danke für deine Nachricht. Ich habe am Samstag Zeit. Wollen wir uns um
15~Uhr im Café treffen?

\medskip
Bis bald

Alex
\end{tcolorbox}

\begin{tcolorbox}[
  colback=referenceGreenBg,
  colframe=genderNeuter,
  title={Final check before sending},
  fonttitle=\bfseries
]
Check the recipient and attachment, use one consistent register, include a
clear subject, answer every question or requested point, and make sure the
message has both an opening and an appropriate closing.
\end{tcolorbox}

\chapter{The German Tense System}

\section{Präsens (Present)}
The Präsens tense is the most frequently used tense in German and covers a wide range of meanings. It is used to describe actions happening now, general truths, habitual actions, and very often future events when a time expression makes the meaning clear. For example, \textit{Ich lerne Deutsch} can mean “I learn German” or “I am learning German,” and \textit{Morgen gehe ich nach Hause} uses the present tense to refer to the future. German does not have a separate continuous tense, so Präsens fulfills both roles.

\section{Perfekt (Present Perfect)}
The Perfekt tense is the main tense for talking about completed actions in the past, especially in spoken German. It is formed using the present tense of \textit{haben} or \textit{sein} plus the past participle of the main verb. For instance, \textit{Ich habe gegessen} means “I ate” or “I have eaten.” In everyday conversation, the Perfekt almost completely replaces the simple past for most verbs, making it essential for understanding and speaking modern German.

\section{Präteritum (Simple Past)}
The Präteritum tense is primarily used in written German, such as novels, newspapers, and formal narratives, to describe past events. In spoken German, it is largely limited to a small group of very common verbs, especially \textit{sein}, \textit{haben}, and the modal verbs, as in \textit{Ich war müde} or \textit{Ich musste gehen}. Although learners often encounter this tense in textbooks, its practical spoken use is quite limited.

\section{Plusquamperfekt (Past Perfect)}
The Plusquamperfekt tense is used to describe an action that had already been completed before another action in the past. It is formed with the simple past of \textit{haben} or \textit{sein} plus the past participle. An example is \textit{Ich hatte gegessen, bevor er kam}, meaning “I had eaten before he arrived.” This tense helps clarify the sequence of past events and is used when the order of actions is important.

\section{Futur I (Future)}
The Futur I tense is formed with \textit{werden} plus the infinitive of the main verb and is used to talk about future actions or make assumptions about the present or future. For example, \textit{Ich werde lernen} means “I will study,” while \textit{Er wird krank sein} expresses a guess, such as “He is probably sick.” However, Germans often prefer the present tense with a time expression instead of Futur I, making this tense less common than learners might expect.

\chapter{Verb Position and the Sentence Bracket}

German word order becomes much easier when the \textbf{finite verb} is found
first. The finite verb is the form that agrees with the subject and carries
the tense: \textit{lernt}, \textit{hat}, \textit{muss}, or \textit{wird}.
Other verb parts often stand at the end of the clause and form a
\textit{Satzklammer} (sentence bracket) with it.

\section{Main clauses: the finite verb is in position two}

In a statement, the finite verb is the second \emph{constituent}, not
necessarily the second word. Any one complete element can occupy position one.

\begin{center}
\small
\rowcolors{2}{referenceBlueBg}{white}
\begin{tabularx}{\textwidth}{@{}>{\bfseries}l l l X@{}}
\toprule
Position 1 & Finite verb & Middle field & End field \\
\midrule
Ich & \textbf{lerne} & heute Abend & Deutsch. \\
Heute Abend & \textbf{lerne} & ich & Deutsch. \\
Nach dem Abendessen & \textbf{lerne} & ich zu Hause & Deutsch. \\
\bottomrule
\end{tabularx}
\end{center}

Only one element comes before the finite verb. If something other than the
subject occupies position one, the subject normally follows the verb:
\textit{Morgen \textbf{fahre} ich nach Berlin}, not
\textit{Morgen ich fahre nach Berlin}.

\section{The Satzklammer}

With separable verbs, compound tenses, modal verbs, the future, and the
passive, the finite verb opens the bracket and the remaining verb material
closes it.

\begin{center}
\small
\rowcolors{2}{referenceGreenBg}{white}
\begin{tabularx}{\textwidth}{@{}>{\bfseries}l l X l@{}}
\toprule
Construction & Left bracket & Middle field & Right bracket \\
\midrule
Separable verb & Ich \textbf{rufe} & dich heute Abend & \textbf{an}. \\
Perfekt & Ich \textbf{habe} & dich gestern & \textbf{angerufen}. \\
Modal verb & Ich \textbf{muss} & dich heute Abend & \textbf{anrufen}. \\
Futur I & Ich \textbf{werde} & dich morgen & \textbf{anrufen}. \\
Passive & Das Formular \textbf{wird} & heute & \textbf{ausgefüllt}. \\
\bottomrule
\end{tabularx}
\end{center}

Material such as objects, adverbs, and negation usually stands inside the
bracket. The right bracket should not be forgotten merely because the middle
field is long.

\section{Placement of \textit{gern(e)}, \textit{nicht}, and \textit{nichts}}

\subsection{\textit{gern} and \textit{gerne}}

\textit{Gern} and \textit{gerne} have the same meaning and are both standard.
They show that somebody likes doing an activity. In the neutral word order,
\textit{gern(e)} stands in the middle field, usually after the finite verb and
pronouns but before the final verb element.

\begin{center}
\small
\rowcolors{2}{referenceBlueBg}{white}
\begin{tabularx}{\textwidth}{@{}>{\bfseries}l X l@{}}
\toprule
Construction & Example & Meaning \\
\midrule
Simple verb & Ich lese \textbf{gern} Krimis. & I like reading crime novels. \\
Modal verb & Ich möchte \textbf{gerne} mitkommen. & I would like to come along. \\
Separable verb & Ich komme \textbf{gern} mit. & I am happy to come along. \\
Perfekt & Ich habe dort \textbf{gern} gearbeitet. & I liked working there. \\
Short answer & Kommst du mit? -- Ja, \textbf{gern}! & Yes, gladly. \\
\bottomrule
\end{tabularx}
\end{center}

The comparison forms are \textit{lieber} and \textit{am liebsten}:
\textit{Ich fahre lieber mit dem Zug};
\textit{Am liebsten bleibe ich zu Hause}.

\subsection{Choosing \textit{nicht}, \textit{nichts}, or \textit{kein}}

These forms do different grammatical jobs.

\begin{center}
\small
\rowcolors{2}{referenceGreenBg}{white}
\begin{tabularx}{\textwidth}{@{}>{\bfseries}l l X@{}}
\toprule
Form & Meaning and function & Example \\
\midrule
nicht & not; negates an action, quality, or particular element
      & \textit{Das ist nicht teuer.} \\
nichts & nothing; stands in place of a noun or object
       & \textit{Ich sehe nichts.} \\
kein & no/not a; accompanies a noun and declines like \textit{ein}
     & \textit{Ich habe keine Zeit.} \\
\bottomrule
\end{tabularx}
\end{center}

To negate one particular element, put \textit{nicht} directly before it:
\textit{nicht heute}, \textit{nicht schnell}, \textit{nicht mit dem Bus},
\textit{nicht nach Berlin}. This position often creates a contrast:
\textit{Ich komme nicht heute, sondern morgen}.

To negate the action or clause as a whole, \textit{nicht} normally stands as
late as possible. It comes before the closing part of a sentence bracket.

\begin{center}
\small
\rowcolors{2}{referenceYellowBg}{white}
\begin{tabularx}{\textwidth}{@{}>{\bfseries}l X@{}}
\toprule
Construction & Example \\
\midrule
Simple verb & \textit{Ich arbeite heute \textbf{nicht}.} \\
Definite object & \textit{Ich kaufe das Auto \textbf{nicht}.} \\
Separable verb & \textit{Ich rufe ihn heute \textbf{nicht} an.} \\
Modal verb & \textit{Ich kann heute \textbf{nicht} kommen.} \\
Perfekt & \textit{Ich habe den Film \textbf{nicht} gesehen.} \\
\bottomrule
\end{tabularx}
\end{center}

\textit{Nichts} occupies an ordinary object position:
\textit{Ich weiß \textbf{nichts} darüber};
\textit{Ich habe \textbf{nichts} verstanden}. Use \textit{kein} with an
indefinite or article-free noun: \textit{Ich trinke \textbf{keinen} Kaffee};
\textit{Sie hat \textbf{keine} Kinder}.

The position changes the meaning:
\textit{Ich koche \textbf{nicht gern}} means \enquote{I do not like cooking}, while
\textit{Ich koche heute \textbf{nicht}} means \enquote{I am not cooking today}.
Other useful combinations are \textit{noch nicht} (not yet),
\textit{nicht mehr} (not any more), and \textit{gar nicht} (not at all).

\section{Subordinate clauses: the finite verb moves to the end}

After a subordinating conjunction such as \textit{weil, dass, obwohl}, or
\textit{wenn}, the finite verb normally stands at the end of its clause. A
separable verb is written as one word there.

\begin{center}
\small
\rowcolors{2}{referenceYellowBg}{white}
\begin{tabularx}{\textwidth}{@{}>{\bfseries}l X@{}}
\toprule
Construction & Subordinate clause \\
\midrule
Simple verb & \ldots{}, weil er heute lange \textbf{arbeitet}. \\
Separable verb & \ldots{}, weil er mich später \textbf{anruft}. \\
Perfekt & \ldots{}, weil er mich gestern \textbf{angerufen hat}. \\
Modal verb & \ldots{}, weil er morgen lange \textbf{arbeiten muss}. \\
Passive & \ldots{}, weil das Formular heute \textbf{ausgefüllt wird}. \\
\bottomrule
\end{tabularx}
\end{center}

Relative clauses follow the same verb-final pattern:
\textit{Das ist die Frau, die mir gestern \textbf{geholfen hat}.}

When the subordinate clause comes first, the complete clause occupies
position one. The finite verb of the following main clause therefore comes
immediately after the comma:

\begin{center}
\textit{Wenn ich Zeit \textbf{habe}, \textbf{rufe} ich dich an.}
\end{center}

\section{Questions and commands}

Yes/no questions and commands normally begin with the finite verb. A question
word, however, occupies position one in an information question.

\begin{center}
\small
\rowcolors{2}{referenceRedBg}{white}
\begin{tabularx}{\textwidth}{@{}>{\bfseries}l l X@{}}
\toprule
Type & Pattern & Example \\
\midrule
Yes/no question & finite verb + subject & \textbf{Kommst} du morgen? \\
Information question & question word + finite verb & Wann \textbf{kommst} du? \\
Command & imperative + remaining information & \textbf{Ruf} mich morgen \textbf{an}! \\
Formal command & verb + \textit{Sie} & \textbf{Rufen Sie} mich morgen \textbf{an}! \\
\bottomrule
\end{tabularx}
\end{center}

\begin{tcolorbox}[
  colback=referenceBlueBg,
  colframe=genderMasculine,
  title={Quick verb-position check},
  fonttitle=\bfseries
]
\begin{enumerate}[leftmargin=*,nosep]
  \item Find the finite verb.
  \item Decide whether the clause is a main clause, question, command, or
        subordinate clause.
  \item Check whether another verb part must close the sentence bracket.
  \item In a main-clause statement, make sure exactly one element precedes the
        finite verb.
\end{enumerate}
\end{tcolorbox}

\chapter{Conjunctions and Word Order}

German connectors do not all affect word order in the same way. The most
useful distinction is between coordinating conjunctions, subordinating
conjunctions, and connecting adverbs.

\begin{tcolorbox}[
  colback=black!3,
  colframe=black!55,
  title={Three patterns to remember},
  fonttitle=\bfseries
]
\begin{tabularx}{\textwidth}{@{}>{\bfseries}l X@{}}
Hauptsatz & subject $+$ \textbf{finite verb} $+$ remaining information \\
Nebensatz & conjunction $+$ subject $+$ remaining information $+$ \textbf{finite verb} \\
Connector adverb & adverb $+$ \textbf{finite verb} $+$ subject $+$ remaining information
\end{tabularx}
\end{tcolorbox}

\section{Coordinating conjunctions: Hauptsatz + Hauptsatz}

Coordinating conjunctions connect elements of equal rank. When they connect
two main clauses, they occupy ``position zero'': the normal verb-second word
order of the following Hauptsatz does not change.

\begin{center}
\small
\rowcolors{2}{referenceBlueBg}{white}
\begin{tabularx}{\textwidth}{@{}>{\bfseries}l l l X@{}}
\toprule
German & English & Function & Example \\
\midrule
und & and & addition & Ich koche, und er \textbf{deckt} den Tisch. \\
oder & or & alternative & Wir fahren heute, oder wir \textbf{fahren} morgen. \\
aber & but & contrast & Sie ist müde, aber sie \textbf{arbeitet} weiter. \\
denn & because/for & reason & Ich bleibe zu Hause, denn ich \textbf{bin} krank. \\
sondern & but rather & correction after a negative & Er fährt nicht, sondern er \textbf{nimmt} den Zug. \\
sowie & as well as & addition & Sie spricht Deutsch sowie Englisch. \\
\bottomrule
\end{tabularx}
\end{center}

\noindent
\textbf{Pattern:}
\quad Hauptsatz, $+$ conjunction $+$ subject $+$ \textbf{finite verb} $+$ \ldots

\section{Subordinating conjunctions: Hauptsatz + Nebensatz}

A subordinating conjunction introduces a dependent clause. The finite verb
moves to the end of that Nebensatz, and a comma separates it from the main
clause.

\subsection{Content, reason, condition, and contrast}

\begin{center}
\small
\rowcolors{2}{referenceGreenBg}{white}
\begin{tabularx}{\textwidth}{@{}>{\bfseries}l l l X@{}}
\toprule
German & English & Function & Example \\
\midrule
dass & that & statement/content & Ich weiß, dass er heute \textbf{kommt}. \\
ob & whether/if & indirect yes/no question & Sie fragt, ob der Laden offen \textbf{ist}. \\
weil & because & reason & Wir gehen nicht, weil es stark \textbf{regnet}. \\
da & since/because & reason, often already known & Da es spät \textbf{ist}, gehen wir nach Hause. \\
wenn & if/when & condition; repeated/present/future time & Ich komme, wenn ich Zeit \textbf{habe}. \\
falls & if/in case & possible condition & Ruf mich an, falls du Hilfe \textbf{brauchst}. \\
obwohl & although & concession/contrast & Er geht spazieren, obwohl es \textbf{regnet}. \\
während & while/whereas & simultaneous time or contrast & Sie liest, während er Musik \textbf{hört}. \\
\bottomrule
\end{tabularx}
\end{center}

\subsection{Time, purpose, result, and manner}

\begin{center}
\small
\rowcolors{2}{referenceGreenBg}{white}
\begin{tabularx}{\textwidth}{@{}>{\bfseries}l l l X@{}}
\toprule
German & English & Function & Example \\
\midrule
als & when & one event/period in the past & Als ich klein \textbf{war}, wohnte ich in Bonn. \\
bevor/ehe & before & earlier event & Wasch dir die Hände, bevor du \textbf{isst}. \\
nachdem & after & earlier completed event & Nachdem wir gegessen \textbf{hatten}, gingen wir spazieren. \\
seit/seitdem & since & starting point continuing until now & Seitdem sie hier \textbf{wohnt}, spricht sie mehr Deutsch. \\
bis & until & end point & Warte, bis der Bus \textbf{kommt}. \\
sobald & as soon as & immediate sequence & Ich rufe an, sobald ich zu Hause \textbf{bin}. \\
solange & as long as & duration/condition & Bleib hier, solange du \textbf{möchtest}. \\
damit & so that & purpose & Ich spreche langsam, damit alle mich \textbf{verstehen}. \\
sodass/so dass & so that/as a result & consequence & Es schneite, sodass die Straßen glatt \textbf{wurden}. \\
indem & by/by means of & manner or method & Man lernt, indem man regelmäßig \textbf{übt}. \\
ohne dass & without & absent accompanying action & Er ging, ohne dass er sich \textbf{verabschiedete}. \\
\bottomrule
\end{tabularx}
\end{center}

\subsection{When the Nebensatz comes first}

The complete Nebensatz occupies position one of the following Hauptsatz.
Consequently, the finite verb of the Hauptsatz comes immediately after the
comma.

\begin{center}
\rowcolors{2}{referenceGreenBg}{white}
\begin{tabularx}{\textwidth}{@{}l X@{}}
\toprule
Order & Example \\
\midrule
HS + NS & Ich bleibe heute zu Hause, weil ich krank \textbf{bin}. \\
NS + HS & Weil ich krank \textbf{bin}, \textbf{bleibe} ich heute zu Hause. \\
\bottomrule
\end{tabularx}
\end{center}

\section{Connecting adverbs: Hauptsatz with inversion}

Words such as \textit{deshalb} and \textit{trotzdem} are adverbs rather than
conjunctions. The following clause remains a Hauptsatz. When the adverb stands
in position one, the finite verb must stand directly after it in position two,
before the subject.

\begin{center}
\small
\rowcolors{2}{referenceYellowBg}{white}
\begin{tabularx}{\textwidth}{@{}>{\bfseries}l l l X@{}}
\toprule
German & English & Function & Example \\
\midrule
deshalb/deswegen & therefore/for that reason & result & Ich bin krank; deshalb \textbf{bleibe} ich zu Hause. \\
trotzdem & nevertheless/even so & contrast & Es regnet; trotzdem \textbf{gehen} wir spazieren. \\
außerdem & moreover/in addition & addition & Sie spricht Deutsch; außerdem \textbf{lernt} sie Spanisch. \\
dann & then & sequence/result & Wir essen; dann \textbf{gehen} wir ins Kino. \\
danach & afterwards & sequence & Der Kurs endet; danach \textbf{fahre} ich nach Hause. \\
sonst & otherwise & alternative/consequence & Beeil dich; sonst \textbf{verpassen} wir den Bus. \\
jedoch & however & contrast & Er war müde; jedoch \textbf{arbeitete} er weiter. \\
\bottomrule
\end{tabularx}
\end{center}

\noindent
\textbf{Compare:}
\quad \textit{weil ich krank \textbf{bin}} (Nebensatz)
\quad but \quad
\textit{deshalb \textbf{bleibe} ich zu Hause} (Hauptsatz).

\section{Paired conjunctions}

\begin{center}
\small
\rowcolors{2}{referenceRedBg}{white}
\begin{tabularx}{\textwidth}{@{}>{\bfseries}l l X@{}}
\toprule
German & English & Example \\
\midrule
sowohl \ldots{} als auch & both \ldots{} and & Sie spricht sowohl Deutsch als auch Englisch. \\
entweder \ldots{} oder & either \ldots{} or & Entweder kochen wir, oder wir bestellen eine Pizza. \\
weder \ldots{} noch & neither \ldots{} nor & Er trinkt weder Kaffee noch Tee. \\
nicht nur \ldots{} sondern auch & not only \ldots{} but also & Sie ist nicht nur klug, sondern auch geduldig. \\
zwar \ldots{} aber & admittedly \ldots{} but & Der Weg ist zwar lang, aber er ist schön. \\
je \ldots{} desto/umso & the \ldots{} the & Je mehr du übst, desto leichter \textbf{wird} es. \\
\bottomrule
\end{tabularx}
\end{center}

With \textit{je \ldots{} desto/umso}, the \textit{je}-clause has subordinate
word order, while the \textit{desto/umso}-clause has main-clause word order:
\textit{Je mehr du \textbf{übst}, desto leichter \textbf{wird} es.}

\section{Quick decision guide}

\begin{enumerate}[leftmargin=*]
  \item If the connector is \textit{und, oder, aber, denn}, or \textit{sondern},
        keep normal Hauptsatz word order after it.
  \item If it is a subordinating conjunction such as \textit{weil, dass,
        obwohl}, or \textit{wenn}, put the finite verb at the end of its clause.
  \item If it is a connecting adverb such as \textit{deshalb, trotzdem}, or
        \textit{außerdem}, put the finite verb immediately after the connector.
  \item Separate a Nebensatz from its Hauptsatz with a comma.
\end{enumerate}

\chapter{Partizip II and the Perfect Tenses}

The Perfekt combines a present-tense form of \textit{haben} or \textit{sein}
with the Partizip II. The auxiliary is finite and normally occupies position
two; the participle closes the sentence bracket.

\begin{center}
\textit{Ich \textbf{habe} den Brief gestern \textbf{geschrieben}.}
\qquad
\textit{Wir \textbf{sind} früh \textbf{abgefahren}.}
\end{center}

\section{Forming the Partizip II}

The following patterns cover most verbs. Strong and mixed forms must still be
learned as principal parts because their vowel or stem may change.

\begin{center}
\small
\rowcolors{2}{referenceBlueBg}{white}
\begin{tabularx}{\textwidth}{@{}>{\bfseries}l l l X@{}}
\toprule
Verb type & Pattern & Infinitive & Partizip II \\
\midrule
Weak & \textit{ge-} + stem + \textit{-t} & machen & gemacht \\
Weak in \textit{-d/-t} & \textit{ge-} + stem + \textit{-et} & arbeiten & gearbeitet \\
Strong & often \textit{ge-} + changed stem + \textit{-en} & schreiben & geschrieben \\
Mixed & stem change + weak ending & bringen & gebracht \\
Separable & prefix + \textit{ge-} + verb & anrufen & angerufen \\
Inseparable prefix & no \textit{ge-} & verstehen & verstanden \\
Verb in \textit{-ieren} & no \textit{ge-}, ending \textit{-t} & studieren & studiert \\
\bottomrule
\end{tabularx}
\end{center}

The common inseparable prefixes are \textit{be-, emp-, ent-, er-, ge-, miss-,
ver-}, and \textit{zer-}: \textit{bezahlt, empfohlen, entdeckt, erzählt,
gefallen, missverstanden, verkauft, zerstört}.

\section{Choosing \textit{haben} or \textit{sein}}

Most verbs form the Perfekt with \textit{haben}. This includes transitive
verbs with an accusative object, reflexive verbs, and verbs describing an
activity or state without a change of location or condition.

\begin{center}
\small
\rowcolors{2}{referenceGreenBg}{white}
\begin{tabularx}{\textwidth}{@{}>{\bfseries}l X X@{}}
\toprule
Auxiliary & Typical use & Examples \\
\midrule
haben & most verbs; transitive, reflexive, and modal verbs
      & \textit{Ich habe gelesen. Sie hat sich beeilt.} \\
sein & change of location or state; also \textit{sein, werden}, and
       \textit{bleiben}
     & \textit{Er ist angekommen. Es ist eingeschlafen.} \\
\bottomrule
\end{tabularx}
\end{center}

Meaning can affect the auxiliary. Compare
\textit{Ich \textbf{bin} nach Berlin gefahren} (movement to a destination)
with \textit{Ich \textbf{habe} den Wagen gefahren} (transitive: driving the
vehicle). Where accepted usage varies, the individual verb table shows both
auxiliaries.

\section{Word order in the Perfekt}

\begin{center}
\small
\rowcolors{2}{referenceYellowBg}{white}
\begin{tabularx}{\textwidth}{@{}>{\bfseries}l X@{}}
\toprule
Clause type & Example \\
\midrule
Main clause & Ich \textbf{habe} das Buch \textbf{gelesen}. \\
Question & \textbf{Hast} du das Buch \textbf{gelesen}? \\
Subordinate clause & \ldots{}, weil ich das Buch \textbf{gelesen habe}. \\
With negation & Ich \textbf{habe} das Buch nicht \textbf{gelesen}. \\
\bottomrule
\end{tabularx}
\end{center}

\section{The Plusquamperfekt}

The Plusquamperfekt places an event before another past event. It uses the
Präteritum of the auxiliary, \textit{hatte} or \textit{war}, plus the same
Partizip II used in the Perfekt.

\begin{center}
\textit{Nachdem ich \textbf{gegessen hatte}, ging ich spazieren.}

\textit{Sie \textbf{war schon abgefahren}, als ich anrief.}
\end{center}

\section{Modal verbs in the Perfekt}

When a modal verb accompanies another infinitive, German normally uses two
infinitives instead of the modal participle. This is called the
\textit{Ersatzinfinitiv}.

\begin{center}
\textit{Ich habe lange \textbf{arbeiten müssen}.}
\qquad
not \quad \textit{Ich habe lange arbeiten gemusst.}
\end{center}

When the modal stands alone, its normal participle is possible:
\textit{Ich habe das nicht gewollt.} In everyday speech, the Präteritum is
often simpler and more natural with modal verbs:
\textit{Ich musste lange arbeiten.}

The Ersatzinfinitiv also creates an important subordinate-clause exception:
the finite auxiliary normally precedes the two infinitives:
\textit{\ldots{}, weil ich lange \textbf{habe arbeiten müssen}}.

\begin{tcolorbox}[
  colback=referenceRedBg,
  colframe=genderFeminine,
  title={Three forms to learn together},
  fonttitle=\bfseries
]
For an irregular verb, learn the infinitive, third-person Präteritum, and
complete Perfekt form together:

\begin{center}
\textbf{gehen -- ging -- ist gegangen}
\qquad
\textbf{sehen -- sah -- hat gesehen}
\end{center}
\end{tcolorbox}

\chapter{Konjunktiv II}

Konjunktiv II presents something as hypothetical, unreal, wished for, or
politely distanced. It is especially common in requests, advice, wishes, and
conditional sentences.

\section{The most important forms}

The forms of \textit{sein}, \textit{haben}, and the modal verbs should be
learned directly. They are generally preferred to constructions with
\textit{würde}.

\begin{center}
\small
\rowcolors{2}{referenceBlueBg}{white}
\begin{tabular}{@{}lccc@{}}
\toprule
 & sein & haben & werden \\
\midrule
ich & wäre & hätte & würde \\
du & wärst & hättest & würdest \\
er/sie/es & wäre & hätte & würde \\
wir & wären & hätten & würden \\
ihr & wärt & hättet & würdet \\
sie/Sie & wären & hätten & würden \\
\bottomrule
\end{tabular}
\end{center}

\begin{center}
\small
\rowcolors{2}{referenceGreenBg}{white}
\begin{tabularx}{\textwidth}{@{}>{\bfseries}l l X@{}}
\toprule
Infinitive & Konjunktiv II & Typical meaning \\
\midrule
können & könnte & could / would be able to \\
dürfen & dürfte & might / would be allowed to \\
müssen & müsste & would have to / should probably \\
sollen & sollte & should / ought to \\
wollen & wollte & wanted / would want to \\
mögen & möchte & would like \\
\bottomrule
\end{tabularx}
\end{center}

Many strong verbs form Konjunktiv II from the Präteritum stem, often with an
umlaut and an ending: \textit{kam -- käme}, \textit{fand -- fände},
\textit{ging -- ginge}, \textit{wusste -- wüsste}. In ordinary speech,
\textit{würde + infinitive} is common instead: \textit{ich würde kommen}.
Weak verbs have forms identical to their Präteritum, so \textit{würde} usually
makes the hypothetical meaning clearer: \textit{ich würde arbeiten}.

\section{Polite requests and suggestions}

Konjunktiv II makes a request less direct.

\begin{center}
\small
\rowcolors{2}{referenceYellowBg}{white}
\begin{tabularx}{\textwidth}{@{}>{\bfseries}l X@{}}
\toprule
Function & Example \\
\midrule
Request & \textit{Könnten Sie mir bitte helfen?} \\
Wish & \textit{Ich hätte gern einen Kaffee.} \\
Preference & \textit{Ich möchte lieber zu Hause bleiben.} \\
Advice & \textit{Du solltest früher schlafen gehen.} \\
Suggestion & \textit{Wir könnten mit dem Zug fahren.} \\
\bottomrule
\end{tabularx}
\end{center}

\section{Unreal conditions in the present or future}

Both clauses normally use Konjunktiv II\@. If the \textit{wenn}-clause comes
first, the finite verb of the main clause follows the comma immediately.

\begin{center}
\textit{Wenn ich mehr Zeit \textbf{hätte}, \textbf{würde} ich öfter lesen.}

\textit{Ich \textbf{käme} mit, wenn ich nicht arbeiten \textbf{müsste}.}
\end{center}

The \textit{wenn}-clause can sometimes be omitted when the condition is clear:
\textit{An deiner Stelle würde ich mit dem Arzt sprechen.}

\section{Unreal situations in the past}

Past Konjunktiv II uses \textit{hätte} or \textit{wäre} plus Partizip II\@. The
choice of auxiliary is the same as in the Perfekt.

\begin{center}
\small
\rowcolors{2}{referenceRedBg}{white}
\begin{tabularx}{\textwidth}{@{}>{\bfseries}l X@{}}
\toprule
Auxiliary & Example \\
\midrule
hätte + Partizip II
  & \textit{Wenn ich gelernt hätte, hätte ich die Prüfung bestanden.} \\
wäre + Partizip II
  & \textit{Wenn wir früher losgefahren wären, wären wir pünktlich angekommen.} \\
\bottomrule
\end{tabularx}
\end{center}

This construction describes a past condition or result that did not occur.
At B1, form the ordinary unreal past with
\textit{hätte/wäre + Partizip II}. A sequence such as
\textit{würde gemacht} is incomplete; more complex forms such as
\textit{würde gemacht haben} exist but are not needed for this pattern.

\section{Wishes and comparisons}

\begin{itemize}[leftmargin=*]
  \item Present wish: \textit{Wenn ich doch mehr Zeit hätte!}
  \item Past regret: \textit{Wäre ich doch früher gekommen!}
  \item Unreal comparison: \textit{Er tut so, als ob er alles wüsste.}
\end{itemize}

\begin{tcolorbox}[
  colback=referenceBlueBg,
  colframe=genderMasculine,
  title={Choosing the form},
  fonttitle=\bfseries
]
Use \textit{wäre}, \textit{hätte}, and the common modal forms directly. Use a
familiar one-word form such as \textit{käme} or \textit{ginge} when it sounds
natural; otherwise use \textit{würde + infinitive}. For an unreal past, use
\textit{hätte/wäre + Partizip II}.
\end{tcolorbox}

\chapter{The Passive Voice}

The active voice focuses on the person or thing performing an action. The
passive focuses on the action or its result. The performer can be omitted when
it is unknown, obvious, or unimportant.

\begin{center}
\textit{Die Mechanikerin repariert das Auto.}
\quad $\longrightarrow$ \quad
\textit{Das Auto wird repariert.}
\end{center}

\section{Process passive: \textit{werden + Partizip II}}

The process passive (\textit{Vorgangspassiv}) describes an action or process.
The finite form of \textit{werden} carries the tense, and the Partizip II
stands at the end.

\begin{center}
\small
\rowcolors{2}{referenceBlueBg}{white}
\begin{tabularx}{\textwidth}{@{}>{\bfseries}l l X@{}}
\toprule
Tense & Formation & Example \\
\midrule
Präsens & wird + Partizip II & Das Auto \textbf{wird repariert}. \\
Präteritum & wurde + Partizip II & Das Auto \textbf{wurde repariert}. \\
Perfekt & ist + Partizip II + worden & Das Auto \textbf{ist repariert worden}. \\
Plusquamperfekt & war + Partizip II + worden & Das Auto \textbf{war repariert worden}. \\
Futur I & wird + Partizip II + werden & Das Auto \textbf{wird repariert werden}. \\
Modal verb & modal + Partizip II + werden & Das Auto \textbf{muss repariert werden}. \\
\bottomrule
\end{tabularx}
\end{center}

In the passive Perfekt, use \textbf{worden}, not \textit{geworden}:
\textit{Das Haus ist gebaut worden}. By contrast, \textit{geworden} is the
Partizip II of \textit{werden} meaning ``to become'':
\textit{Er ist müde geworden}.

\section{What happens to the objects?}

An accusative object in the active sentence becomes the nominative subject of
the passive sentence. A dative object remains dative.

\begin{center}
\small
\rowcolors{2}{referenceGreenBg}{white}
\begin{tabularx}{\textwidth}{@{}>{\bfseries}l X@{}}
\toprule
Active & Passive \\
\midrule
Man repariert \textbf{das Auto}.
  & \textbf{Das Auto} wird repariert. \\
Man gibt \textbf{dem Gast} \textbf{den Schlüssel}.
  & \textbf{Der Schlüssel} wird \textbf{dem Gast} gegeben. \\
Man hilft \textbf{dem Mann}.
  & \textbf{Dem Mann} wird geholfen. \\
\bottomrule
\end{tabularx}
\end{center}

The final example has no nominative subject because \textit{helfen} governs
the dative. This subjectless passive is normal German.

\section{Naming the performer or cause}

Use \textit{von + Dativ} for a person, institution, or originator. Use
\textit{durch + Akkusativ} especially for a means, cause, or process.

\begin{itemize}[leftmargin=*]
  \item \textit{Der Brief wurde \textbf{von der Chefin} geschrieben.}
  \item \textit{Die Straße wurde \textbf{durch den Sturm} blockiert.}
\end{itemize}

\section{State passive: \textit{sein + Partizip II}}

The state passive (\textit{Zustandspassiv}) describes the condition resulting
from an action, not the action itself.

\begin{center}
\small
\rowcolors{2}{referenceYellowBg}{white}
\begin{tabularx}{\textwidth}{@{}>{\bfseries}l l X@{}}
\toprule
Meaning & Formation & Example \\
\midrule
Process & werden + Partizip II & Die Tür \textbf{wird geschlossen}. \\
Resulting state & sein + Partizip II & Die Tür \textbf{ist geschlossen}. \\
Past process & wurde + Partizip II & Die Tür \textbf{wurde geschlossen}. \\
Past state & war + Partizip II & Die Tür \textbf{war geschlossen}. \\
\bottomrule
\end{tabularx}
\end{center}

\section{Impersonal passive}

An activity can be presented without naming any performer. Initial
\textit{es} disappears when another element occupies position one.

\begin{center}
\textit{Es wird heute getanzt.}
\qquad
\textit{Heute wird getanzt.}

\textit{Im Unterricht wird viel gesprochen.}
\end{center}

\section{Alternatives to the passive}

German often uses \textit{man} when the performer is general, or
\textit{sich lassen + infinitive} to express possibility.

\begin{itemize}[leftmargin=*]
  \item Passive: \textit{Das Problem kann gelöst werden.}
  \item With \textit{man}: \textit{Man kann das Problem lösen.}
  \item With \textit{sich lassen}: \textit{Das Problem lässt sich lösen.}
\end{itemize}

\begin{tcolorbox}[
  colback=referenceRedBg,
  colframe=genderFeminine,
  title={Word-order reminder},
  fonttitle=\bfseries
]
In a subordinate clause, the finite passive auxiliary moves to the end:
\textit{Ich weiß, dass das Auto heute \textbf{repariert wird}.}
With a modal verb: \textit{Ich weiß, dass das Auto heute
\textbf{repariert werden muss}.}
\end{tcolorbox}

\chapter{Infinitive Constructions}

An infinitive can appear without \textit{zu} or as part of a construction with
\textit{zu}. The position of \textit{zu}, the understood subject, and the
comma are the most important points for learners.

\section{The infinitive without \textit{zu}}

No \textit{zu} is used after modal verbs or \textit{werden}, nor in several
common constructions with verbs of perception, movement, and causation.

\begin{center}
\small
\rowcolors{2}{referenceBlueBg}{white}
\begin{tabularx}{\textwidth}{@{}>{\bfseries}l X@{}}
\toprule
Construction & Example \\
\midrule
Modal verb & \textit{Ich muss heute arbeiten.} \\
Futur with \textit{werden} & \textit{Wir werden morgen kommen.} \\
\textit{lassen} & \textit{Lass mich bitte ausreden.} \\
Perception & \textit{Ich höre die Kinder spielen.} \\
Movement + activity & \textit{Wir gehen heute schwimmen.} \\
\textit{lernen} & \textit{Das Kind lernt lesen.} \\
\bottomrule
\end{tabularx}
\end{center}

\section{The infinitive with \textit{zu}}

Many verbs, adjectives, nouns, and impersonal expressions introduce a
\textit{zu}-infinitive.

\begin{center}
\small
\rowcolors{2}{referenceGreenBg}{white}
\begin{tabularx}{\textwidth}{@{}>{\bfseries}l X@{}}
\toprule
Trigger & Example \\
\midrule
Verb & \textit{Ich versuche, jeden Tag Deutsch zu lernen.} \\
Adjective & \textit{Es ist wichtig, genug zu schlafen.} \\
Noun & \textit{Ich habe den Wunsch, in Berlin zu studieren.} \\
Expression with anticipatory \textit{es}
  & \textit{Ich finde es schwer, früh aufzustehen.} \\
\bottomrule
\end{tabularx}
\end{center}

The understood subject of the infinitive is usually supplied by the main
clause. Compare:

\begin{itemize}[leftmargin=*]
  \item \textit{Ich hoffe, bald zu reisen.} --- \textit{ich} will travel.
  \item \textit{Ich bitte Paul, bald zu kommen.} --- \textit{Paul} should come.
\end{itemize}

If the two actions have different explicit subjects, use a finite subordinate
clause instead: \textit{Ich hoffe, dass Paul bald kommt.}

\section{Where \textit{zu} goes}

\begin{center}
\small
\rowcolors{2}{referenceYellowBg}{white}
\begin{tabularx}{\textwidth}{@{}>{\bfseries}l l l X@{}}
\toprule
Verb type & Infinitive & With \textit{zu} & Example \\
\midrule
Simple & lernen & zu lernen & \textit{Ich versuche, Deutsch zu lernen.} \\
Inseparable & verstehen & zu verstehen & \textit{Es ist schwer, das zu verstehen.} \\
Verb in \textit{-ieren} & reservieren & zu reservieren & \textit{Ich vergaß, zu reservieren.} \\
Separable & anrufen & anzurufen & \textit{Ich verspreche, dich anzurufen.} \\
Separable & aufstehen & aufzustehen & \textit{Ich versuche, früh aufzustehen.} \\
\bottomrule
\end{tabularx}
\end{center}

With a separable verb, \textit{zu} is inserted between the prefix and the verb
stem: \textit{an + zu + rufen = anzurufen}.

\section{Purpose, absence, and alternatives}

\begin{center}
\small
\rowcolors{2}{referenceRedBg}{white}
\begin{tabularx}{\textwidth}{@{}>{\bfseries}l l X@{}}
\toprule
Construction & Meaning & Example \\
\midrule
um \ldots{} zu & in order to & \textit{Ich lerne viel, um die Prüfung zu bestehen.} \\
ohne \ldots{} zu & without doing & \textit{Er ging, ohne sich zu verabschieden.} \\
statt/anstatt \ldots{} zu & instead of doing
  & \textit{Sie fährt Rad, statt das Auto zu nehmen.} \\
\bottomrule
\end{tabularx}
\end{center}

Use \textit{um \ldots{} zu} when the understood subject is the same in both
parts. Use \textit{damit} plus a finite clause when the subjects differ:

\begin{center}
\textit{Ich spreche langsam, um verständlich zu sein.}

\textit{Ich spreche langsam, damit die Kinder mich verstehen.}
\end{center}

\section{\textit{brauchen + zu}}

In standard usage, \textit{brauchen + zu} normally occurs with a negative or
limiting expression and has a modal meaning:

\begin{center}
\textit{Du brauchst heute nicht zu kommen.}
\qquad
\textit{Du brauchst nur anzurufen.}
\end{center}

\section{Perfect infinitive}

To show that the infinitive action happened earlier, use Partizip II plus
\textit{zu haben} or \textit{zu sein}. The auxiliary follows the same choice
as in the Perfekt.

\begin{itemize}[leftmargin=*]
  \item \textit{Sie behauptet, den Brief geschrieben \textbf{zu haben}.}
  \item \textit{Er scheint früh angekommen \textbf{zu sein}.}
\end{itemize}

\section{Commas with infinitive groups}

Under the current official spelling rules, an expanded, clause-like
infinitive group is separated by a comma:
\textit{Sie verspricht, morgen pünktlich zu kommen.} A comma is also required
with \textit{um, ohne, statt, anstatt, außer}, and \textit{als}, and when the
group is inserted into a sentence it is enclosed by a pair of commas.

A bare, unexpanded \textit{zu}-infinitive can be written without a comma:
\textit{Sie verspricht zu kommen.} No comma separates an integrated infinitive
construction with verbs such as \textit{brauchen} or \textit{scheinen}:
\textit{Du brauchst nicht zu warten}; \textit{Er scheint zu schlafen}.

\begin{tcolorbox}[
  colback=referenceBlueBg,
  colframe=genderMasculine,
  title={Quick choice},
  fonttitle=\bfseries
]
\begin{itemize}[leftmargin=*,nosep]
  \item After a modal verb: infinitive without \textit{zu}.
  \item Same subject and a purpose: \textit{um \ldots{} zu}.
  \item Different subjects: usually \textit{damit} or another finite clause.
  \item Separable verb: put \textit{zu} between prefix and stem.
  \item Expanded clause-like infinitive group: separate it with a comma.
\end{itemize}
\end{tcolorbox}

\chapter{Strong-Verb Stem Changes}

\section{What do regular, irregular, weak, and strong mean?}

The terms describe two related ideas. \textbf{Weak} and \textbf{strong}
describe how a verb forms its past tenses, whereas \textbf{regular} and
\textbf{irregular} describe how predictable its forms are. In this book,
\textit{regular} normally means a weak verb that follows the expected pattern.

\begin{center}
\small
\rowcolors{2}{referenceBlueBg}{white}
\begin{tabularx}{\textwidth}{@{}>{\bfseries}p{0.17\textwidth} X p{0.31\textwidth}@{}}
\toprule
Label & Meaning & Example forms \\
\midrule
Weak / regular & Adds \textit{-te} in the Präteritum and normally
ends the Partizip II in \textit{-t}; there is no unpredictable stem-vowel
change. & \textit{machen -- machte -- gemacht} \\
Strong & Changes its stem vowel in the Präteritum and normally ends the
Partizip II in \textit{-en}. Its key forms must be learnt. &
\textit{sprechen -- sprach -- gesprochen} \\
Mixed & Combines a stem change with the weak endings \textit{-te} and
\textit{-t}. & \textit{denken -- dachte -- gedacht} \\
Irregular & An umbrella label for a verb with forms that cannot all be
predicted from the normal weak pattern. Strong and mixed verbs are therefore
irregular, as are special verbs such as \textit{sein}. &
\textit{sein -- war -- gewesen} \\
\bottomrule
\end{tabularx}
\end{center}

\begin{tcolorbox}[
  colback=referenceYellowBg,
  colframe=genderPlural,
  title={How to read the labels in the verb tables},
  fonttitle=\bfseries
]
\textbf{Regular and weak usually mean the same thing in this book.}
\textbf{Strong verbs are irregular}, but not every irregular verb is strong:
mixed verbs and highly irregular verbs form other groups. A label such as
\textit{reflexive}, \textit{separable}, or \textit{+ Dativ} gives different
information and does not determine whether the verb is weak or strong.

Some regular patterns have predictable details: verbs ending in
\textit{-ieren} form the Partizip II without \textit{ge-}
(\textit{studiert}), and separable verbs place \textit{ge-} after the prefix
(\textit{aufgemacht}).
\end{tcolorbox}

Strong verbs form the Präteritum by changing the stem vowel rather than by
adding \textit{-te}. Their Partizip II normally ends in \textit{-en}. Some also
change their vowel in the second- and third-person singular Präsens. Because
the changes are not completely predictable, learn each verb through its key
forms:

\begin{tcolorbox}[
  colback=referenceBlueBg,
  colframe=genderMasculine,
  title={Recommended learning pattern},
  fonttitle=\bfseries
]
\centering
\textbf{Infinitiv} $\longrightarrow$ \textbf{er/sie/es im Präsens}
$\longrightarrow$ \textbf{Präteritum} $\longrightarrow$ \textbf{Perfekt}

\medskip
\textit{sprechen $\longrightarrow$ spricht $\longrightarrow$ sprach
$\longrightarrow$ hat gesprochen}
\end{tcolorbox}

\section{Learning shortcuts}

\begin{itemize}[leftmargin=*]
  \item Learn the third-person singular Präsens whenever it changes:
        \textit{nehmen -- er nimmt}, not merely \textit{nehmen}.
  \item Learn the Präteritum and Perfekt together:
        \textit{sprechen -- sprach -- hat gesprochen}.
  \item Mark verbs that normally use \textit{sein}, especially movement and
        change-of-state verbs such as \textit{gehen, kommen, bleiben}, and
        \textit{fallen}.
  \item Learn a prefixed verb with its family, but check whether the prefix is
        separable: \textit{aufstehen -- stand auf -- ist aufgestanden}, but
        \textit{verstehen -- verstand -- hat verstanden}.
\end{itemize}

\section{Strong verbs with a Präsens stem change}

The vowel change normally appears only with \textit{du} and
\textit{er/sie/es}. The plural forms retain the infinitive vowel:
\textit{ich fahre, du fährst, er fährt, wir fahren}.

\begin{center}
\small
\rowcolors{2}{referenceBlueBg}{white}
\begin{tabularx}{\textwidth}{@{}>{\bfseries}l l l l X@{}}
\toprule
Infinitive & English & er/sie/es & Präteritum & Perfekt \\
\midrule
empfehlen & recommend & empfiehlt & empfahl & hat empfohlen \\
essen & eat & isst & aß & hat gegessen \\
fahren & drive/travel & fährt & fuhr & ist/hat gefahren \\
fallen & fall & fällt & fiel & ist gefallen \\
fangen & catch & fängt & fing & hat gefangen \\
fressen & eat (animals) & frisst & fraß & hat gefressen \\
geben & give & gibt & gab & hat gegeben \\
gefallen & please & gefällt & gefiel & hat gefallen \\
gelten & be valid & gilt & galt & hat gegolten \\
halten & hold/stop & hält & hielt & hat gehalten \\
helfen & help & hilft & half & hat geholfen \\
laden & load/invite & lädt & lud & hat geladen \\
lassen & let/leave & lässt & ließ & hat gelassen \\
laufen & run/walk & läuft & lief & ist gelaufen \\
lesen & read & liest & las & hat gelesen \\
nehmen & take & nimmt & nahm & hat genommen \\
raten & advise/guess & rät & riet & hat geraten \\
schlafen & sleep & schläft & schlief & hat geschlafen \\
schlagen & hit & schlägt & schlug & hat geschlagen \\
sehen & see & sieht & sah & hat gesehen \\
sprechen & speak & spricht & sprach & hat gesprochen \\
stechen & sting/stab & sticht & stach & hat gestochen \\
tragen & carry/wear & trägt & trug & hat getragen \\
treffen & meet/hit & trifft & traf & hat getroffen \\
treten & step/kick & tritt & trat & ist/hat getreten \\
vergessen & forget & vergisst & vergaß & hat vergessen \\
waschen & wash & wäscht & wusch & hat gewaschen \\
werfen & throw & wirft & warf & hat geworfen \\
werden & become & wird & wurde & ist geworden \\
\bottomrule
\end{tabularx}
\end{center}

\section{Strong verbs without a Präsens stem change}

These verbs look regular in the Präsens but change their stem in the
Präteritum, the Partizip II, or both.

\begin{center}
\small
\rowcolors{2}{referenceGreenBg}{white}
\begin{tabularx}{\textwidth}{@{}>{\bfseries}l l l l X@{}}
\toprule
Infinitive & English & er/sie/es & Präteritum & Perfekt \\
\midrule
beginnen & begin & beginnt & begann & hat begonnen \\
beißen & bite & beißt & biss & hat gebissen \\
bieten & offer & bietet & bot & hat geboten \\
bitten & ask/request & bittet & bat & hat gebeten \\
bleiben & remain & bleibt & blieb & ist geblieben \\
finden & find & findet & fand & hat gefunden \\
fliegen & fly & fliegt & flog & ist/hat geflogen \\
fliehen & flee & flieht & floh & ist geflohen \\
gehen & go & geht & ging & ist gegangen \\
gießen & pour & gießt & goss & hat gegossen \\
genießen & enjoy & genießt & genoss & hat genossen \\
gewinnen & win & gewinnt & gewann & hat gewonnen \\
heben & lift & hebt & hob & hat gehoben \\
heißen & be called & heißt & hieß & hat geheißen \\
kommen & come & kommt & kam & ist gekommen \\
leihen & lend/borrow & leiht & lieh & hat geliehen \\
rufen & call & ruft & rief & hat gerufen \\
scheinen & shine/seem & scheint & schien & hat geschienen \\
schließen & close & schließt & schloss & hat geschlossen \\
schreiben & write & schreibt & schrieb & hat geschrieben \\
sitzen & sit & sitzt & saß & hat/ist gesessen \\
stehen & stand & steht & stand & hat/ist gestanden \\
steigen & climb/rise & steigt & stieg & ist gestiegen \\
streichen & paint/cancel & streicht & strich & hat gestrichen \\
trinken & drink & trinkt & trank & hat getrunken \\
tun & do & tut & tat & hat getan \\
verzeihen & forgive & verzeiht & verzieh & hat verziehen \\
ziehen & pull/move & zieht & zog & ist/hat gezogen \\
\bottomrule
\end{tabularx}
\end{center}

\section{Common prefixed families}

Many prefixed verbs preserve the stem change of their base verb. Inseparable
prefixes such as \textit{be-, emp-, ent-, er-, ver-}, and \textit{zer-} do not
add \textit{ge-} in the Partizip II.

\begin{center}
\small
\rowcolors{2}{referenceYellowBg}{white}
\begin{tabularx}{\textwidth}{@{}>{\bfseries}l l X@{}}
\toprule
Family & Pattern & Examples \\
\midrule
kommen & kam -- gekommen & bekommen: bekam -- bekommen \\
lassen & ließ -- gelassen & entlassen: entließ -- entlassen \\
leihen & lieh -- geliehen & verleihen: verlieh -- verliehen \\
nehmen & nahm -- genommen & übernehmen: übernahm -- übernommen; unternehmen: unternahm -- unternommen \\
schreiben & schrieb -- geschrieben & beschreiben: beschrieb -- beschrieben; unterschreiben: unterschrieb -- unterschrieben \\
sprechen & sprach -- gesprochen & widersprechen: widersprach -- widersprochen \\
stehen & stand -- gestanden & bestehen: bestand -- bestanden; entstehen: entstand -- entstanden; verstehen: verstand -- verstanden \\
halten & hielt -- gehalten & erhalten: erhielt -- erhalten; unterhalten: unterhielt -- unterhalten \\
ziehen & zog -- gezogen & beziehen: bezog -- bezogen \\
\bottomrule
\end{tabularx}
\end{center}

\section{Mixed verbs and modal verbs}

Mixed verbs combine a stem change with the weak \textit{-te} Präteritum and
the \textit{-t} Partizip ending. Modal verbs follow a closely related irregular
pattern and lose their umlaut in the Präteritum.

\begin{center}
\small
\rowcolors{2}{referenceRedBg}{white}
\begin{tabularx}{\textwidth}{@{}>{\bfseries}l l l l X@{}}
\toprule
Infinitive & English & er/sie/es & Präteritum & Perfekt \\
\midrule
bringen & bring & bringt & brachte & hat gebracht \\
denken & think & denkt & dachte & hat gedacht \\
kennen & know/be familiar with & kennt & kannte & hat gekannt \\
nennen & name/call & nennt & nannte & hat genannt \\
wissen & know a fact & weiß & wusste & hat gewusst \\
\addlinespace
dürfen & be allowed to & darf & durfte & hat gedurft \\
können & be able to & kann & konnte & hat gekonnt \\
mögen & like & mag & mochte & hat gemocht \\
müssen & have to & muss & musste & hat gemusst \\
sollen & should/be supposed to & soll & sollte & hat gesollt \\
wollen & want & will & wollte & hat gewollt \\
\bottomrule
\end{tabularx}
\end{center}


\part{Verb Reference}
\begin{tcolorbox}[colback=referenceYellowBg,colframe=genderPlural]
\textbf{Table convention.} A single em dash (---) under \textbf{IMPERATIV}
means that no normal imperative is shown. Under \textbf{MIT PRÄPOSITIONEN} or
\textbf{TRENNBARE VERBEN}, it means that the table intentionally lists no item
in that category; it is not an accidentally blank section. Separable compounds
are listed compactly beneath their base verb, with a vertical bar marking the
split point (for example, \textbf{aus$|$stellen}). Each compound remains in the
German and English indexes, which lead to the applicable stem table.
\end{tcolorbox}
% All verb records live in one canonical source and use the same renderer.
% SINGLE CANONICAL VERB SOURCE (consolidated 2026-09-25).
% Edit verb data here; style/germanbook.sty owns the shared layout.
% scripts/build_unified_verb_list.py is retained only as a migration audit.

\chapter{Auxiliary Verbs}

% Migrated from verbs/auxiliary.tex:1
\VerbTable{
  name = {sein},
  english = {to be},
  type = {irregular, + Nominativ},
  auxiliary = {sein},
  style = {nom},
  present = {bin,bist,ist,sind,seid,sind},
  preterite = {war,warst,war,waren,wart,waren},
  subjunctive = {wäre,wärest,wäre,wären,wäret,wären},
  imperative = {sei (du)!,seien wir!,seid (ihr)!,seien Sie!},
  perfect = {sein + \textbf{gewesen}},
  future = {werden + \textbf{sein}},
  pluperfect = {waren + \textbf{gewesen}},
  prepositions = {
    & \textbf{sein in} + Dat & (to be located in)\\
    & \textbf{sein auf} + Dat & (to be on / at)\\
    &\textbf{sein bei} + Dat & (to be with /at)\\
    &\textbf{sein von} + Dat & (to be from / made of)\\
    &\textbf{sein aus} + Dat & (to be from / made of)\\
    &\textbf{sein für} + Akk & (to be for / in favour of)\\
    &\textbf{sein mit} + Dat & (to be with / equipped with)
  },
  separable = {---}
}

% Migrated from verbs/auxiliary.tex:2
\VerbTable{
  name = {haben},
  english = {to have},
  type = {irregular, + Akkusativ},
  auxiliary = {haben},
  style = {acc},
  present = {habe,hast,hat,haben,habt,haben},
  preterite = {hatte,hattest,hatte,hatten,hattet,hatten},
  subjunctive = {hätte,hättest,hätte,hätten,hättet,hätten},
  imperative = {hab(e) (du)!,haben wir!,habt (ihr)!,haben Sie!},
  perfect = {haben + \textbf{gehabt}},
  future = {werden + \textbf{haben}},
  pluperfect = {hatten + \textbf{gehabt}},
  prepositions = {
    & \textbf{haben an} + Dat & (to have a problem / interest /fault)\\
    & \textbf{haben für} + Akk & (to have for / feel towards)\\
    &\textbf{haben von} + Dat & (to have heard / get from)\\
    &\textbf{haben auf} + Akk & (to be in the mood for)\\
    &\textbf{haben vor} + Dat & (to plan / intend)
  },
  separable = {---}
}

% Migrated from verbs/auxiliary.tex:3
\VerbTable{
  name = {werden},
  english = {to become},
  type = {irregular},
  auxiliary = {sein},
  style = {acc},
  present = {werde,wirst,wird,werden,werdet,werden},
  preterite = {wurde,wurdest,wurde,wurden,wurdet,wurden},
  subjunctive = {würde,würdest,würde,würden,würdet,würden},
  imperative = {werde (du)!,werden wir!,werdet (ihr)!,werden Sie!},
  perfect = {sein + \textbf{geworden}},
  future = {werden + \textbf{werden}},
  pluperfect = {waren + \textbf{geworden}},
  prepositions = {---},
  separable = {---}
}

% Migrated from verbs/auxiliary.tex:4
\VerbTable{
  name = {lassen},
  english = {to let},
  type = {irregular},
  auxiliary = {haben},
  style = {default},
  present = {lasse,lässt,lässt,lassen,lasst,lassen},
  preterite = {ließ,ließt,ließ,ließen,ließt,ließen},
  subjunctive = {ließe,ließest,ließe,ließen,ließet,ließen},
  imperative = {lass(e) (du)!,lassen wir!,lasst (ihr)!,lassen Sie!},
  perfect = {haben + \textbf{gelassen}},
  future = {werden + \textbf{lassen}},
  pluperfect = {hatten + \textbf{gelassen}},
  prepositions = {---},
  separable = {---}
}

\chapter{Modal Verbs}

% Migrated from verbs/modal.tex:1
\VerbTable{
  name = {dürfen},
  english = {to be allowed to},
  type = {Modal},
  auxiliary = {haben},
  style = {modal},
  present = {darf,darfst,darf,dürfen,dürft,dürfen},
  preterite = {durfte,durftest,durfte,durften,durftet,durften},
  subjunctive = {dürfte,dürftest,dürfte,dürften,dürftet,dürften},
  imperative = {---},
  perfect = {haben + \textbf{gedurft}},
  future = {werden + \textbf{dürfen}},
  pluperfect = {hatten + \textbf{gedurft}},
  prepositions = {---},
  separable = {---}
}

% Migrated from verbs/modal.tex:2
\VerbTable{
  name = {können},
  english = {to be able to},
  type = {Modal},
  auxiliary = {haben},
  style = {modal},
  present = {kann,kannst,kann,können,könnt,können},
  preterite = {konnte,konntest,konnte,konnten,konntet,konnten},
  subjunctive = {könnte,könntest,könnte,könnten,könntet,könnten},
  imperative = {---},
  perfect = {haben + \textbf{gekonnt}},
  future = {werden + \textbf{können}},
  pluperfect = {hatten + \textbf{gekonnt}},
  prepositions = {---},
  separable = {---}
}

% Migrated from verbs/modal.tex:3
\VerbTable{
  name = {mögen},
  english = {to like},
  type = {Modal},
  auxiliary = {haben},
  style = {modal},
  present = {mag,magst,mag,mögen,mögt,mögen},
  preterite = {mochte,mochtest,mochte,mochten,mochtet,mochten},
  subjunctive = {möchte,möchtest,möchte,möchten,möchtet,möchten},
  imperative = {---},
  perfect = {haben + \textbf{gemocht}},
  future = {werden + \textbf{mögen}},
  pluperfect = {hatten + \textbf{gemocht}},
  prepositions = {---},
  separable = {---}
}

% Migrated from verbs/modal.tex:4
\VerbTable{
  name = {müssen},
  english = {to have to},
  type = {Modal},
  auxiliary = {haben},
  style = {modal},
  present = {muss,musst,muss,müssen,müsst,müssen},
  preterite = {musste,musstest,musste,mussten,musstet,mussten},
  subjunctive = {müsste,müsstest,müsste,müssten,müsstet,müssten},
  imperative = {---},
  perfect = {haben + \textbf{gemusst}},
  future = {werden + \textbf{müssen}},
  pluperfect = {hatten + \textbf{gemusst}},
  prepositions = {---},
  separable = {---}
}

% Migrated from verbs/modal.tex:5
\VerbTable{
  name = {sollen},
  english = {should},
  type = {Modal},
  auxiliary = {haben},
  style = {modal},
  present = {soll,sollst,soll,sollen,sollt,sollen},
  preterite = {sollte,solltest,sollte,sollten,solltet,sollten},
  subjunctive = {sollte,solltest,sollte,sollten,solltet,sollten},
  imperative = {soll(e) du!,sollen wir!,sollt (ihr),sollen Sie!},
  perfect = {haben + \textbf{gesollt}},
  future = {werden + \textbf{sollen}},
  pluperfect = {hatten + \textbf{gesollt}},
  prepositions = {---},
  separable = {---}
}

% Migrated from verbs/modal.tex:6
\VerbTable{
  name = {wollen},
  english = {to want},
  type = {Modal},
  auxiliary = {haben},
  style = {modal},
  present = {will,willst,will,wollen,wollt,wollen},
  preterite = {wollte,wolltest,wollte,wollten,wolltet,wollten},
  subjunctive = {wollte,wolltest,wollte,wollten,wolltet,wollten},
  imperative = {wolle! (du),wollen wir!,wollt! (ihr),wollen Sie!},
  perfect = {haben + \textbf{gewollt}},
  future = {werden + \textbf{wollen}},
  pluperfect = {hatten + \textbf{gewollt}},
  prepositions = {---},
  separable = {---}
}

\chapter{Reflexive Verbs}

% Migrated from verbs/reflexive.tex:1
\VerbTable{
  name = {sich ändern},
  english = {to change},
  type = {regular, + Akkusativ},
  auxiliary = {haben},
  style = {acc},
  present = {ändere mich,änderst dich,ändert sich,ändern uns,ändert euch,ändern sich},
  preterite = {änderte mich,ändertest dich,änderte sich,änderten uns,ändertet euch,änderten sich},
  subjunctive = {},
  imperative = {änder(e) (du) dich!,ändern wir uns!,ändert (ihr) euch!,ändern Sie sich!},
  perfect = {haben + sich + \textbf{geändert}},
  future = {werden + sich + \textbf{ändern}},
  pluperfect = {hatten + sich + \textbf{geändert}},
  prepositions = {---},
  separable = {---}
}

% Migrated from verbs/reflexive.tex:2
\VerbTable{
  name = {sich ärgern},
  english = {to be annoyed},
  type = {regular, + Akkusativ},
  auxiliary = {haben},
  style = {acc},
  present = {ärgere mich,ärgerst dich,ärgert sich,ärgern uns,ärgert euch,ärgern sich},
  preterite = {ärgerte mich,ärgertest dich,ärgerte sich,ärgerten uns,ärgertet euch,ärgerten sich},
  subjunctive = {},
  imperative = {ärger(e) (du) dich!,ärgern wir uns!,ärgert (ihr) euch!,ärgern Sie sich!},
  perfect = {haben + sich + \textbf{geärgert}},
  future = {werden + sich + \textbf{ärgern}},
  pluperfect = {hatten + sich + \textbf{geärgert}},
  prepositions = {& \textbf{sich ärgern über} + Akk & (to be annoyed about)},
  separable = {---}
}

% Migrated from verbs/b1/a.tex:15
\VerbTable{
  name = {sich amüsieren},
  english = {to amuse},
  type = {regular},
  auxiliary = {haben},
  style = {default},
  present = {amüsiere mich,amüsierst dich,amüsiert sich,amüsieren uns,amüsiert euch,amüsieren sich},
  preterite = {amüsierte mich,amüsiertest dich,amüsierte sich,amüsierten uns,amüsiertet euch,amüsierten sich},
  subjunctive = {},
  imperative = {amüsier dich!,amüsieren wir uns!,amüsiert euch!,amüsieren Sie sich!},
  perfect = {haben + \textbf{sich amüsiert}},
  future = {werden + \textbf{sich amüsieren}},
  pluperfect = {hatten + \textbf{sich amüsiert}},
  prepositions = {---},
  separable = {---}
}

% Migrated from verbs/b1/a.tex:34
\VerbTable{
  name = {sich anstrengen},
  english = {to exert oneself},
  type = {regular},
  auxiliary = {haben},
  style = {default},
  present = {strenge mich an,strengst dich an,strengt sich an,strengen uns an,strengt euch an,strengen sich an},
  preterite = {strengte mich an,strengtest dich an,strengte sich an,strengten uns an,strengtet euch an,strengten sich an},
  subjunctive = {},
  imperative = {streng dich an!,strengen wir uns an!,strengt euch an!,strengen Sie sich an!},
  perfect = {haben + \textbf{sich angestrengt}},
  future = {werden + \textbf{sich anstrengen}},
  pluperfect = {hatten + \textbf{sich angestrengt}},
  prepositions = {---},
  separable = {---}
}

% Migrated from verbs/reflexive.tex:3
\VerbTable{
  name = {sich bedanken},
  english = {to thank},
  type = {regular},
  auxiliary = {haben},
  style = {default},
  present = {bedanke mich,bedankst dich,bedankt sich,bedanken uns,bedankt euch,bedanken sich},
  preterite = {bedankte mich,bedanktest dich,vbedankte sich,bedankten uns,bedanktet euch,bedankten sich},
  subjunctive = {},
  imperative = {bedank(e) (du) dich!,bedanken wir uns!,bedankt (ihr) euch!,bedanken Sie sich!},
  perfect = {haben + sich + \textbf{bedankt}},
  future = {werden + sich + \textbf{bedanken}},
  pluperfect = {hatten + sich + \textbf{bedankt}},
  prepositions = {
    & \textbf{sich bedanken für} + Akk & (to thank for)\\
    & \textbf{sich bedanken bei} + Dat & (to be thankful to someone)
  },
  separable = {---}
}

% Migrated from verbs/reflexive.tex:4
\VerbTable{
  name = {sich beeilen},
  english = {to hurry},
  type = {regular},
  auxiliary = {haben},
  style = {default},
  present = {beeile mich,beeilst dich,beeilt sich,beeilen uns,beeilt euch,beeilen sich},
  preterite = {beeilte mich,beeiltest dich,beeilte sich,beeilten uns,beeiltet euch,beeilten sich},
  subjunctive = {},
  imperative = {beeil(e) (du) dich!,beeilen wir uns!,beeilt (ihr) euch!,beeilen Sie sich!},
  perfect = {haben + sich + \textbf{beeilt}},
  future = {werden + sich + \textbf{beeilen}},
  pluperfect = {hatten + sich + \textbf{beeilt}},
  prepositions = {& \textbf{sich beeilen mit} + Dat & (to hurry up with)},
  separable = {---}
}

% Migrated from verbs/reflexive.tex:5
\VerbTable{
  name = {sich befinden},
  english = {to be located},
  type = {irregular},
  auxiliary = {haben},
  style = {default},
  present = {befinde mich,befindest dich,befindet sich,befinden uns,befindet euch,befinden sich},
  preterite = {befand mich,befandest dich,befand sich,befanden uns,befandet euch,befanden sich},
  subjunctive = {},
  imperative = {befind(e) (du) dich!,befinden wir uns!,befindet (ihr) euch!,befinden Sie sich!},
  perfect = {haben + sich + \textbf{befunden}},
  future = {werden + sich + \textbf{befinden}},
  pluperfect = {hatten + sich + \textbf{befunden}},
  prepositions = {& \textbf{sich befinden in} + Akk & (to be located in)},
  separable = {---}
}

% Migrated from verbs/reflexive.tex:6
\VerbTable{
  name = {sich bemühen},
  english = {to make an effort},
  type = {regular, reflexive},
  auxiliary = {haben},
  style = {default},
  present = {bemühe mich,bemühst dich,bemüht sich,bemühen uns,bemüht euch,bemühen sich},
  preterite = {bemühte mich,bemühtest dich,bemühte sich,bemühten uns,bemühtet euch,bemühten sich},
  subjunctive = {},
  imperative = {bemüh(e) dich!,bemühen wir uns!,bemüht euch!,bemühen Sie sich!},
  perfect = {haben + sich + \textbf{bemüht}},
  future = {werden + sich + \textbf{bemühen}},
  pluperfect = {hatten + sich + \textbf{bemüht}},
  prepositions = {
    & \textbf{sich bemühen um} + Akk & (to make an effort for / to seek)\\
    & \textbf{sich bemühen zu} + Inf & (to try to)\\
  },
  separable = {---}
}

% Migrated from verbs/reflexive.tex:7
\VerbTable{
  name = {sich beschäftigen},
  english = {to occupy oneself / to deal with},
  type = {regular},
  auxiliary = {haben},
  style = {default},
  present = {beschäftige mich,beschäftigst dich,beschäftigt sich,beschäftigen uns,beschäftigt euch,beschäftigen sich},
  preterite = {beschäftigte mich,beschäftigtest dich,beschäftigte sich,beschäftigten uns,beschäftigtet euch,beschäftigten sich},
  subjunctive = {},
  imperative = {beschäftige dich!,beschäftigen wir uns!,beschäftigt euch!,beschäftigen Sie sich!},
  perfect = {haben + sich + \textbf{beschäftigt}},
  future = {werden + sich + \textbf{beschäftigen}},
  pluperfect = {hatten + sich + \textbf{beschäftigt}},
  prepositions = {& \textbf{sich beschäftigen mit} + Dat & (to deal with / be occupied with)\\},
  separable = {---}
}

% Migrated from verbs/reflexive.tex:8
\VerbTable{
  name = {sich beschweren},
  english = {to complain},
  type = {regular},
  auxiliary = {haben},
  style = {default},
  present = {beschwere mich,beschwerst dich,beschwert sich,beschweren uns,beschwert euch,beschweren sich},
  preterite = {beschwerte mich,beschwertest dich,beschwerte sich,beschwerten uns,beschwertet euch,beschwerten sich},
  subjunctive = {},
  imperative = {beschwer(e) (du) dich!,beschweren wir uns!,beschwert (ihr) euch!,beschweren Sie sich!},
  perfect = {haben + sich + \textbf{beschwert}},
  future = {werden + sich + \textbf{beschweren}},
  pluperfect = {hatten + sich + \textbf{beschwert}},
  prepositions = {& \textbf{sich beschweren über} + Akk & (to complain about something)},
  separable = {---}
}

% Migrated from verbs/b1/b.tex:27
\VerbTable{
  name = {sich beteiligen},
  english = {to participate},
  type = {regular},
  auxiliary = {haben},
  style = {default},
  present = {beteilige mich,beteiligst dich,beteiligt sich,beteiligen uns,beteiligt euch,beteiligen sich},
  preterite = {beteiligte mich,beteiligtest dich,beteiligte sich,beteiligten uns,beteiligtet euch,beteiligten sich},
  subjunctive = {},
  imperative = {beteilig dich!,beteiligen wir uns!,beteiligt euch!,beteiligen Sie sich!},
  perfect = {haben + \textbf{sich beteiligt}},
  future = {werden + \textbf{sich beteiligen}},
  pluperfect = {hatten + \textbf{sich beteiligt}},
  prepositions = {---},
  separable = {---}
}

% Migrated from verbs/reflexive.tex:9
\VerbTable{
  name = {sich bewegen},
  english = {to move oneself},
  type = {irregular},
  auxiliary = {haben / sein},
  style = {default},
  present = {bewege mich,bewegst dich,bewegt sich,bewegen uns,bewegt euch,bewegen sich},
  preterite = {bewegte mich,bewegtest dich,bewegte sich,bewegten uns,bewegtet euch,bewegten sich},
  subjunctive = {},
  imperative = {beweg dich (du)!,bewegen wir uns!,bewegt euch (ihr)!,bewegen Sie sich!},
  perfect = {haben / sein + sich + \textbf{bewegt}},
  future = {werden + sich + \textbf{bewegen}},
  pluperfect = {haben / waren + sich + \textbf{bewegt}},
  prepositions = {
    & \textbf{sich bewegen in} + Dat & (to move in)\\
    & \textbf{sich bewegen auf} + Akk & (to move toward)\\
    & \textbf{sich bewegen durch} + Akk & (to move through)\\
  },
  separable = {---}
}

% Migrated from verbs/reflexive.tex:10
\VerbTable{
  name = {sich drehen},
  english = {to turn / to rotate},
  type = {regular},
  auxiliary = {haben},
  style = {default},
  present = {drehe mich,drehst dich,dreht sich,drehen uns,dreht euch,drehen sich},
  preterite = {drehte mich,drehtest dich,drehte sich,drehten uns,drehtet euch,drehten sich},
  subjunctive = {},
  imperative = {dreh(e) dich!,drehen wir uns!,dreht euch!,drehen Sie sich!},
  perfect = {haben + sich + \textbf{gedreht}},
  future = {werden + sich + \textbf{drehen}},
  pluperfect = {hatten + sich + \textbf{gedreht}},
  prepositions = {
    & \textbf{sich drehen um} + Akk & (to revolve around / be about)\\
    & \textbf{sich drehen nach} + Dat & (to turn toward)\\
  },
  separable = {
    & \textbf{um$|$drehen} & (to turn around)\\
    & \textbf{ab$|$drehen} & (to turn off / shoot a film)\\
  }
}

% Migrated from verbs/reflexive.tex:11
\VerbTable{
  name = {sich eignen},
  english = {to be suitable},
  type = {regular},
  auxiliary = {haben},
  style = {default},
  present = {eigne mich,eignest dich,eignet sich,eignen uns,eignet euch,eignen sich},
  preterite = {eignete mich,eignetest dich,eignete sich,eigneten uns,eignetet euch,eigneten sich},
  subjunctive = {},
  imperative = {eigne dich!,eignen wir uns!,eignet euch!,eignen Sie sich!},
  perfect = {haben + sich + \textbf{geeignet}},
  future = {werden + sich + \textbf{eignen}},
  pluperfect = {hatten + sich + \textbf{geeignet}},
  prepositions = {
    & \textbf{sich eignen für} + Akk & (to be suitable for)\\
    & \textbf{sich eignen als} + Nom & (to be suitable as)\\
  },
  separable = {---}
}

% Migrated from verbs/b1/e.tex:6
\VerbTable{
  name = {sich einigen},
  english = {to agree},
  type = {regular},
  auxiliary = {haben},
  style = {default},
  present = {einige mich,einigst dich,einigt sich,einigen uns,einigt euch,einigen sich},
  preterite = {einigte mich,einigtest dich,einigte sich,einigten uns,einigtet euch,einigten sich},
  subjunctive = {},
  imperative = {einig dich!,einigen wir uns!,einigt euch!,einigen Sie sich!},
  perfect = {haben + \textbf{sich geeinigt}},
  future = {werden + \textbf{sich einigen}},
  pluperfect = {hatten + \textbf{sich geeinigt}},
  prepositions = {---},
  separable = {---}
}

% Migrated from verbs/reflexive.tex:15
\VerbTable{
  name = {sich engagieren},
  english = {to get involved / commit oneself},
  type = {regular},
  auxiliary = {haben},
  style = {default},
  present = {engagiere mich,engagierst dich,engagiert sich,engagieren uns,engagiert euch,engagieren sich},
  preterite = {engagierte mich,engagiertest dich,engagierte sich,engagierten uns,engagiertet euch,engagierten sich},
  subjunctive = {},
  imperative = {engagiere dich!,engagieren wir uns!,engagiert euch!,engagieren Sie sich!},
  perfect = {haben + sich + \textbf{engagiert}},
  future = {werden + sich + \textbf{engagieren}},
  pluperfect = {hatten + sich + \textbf{engagiert}},
  prepositions = {
    & \textbf{sich engagieren für} + Akk & (to commit oneself to / support)\\
    & \textbf{sich engagieren bei} + Dat & (to get involved at / with)\\
  },
  separable = {---}
}

% Migrated from verbs/reflexive.tex:16
\VerbTable{
  name = {sich entscheiden},
  english = {to decide},
  type = {regular},
  auxiliary = {haben},
  style = {default},
  present = {entscheide mich,entscheidest dich,entscheidet sich,entscheiden uns,entscheidet euch,entscheiden sich},
  preterite = {entschied mich,entschiedest dich,entschied sich,entschieden uns,entschiedet euch,entschieden sich},
  subjunctive = {},
  imperative = {entscheid(e) (du) dich!,entscheiden wir uns!,entscheidet (ihr) euch!,entscheiden Sie sich!},
  perfect = {haben + sich + \textbf{entschieden}},
  future = {werden + sich + \textbf{entscheiden}},
  pluperfect = {hatten + sich + \textbf{entschieden}},
  prepositions = {
    & \textbf{sich entscheiden für} + Akk & (to decide for something)\\
    & \textbf{sich entscheiden gegen} + Akk & (to decide against something)
  },
  separable = {---}
}

% Migrated from verbs/b1/e.tex:23
\VerbTable{
  name = {sich entschließen},
  english = {to decide},
  type = {strong / irregular},
  auxiliary = {haben},
  style = {default},
  present = {entschließe mich,entschließt dich,entschließt sich,entschließen uns,entschließt euch,entschließen sich},
  preterite = {entschloss mich,entschlossest dich,entschloss sich,entschlossen uns,entschlosst euch,entschlossen sich},
  subjunctive = {},
  imperative = {entschließ dich!,entschließen wir uns!,entschließt euch!,entschließen Sie sich!},
  perfect = {haben + \textbf{sich entschlossen}},
  future = {werden + \textbf{sich entschließen}},
  pluperfect = {hatten + \textbf{sich entschlossen}},
  prepositions = {---},
  separable = {---}
}

% Migrated from verbs/reflexive.tex:17
\VerbTable{
  name = {sich entschuldigen},
  english = {to apologise},
  type = {regular, + Akkusativ},
  auxiliary = {haben},
  style = {acc},
  present = {entschuldige mich,entschuldigst dich,entschuldigt sich,entschuldigen uns,entschuldigt euch,entschuldigen sich},
  preterite = {entschuldigte mich,entschuldigtest dich,entschuldigte sich,entschuldigten uns,entschuldigtet euch,entschuldigten sich},
  subjunctive = {},
  imperative = {entschuldig(e) (du) dich!,entschuldigen wir uns!,entschuldigt (ihr) euch!,entschuldigen Sie sich!},
  perfect = {haben + sich + \textbf{entschuldigt}},
  future = {werden + sich + \textbf{entschuldigen}},
  pluperfect = {hatten + sich + \textbf{entschuldigt}},
  prepositions = {
    & \textbf{sich entschuldigen für} + Akk & (to apologise for)\\
    & \textbf{sich entschuldigen bei} + Dat & (to apologise to)\\
    & \textbf{sich entschuldigen mit} + Dat & (to plead something)
  },
  separable = {---}
}

% Migrated from verbs/reflexive.tex:18
\VerbTable{
  name = {sich entspannen},
  english = {to relax},
  type = {regular},
  auxiliary = {haben},
  style = {default},
  present = {entspanne mich,entspannst dich,entspannt sich,entspannen uns,entspannt euch,entspannen sich},
  preterite = {entspannte mich,entspanntest dich,entspannte sich,entspannten uns,entspanntet euch,entspannten sich},
  subjunctive = {},
  imperative = {entspann dich (du)!,entspannen wir uns!,entspannt euch (ihr)!,entspannen Sie sich!},
  perfect = {haben + sich + \textbf{entspannt}},
  future = {werden + sich + \textbf{entspannen}},
  pluperfect = {hatten + sich + \textbf{entspannt}},
  prepositions = {
    & \textbf{sich entspannen bei} + Dat & (to relax while / during)\\
    & \textbf{sich entspannen mit} + Dat & (to relax with)\\
    & \textbf{sich entspannen nach} + Dat & (to relax after)\\
  },
  separable = {---}
}

% Migrated from verbs/reflexive.tex:19
\VerbTable{
  name = {sich entwickeln},
  english = {to develop},
  type = {regular},
  auxiliary = {haben},
  style = {default},
  present = {entwickele mich,entwickelst dich,entwickelt sich,entwickeln uns,entwickelt euch,entwickeln sich},
  preterite = {entwickelte mich,ventwickeltest dich,entwickelte sich,entwickelten uns,entwickeltet euch,entwickelten sich},
  subjunctive = {},
  imperative = {entwickel(e) (du) dich!,entwickelen wir uns!,entwickelt (ihr) euch!,entwickelen Sie sich!},
  perfect = {haben + sich + \textbf{entwickelt}},
  future = {werden + sich + \textbf{entwickeln}},
  pluperfect = {hatten + sich + \textbf{entwickelt}},
  prepositions = {
    & \textbf{sich entwickeln aus} + Dat & (to evolve from something)\\
    & \textbf{sich entwickeln zu} + Dat & (to evolve into something)
  },
  separable = {---}
}

% Migrated from verbs/b1/e.tex:25
\VerbTable{
  name = {sich ereignen},
  english = {to happen / occur},
  type = {regular},
  auxiliary = {haben},
  style = {default},
  present = {ereigne mich,ereignest dich,ereignet sich,ereignen uns,ereignet euch,ereignen sich},
  preterite = {ereignete mich,ereignetest dich,ereignete sich,ereigneten uns,ereignetet euch,ereigneten sich},
  subjunctive = {},
  imperative = {ereigne dich!,ereignen wir uns!,ereignet euch!,ereignen Sie sich!},
  perfect = {haben + \textbf{sich ereignet}},
  future = {werden + \textbf{sich ereignen}},
  pluperfect = {hatten + \textbf{sich ereignet}},
  prepositions = {---},
  separable = {---}
}

% Migrated from verbs/reflexive.tex:20
\VerbTable{
  name = {sich erholen},
  english = {to recover / rest},
  type = {regular},
  auxiliary = {haben},
  style = {default},
  present = {erhole mich,erholst dich,erholt sich,erholen uns,erholt euch,erholen sich},
  preterite = {erholte mich,erholtest dich,erholte sich,erholten uns,erholtet euch,erholten sich},
  subjunctive = {},
  imperative = {erhole dich!,erholen wir uns!,erholt euch!,erholen Sie sich!},
  perfect = {haben + sich + \textbf{erholt}},
  future = {werden + sich + \textbf{erholen}},
  pluperfect = {hatten + sich + \textbf{erholt}},
  prepositions = {
    & \textbf{sich erholen von} + Dat & (to recover from)\\
    & \textbf{sich erholen nach} + Dat & (to recover after)\\
  },
  separable = {---}
}

% Migrated from verbs/reflexive.tex:12
\VerbTable{
  name = {sich erinnern},
  english = {to remember},
  type = {regular},
  auxiliary = {haben},
  style = {default},
  present = {erinnere mich,erinnerst dich,erinnert sich,erinnern uns,erinnert euch,erinnern sich},
  preterite = {erinnerte mich,erinnertest dich,erinnerte sich,erinnerten uns,erinnertet euch,erinnerten sich},
  subjunctive = {},
  imperative = {erinner(e) (du) dich!,erinnern wir uns!,erinnert (ihr) euch!,erinnern Sie sich!},
  perfect = {haben + sich + \textbf{erinnert}},
  future = {werden + sich + \textbf{erinnern}},
  pluperfect = {hatten + sich + \textbf{erinnert}},
  prepositions = {& \textbf{sich erinnern an} + Akk & (to recall something)},
  separable = {---}
}

% Migrated from verbs/b1/e.tex:30
\VerbTable{
  name = {sich erkälten},
  english = {to catch a cold},
  type = {regular},
  auxiliary = {haben},
  style = {default},
  present = {erkälte mich,erkältest dich,erkältet sich,erkälten uns,erkältet euch,erkälten sich},
  preterite = {erkältete mich,erkältetest dich,erkältete sich,erkälteten uns,erkältetet euch,erkälteten sich},
  subjunctive = {},
  imperative = {erkälte dich!,erkälten wir uns!,erkältet euch!,erkälten Sie sich!},
  perfect = {haben + \textbf{sich erkältet}},
  future = {werden + \textbf{sich erkälten}},
  pluperfect = {hatten + \textbf{sich erkältet}},
  prepositions = {---},
  separable = {---}
}

% Migrated from verbs/reflexive.tex:13
\VerbTable{
  name = {sich erkundigen},
  english = {to inquire},
  type = {regular},
  auxiliary = {haben},
  style = {default},
  present = {erkundige mich,erkundigst dich,erkundigt sich,erkundigen uns,erkundigt euch,erkundigen sich},
  preterite = {erkundigte mich,erkundigtest dich,erkundigte sich,erkundigten uns,erkundigtet euch,erkundigten sich},
  subjunctive = {},
  imperative = {erkundige dich!,erkundigen wir uns!,erkundigt euch!,erkundigen Sie sich!},
  perfect = {haben + sich + \textbf{erkundigt}},
  future = {werden + sich + \textbf{erkundigen}},
  pluperfect = {hatten + sich + \textbf{erkundigt}},
  prepositions = {
    & \textbf{sich erkundigen nach} + Dat & (to inquire about)\\
    & \textbf{sich erkundigen bei} + Dat & (to inquire with / ask at)\\
  },
  separable = {---}
}

% Migrated from verbs/reflexive.tex:14
\VerbTable{
  name = {sich ernähren},
  english = {to eat / nourish oneself},
  type = {regular},
  auxiliary = {haben},
  style = {default},
  present = {ernähre mich,ernährst dich,ernährt sich,ernähren uns,ernährt euch,ernähren sich},
  preterite = {ernährte mich,ernährtest dich,ernährte sich,ernährten uns,ernährtet euch,ernährten sich},
  subjunctive = {},
  imperative = {ernähr dich (du)!,ernähren wir uns!,ernährt euch (ihr)!,ernähren Sie sich!},
  perfect = {haben + sich + \textbf{ernährt}},
  future = {werden + sich + \textbf{ernähren}},
  pluperfect = {hatten + sich + \textbf{ernährt}},
  prepositions = {
    & \textbf{sich ernähren von} + Dat & (to live on / eat mainly)\\
    & \textbf{sich ernähren mit} + Dat & (to nourish oneself with)\\
  },
  separable = {---}
}

% Migrated from verbs/reflexive.tex:21
\VerbTable{
  name = {sich freuen},
  english = {to be pleased about},
  type = {regular},
  auxiliary = {haben},
  style = {default},
  present = {freue mich,freust dich,freut sich,freuen uns,freut euch,freuen sich},
  preterite = {freute mich,freutest dich,freute sich,freuten uns,freutet euch,freiten sich},
  subjunctive = {},
  imperative = {freu(e) (du) dich!,freuen wir uns!,freut (ihr) euch!,freuen Sie sich!},
  perfect = {haben + sich + \textbf{gefreut}},
  future = {werden + sich + \textbf{freuen}},
  pluperfect = {hatten + sich + \textbf{gefreut}},
  prepositions = {
    & \textbf{sich freuen über} + Akk & (to be please about something)\\
    & \textbf{sich freuen auf} + Akk & (to look forward to something)
  },
  separable = {---}
}

% Migrated from verbs/reflexive.tex:22
\VerbTable{
  name = {sich fühlen},
  english = {to feel},
  type = {regular, + Akkusativ},
  auxiliary = {haben},
  style = {acc},
  present = {fühle mich,fühlst dich,fühlt sich,fühlen uns,fühlt euch,fühlen sich},
  preterite = {fühlte mich,fühltest dich,fühlte sich,fühlten uns,fühltet euch,fühlten sich},
  subjunctive = {},
  imperative = {fühl(e) (du) dich!,fühlen wir uns!,fühlt (ihr) euch!,fühlen Sie sich!},
  perfect = {haben + sich + \textbf{gefühlt}},
  future = {werden + sich + \textbf{fühlen}},
  pluperfect = {hatten + sich + \textbf{gefühlt}},
  prepositions = {---},
  separable = {& \textbf{sich wohl$|$fühlen} & (to be comfortable)}
}

% Migrated from verbs/reflexive.tex:23
\VerbTable{
  name = {sich fürchten},
  english = {to be afraid},
  type = {regular},
  auxiliary = {haben},
  style = {default},
  present = {fürchte mich,fürchtest dich,fürchtet sich,fürchten uns,fürchtet euch,fürchten sich},
  preterite = {fürchtete mich,fürchtetest dich,fürchtete sich,fürchteten uns,fürchtetet euch,fürchteten sich},
  subjunctive = {},
  imperative = {fürchte dich!,fürchten wir uns!,fürchtet euch!,fürchten Sie sich!},
  perfect = {haben + sich + \textbf{gefürchtet}},
  future = {werden + sich + \textbf{fürchten}},
  pluperfect = {hatten + sich + \textbf{gefürchtet}},
  prepositions = {
    & \textbf{sich fürchten vor} + Dat & (to be afraid of)\\
    & \textbf{sich fürchten um} + Akk & (to fear for)\\
  },
  separable = {---}
}

% Migrated from verbs/reflexive.tex:24
\VerbTable{
  name = {sich gewöhnen},
  english = {to get used to},
  type = {regular},
  auxiliary = {haben},
  style = {default},
  present = {gewöhne mich,gewöhnst dich,gewöhnt sich,gewöhnen uns,gewöhnt euch,gewöhnen sich},
  preterite = {gewöhnte mich,gewöhntest dich,gewöhnte sich,gewöhnten uns,gewöhntet euch,gewöhnten sich},
  subjunctive = {},
  imperative = {gewöhn(e) dich!,gewöhnen wir uns!,gewöhnt euch!,gewöhnen Sie sich!},
  perfect = {haben + sich + \textbf{gewöhnt}},
  future = {werden + sich + \textbf{gewöhnen}},
  pluperfect = {hatten + sich + \textbf{gewöhnt}},
  prepositions = {& \textbf{sich gewöhnen an} + Akk & (to get used to)\\},
  separable = {---}
}

% Migrated from verbs/reflexive.tex:25
\VerbTable{
  name = {sich hüten},
  english = {to beware / to be careful},
  type = {regular},
  auxiliary = {haben},
  style = {default},
  present = {hüte mich,hütest dich,hütet sich,hüten uns,hütet euch,hüten sich},
  preterite = {hütete mich,hütetest dich,hütete sich,hüteten uns,hütetet euch,hüteten sich},
  subjunctive = {},
  imperative = {hüte dich!,hüten wir uns!,hütet euch!,hüten Sie sich!},
  perfect = {haben + sich + \textbf{gehütet}},
  future = {werden + sich + \textbf{hüten}},
  pluperfect = {hatten + sich + \textbf{gehütet}},
  prepositions = {
    & \textbf{sich hüten vor} + Dat & (to beware of)\\
    & \textbf{sich hüten, zu} + Inf & (to be careful not to)\\
  },
  separable = {---}
}

% Migrated from verbs/reflexive.tex:26
\VerbTable{
  name = {sich interessieren},
  english = {to be interest in},
  type = {regular},
  auxiliary = {haben},
  style = {default},
  present = {interessiere mich,interessierst dich,interessiert sich,interessieren uns,interessiert euch,interessieren sich},
  preterite = {interessierte mich,interessiertest dich,interessierte sich,interessierten uns,interessiertet euch,interessierten sich},
  subjunctive = {},
  imperative = {interessier(e) (du) dich!,interessieren wir uns!,interessiert (ihr) euch!,interessieren Sie sich!},
  perfect = {haben + sich + \textbf{interessiert}},
  future = {werden + sich + \textbf{interessieren}},
  pluperfect = {hatten + sich + \textbf{interessiert}},
  prepositions = {& \textbf{sich interessieren für} + Akk & (to be interested in something)\\},
  separable = {---}
}

% Migrated from verbs/reflexive.tex:27
\VerbTable{
  name = {sich irren},
  english = {to be mistaken},
  type = {regular},
  auxiliary = {haben},
  style = {default},
  present = {irre mich,irrst dich,irrt sich,irren uns,irrt euch,irren sich},
  preterite = {irrte mich,irrtest dich,irrte sich,irrten uns,irrtet euch,irrten sich},
  subjunctive = {},
  imperative = {irr(e) dich!,irren wir uns!,irrt euch!,irren Sie sich!},
  perfect = {haben + sich + \textbf{geirrt}},
  future = {werden + sich + \textbf{irren}},
  pluperfect = {hatten + sich + \textbf{geirrt}},
  prepositions = {
    & \textbf{sich irren in} + Dat & (to be mistaken about)\\
    & \textbf{sich irren bei} + Dat & (to make a mistake in/with)\\
  },
  separable = {---}
}

% Migrated from verbs/reflexive.tex:28
\VerbTable{
  name = {sich konzentrieren},
  english = {to concentrate},
  type = {regular},
  auxiliary = {haben},
  style = {default},
  present = {konzentriere mich,konzentrierst dich,konzentriert sich,konzentrieren uns,konzentriert euch,konzentrieren sich},
  preterite = {konzentrierte mich,konzentriertest dich,konzentrierte sich,konzentrierten uns,konzentriertet euch,konzentrierten sich},
  subjunctive = {},
  imperative = {konzentrier(e) (du) dich!,konzentrieren wir uns!,konzentriert (ihr) euch!,konzentrieren Sie sich!},
  perfect = {haben + sich + \textbf{konzentriert}},
  future = {werden + sich + \textbf{konzentrieren}},
  pluperfect = {hatten + sich + \textbf{konzentriert}},
  prepositions = {& \textbf{sich konzentrieren auf} + Akk & (to concentrate on something)},
  separable = {---}
}

% Migrated from verbs/reflexive.tex:29
\VerbTable{
  name = {sich kümmern},
  english = {to take care of},
  type = {regular},
  auxiliary = {haben},
  style = {default},
  present = {kümmere mich,kümmerst dich,kümmert sich,kümmern uns,kümmert euch,kümmern sich},
  preterite = {kümmerte mich,kümmertest dich,kümmerte sich,kümmerten uns,kümmertet euch,kümmerten sich},
  subjunctive = {},
  imperative = {kümmer(e) (du) dich!,kümmern wir uns!,kümmert (ihr) euch!,kümmern Sie sich!},
  perfect = {haben + sich + \textbf{gekümmert}},
  future = {werden + sich + \textbf{kümmern}},
  pluperfect = {hatten + sich + \textbf{gekümmert}},
  prepositions = {& \textbf{sich kümmern um} + Akk & (to take care of)},
  separable = {---}
}

% Migrated from verbs/b1/l.tex:1
\VerbTable{
  name = {sich langweilen},
  english = {to be bored / bore},
  type = {regular},
  auxiliary = {haben},
  style = {default},
  present = {langweile mich,langweilst dich,langweilt sich,langweilen uns,langweilt euch,langweilen sich},
  preterite = {langweilte mich,langweiltest dich,langweilte sich,langweilten uns,langweiltet euch,langweilten sich},
  subjunctive = {},
  imperative = {langweil dich!,langweilen wir uns!,langweilt euch!,langweilen Sie sich!},
  perfect = {haben + \textbf{sich gelangweilt}},
  future = {werden + \textbf{sich langweilen}},
  pluperfect = {hatten + \textbf{sich gelangweilt}},
  prepositions = {---},
  separable = {---}
}

% Migrated from verbs/b1/l.tex:8
\VerbTable{
  name = {sich lohnen},
  english = {to be worthwhile},
  type = {regular},
  auxiliary = {haben},
  style = {default},
  present = {lohne mich,lohnst dich,lohnt sich,lohnen uns,lohnt euch,lohnen sich},
  preterite = {lohnte mich,lohntest dich,lohnte sich,lohnten uns,lohntet euch,lohnten sich},
  subjunctive = {},
  imperative = {lohn dich!,lohnen wir uns!,lohnt euch!,lohnen Sie sich!},
  perfect = {haben + \textbf{sich gelohnt}},
  future = {werden + \textbf{sich lohnen}},
  pluperfect = {hatten + \textbf{sich gelohnt}},
  prepositions = {---},
  separable = {---}
}

% Migrated from verbs/reflexive.tex:30
\VerbTable{
  name = {sich melden},
  english = {to get in touch},
  type = {regular},
  auxiliary = {haben},
  style = {default},
  present = {melde mich,meldest dich,meldet sich,melden uns,meldet euch,melden sich},
  preterite = {meldete mich,meldetest dich,meldete sich,meldeten uns,meldetet euch,meldeten sich},
  subjunctive = {},
  imperative = {melde dich (du)!,melden wir uns!,meldet euch (ihr)!,melden Sie sich!},
  perfect = {haben + sich + \textbf{gemeldet}},
  future = {werden + sich +\textbf{melden}},
  pluperfect = {hatten + sich + \textbf{gemeldet}},
  prepositions = {
    & \textbf{sich melden bei} + Dat & (to contact)\\
    & \textbf{sich melden für} + Akk & (to register for)\\
  },
  separable = {
    & \textbf{an$|$melden} & (to register)\\
    & \textbf{ab$|$melden} & (to deregister)\\
  }
}

% Migrated from verbs/b1/n.tex:3
\VerbTable{
  name = {sich nähern},
  english = {to approach},
  type = {regular},
  auxiliary = {haben},
  style = {default},
  present = {nähre mich,näherst dich,nähert sich,nähern uns,nähert euch,nähern sich},
  preterite = {näherte mich,nähertest dich,näherte sich,näherten uns,nähertet euch,näherten sich},
  subjunctive = {},
  imperative = {nähre dich!,nähern wir uns!,nähert euch!,nähern Sie sich!},
  perfect = {haben + \textbf{sich genähert}},
  future = {werden + \textbf{sich nähern}},
  pluperfect = {hatten + \textbf{sich genähert}},
  prepositions = {---},
  separable = {---}
}

% Migrated from verbs/reflexive.tex:31
\VerbTable{
  name = {sich rächen},
  english = {to take revenge},
  type = {regular},
  auxiliary = {haben},
  style = {default},
  present = {räche mich,rächst dich,rächt sich,rächen uns,rächt euch,rächen sich},
  preterite = {rächte mich,rächtest dich,rächte sich,rächten uns,rächtet euch,rächten sich},
  subjunctive = {},
  imperative = {räch(e) dich!,rächen wir uns!,rächt euch!,rächen Sie sich!},
  perfect = {haben + sich + \textbf{gerächt}},
  future = {werden + sich + \textbf{rächen}},
  pluperfect = {hatten + sich + \textbf{gerächt}},
  prepositions = {
    & \textbf{sich rächen an} + Dat & (to take revenge on)\\
    & \textbf{sich rächen für} + Akk & (to take revenge for)\\
  },
  separable = {---}
}

% Migrated from verbs/reflexive.tex:32
\VerbTable{
  name = {sich regen},
  english = {to stir / to move},
  type = {regular},
  auxiliary = {haben},
  style = {default},
  present = {rege mich,regst dich,regt sich,regen uns,regt euch,regen sich},
  preterite = {regte mich,regtest dich,regte sich,regten uns,regtet euch,regten sich},
  subjunctive = {},
  imperative = {reg(e) dich!,regen wir uns!,regt euch!,regen Sie sich!},
  perfect = {haben + sich + \textbf{geregt}},
  future = {werden + sich + \textbf{regen}},
  pluperfect = {hatten + sich + \textbf{geregt}},
  prepositions = {
    & \textbf{sich regen in} + Dat & (to stir in / move in)\\
    & \textbf{sich regen bei} + Dat & (to stir at / during)\\
  },
  separable = {
    & \textbf{an$|$regen} & (to stimulate / suggest)\\
    & \textbf{auf$|$regen} & (to upset / excite)\\
  }
}

% Migrated from verbs/reflexive.tex:33
\VerbTable{
  name = {sich schämen},
  english = {to be ashamed},
  type = {regular, + Akkusativ},
  auxiliary = {haben},
  style = {acc},
  present = {schäme mich,schämst dich,schämt sich,schämen uns,schämt euch,schämen sich},
  preterite = {schämte mich,schämtest dich,schämte sich,schämten uns,schämtet euch,schämten sich},
  subjunctive = {},
  imperative = {schäm(e) (du) dich!,schämen wir uns!,schämt (ihr) euch!,schämen Sie sich!},
  perfect = {haben + sich + \textbf{geschämt}},
  future = {werden + sich + \textbf{schämen}},
  pluperfect = {hatten + sich + \textbf{geschämt}},
  prepositions = {& \textbf{sich schämen für} + Akk & (to be ashamed of)},
  separable = {---}
}

% Migrated from verbs/reflexive.tex:34
\VerbTable{
  name = {sich schminken},
  english = {to put on makeup},
  type = {regular},
  auxiliary = {haben},
  style = {default},
  present = {schminke mich,schminkst dich,schminkt sich,schminken uns,schminkt euch,schminken sich},
  preterite = {schminkte mich,schminktest dich,schminkte sich,schminkten uns,schminktet euch,schminkten sich},
  subjunctive = {},
  imperative = {schminke dich!,schminken wir uns!,schminkt euch!,schminken Sie sich!},
  perfect = {haben + sich + \textbf{geschminkt}},
  future = {werden + sich + \textbf{schminken}},
  pluperfect = {hatten + sich + \textbf{geschminkt}},
  prepositions = {
    & \textbf{sich schminken für} + Akk & (to put on makeup for)\\
    & \textbf{sich schminken mit} + Dat & (to put makeup on with)\\
  },
  separable = {
    & \textbf{ab$|$schminken} & (to remove makeup)\\
    & \textbf{an$|$schminken} & (to apply makeup)\\
  }
}

% Migrated from verbs/reflexive.tex:35
\VerbTable{
  name = {sich schützen},
  english = {to protect oneself},
  type = {regular},
  auxiliary = {haben},
  style = {default},
  present = {schütze mich,schützt dich,schützt sich,schützen uns,schützt euch,schützen sich},
  preterite = {schützte mich,schütztest dich,schützte sich,schützten uns,schütztet euch,schützten sich},
  subjunctive = {},
  imperative = {schütz(e) dich!,schützen wir uns!,schützt euch!,schützen Sie sich!},
  perfect = {haben + sich + \textbf{geschützt}},
  future = {werden + sich + \textbf{schützen}},
  pluperfect = {hatten + sich + \textbf{geschützt}},
  prepositions = {
    & \textbf{sich schützen vor} + Dat & (to protect oneself from)\\
    & \textbf{sich schützen gegen} + Akk & (to protect oneself against)\\
  },
  separable = {---}
}

% Migrated from verbs/reflexive.tex:36
\VerbTable{
  name = {sich sehnen},
  english = {to long for},
  type = {regular, reflexive},
  auxiliary = {haben},
  style = {default},
  present = {sehne mich,sehnst dich,sehnt sich,sehnen uns,sehnt euch,sehnen sich},
  preterite = {sehnte mich,sehntest dich,sehnte sich,sehnten uns,sehntet euch,sehnten sich},
  subjunctive = {},
  imperative = {sehn(e) dich!,sehnen wir uns!,sehnt euch!,sehnen Sie sich!},
  perfect = {haben + sich + \textbf{gesehnt}},
  future = {werden + sich + \textbf{sehnen}},
  pluperfect = {hatten + sich + \textbf{gesehnt}},
  prepositions = {& \textbf{sich sehnen nach} + Dat & (to long for)\\},
  separable = {---}
}

% Migrated from verbs/reflexive.tex:37
\VerbTable{
  name = {sich setzen},
  english = {to sit down / take a seat},
  type = {regular},
  auxiliary = {haben},
  style = {default},
  present = {setze mich,setzt dich,setzt sich,setzen uns,setzt euch,setzen sich},
  preterite = {setzte mich,setztest dich,setzte sich,setzten uns,setztet euch,setzten sich},
  subjunctive = {},
  imperative = {setze dich!,setzen wir uns!,setzt euch!,setzen Sie sich!},
  perfect = {haben + sich + \textbf{gesetzt}},
  future = {werden + sich + \textbf{setzen}},
  pluperfect = {hatten + sich + \textbf{gesetzt}},
  prepositions = {
    & \textbf{sich setzen auf} + Akk & (to sit down on)\\
    & \textbf{sich setzen neben} + Akk & (to sit next to)\\
  },
  separable = {& \textbf{hin$|$setzen} & (to sit down)\\}
}

% Migrated from verbs/reflexive.tex:38
\VerbTable{
  name = {sich sorgen},
  english = {to worry},
  type = {regular},
  auxiliary = {haben},
  style = {default},
  present = {sorge mich,sorgst dich,sorgt sich,sorgen uns,sorgt euch,sorgen sich},
  preterite = {sorgte mich,sorgtest dich,sorgte sich,sorgten uns,sorgtet euch,sorgten sich},
  subjunctive = {},
  imperative = {sorg(e) (du) dich!,sorgen wir uns!,sorgt (ihr) euch!,sorgen Sie sich!},
  perfect = {haben + sich + \textbf{gesorgt}},
  future = {werden + sich + \textbf{sorgen}},
  pluperfect = {hatten + sich + \textbf{gesorgt}},
  prepositions = {& \textbf{sich sorgen um} + Akk & (to worry about something)},
  separable = {---}
}

% Migrated from verbs/reflexive.tex:39
\VerbTable{
  name = {sich stellen},
  english = {to surrender},
  type = {regular, + Akkusativ},
  auxiliary = {haben},
  style = {acc},
  present = {stelle mich,stellst dich,stellt sich,stellen uns,stellt euch,stellen sich},
  preterite = {stellte mich,stelltest dich,stellte sich,stellten uns,stelltet euch,stellten sich},
  subjunctive = {},
  imperative = {stell(e) (du) dich!,stellen wir uns!,stellt (ihr) euch!,stellen Sie sich!},
  perfect = {haben + sich + \textbf{gestellt}},
  future = {werden + sich + \textbf{stellen}},
  pluperfect = {hatten + sich + \textbf{gestellt}},
  prepositions = {---},
  separable = {& \textbf{sich vor$|$stellen} & (to introduce oneself)}
}

% Migrated from verbs/reflexive.tex:40
\VerbTable{
  name = {sich teilen},
  english = {to share},
  type = {regular, + Dativ},
  auxiliary = {haben},
  style = {dat},
  present = {teile mir,teilst dir,teilt sich,teilen uns,teilt euch,teilen sich},
  preterite = {teilte mir,teiltest dir,teilte sich,teilten uns,teiltet euch,teilten sich},
  subjunctive = {},
  imperative = {teil(e) (du) dir!,teilen wir uns!,teilt (ihr) euch!,teilen Sie sich!},
  perfect = {haben + sich + \textbf{geteilt}},
  future = {werden + sich + \textbf{teilen}},
  pluperfect = {hatten + sich + \textbf{geteilt}},
  prepositions = {
    & \textbf{sich teilen in} + Akk & (to split into parts)\\
    & \textbf{sich teilen mit} + Dat & (to share with someone)
  },
  separable = {---}
}

% Migrated from verbs/reflexive.tex:41
\VerbTable{
  name = {sich treffen},
  english = {to come forward},
  type = {irregular, + Akkusativ},
  auxiliary = {haben},
  style = {acc},
  present = {treffe mich,triffst dich,trifft sich,treffen uns,trefft euch,treffen sich},
  preterite = {traf mich,trafst dich,traf sich,trafen uns,traft euch,trafen sich},
  subjunctive = {},
  imperative = {triff(e) (du) dich!,treffen wir uns!,trefft (ihr) euch!,treffen Sie sich!},
  perfect = {haben + sich + \textbf{getroffen}},
  future = {werden + sich + \textbf{treffen}},
  pluperfect = {hatten + sich + \textbf{getroffen}},
  prepositions = {
    & \textbf{sich treffen mit} + Dat & (to meet someone)\\
    & \textbf{sich treffen zu} + Dat & (to meet for a purpose)\\
  },
  separable = {---}
}

% Migrated from verbs/reflexive.tex:42
\VerbTable{
  name = {sich trennen},
  english = {to separate},
  type = {regular, + Akkusativ},
  auxiliary = {haben},
  style = {acc},
  present = {trenne mich,trennst dich,trennt sich,trennen uns,trennt euch,trennen sich},
  preterite = {trennte mich,trenntest dich,trennte sich,trennten uns,trenntet euch,trennten sich},
  subjunctive = {},
  imperative = {trenn(e) (du) dich!,trennen wir uns!,trennt (ihr) euch!,trennen Sie sich!},
  perfect = {haben + sich + \textbf{getrennt}},
  future = {werden + sich + \textbf{trennen}},
  pluperfect = {hatten + sich + \textbf{getrennt}},
  prepositions = {& \textbf{sich trennen von} + Dat & (to separate from someone / something)},
  separable = {---}
}

% Migrated from verbs/b1/u.tex:16
\VerbTable{
  name = {sich umziehen},
  english = {to change clothes},
  type = {reflexive, separable, strong / irregular},
  auxiliary = {haben},
  style = {default},
  present = {ziehe mich um,ziehst dich um,zieht sich um,ziehen uns um,zieht euch um,ziehen sich um},
  preterite = {zog mich um,zogst dich um,zog sich um,zogen uns um,zogt euch um,zogen sich um},
  subjunctive = {},
  imperative = {zieh dich um!,ziehen wir uns um!,zieht euch um!,ziehen Sie sich um!},
  perfect = {haben + \textbf{sich umgezogen}},
  future = {werden + \textbf{sich umziehen}},
  pluperfect = {hatten + \textbf{sich umgezogen}},
  prepositions = {---},
  separable = {---}
}

% Migrated from verbs/reflexive.tex:43
\VerbTable{
  name = {sich unterhalten},
  english = {to chat},
  type = {irregular},
  auxiliary = {haben},
  style = {default},
  present = {unterhalte mich,unterhältst dich,unterhält sich,unterhalten uns,unterhaltet euch,unterhalten sich},
  preterite = {underhielt mich,underhieltst dich,underhielt sich,underhielten uns,underhieltet euch,underhielten sich},
  subjunctive = {},
  imperative = {unterhalt(e) (du) dich!,unterhalten wir uns!,unterhaltt (ihr) euch!,unterhalten Sie sich!},
  perfect = {haben + sich + \textbf{unterhalten}},
  future = {werden + sich + \textbf{unterhalten}},
  pluperfect = {hatten + sich + \textbf{unterhalten}},
  prepositions = {
    & \textbf{sich unterhalten über} + Akk & (to chat about)\\
    & \textbf{sich unterhalten mit} + Dat & (to chat mit)\\
  },
  separable = {---}
}

% Migrated from verbs/reflexive.tex:44
\VerbTable{
  name = {sich verabreden},
  english = {to arrange / make an appointment},
  type = {regular},
  auxiliary = {haben},
  style = {default},
  present = {verabrede mich,verabredest dich,verabredet sich,verabreden uns,verabredet euch,verabreden sich},
  preterite = {verabredete mich,verabredetest dich,verabredete sich,verabredeten uns,verabredetet euch,verabredeten sich},
  subjunctive = {},
  imperative = {verabrede dich!,verabreden wir uns!,verabredet euch!,verabreden Sie sich!},
  perfect = {haben + sich + \textbf{verabredet}},
  future = {werden + sich + \textbf{verabreden}},
  pluperfect = {hatten + sich + \textbf{verabredet}},
  prepositions = {& \textbf{sich verabreden mit} + Dat & (to arrange with)\\},
  separable = {---}
}

% Migrated from verbs/b1/v.tex:7
\VerbTable{
  name = {sich vergnügen},
  english = {to enjoy oneself},
  type = {regular},
  auxiliary = {haben},
  style = {default},
  present = {vergnüge mich,vergnügst dich,vergnügt sich,vergnügen uns,vergnügt euch,vergnügen sich},
  preterite = {vergnügte mich,vergnügtest dich,vergnügte sich,vergnügten uns,vergnügtet euch,vergnügten sich},
  subjunctive = {},
  imperative = {vergnüg dich!,vergnügen wir uns!,vergnügt euch!,vergnügen Sie sich!},
  perfect = {haben + \textbf{sich vergnügt}},
  future = {werden + \textbf{sich vergnügen}},
  pluperfect = {hatten + \textbf{sich vergnügt}},
  prepositions = {---},
  separable = {---}
}

% Migrated from verbs/b1/v.tex:10
\VerbTable{
  name = {sich verhalten},
  english = {to behave},
  type = {strong / irregular},
  auxiliary = {haben},
  style = {default},
  present = {verhalte mich,verhältst dich,verhält sich,verhalten uns,verhaltet euch,verhalten sich},
  preterite = {verhielt mich,verhieltest dich,verhielt sich,verhielten uns,verhieltet euch,verhielten sich},
  subjunctive = {},
  imperative = {verhalte dich!,verhalten wir uns!,verhaltet euch!,verhalten Sie sich!},
  perfect = {haben + \textbf{sich verhalten}},
  future = {werden + \textbf{sich verhalten}},
  pluperfect = {hatten + \textbf{sich verhalten}},
  prepositions = {---},
  separable = {---}
}

% Migrated from verbs/reflexive.tex:45
\VerbTable{
  name = {sich verlassen},
  english = {to rely},
  type = {irregular},
  auxiliary = {haben},
  style = {default},
  present = {verlasse mich,verlässt dich,verlässt sich,verlassen uns,verlasst euch,verlassen sich},
  preterite = {verließ mich,verließest dich,verließ sich,verließen uns,verließt euch,verließen sich},
  subjunctive = {},
  imperative = {verlass(e) dich!,verlassen wir uns!,verlasst euch!,verlassen Sie sich!},
  perfect = {haben + sich + \textbf{verlassen}},
  future = {werden + sich + \textbf{verlassen}},
  pluperfect = {hatten + sich + \textbf{verlassen}},
  prepositions = {& \textbf{sich verlassen auf} + Akk & (to rely on)\\},
  separable = {---}
}

% Migrated from verbs/b1/v.tex:14
\VerbTable{
  name = {sich verlaufen},
  english = {to get lost / run},
  type = {strong / irregular},
  auxiliary = {haben},
  style = {default},
  present = {verlaufe mich,verläufst dich,verläuft sich,verlaufen uns,verlauft euch,verlaufen sich},
  preterite = {verlief mich,verliefst dich,verlief sich,verliefen uns,verlieft euch,verliefen sich},
  subjunctive = {},
  imperative = {verlauf dich!,verlaufen wir uns!,verlauft euch!,verlaufen Sie sich!},
  perfect = {haben + \textbf{sich verlaufen}},
  future = {werden + \textbf{sich verlaufen}},
  pluperfect = {hatten + \textbf{sich verlaufen}},
  prepositions = {---},
  separable = {---}
}

% Migrated from verbs/reflexive.tex:46
\VerbTable{
  name = {sich verletzen},
  english = {to injure oneself},
  type = {regular},
  auxiliary = {haben},
  style = {default},
  present = {verletze mich,verletzt dich,verletzt sich,verletzen uns,verletzt euch,verletzen sich},
  preterite = {verletzte mich,verletztest dich,verletzte sich,verletzten uns,verletz tet euch,verletzten sich},
  subjunctive = {},
  imperative = {verletze dich (du) nicht!,verletzen wir uns nicht!,verletzt euch (ihr) nicht!,verletzen Sie sich nicht!},
  perfect = {haben + sich + \textbf{verletzt}},
  future = {werden + sich + \textbf{verletzen}},
  pluperfect = {hatten + sich + \textbf{verletzt}},
  prepositions = {
    & \textbf{sich verletzen an} + Dat & (to injure oneself on)\\
    & \textbf{sich verletzen bei} + Dat & (to injure oneself while doing)\\
  },
  separable = {---}
}

% Migrated from verbs/reflexive.tex:47
\VerbTable{
  name = {sich verlieben},
  english = {to fall in love},
  type = {regular},
  auxiliary = {haben},
  style = {default},
  present = {verliebe mir,verliebst dir,verliebt sich,verlieben uns,verliebt euch,verlieben sich},
  preterite = {verliebte mir,verliebtest dir,verliebte sich,verliebten uns,verliebtet euch,verliebten sich},
  subjunctive = {},
  imperative = {verlieb(e) (du) dir!,verlieben wir uns!,verliebt (ihr) euch!,verlieben Sie sich!},
  perfect = {haben + sich + \textbf{verliebt}},
  future = {werden + sich + \textbf{verlieben}},
  pluperfect = {hatten + sich + \textbf{verliebt}},
  prepositions = {& \textbf{sich verlieben in} + Akk & (to fall in love with someone)},
  separable = {---}
}

% Migrated from verbs/reflexive.tex:48
\VerbTable{
  name = {sich verschreiben},
  english = {to commit / write incorrectly},
  type = {irregular},
  auxiliary = {haben},
  style = {default},
  present = {verschreibe mich,verschreibst dich,verschreibt sich,verschreiben uns,verschreibt euch,verschreiben sich},
  preterite = {verschrieb mich,verschriebst dich,verschrieb sich,verschrieben uns,verschriebt euch,verschrieben sich},
  subjunctive = {},
  imperative = {verschreibe (du) dich!,verschreiben wir uns!,verschreibt (ihr) euch!,verschreiben Sie sich!},
  perfect = {haben + sich + \textbf{verschrieben}},
  future = {werden + sich + \textbf{verschreiben}},
  pluperfect = {hatten + sich + \textbf{verschrieben}},
  prepositions = {---},
  separable = {---}
}

% Migrated from verbs/reflexive.tex:49
\VerbTable{
  name = {sich verstehen},
  english = {to get along},
  type = {irregular, + Akkusativ},
  auxiliary = {haben},
  style = {acc},
  present = {verstehe mich,verstehst dich,versteht sich,verstehen uns,versteht euch,verstehen sich},
  preterite = {verstand mich,verstandest dich,verstand sich,verstanden uns,verstandet euch,verstanden sich},
  subjunctive = {},
  imperative = {versteh(e) (du) dich!,verstehen wir uns!,versteht (ihr) euch!,verstehen Sie sich!},
  perfect = {haben + sich + \textbf{verstanden}},
  future = {werden + sich + \textbf{verstehen}},
  pluperfect = {hatten + sich + \textbf{verstanden}},
  prepositions = {& \textbf{sich verlieben mit} + Dat & (to get along with)},
  separable = {---}
}

% Migrated from verbs/reflexive.tex:50
\VerbTable{
  name = {sich verwählen},
  english = {to misdial / miscall},
  type = {regular},
  auxiliary = {haben},
  style = {default},
  present = {verwähle mich,verwählst dich,verwählt sich,verwählen uns,verwählt euch,verwählen sich},
  preterite = {verwählte mich,verwähltest dich,verwählte sich,verwählten uns,verwähltet euch,verwählten sich},
  subjunctive = {},
  imperative = {verwähle dich!,verwählen wir uns!,verwählt euch!,verwählen Sie sich!},
  perfect = {haben + sich + \textbf{verwählt}},
  future = {werden + sich + \textbf{verwählen}},
  pluperfect = {hatten + sich + \textbf{verwählt}},
  prepositions = {---},
  separable = {---}
}

% Migrated from verbs/reflexive.tex:51
\VerbTable{
  name = {sich waschen},
  english = {to wash},
  type = {irregular, + Akkusativ},
  auxiliary = {haben},
  style = {acc},
  present = {wasche mich,wäschst dich,wäscht sich,waschen uns,wascht euch,waschen sich},
  preterite = {wunch mich,wunchst dich,wunch sich,wunchen uns,wuncht euch,wunchen sich},
  subjunctive = {},
  imperative = {wasch(e) (du) dich!,waschen wir uns!,wascht (ihr) euch!,waschen Sie sich!},
  perfect = {haben + sich + \textbf{gewaschen}},
  future = {werden + sich + \textbf{waschen}},
  pluperfect = {hatten + sich + \textbf{gewaschen}},
  prepositions = {---},
  separable = {---}
}

% Migrated from verbs/b1/w.tex:6
\VerbTable{
  name = {sich weigern},
  english = {to refuse},
  type = {regular},
  auxiliary = {haben},
  style = {default},
  present = {weigre mich,weigerst dich,weigert sich,weigern uns,weigert euch,weigern sich},
  preterite = {weigerte mich,weigertest dich,weigerte sich,weigerten uns,weigertet euch,weigerten sich},
  subjunctive = {},
  imperative = {weigre dich!,weigern wir uns!,weigert euch!,weigern Sie sich!},
  perfect = {haben + \textbf{sich geweigert}},
  future = {werden + \textbf{sich weigern}},
  pluperfect = {hatten + \textbf{sich geweigert}},
  prepositions = {---},
  separable = {---}
}

% Migrated from verbs/reflexive.tex:52
\VerbTable{
  name = {sich wundern},
  english = {to wonder / be suprised},
  type = {regular},
  auxiliary = {haben},
  style = {default},
  present = {wundere mich,wunderst dich,wundert sich,wundern uns,wundert euch,wundern sich},
  preterite = {wunderte mich,wundertest dich,wunderte sich,wunderten uns,wundertet euch,wunderten sich},
  subjunctive = {},
  imperative = {wunder(e) (du) dich!,wundern wir uns!,wundert (ihr) euch!,wundern Sie sich!},
  perfect = {haben + sich + \textbf{gewundert}},
  future = {werden + sich + \textbf{wundern}},
  pluperfect = {hatten + sich + \textbf{gewundert}},
  prepositions = {& \textbf{sich wundern über} + Akk & (to wonder / be suprised about)},
  separable = {---}
}

% Migrated from verbs/reflexive.tex:53
\VerbTable{
  name = {sich ziehen},
  english = {to drag},
  type = {irregular, + Akkusativ},
  auxiliary = {haben},
  style = {acc},
  present = {ziehe mich,ziehst dich,zieht sich,ziehen uns,zieht euch,ziehen sich},
  preterite = {zog mich,zogst dich,zog sich,zogen uns,zogt euch,zogen sich},
  subjunctive = {},
  imperative = {zieh(e) (du) dich!,ziehen wir uns!,zieht (ihr) euch!,ziehen Sie sich!},
  perfect = {haben + sich + \textbf{gezogen}},
  future = {werden + sich + \textbf{ziehen}},
  pluperfect = {hatten + sich + \textbf{gezogen}},
  prepositions = {---},
  separable = {& \textbf{sich an$|$ziehan} + Akk & (to get dressed)}
}

\chapter{Alphabetical Verb Tables}

\begin{tcolorbox}[colback=referenceBlueBg,colframe=genderMasculine,title={B1 coverage note}]
This reference contains 715 complete verb entries. The expanded selection targets the verbs most useful through CEFR B1; CEFR does not prescribe one official exhaustive verb list. Each entry includes all six present and simple-past forms, four imperative forms where they exist, and the perfect, future and pluperfect constructions.
\end{tcolorbox}

\section{A}

% Migrated from verbs/b1/a.tex:1
\VerbTable{
  name = {abbiegen},
  english = {to turn},
  type = {separable, strong / irregular},
  auxiliary = {sein},
  style = {default},
  present = {biege ab,biegst ab,biegt ab,biegen ab,biegt ab,biegen ab},
  preterite = {bog ab,bogst ab,bog ab,bogen ab,bogt ab,bogen ab},
  subjunctive = {},
  imperative = {bieg (du) ab!,biegen wir ab!,biegt (ihr) ab!,biegen Sie ab!},
  perfect = {sein + \textbf{abgebogen}},
  future = {werden + \textbf{abbiegen}},
  pluperfect = {waren + \textbf{abgebogen}},
  prepositions = {---},
  separable = {---}
}

% Migrated from verbs/b1/a.tex:2
\VerbTable{
  name = {abfahren},
  english = {to depart / leave},
  type = {separable, strong / irregular},
  auxiliary = {sein},
  style = {default},
  present = {fahre ab,fährst ab,fährt ab,fahren ab,fahrt ab,fahren ab},
  preterite = {fuhr ab,fuhrst ab,fuhr ab,fuhren ab,fuhrt ab,fuhren ab},
  subjunctive = {},
  imperative = {fahr (du) ab!,fahren wir ab!,fahrt (ihr) ab!,fahren Sie ab!},
  perfect = {sein + \textbf{abgefahren}},
  future = {werden + \textbf{abfahren}},
  pluperfect = {waren + \textbf{abgefahren}},
  prepositions = {---},
  separable = {---}
}

% Migrated from verbs/b1/a.tex:3
\VerbTable{
  name = {abgeben},
  english = {to hand in / give away},
  type = {separable, strong / irregular},
  auxiliary = {haben},
  style = {default},
  present = {gebe ab,gibst ab,gibt ab,geben ab,gebt ab,geben ab},
  preterite = {gab ab,gabst ab,gab ab,gaben ab,gabt ab,gaben ab},
  subjunctive = {},
  imperative = {gib (du) ab!,geben wir ab!,gebt (ihr) ab!,geben Sie ab!},
  perfect = {haben + \textbf{abgegeben}},
  future = {werden + \textbf{abgeben}},
  pluperfect = {hatten + \textbf{abgegeben}},
  prepositions = {---},
  separable = {---}
}

% Migrated from verbs/b1/a.tex:4
\VerbTable{
  name = {abhängen},
  english = {to depend / hang},
  type = {separable, strong / irregular},
  auxiliary = {haben},
  style = {default},
  present = {hänge ab,hängst ab,hängt ab,hängen ab,hängt ab,hängen ab},
  preterite = {hing ab,hingst ab,hing ab,hingen ab,hingt ab,hingen ab},
  subjunctive = {},
  imperative = {häng (du) ab!,hängen wir ab!,hängt (ihr) ab!,hängen Sie ab!},
  perfect = {haben + \textbf{abgehangen}},
  future = {werden + \textbf{abhängen}},
  pluperfect = {hatten + \textbf{abgehangen}},
  prepositions = {---},
  separable = {---}
}

% Migrated from verbs/b1/a.tex:5
\VerbTable{
  name = {abheben},
  english = {to withdraw / take off},
  type = {separable, strong / irregular},
  auxiliary = {haben},
  style = {default},
  present = {hebe ab,hebst ab,hebt ab,heben ab,hebt ab,heben ab},
  preterite = {hob ab,hobst ab,hob ab,hoben ab,hobt ab,hoben ab},
  subjunctive = {},
  imperative = {heb (du) ab!,heben wir ab!,hebt (ihr) ab!,heben Sie ab!},
  perfect = {haben + \textbf{abgehoben}},
  future = {werden + \textbf{abheben}},
  pluperfect = {hatten + \textbf{abgehoben}},
  prepositions = {---},
  separable = {---}
}

% Migrated from verbs/b1/a.tex:6
\VerbTable{
  name = {abholen},
  english = {to pick up},
  type = {separable, regular},
  auxiliary = {haben},
  style = {default},
  present = {hole ab,holst ab,holt ab,holen ab,holt ab,holen ab},
  preterite = {holte ab,holtest ab,holte ab,holten ab,holtet ab,holten ab},
  subjunctive = {},
  imperative = {hol (du) ab!,holen wir ab!,holt (ihr) ab!,holen Sie ab!},
  perfect = {haben + \textbf{abgeholt}},
  future = {werden + \textbf{abholen}},
  pluperfect = {hatten + \textbf{abgeholt}},
  prepositions = {---},
  separable = {---}
}

% Migrated from verbs/b1/a.tex:7
\VerbTable{
  name = {ablehnen},
  english = {to reject},
  type = {separable, regular},
  auxiliary = {haben},
  style = {default},
  present = {lehne ab,lehnst ab,lehnt ab,lehnen ab,lehnt ab,lehnen ab},
  preterite = {lehnte ab,lehntest ab,lehnte ab,lehnten ab,lehntet ab,lehnten ab},
  subjunctive = {},
  imperative = {lehn (du) ab!,lehnen wir ab!,lehnt (ihr) ab!,lehnen Sie ab!},
  perfect = {haben + \textbf{abgelehnt}},
  future = {werden + \textbf{ablehnen}},
  pluperfect = {hatten + \textbf{abgelehnt}},
  prepositions = {---},
  separable = {---}
}

% Migrated from verbs/b1/a.tex:8
\VerbTable{
  name = {abmachen},
  english = {to arrange / agree},
  type = {separable, regular},
  auxiliary = {haben},
  style = {default},
  present = {mache ab,machst ab,macht ab,machen ab,macht ab,machen ab},
  preterite = {machte ab,machtest ab,machte ab,machten ab,machtet ab,machten ab},
  subjunctive = {},
  imperative = {mach (du) ab!,machen wir ab!,macht (ihr) ab!,machen Sie ab!},
  perfect = {haben + \textbf{abgemacht}},
  future = {werden + \textbf{abmachen}},
  pluperfect = {hatten + \textbf{abgemacht}},
  prepositions = {---},
  separable = {---}
}

% Migrated from verbs/b1/a.tex:9
\VerbTable{
  name = {abnehmen},
  english = {to take off / lose weight / decrease},
  type = {separable, strong / irregular},
  auxiliary = {haben},
  style = {default},
  present = {nehme ab,nimmst ab,nimmt ab,nehmen ab,nehmt ab,nehmen ab},
  preterite = {nahm ab,nahmst ab,nahm ab,nahmen ab,nahmt ab,nahmen ab},
  subjunctive = {},
  imperative = {nimm (du) ab!,nehmen wir ab!,nehmt (ihr) ab!,nehmen Sie ab!},
  perfect = {haben + \textbf{abgenommen}},
  future = {werden + \textbf{abnehmen}},
  pluperfect = {hatten + \textbf{abgenommen}},
  prepositions = {---},
  separable = {---}
}

% Migrated from verbs/b1/a.tex:10
\VerbTable{
  name = {abonnieren},
  english = {to subscribe},
  type = {regular},
  auxiliary = {haben},
  style = {default},
  present = {abonniere,abonnierst,abonniert,abonnieren,abonniert,abonnieren},
  preterite = {abonnierte,abonniertest,abonnierte,abonnierten,abonniertet,abonnierten},
  subjunctive = {},
  imperative = {abonnier (du)!,abonnieren wir!,abonniert (ihr)!,abonnieren Sie!},
  perfect = {haben + \textbf{abonniert}},
  future = {werden + \textbf{abonnieren}},
  pluperfect = {hatten + \textbf{abonniert}},
  prepositions = {---},
  separable = {---}
}

% Migrated from verbs/b1/a.tex:11
\VerbTable{
  name = {absagen},
  english = {to cancel},
  type = {separable, regular},
  auxiliary = {haben},
  style = {default},
  present = {sage ab,sagst ab,sagt ab,sagen ab,sagt ab,sagen ab},
  preterite = {sagte ab,sagtest ab,sagte ab,sagten ab,sagtet ab,sagten ab},
  subjunctive = {},
  imperative = {sag (du) ab!,sagen wir ab!,sagt (ihr) ab!,sagen Sie ab!},
  perfect = {haben + \textbf{abgesagt}},
  future = {werden + \textbf{absagen}},
  pluperfect = {hatten + \textbf{abgesagt}},
  prepositions = {---},
  separable = {---}
}

% Migrated from verbs/b1/a.tex:12
\VerbTable{
  name = {abschreiben},
  english = {to copy},
  type = {separable, strong / irregular},
  auxiliary = {haben},
  style = {default},
  present = {schreibe ab,schreibst ab,schreibt ab,schreiben ab,schreibt ab,schreiben ab},
  preterite = {schrieb ab,schriebst ab,schrieb ab,schrieben ab,schriebt ab,schrieben ab},
  subjunctive = {},
  imperative = {schreib (du) ab!,schreiben wir ab!,schreibt (ihr) ab!,schreiben Sie ab!},
  perfect = {haben + \textbf{abgeschrieben}},
  future = {werden + \textbf{abschreiben}},
  pluperfect = {hatten + \textbf{abgeschrieben}},
  prepositions = {---},
  separable = {---}
}

% Migrated from verbs/b1/a.tex:13
\VerbTable{
  name = {abstimmen},
  english = {to vote},
  type = {separable, regular},
  auxiliary = {haben},
  style = {default},
  present = {stimme ab,stimmst ab,stimmt ab,stimmen ab,stimmt ab,stimmen ab},
  preterite = {stimmte ab,stimmtest ab,stimmte ab,stimmten ab,stimmtet ab,stimmten ab},
  subjunctive = {},
  imperative = {stimm (du) ab!,stimmen wir ab!,stimmt (ihr) ab!,stimmen Sie ab!},
  perfect = {haben + \textbf{abgestimmt}},
  future = {werden + \textbf{abstimmen}},
  pluperfect = {hatten + \textbf{abgestimmt}},
  prepositions = {---},
  separable = {---}
}

% Migrated from verbs/b1/a.tex:14
\VerbTable{
  name = {abwaschen},
  english = {to wash up},
  type = {separable, strong / irregular},
  auxiliary = {haben},
  style = {default},
  present = {wasche ab,wäschst ab,wäscht ab,waschen ab,wascht ab,waschen ab},
  preterite = {wusch ab,wuschst ab,wusch ab,wuschen ab,wuscht ab,wuschen ab},
  subjunctive = {},
  imperative = {wasch (du) ab!,waschen wir ab!,wascht (ihr) ab!,waschen Sie ab!},
  perfect = {haben + \textbf{abgewaschen}},
  future = {werden + \textbf{abwaschen}},
  pluperfect = {hatten + \textbf{abgewaschen}},
  prepositions = {---},
  separable = {---}
}

% Migrated from verbs/a.tex:1
\VerbTable{
  name = {achten},
  english = {to pay attention to / to respect},
  type = {regular},
  auxiliary = {haben},
  style = {default},
  present = {achte,achtest,achtet,achten,achtet,achten},
  preterite = {achtete,achtetest,achtete,achteten,achtetet,achteten},
  subjunctive = {},
  imperative = {acht(e) (du)!,achten wir!,achtet (ihr)!,achten Sie!},
  perfect = {haben + \textbf{geachtet}},
  future = {werden + \textbf{achten}},
  pluperfect = {hatten + \textbf{geachtet}},
  prepositions = {
    & \textbf{achten auf} + Akk & (to pay attention to)\\
    & \textbf{achten} + Akk & (to respect / value)\\
  },
  separable = {---}
}

% Migrated from verbs/a.tex:4
\VerbTable{
  name = {ändern},
  english = {to change},
  type = {regular, + Akkusativ},
  auxiliary = {haben},
  style = {acc},
  present = {ändere,änderst,ändert,ändern,ändert,ändern},
  preterite = {änderte,ändertest,änderte,änderten,ändertet,änderten},
  subjunctive = {},
  imperative = {änder(e) (du)!,ändern wir!,ändert (ihr)!,ändern Sie!},
  perfect = {haben + \textbf{geändert}},
  future = {werden + \textbf{ändern}},
  pluperfect = {hatten + \textbf{geändert}},
  prepositions = {---},
  separable = {---}
}

% Migrated from verbs/a.tex:7
\VerbTable{
  name = {ärgern},
  english = {to annoy},
  type = {regular},
  auxiliary = {haben},
  style = {default},
  present = {ärgere,ärgerst,ärgert,ärgern,ärgert,ärgern},
  preterite = {ärgerte,ärgertest,ärgerte,ärgerten,ärgertet,ärgerten},
  subjunctive = {},
  imperative = {ärger(e) (du)!,ärgern wir!,ärgert (ihr)!,ärgern Sie!},
  perfect = {haben + \textbf{geärgert}},
  future = {werden + \textbf{ärgern}},
  pluperfect = {hatten + \textbf{geärgert}},
  prepositions = {& \textbf{ärgern über} + Akk & (to be annoyed about)\\},
  separable = {---}
}

% Migrated from verbs/a.tex:2
\VerbTable{
  name = {ahmen},
  english = {to imitate},
  type = {regular},
  auxiliary = {haben},
  style = {default},
  present = {ahme,ahmst,ahmt,ahmen,ahmt,ahmen},
  preterite = {ahmte,ahmtest,ahmte,ahmten,ahmtet,ahmten},
  subjunctive = {},
  imperative = {ahme (du)!,ahmen wir!,ahmt (ihr)!,ahmen Sie!},
  perfect = {haben + \textbf{geahmt}},
  future = {werden + \textbf{ahmen}},
  pluperfect = {hatten + \textbf{geahmt}},
  prepositions = {---},
  separable = {& \textbf{nach$|$ahmen} & (to imitate / emulate)\\}
}

% Migrated from verbs/a.tex:3
\VerbTable{
  name = {akzeptieren},
  english = {to accept},
  type = {regular},
  auxiliary = {haben},
  style = {default},
  present = {akzeptiere,akzeptierst,akzeptiert,akzeptieren,akzeptiert,akzeptieren},
  preterite = {akzeptierte,akzeptiertest,akzeptierte,akzeptierten,akzeptiertet,akzeptierten},
  subjunctive = {},
  imperative = {akzeptiere (du)!,akzeptieren wir!,akzeptiert (ihr)!,akzeptieren Sie!},
  perfect = {haben + \textbf{akzeptiert}},
  future = {werden + \textbf{akzeptieren}},
  pluperfect = {hatten + \textbf{akzeptiert}},
  prepositions = {& \textbf{akzeptieren} + Akk & (to accept something)\\},
  separable = {---}
}

% Migrated from verbs/b1/a.tex:16
\VerbTable{
  name = {analysieren},
  english = {to analyze},
  type = {regular},
  auxiliary = {haben},
  style = {default},
  present = {analysiere,analysierst,analysiert,analysieren,analysiert,analysieren},
  preterite = {analysierte,analysiertest,analysierte,analysierten,analysiertet,analysierten},
  subjunctive = {},
  imperative = {analysier (du)!,analysieren wir!,analysiert (ihr)!,analysieren Sie!},
  perfect = {haben + \textbf{analysiert}},
  future = {werden + \textbf{analysieren}},
  pluperfect = {hatten + \textbf{analysiert}},
  prepositions = {---},
  separable = {---}
}

% Migrated from verbs/b1/a.tex:17
\VerbTable{
  name = {anbieten},
  english = {to offer},
  type = {separable, strong / irregular},
  auxiliary = {haben},
  style = {default},
  present = {biete an,bietest an,bietet an,bieten an,bietet an,bieten an},
  preterite = {bot an,botest an,bot an,boten an,botet an,boten an},
  subjunctive = {},
  imperative = {biete (du) an!,bieten wir an!,bietet (ihr) an!,bieten Sie an!},
  perfect = {haben + \textbf{angeboten}},
  future = {werden + \textbf{anbieten}},
  pluperfect = {hatten + \textbf{angeboten}},
  prepositions = {---},
  separable = {---}
}

% Migrated from verbs/b1/a.tex:18
\VerbTable{
  name = {anerkennen},
  english = {to acknowledge},
  type = {separable, regular},
  auxiliary = {haben},
  style = {default},
  present = {erkenne an,erkennst an,erkennt an,erkennen an,erkennt an,erkennen an},
  preterite = {erkannte an,erkanntest an,erkannte an,erkannten an,erkanntet an,erkannten an},
  subjunctive = {},
  imperative = {erkenn (du) an!,erkennen wir an!,erkennt (ihr) an!,erkennen Sie an!},
  perfect = {haben + \textbf{anerkannt}},
  future = {werden + \textbf{anerkennen}},
  pluperfect = {hatten + \textbf{anerkannt}},
  prepositions = {---},
  separable = {---}
}

% Migrated from verbs/b1/a.tex:19
\VerbTable{
  name = {anfangen},
  english = {to get started},
  type = {separable, strong / irregular},
  auxiliary = {haben},
  style = {default},
  present = {fange an,fängst an,fängt an,fangen an,fangt an,fangen an},
  preterite = {fing an,fingst an,fing an,fingen an,fingt an,fingen an},
  subjunctive = {},
  imperative = {fang (du) an!,fangen wir an!,fangt (ihr) an!,fangen Sie an!},
  perfect = {haben + \textbf{angefangen}},
  future = {werden + \textbf{anfangen}},
  pluperfect = {hatten + \textbf{angefangen}},
  prepositions = {---},
  separable = {---}
}

% Migrated from verbs/b1/a.tex:20
\VerbTable{
  name = {angeben},
  english = {to specify},
  type = {separable, strong / irregular},
  auxiliary = {haben},
  style = {default},
  present = {gebe an,gibst an,gibt an,geben an,gebt an,geben an},
  preterite = {gab an,gabst an,gab an,gaben an,gabt an,gaben an},
  subjunctive = {},
  imperative = {gib (du) an!,geben wir an!,gebt (ihr) an!,geben Sie an!},
  perfect = {haben + \textbf{angegeben}},
  future = {werden + \textbf{angeben}},
  pluperfect = {hatten + \textbf{angegeben}},
  prepositions = {---},
  separable = {---}
}

% Migrated from verbs/b1/a.tex:21
\VerbTable{
  name = {anhaben},
  english = {to have on},
  type = {separable, regular},
  auxiliary = {haben},
  style = {default},
  present = {habe an,hast an,hat an,haben an,habt an,haben an},
  preterite = {hatte an,hattest an,hatte an,hatten an,hattet an,hatten an},
  subjunctive = {},
  imperative = {hab (du) an!,haben wir an!,habt (ihr) an!,haben Sie an!},
  perfect = {haben + \textbf{angehabt}},
  future = {werden + \textbf{anhaben}},
  pluperfect = {hatten + \textbf{angehabt}},
  prepositions = {---},
  separable = {---}
}

% Migrated from verbs/b1/a.tex:22
\VerbTable{
  name = {anklicken},
  english = {to click},
  type = {separable, regular},
  auxiliary = {haben},
  style = {default},
  present = {klicke an,klickst an,klickt an,klicken an,klickt an,klicken an},
  preterite = {klickte an,klicktest an,klickte an,klickten an,klicktet an,klickten an},
  subjunctive = {},
  imperative = {klick (du) an!,klicken wir an!,klickt (ihr) an!,klicken Sie an!},
  perfect = {haben + \textbf{angeklickt}},
  future = {werden + \textbf{anklicken}},
  pluperfect = {hatten + \textbf{angeklickt}},
  prepositions = {---},
  separable = {---}
}

% Migrated from verbs/b1/a.tex:23
\VerbTable{
  name = {ankommen},
  english = {to arrive},
  type = {separable, strong / irregular},
  auxiliary = {sein},
  style = {default},
  present = {komme an,kommst an,kommt an,kommen an,kommt an,kommen an},
  preterite = {kam an,kamst an,kam an,kamen an,kamt an,kamen an},
  subjunctive = {},
  imperative = {komm (du) an!,kommen wir an!,kommt (ihr) an!,kommen Sie an!},
  perfect = {sein + \textbf{angekommen}},
  future = {werden + \textbf{ankommen}},
  pluperfect = {waren + \textbf{angekommen}},
  prepositions = {---},
  separable = {---}
}

% Migrated from verbs/b1/a.tex:24
\VerbTable{
  name = {ankündigen},
  english = {to announce},
  type = {separable, regular},
  auxiliary = {haben},
  style = {default},
  present = {kündige an,kündigst an,kündigt an,kündigen an,kündigt an,kündigen an},
  preterite = {kündigte an,kündigtest an,kündigte an,kündigten an,kündigtet an,kündigten an},
  subjunctive = {},
  imperative = {kündig (du) an!,kündigen wir an!,kündigt (ihr) an!,kündigen Sie an!},
  perfect = {haben + \textbf{angekündigt}},
  future = {werden + \textbf{ankündigen}},
  pluperfect = {hatten + \textbf{angekündigt}},
  prepositions = {---},
  separable = {---}
}

% Migrated from verbs/b1/a.tex:25
\VerbTable{
  name = {anmelden},
  english = {to register / sign up},
  type = {separable, regular},
  auxiliary = {haben},
  style = {default},
  present = {melde an,meldest an,meldet an,melden an,meldet an,melden an},
  preterite = {meldete an,meldetest an,meldete an,meldeten an,meldetet an,meldeten an},
  subjunctive = {},
  imperative = {melde (du) an!,melden wir an!,meldet (ihr) an!,melden Sie an!},
  perfect = {haben + \textbf{angemeldet}},
  future = {werden + \textbf{anmelden}},
  pluperfect = {hatten + \textbf{angemeldet}},
  prepositions = {---},
  separable = {---}
}

% Migrated from verbs/b1/a.tex:26
\VerbTable{
  name = {annehmen},
  english = {to accept},
  type = {separable, strong / irregular},
  auxiliary = {haben},
  style = {default},
  present = {nehme an,nimmst an,nimmt an,nehmen an,nehmt an,nehmen an},
  preterite = {nahm an,nahmst an,nahm an,nahmen an,nahmt an,nahmen an},
  subjunctive = {},
  imperative = {nimm (du) an!,nehmen wir an!,nehmt (ihr) an!,nehmen Sie an!},
  perfect = {haben + \textbf{angenommen}},
  future = {werden + \textbf{annehmen}},
  pluperfect = {hatten + \textbf{angenommen}},
  prepositions = {---},
  separable = {---}
}

% Migrated from verbs/b1/a.tex:27
\VerbTable{
  name = {anrufen},
  english = {to call},
  type = {separable, strong / irregular},
  auxiliary = {haben},
  style = {default},
  present = {rufe an,rufst an,ruft an,rufen an,ruft an,rufen an},
  preterite = {rief an,riefst an,rief an,riefen an,rieft an,riefen an},
  subjunctive = {},
  imperative = {ruf (du) an!,rufen wir an!,ruft (ihr) an!,rufen Sie an!},
  perfect = {haben + \textbf{angerufen}},
  future = {werden + \textbf{anrufen}},
  pluperfect = {hatten + \textbf{angerufen}},
  prepositions = {---},
  separable = {---}
}

% Migrated from verbs/b1/a.tex:28
\VerbTable{
  name = {anschaffen},
  english = {to purchase / acquire},
  type = {separable, regular},
  auxiliary = {haben},
  style = {default},
  present = {schaffe an,schaffst an,schafft an,schaffen an,schafft an,schaffen an},
  preterite = {schaffte an,schafftest an,schaffte an,schafften an,schafftet an,schafften an},
  subjunctive = {},
  imperative = {schaff (du) an!,schaffen wir an!,schafft (ihr) an!,schaffen Sie an!},
  perfect = {haben + \textbf{angeschafft}},
  future = {werden + \textbf{anschaffen}},
  pluperfect = {hatten + \textbf{angeschafft}},
  prepositions = {---},
  separable = {---}
}

% Migrated from verbs/b1/a.tex:29
\VerbTable{
  name = {anschließen},
  english = {to connect / join},
  type = {separable, strong / irregular},
  auxiliary = {haben},
  style = {default},
  present = {schließe an,schließt an,schließt an,schließen an,schließt an,schließen an},
  preterite = {schloss an,schlossest an,schloss an,schlossen an,schlosst an,schlossen an},
  subjunctive = {},
  imperative = {schließ (du) an!,schließen wir an!,schließt (ihr) an!,schließen Sie an!},
  perfect = {haben + \textbf{angeschlossen}},
  future = {werden + \textbf{anschließen}},
  pluperfect = {hatten + \textbf{angeschlossen}},
  prepositions = {---},
  separable = {---}
}

% Migrated from verbs/b1/a.tex:30
\VerbTable{
  name = {anschnallen},
  english = {to buckle up},
  type = {separable, regular},
  auxiliary = {haben},
  style = {default},
  present = {schnalle an,schnallst an,schnallt an,schnallen an,schnallt an,schnallen an},
  preterite = {schnallte an,schnalltest an,schnallte an,schnallten an,schnalltet an,schnallten an},
  subjunctive = {},
  imperative = {schnall (du) an!,schnallen wir an!,schnallt (ihr) an!,schnallen Sie an!},
  perfect = {haben + \textbf{angeschnallt}},
  future = {werden + \textbf{anschnallen}},
  pluperfect = {hatten + \textbf{angeschnallt}},
  prepositions = {---},
  separable = {---}
}

% Migrated from verbs/b1/a.tex:31
\VerbTable{
  name = {ansehen},
  english = {to look at / watch},
  type = {separable, strong / irregular},
  auxiliary = {haben},
  style = {default},
  present = {sehe an,siehst an,sieht an,sehen an,seht an,sehen an},
  preterite = {sah an,sahst an,sah an,sahen an,saht an,sahen an},
  subjunctive = {},
  imperative = {sieh (du) an!,sehen wir an!,seht (ihr) an!,sehen Sie an!},
  perfect = {haben + \textbf{angesehen}},
  future = {werden + \textbf{ansehen}},
  pluperfect = {hatten + \textbf{angesehen}},
  prepositions = {---},
  separable = {---}
}

% Migrated from verbs/b1/a.tex:32
\VerbTable{
  name = {ansprechen},
  english = {to address / speak to},
  type = {separable, strong / irregular},
  auxiliary = {haben},
  style = {default},
  present = {spreche an,sprichst an,spricht an,sprechen an,sprecht an,sprechen an},
  preterite = {sprach an,sprachst an,sprach an,sprachen an,spracht an,sprachen an},
  subjunctive = {},
  imperative = {sprich (du) an!,sprechen wir an!,sprecht (ihr) an!,sprechen Sie an!},
  perfect = {haben + \textbf{angesprochen}},
  future = {werden + \textbf{ansprechen}},
  pluperfect = {hatten + \textbf{angesprochen}},
  prepositions = {---},
  separable = {---}
}

% Migrated from verbs/b1/a.tex:33
\VerbTable{
  name = {anstellen},
  english = {to employ / switch on},
  type = {separable, regular},
  auxiliary = {haben},
  style = {default},
  present = {stelle an,stellst an,stellt an,stellen an,stellt an,stellen an},
  preterite = {stellte an,stelltest an,stellte an,stellten an,stelltet an,stellten an},
  subjunctive = {},
  imperative = {stell (du) an!,stellen wir an!,stellt (ihr) an!,stellen Sie an!},
  perfect = {haben + \textbf{angestellt}},
  future = {werden + \textbf{anstellen}},
  pluperfect = {hatten + \textbf{angestellt}},
  prepositions = {---},
  separable = {---}
}

% Migrated from verbs/a.tex:5
\VerbTable{
  name = {antworten},
  english = {to answer},
  type = {regular, + Dativ},
  auxiliary = {haben},
  style = {dat},
  present = {antworte,antwortest,antwortet,antworten,antwortet,antworten},
  preterite = {antwortete,antwortetest,antwortete,antworteten,antwortetet,antworteten},
  subjunctive = {},
  imperative = {antwort(e) (du)!,antworten wir!,antwortet (ihr)!,antworten Sie!},
  perfect = {haben + \textbf{geantwortet}},
  future = {werden + \textbf{antworten}},
  pluperfect = {hatten + \textbf{geantwortet}},
  prepositions = {& \textbf{antworten auf} + Akk & (to answer something)\\},
  separable = {---}
}

% Migrated from verbs/b1/a.tex:35
\VerbTable{
  name = {anwenden},
  english = {to apply},
  type = {separable, mixed / irregular},
  auxiliary = {haben},
  style = {default},
  present = {wende an,wendest an,wendet an,wenden an,wendet an,wenden an},
  preterite = {wandte an,wandtest an,wandte an,wandten an,wandtet an,wandten an},
  subjunctive = {},
  imperative = {wend (du) an!,wenden wir an!,wendet (ihr) an!,wenden Sie an!},
  perfect = {haben + \textbf{angewandt / angewendet}},
  future = {werden + \textbf{anwenden}},
  pluperfect = {hatten + \textbf{angewandt / angewendet}},
  prepositions = {---},
  separable = {---}
}

% Migrated from verbs/b1/a.tex:36
\VerbTable{
  name = {anzeigen},
  english = {to display / report},
  type = {separable, regular},
  auxiliary = {haben},
  style = {default},
  present = {zeige an,zeigst an,zeigt an,zeigen an,zeigt an,zeigen an},
  preterite = {zeigte an,zeigtest an,zeigte an,zeigten an,zeigtet an,zeigten an},
  subjunctive = {},
  imperative = {zeig (du) an!,zeigen wir an!,zeigt (ihr) an!,zeigen Sie an!},
  perfect = {haben + \textbf{angezeigt}},
  future = {werden + \textbf{anzeigen}},
  pluperfect = {hatten + \textbf{angezeigt}},
  prepositions = {---},
  separable = {---}
}

% Migrated from verbs/b1/a.tex:37
\VerbTable{
  name = {anziehen},
  english = {to put on / get dressed},
  type = {separable, strong / irregular},
  auxiliary = {haben},
  style = {default},
  present = {ziehe an,ziehst an,zieht an,ziehen an,zieht an,ziehen an},
  preterite = {zog an,zogst an,zog an,zogen an,zogt an,zogen an},
  subjunctive = {},
  imperative = {zieh (du) an!,ziehen wir an!,zieht (ihr) an!,ziehen Sie an!},
  perfect = {haben + \textbf{angezogen}},
  future = {werden + \textbf{anziehen}},
  pluperfect = {hatten + \textbf{angezogen}},
  prepositions = {---},
  separable = {---}
}

% Migrated from verbs/a.tex:6
\VerbTable{
  name = {arbeiten},
  english = {to work},
  type = {regular},
  auxiliary = {haben},
  style = {default},
  present = {arbeite,arbeitest,arbeitet,arbeiten,arbeitet,arbeiten},
  preterite = {arbeitete,arbeitetest,arbeitete,arbeiteten,arbeitetet,arbeiteten},
  subjunctive = {},
  imperative = {arbeit(e) (du)!,arbeiten wir!,arbeitet (ihr)!,arbeiten Sie!},
  perfect = {haben + \textbf{gearbeitet}},
  future = {werden + \textbf{arbeiten}},
  pluperfect = {hatten + \textbf{gearbeitet}},
  prepositions = {
    & \textbf{arbeiten an} + Dat & (to work on)\\
    & \textbf{arbeiten für} + Akk & (to work for)\\
  },
  separable = {---}
}

% Migrated from verbs/a.tex:8
\VerbTable{
  name = {atmen},
  english = {to breathe},
  type = {regular},
  auxiliary = {haben},
  style = {default},
  present = {atme,atmest,atmet,atmen,atmet,atmen},
  preterite = {atmete,atmetest,atmete,atmeten,atmetet,atmeten},
  subjunctive = {},
  imperative = {atme (du)!,atmen wir!,atmet (ihr)!,atmen Sie!},
  perfect = {haben + \textbf{geatmet}},
  future = {werden + \textbf{atmen}},
  pluperfect = {hatten + \textbf{geatmet}},
  prepositions = {
    & \textbf{atmen in} + Akk & (to breathe into)\\
    & \textbf{atmen durch} + Akk & (to breathe through)\\
    & \textbf{atmen mit} + Dat & (to breathe with)\\
  },
  separable = {
    & \textbf{ein$|$atmen} & (to inhale)\\
    & \textbf{aus$|$atmen} & (to exhale)\\
    & \textbf{auf$|$atmen} & (to breathe again / breathe easier)\\
  }
}

% Migrated from verbs/b1/a.tex:38
\VerbTable{
  name = {auffallen},
  english = {to stand out / be noticed},
  type = {separable, strong / irregular},
  auxiliary = {sein},
  style = {default},
  present = {falle auf,fällst auf,fällt auf,fallen auf,fallt auf,fallen auf},
  preterite = {fiel auf,fielst auf,fiel auf,fielen auf,fielt auf,fielen auf},
  subjunctive = {},
  imperative = {fall (du) auf!,fallen wir auf!,fallt (ihr) auf!,fallen Sie auf!},
  perfect = {sein + \textbf{aufgefallen}},
  future = {werden + \textbf{auffallen}},
  pluperfect = {waren + \textbf{aufgefallen}},
  prepositions = {---},
  separable = {---}
}

% Migrated from verbs/b1/a.tex:39
\VerbTable{
  name = {auffordern},
  english = {to prompt},
  type = {separable, regular},
  auxiliary = {haben},
  style = {default},
  present = {fordre auf,forderst auf,fordert auf,fordern auf,fordert auf,fordern auf},
  preterite = {forderte auf,fordertest auf,forderte auf,forderten auf,fordertet auf,forderten auf},
  subjunctive = {},
  imperative = {fordre (du) auf!,fordern wir auf!,fordert (ihr) auf!,fordern Sie auf!},
  perfect = {haben + \textbf{aufgefordert}},
  future = {werden + \textbf{auffordern}},
  pluperfect = {hatten + \textbf{aufgefordert}},
  prepositions = {---},
  separable = {---}
}

% Migrated from verbs/b1/a.tex:40
\VerbTable{
  name = {aufführen},
  english = {to perform / list},
  type = {separable, regular},
  auxiliary = {haben},
  style = {default},
  present = {führe auf,führst auf,führt auf,führen auf,führt auf,führen auf},
  preterite = {führte auf,führtest auf,führte auf,führten auf,führtet auf,führten auf},
  subjunctive = {},
  imperative = {führ (du) auf!,führen wir auf!,führt (ihr) auf!,führen Sie auf!},
  perfect = {haben + \textbf{aufgeführt}},
  future = {werden + \textbf{aufführen}},
  pluperfect = {hatten + \textbf{aufgeführt}},
  prepositions = {---},
  separable = {---}
}

% Migrated from verbs/b1/a.tex:41
\VerbTable{
  name = {aufgeben},
  english = {to give up / post},
  type = {separable, strong / irregular},
  auxiliary = {haben},
  style = {default},
  present = {gebe auf,gibst auf,gibt auf,geben auf,gebt auf,geben auf},
  preterite = {gab auf,gabst auf,gab auf,gaben auf,gabt auf,gaben auf},
  subjunctive = {},
  imperative = {gib (du) auf!,geben wir auf!,gebt (ihr) auf!,geben Sie auf!},
  perfect = {haben + \textbf{aufgegeben}},
  future = {werden + \textbf{aufgeben}},
  pluperfect = {hatten + \textbf{aufgegeben}},
  prepositions = {---},
  separable = {---}
}

% Migrated from verbs/b1/a.tex:42
\VerbTable{
  name = {aufhalten},
  english = {to stop},
  type = {separable, strong / irregular},
  auxiliary = {haben},
  style = {default},
  present = {halte auf,hältst auf,hält auf,halten auf,haltet auf,halten auf},
  preterite = {hielt auf,hieltest auf,hielt auf,hielten auf,hieltet auf,hielten auf},
  subjunctive = {},
  imperative = {halte (du) auf!,halten wir auf!,haltet (ihr) auf!,halten Sie auf!},
  perfect = {haben + \textbf{aufgehalten}},
  future = {werden + \textbf{aufhalten}},
  pluperfect = {hatten + \textbf{aufgehalten}},
  prepositions = {---},
  separable = {---}
}

% Migrated from verbs/b1/a.tex:43
\VerbTable{
  name = {aufheben},
  english = {to pick up / keep / cancel},
  type = {separable, strong / irregular},
  auxiliary = {haben},
  style = {default},
  present = {hebe auf,hebst auf,hebt auf,heben auf,hebt auf,heben auf},
  preterite = {hob auf,hobst auf,hob auf,hoben auf,hobt auf,hoben auf},
  subjunctive = {},
  imperative = {heb (du) auf!,heben wir auf!,hebt (ihr) auf!,heben Sie auf!},
  perfect = {haben + \textbf{aufgehoben}},
  future = {werden + \textbf{aufheben}},
  pluperfect = {hatten + \textbf{aufgehoben}},
  prepositions = {---},
  separable = {---}
}

% Migrated from verbs/b1/a.tex:44
\VerbTable{
  name = {aufhören},
  english = {to stop},
  type = {separable, regular},
  auxiliary = {haben},
  style = {default},
  present = {höre auf,hörst auf,hört auf,hören auf,hört auf,hören auf},
  preterite = {hörte auf,hörtest auf,hörte auf,hörten auf,hörtet auf,hörten auf},
  subjunctive = {},
  imperative = {hör (du) auf!,hören wir auf!,hört (ihr) auf!,hören Sie auf!},
  perfect = {haben + \textbf{aufgehört}},
  future = {werden + \textbf{aufhören}},
  pluperfect = {hatten + \textbf{aufgehört}},
  prepositions = {---},
  separable = {---}
}

% Migrated from verbs/b1/a.tex:45
\VerbTable{
  name = {aufladen},
  english = {to recharge},
  type = {separable, strong / irregular},
  auxiliary = {haben},
  style = {default},
  present = {lade auf,lädst auf,lädt auf,laden auf,ladet auf,laden auf},
  preterite = {lud auf,ludest auf,lud auf,luden auf,ludet auf,luden auf},
  subjunctive = {},
  imperative = {lad (du) auf!,laden wir auf!,ladet (ihr) auf!,laden Sie auf!},
  perfect = {haben + \textbf{aufgeladen}},
  future = {werden + \textbf{aufladen}},
  pluperfect = {hatten + \textbf{aufgeladen}},
  prepositions = {---},
  separable = {---}
}

% Migrated from verbs/b1/a.tex:46
\VerbTable{
  name = {auflösen},
  english = {to dissolve},
  type = {separable, regular},
  auxiliary = {haben},
  style = {default},
  present = {löse auf,löst auf,löst auf,lösen auf,löst auf,lösen auf},
  preterite = {löste auf,löstest auf,löste auf,lösten auf,löstet auf,lösten auf},
  subjunctive = {},
  imperative = {lös (du) auf!,lösen wir auf!,löst (ihr) auf!,lösen Sie auf!},
  perfect = {haben + \textbf{aufgelöst}},
  future = {werden + \textbf{auflösen}},
  pluperfect = {hatten + \textbf{aufgelöst}},
  prepositions = {---},
  separable = {---}
}

% Migrated from verbs/b1/a.tex:47
\VerbTable{
  name = {aufnehmen},
  english = {to record / admit},
  type = {separable, strong / irregular},
  auxiliary = {haben},
  style = {default},
  present = {nehme auf,nimmst auf,nimmt auf,nehmen auf,nehmt auf,nehmen auf},
  preterite = {nahm auf,nahmst auf,nahm auf,nahmen auf,nahmt auf,nahmen auf},
  subjunctive = {},
  imperative = {nimm (du) auf!,nehmen wir auf!,nehmt (ihr) auf!,nehmen Sie auf!},
  perfect = {haben + \textbf{aufgenommen}},
  future = {werden + \textbf{aufnehmen}},
  pluperfect = {hatten + \textbf{aufgenommen}},
  prepositions = {---},
  separable = {---}
}

% Migrated from verbs/b1/a.tex:48
\VerbTable{
  name = {aufpassen},
  english = {to pay attention / watch out},
  type = {separable, regular},
  auxiliary = {haben},
  style = {default},
  present = {passe auf,passt auf,passt auf,passen auf,passt auf,passen auf},
  preterite = {passte auf,passtest auf,passte auf,passten auf,passtet auf,passten auf},
  subjunctive = {},
  imperative = {pass (du) auf!,passen wir auf!,passt (ihr) auf!,passen Sie auf!},
  perfect = {haben + \textbf{aufgepasst}},
  future = {werden + \textbf{aufpassen}},
  pluperfect = {hatten + \textbf{aufgepasst}},
  prepositions = {---},
  separable = {---}
}

% Migrated from verbs/b1/a.tex:49
\VerbTable{
  name = {aufräumen},
  english = {to clean up},
  type = {separable, regular},
  auxiliary = {haben},
  style = {default},
  present = {räume auf,räumst auf,räumt auf,räumen auf,räumt auf,räumen auf},
  preterite = {räumte auf,räumtest auf,räumte auf,räumten auf,räumtet auf,räumten auf},
  subjunctive = {},
  imperative = {räum (du) auf!,räumen wir auf!,räumt (ihr) auf!,räumen Sie auf!},
  perfect = {haben + \textbf{aufgeräumt}},
  future = {werden + \textbf{aufräumen}},
  pluperfect = {hatten + \textbf{aufgeräumt}},
  prepositions = {---},
  separable = {---}
}

% Migrated from verbs/b1/a.tex:50
\VerbTable{
  name = {aufregen},
  english = {to upset / get worked up},
  type = {separable, regular},
  auxiliary = {haben},
  style = {default},
  present = {rege auf,regst auf,regt auf,regen auf,regt auf,regen auf},
  preterite = {regte auf,regtest auf,regte auf,regten auf,regtet auf,regten auf},
  subjunctive = {},
  imperative = {reg (du) auf!,regen wir auf!,regt (ihr) auf!,regen Sie auf!},
  perfect = {haben + \textbf{aufgeregt}},
  future = {werden + \textbf{aufregen}},
  pluperfect = {hatten + \textbf{aufgeregt}},
  prepositions = {---},
  separable = {---}
}

% Migrated from verbs/b1/a.tex:51
\VerbTable{
  name = {aufstehen},
  english = {to stand up},
  type = {separable, strong / irregular},
  auxiliary = {sein},
  style = {default},
  present = {stehe auf,stehst auf,steht auf,stehen auf,steht auf,stehen auf},
  preterite = {stand auf,standest auf,stand auf,standen auf,standet auf,standen auf},
  subjunctive = {},
  imperative = {steh (du) auf!,stehen wir auf!,steht (ihr) auf!,stehen Sie auf!},
  perfect = {sein + \textbf{aufgestanden}},
  future = {werden + \textbf{aufstehen}},
  pluperfect = {waren + \textbf{aufgestanden}},
  prepositions = {---},
  separable = {---}
}

% Migrated from verbs/b1/a.tex:52
\VerbTable{
  name = {auftreten},
  english = {to appear / occur},
  type = {separable, strong / irregular},
  auxiliary = {sein},
  style = {default},
  present = {trete auf,trittst auf,tritt auf,treten auf,tretet auf,treten auf},
  preterite = {trat auf,tratest auf,trat auf,traten auf,tratet auf,traten auf},
  subjunctive = {},
  imperative = {tritt (du) auf!,treten wir auf!,tretet (ihr) auf!,treten Sie auf!},
  perfect = {sein + \textbf{aufgetreten}},
  future = {werden + \textbf{auftreten}},
  pluperfect = {waren + \textbf{aufgetreten}},
  prepositions = {---},
  separable = {---}
}

% Migrated from verbs/b1/a.tex:53
\VerbTable{
  name = {aufwachen},
  english = {to wake up},
  type = {separable, regular},
  auxiliary = {sein},
  style = {default},
  present = {wache auf,wachst auf,wacht auf,wachen auf,wacht auf,wachen auf},
  preterite = {wachte auf,wachtest auf,wachte auf,wachten auf,wachtet auf,wachten auf},
  subjunctive = {},
  imperative = {wach (du) auf!,wachen wir auf!,wacht (ihr) auf!,wachen Sie auf!},
  perfect = {sein + \textbf{aufgewacht}},
  future = {werden + \textbf{aufwachen}},
  pluperfect = {waren + \textbf{aufgewacht}},
  prepositions = {---},
  separable = {---}
}

% Migrated from verbs/b1/a.tex:54
\VerbTable{
  name = {ausdrucken},
  english = {to print},
  type = {separable, regular},
  auxiliary = {haben},
  style = {default},
  present = {drucke aus,druckst aus,druckt aus,drucken aus,druckt aus,drucken aus},
  preterite = {druckte aus,drucktest aus,druckte aus,druckten aus,drucktet aus,druckten aus},
  subjunctive = {},
  imperative = {druck (du) aus!,drucken wir aus!,druckt (ihr) aus!,drucken Sie aus!},
  perfect = {haben + \textbf{ausgedruckt}},
  future = {werden + \textbf{ausdrucken}},
  pluperfect = {hatten + \textbf{ausgedruckt}},
  prepositions = {---},
  separable = {---}
}

% Migrated from verbs/b1/a.tex:55
\VerbTable{
  name = {ausfallen},
  english = {to be cancelled / fail},
  type = {separable, strong / irregular},
  auxiliary = {sein},
  style = {default},
  present = {falle aus,fällst aus,fällt aus,fallen aus,fallt aus,fallen aus},
  preterite = {fiel aus,fielst aus,fiel aus,fielen aus,fielt aus,fielen aus},
  subjunctive = {},
  imperative = {fall (du) aus!,fallen wir aus!,fallt (ihr) aus!,fallen Sie aus!},
  perfect = {sein + \textbf{ausgefallen}},
  future = {werden + \textbf{ausfallen}},
  pluperfect = {waren + \textbf{ausgefallen}},
  prepositions = {---},
  separable = {---}
}

% Migrated from verbs/b1/a.tex:56
\VerbTable{
  name = {ausfüllen},
  english = {to fill in},
  type = {separable, regular},
  auxiliary = {haben},
  style = {default},
  present = {fülle aus,füllst aus,füllt aus,füllen aus,füllt aus,füllen aus},
  preterite = {füllte aus,fülltest aus,füllte aus,füllten aus,fülltet aus,füllten aus},
  subjunctive = {},
  imperative = {füll (du) aus!,füllen wir aus!,füllt (ihr) aus!,füllen Sie aus!},
  perfect = {haben + \textbf{ausgefüllt}},
  future = {werden + \textbf{ausfüllen}},
  pluperfect = {hatten + \textbf{ausgefüllt}},
  prepositions = {---},
  separable = {---}
}

% Migrated from verbs/b1/a.tex:57
\VerbTable{
  name = {ausgeben},
  english = {to spend / hand out},
  type = {separable, strong / irregular},
  auxiliary = {haben},
  style = {default},
  present = {gebe aus,gibst aus,gibt aus,geben aus,gebt aus,geben aus},
  preterite = {gab aus,gabst aus,gab aus,gaben aus,gabt aus,gaben aus},
  subjunctive = {},
  imperative = {gib (du) aus!,geben wir aus!,gebt (ihr) aus!,geben Sie aus!},
  perfect = {haben + \textbf{ausgegeben}},
  future = {werden + \textbf{ausgeben}},
  pluperfect = {hatten + \textbf{ausgegeben}},
  prepositions = {---},
  separable = {---}
}

% Migrated from verbs/b1/a.tex:58
\VerbTable{
  name = {ausgehen},
  english = {to go out / assume},
  type = {separable, strong / irregular},
  auxiliary = {sein},
  style = {default},
  present = {gehe aus,gehst aus,geht aus,gehen aus,geht aus,gehen aus},
  preterite = {ging aus,gingst aus,ging aus,gingen aus,gingt aus,gingen aus},
  subjunctive = {},
  imperative = {geh (du) aus!,gehen wir aus!,geht (ihr) aus!,gehen Sie aus!},
  perfect = {sein + \textbf{ausgegangen}},
  future = {werden + \textbf{ausgehen}},
  pluperfect = {waren + \textbf{ausgegangen}},
  prepositions = {---},
  separable = {---}
}

% Migrated from verbs/b1/a.tex:59
\VerbTable{
  name = {ausmachen},
  english = {to switch off / constitute},
  type = {separable, regular},
  auxiliary = {haben},
  style = {default},
  present = {mache aus,machst aus,macht aus,machen aus,macht aus,machen aus},
  preterite = {machte aus,machtest aus,machte aus,machten aus,machtet aus,machten aus},
  subjunctive = {},
  imperative = {mach (du) aus!,machen wir aus!,macht (ihr) aus!,machen Sie aus!},
  perfect = {haben + \textbf{ausgemacht}},
  future = {werden + \textbf{ausmachen}},
  pluperfect = {hatten + \textbf{ausgemacht}},
  prepositions = {---},
  separable = {---}
}

% Migrated from verbs/b1/a.tex:60
\VerbTable{
  name = {ausreichen},
  english = {to suffice},
  type = {separable, regular},
  auxiliary = {haben},
  style = {default},
  present = {reiche aus,reichst aus,reicht aus,reichen aus,reicht aus,reichen aus},
  preterite = {reichte aus,reichtest aus,reichte aus,reichten aus,reichtet aus,reichten aus},
  subjunctive = {},
  imperative = {reich (du) aus!,reichen wir aus!,reicht (ihr) aus!,reichen Sie aus!},
  perfect = {haben + \textbf{ausgereicht}},
  future = {werden + \textbf{ausreichen}},
  pluperfect = {hatten + \textbf{ausgereicht}},
  prepositions = {---},
  separable = {---}
}

% Migrated from verbs/b1/a.tex:61
\VerbTable{
  name = {ausrichten},
  english = {to pass on / align},
  type = {separable, regular},
  auxiliary = {haben},
  style = {default},
  present = {richte aus,richtest aus,richtet aus,richten aus,richtet aus,richten aus},
  preterite = {richtete aus,richtetest aus,richtete aus,richteten aus,richtetet aus,richteten aus},
  subjunctive = {},
  imperative = {richte (du) aus!,richten wir aus!,richtet (ihr) aus!,richten Sie aus!},
  perfect = {haben + \textbf{ausgerichtet}},
  future = {werden + \textbf{ausrichten}},
  pluperfect = {hatten + \textbf{ausgerichtet}},
  prepositions = {---},
  separable = {---}
}

% Migrated from verbs/b1/a.tex:62
\VerbTable{
  name = {ausruhen},
  english = {to rest},
  type = {separable, regular},
  auxiliary = {haben},
  style = {default},
  present = {ruhe aus,ruhst aus,ruht aus,ruhen aus,ruht aus,ruhen aus},
  preterite = {ruhte aus,ruhtest aus,ruhte aus,ruhten aus,ruhtet aus,ruhten aus},
  subjunctive = {},
  imperative = {ruh (du) aus!,ruhen wir aus!,ruht (ihr) aus!,ruhen Sie aus!},
  perfect = {haben + \textbf{ausgeruht}},
  future = {werden + \textbf{ausruhen}},
  pluperfect = {hatten + \textbf{ausgeruht}},
  prepositions = {---},
  separable = {---}
}

% Migrated from verbs/b1/a.tex:63
\VerbTable{
  name = {ausschließen},
  english = {to exclude},
  type = {separable, strong / irregular},
  auxiliary = {haben},
  style = {default},
  present = {schließe aus,schließt aus,schließt aus,schließen aus,schließt aus,schließen aus},
  preterite = {schloss aus,schlossest aus,schloss aus,schlossen aus,schlosst aus,schlossen aus},
  subjunctive = {},
  imperative = {schließ (du) aus!,schließen wir aus!,schließt (ihr) aus!,schließen Sie aus!},
  perfect = {haben + \textbf{ausgeschlossen}},
  future = {werden + \textbf{ausschließen}},
  pluperfect = {hatten + \textbf{ausgeschlossen}},
  prepositions = {---},
  separable = {---}
}

% Migrated from verbs/b1/a.tex:64
\VerbTable{
  name = {aussehen},
  english = {to look},
  type = {separable, strong / irregular},
  auxiliary = {haben},
  style = {default},
  present = {sehe aus,siehst aus,sieht aus,sehen aus,seht aus,sehen aus},
  preterite = {sah aus,sahst aus,sah aus,sahen aus,saht aus,sahen aus},
  subjunctive = {},
  imperative = {sieh (du) aus!,sehen wir aus!,seht (ihr) aus!,sehen Sie aus!},
  perfect = {haben + \textbf{ausgesehen}},
  future = {werden + \textbf{aussehen}},
  pluperfect = {hatten + \textbf{ausgesehen}},
  prepositions = {---},
  separable = {---}
}

% Migrated from verbs/b1/a.tex:65
\VerbTable{
  name = {aussprechen},
  english = {to pronounce},
  type = {separable, strong / irregular},
  auxiliary = {haben},
  style = {default},
  present = {spreche aus,sprichst aus,spricht aus,sprechen aus,sprecht aus,sprechen aus},
  preterite = {sprach aus,sprachst aus,sprach aus,sprachen aus,spracht aus,sprachen aus},
  subjunctive = {},
  imperative = {sprich (du) aus!,sprechen wir aus!,sprecht (ihr) aus!,sprechen Sie aus!},
  perfect = {haben + \textbf{ausgesprochen}},
  future = {werden + \textbf{aussprechen}},
  pluperfect = {hatten + \textbf{ausgesprochen}},
  prepositions = {---},
  separable = {---}
}

% Migrated from verbs/b1/a.tex:66
\VerbTable{
  name = {ausstellen},
  english = {to issue / exhibit},
  type = {separable, regular},
  auxiliary = {haben},
  style = {default},
  present = {stelle aus,stellst aus,stellt aus,stellen aus,stellt aus,stellen aus},
  preterite = {stellte aus,stelltest aus,stellte aus,stellten aus,stelltet aus,stellten aus},
  subjunctive = {},
  imperative = {stell (du) aus!,stellen wir aus!,stellt (ihr) aus!,stellen Sie aus!},
  perfect = {haben + \textbf{ausgestellt}},
  future = {werden + \textbf{ausstellen}},
  pluperfect = {hatten + \textbf{ausgestellt}},
  prepositions = {---},
  separable = {---}
}

% Migrated from verbs/b1/a.tex:67
\VerbTable{
  name = {aussuchen},
  english = {to choose},
  type = {separable, regular},
  auxiliary = {haben},
  style = {default},
  present = {suche aus,suchst aus,sucht aus,suchen aus,sucht aus,suchen aus},
  preterite = {suchte aus,suchtest aus,suchte aus,suchten aus,suchtet aus,suchten aus},
  subjunctive = {},
  imperative = {such (du) aus!,suchen wir aus!,sucht (ihr) aus!,suchen Sie aus!},
  perfect = {haben + \textbf{ausgesucht}},
  future = {werden + \textbf{aussuchen}},
  pluperfect = {hatten + \textbf{ausgesucht}},
  prepositions = {---},
  separable = {---}
}

% Migrated from verbs/b1/a.tex:68
\VerbTable{
  name = {auswählen},
  english = {to select},
  type = {separable, regular},
  auxiliary = {haben},
  style = {default},
  present = {wähle aus,wählst aus,wählt aus,wählen aus,wählt aus,wählen aus},
  preterite = {wählte aus,wähltest aus,wählte aus,wählten aus,wähltet aus,wählten aus},
  subjunctive = {},
  imperative = {wähl (du) aus!,wählen wir aus!,wählt (ihr) aus!,wählen Sie aus!},
  perfect = {haben + \textbf{ausgewählt}},
  future = {werden + \textbf{auswählen}},
  pluperfect = {hatten + \textbf{ausgewählt}},
  prepositions = {---},
  separable = {---}
}

% Migrated from verbs/b1/a.tex:69
\VerbTable{
  name = {ausziehen},
  english = {to take off / move out},
  type = {separable, strong / irregular},
  auxiliary = {haben / sein},
  style = {default},
  present = {ziehe aus,ziehst aus,zieht aus,ziehen aus,zieht aus,ziehen aus},
  preterite = {zog aus,zogst aus,zog aus,zogen aus,zogt aus,zogen aus},
  subjunctive = {},
  imperative = {zieh (du) aus!,ziehen wir aus!,zieht (ihr) aus!,ziehen Sie aus!},
  perfect = {haben / sein + \textbf{ausgezogen}},
  future = {werden + \textbf{ausziehen}},
  pluperfect = {hatten / waren + \textbf{ausgezogen}},
  prepositions = {---},
  separable = {---}
}

\section{B}

% Migrated from verbs/b1/b.tex:1
\VerbTable{
  name = {backen},
  english = {to bake},
  type = {regular},
  auxiliary = {haben},
  style = {default},
  present = {backe,backst,backt,backen,backt,backen},
  preterite = {backte,backtest,backte,backten,backtet,backten},
  subjunctive = {},
  imperative = {back (du)!,backen wir!,backt (ihr)!,backen Sie!},
  perfect = {haben + \textbf{gebacken}},
  future = {werden + \textbf{backen}},
  pluperfect = {hatten + \textbf{gebacken}},
  prepositions = {---},
  separable = {---}
}

% Migrated from verbs/b.tex:1
\VerbTable{
  name = {baden},
  english = {to bathe},
  type = {regular},
  auxiliary = {haben},
  style = {default},
  present = {bade,badest,badet,baden,badet,baden},
  preterite = {badete,badetest,badete,badeten,badetet,badeten},
  subjunctive = {},
  imperative = {bade (du)!,baden wir!,badet (ihr)!,baden Sie!},
  perfect = {haben + \textbf{gebadet}},
  future = {werden + \textbf{baden}},
  pluperfect = {hatten + \textbf{gebadet}},
  prepositions = {& \textbf{baden in} + Dat & (to bathe / swim in)\\},
  separable = {---}
}

% Migrated from verbs/b1/b.tex:2
\VerbTable{
  name = {basteln},
  english = {to make handicrafts},
  type = {regular},
  auxiliary = {haben},
  style = {default},
  present = {bastle,bastelst,bastelt,basteln,bastelt,basteln},
  preterite = {bastelte,basteltest,bastelte,bastelten,basteltet,bastelten},
  subjunctive = {},
  imperative = {bastle (du)!,basteln wir!,bastelt (ihr)!,basteln Sie!},
  perfect = {haben + \textbf{gebastelt}},
  future = {werden + \textbf{basteln}},
  pluperfect = {hatten + \textbf{gebastelt}},
  prepositions = {---},
  separable = {---}
}

% Migrated from verbs/b.tex:2
\VerbTable{
  name = {bauen},
  english = {to build},
  type = {regular},
  auxiliary = {haben},
  style = {default},
  present = {baue,baust,baut,bauen,baut,bauen},
  preterite = {baute,bautest,baute,bauten,bautet,bauten},
  subjunctive = {},
  imperative = {bau(e) (du)!,bauen wir!,baut (ihr)!,bauen Sie!},
  perfect = {haben + \textbf{gebaut}},
  future = {werden + \textbf{bauen}},
  pluperfect = {hatten + \textbf{gebaut}},
  prepositions = {
    & \textbf{bauen auf} + Akk & (to rely on / build on)\\
    & \textbf{bauen an} + Dat & (to be working on building)\\
    & \textbf{bauen in} + Akk & (to install into)\\
  },
  separable = {
    & \textbf{ab$|$bauen} & (to dismantle / reduce)\\
    & \textbf{auf$|$bauen} & (to build up / establish)\\
    & \textbf{ein$|$bauen} & (to install)\\
    & \textbf{um$|$bauen} & (to rebuild / remodel)
  }
}

% Migrated from verbs/b1/b.tex:3
\VerbTable{
  name = {beachten},
  english = {to note},
  type = {regular},
  auxiliary = {haben},
  style = {default},
  present = {beachte,beachtest,beachtet,beachten,beachtet,beachten},
  preterite = {beachtete,beachtetest,beachtete,beachteten,beachtetet,beachteten},
  subjunctive = {},
  imperative = {beachte (du)!,beachten wir!,beachtet (ihr)!,beachten Sie!},
  perfect = {haben + \textbf{beachtet}},
  future = {werden + \textbf{beachten}},
  pluperfect = {hatten + \textbf{beachtet}},
  prepositions = {---},
  separable = {---}
}

% Migrated from verbs/b.tex:3
\VerbTable{
  name = {beantragen},
  english = {to apply for},
  type = {regular, + Akkusativ},
  auxiliary = {haben},
  style = {acc},
  present = {beantrage,beantragst,beantragt,beantragen,beantragt,beantragen},
  preterite = {beantragte,beantragtest,beantragte,beantragten,beantragtet,beantragten},
  subjunctive = {},
  imperative = {beantrage (du)!,beantragen wir!,beantragt (ihr)!,beantragen Sie!},
  perfect = {haben + \textbf{beantragt}},
  future = {werden + \textbf{beantragen}},
  pluperfect = {hatten + \textbf{beantragt}},
  prepositions = {& \textbf{beantragen bei} + Dat & (to apply at)\\},
  separable = {---}
}

% Migrated from verbs/b1/b.tex:4
\VerbTable{
  name = {beantworten},
  english = {to answer},
  type = {regular},
  auxiliary = {haben},
  style = {default},
  present = {beantworte,beantwortest,beantwortet,beantworten,beantwortet,beantworten},
  preterite = {beantwortete,beantwortetest,beantwortete,beantworteten,beantwortetet,beantworteten},
  subjunctive = {},
  imperative = {beantworte (du)!,beantworten wir!,beantwortet (ihr)!,beantworten Sie!},
  perfect = {haben + \textbf{beantwortet}},
  future = {werden + \textbf{beantworten}},
  pluperfect = {hatten + \textbf{beantwortet}},
  prepositions = {---},
  separable = {---}
}

% Migrated from verbs/b.tex:4
\VerbTable{
  name = {bedeuten},
  english = {to mean something},
  type = {regular, + Akkusativ},
  auxiliary = {haben},
  style = {acc},
  present = {bedeute,bedeutest,bedeutet,bedeuten,bedeutet,bedeuten},
  preterite = {bedeutete,bedeutetest,bedeutete,bedeuteten,bedeutetet,bedeuteten},
  subjunctive = {},
  imperative = {bedeut(e) (du)!,bedeuten wir!,bedeutet (ihr)!,bedeuten Sie!},
  perfect = {haben + \textbf{bedeutet}},
  future = {werden + \textbf{bedeuten}},
  pluperfect = {hatten + \textbf{bedeutet}},
  prepositions = {---},
  separable = {---}
}

% Migrated from verbs/b1/b.tex:5
\VerbTable{
  name = {bedienen},
  english = {to serve / operate},
  type = {regular},
  auxiliary = {haben},
  style = {default},
  present = {bediene,bedienst,bedient,bedienen,bedient,bedienen},
  preterite = {bediente,bedientest,bediente,bedienten,bedientet,bedienten},
  subjunctive = {},
  imperative = {bedien (du)!,bedienen wir!,bedient (ihr)!,bedienen Sie!},
  perfect = {haben + \textbf{bedient}},
  future = {werden + \textbf{bedienen}},
  pluperfect = {hatten + \textbf{bedient}},
  prepositions = {---},
  separable = {---}
}

% Migrated from verbs/b1/b.tex:6
\VerbTable{
  name = {beeinflussen},
  english = {to influence},
  type = {regular},
  auxiliary = {haben},
  style = {default},
  present = {beeinflusse,beeinflusst,beeinflusst,beeinflussen,beeinflusst,beeinflussen},
  preterite = {beeinflusste,beeinflusstest,beeinflusste,beeinflussten,beeinflusstet,beeinflussten},
  subjunctive = {},
  imperative = {beeinfluss (du)!,beeinflussen wir!,beeinflusst (ihr)!,beeinflussen Sie!},
  perfect = {haben + \textbf{beeinflusst}},
  future = {werden + \textbf{beeinflussen}},
  pluperfect = {hatten + \textbf{beeinflusst}},
  prepositions = {---},
  separable = {---}
}

% Migrated from verbs/b1/b.tex:7
\VerbTable{
  name = {beenden},
  english = {to terminate},
  type = {regular},
  auxiliary = {haben},
  style = {default},
  present = {beende,beendest,beendet,beenden,beendet,beenden},
  preterite = {beendete,beendetest,beendete,beendeten,beendetet,beendeten},
  subjunctive = {},
  imperative = {beende (du)!,beenden wir!,beendet (ihr)!,beenden Sie!},
  perfect = {haben + \textbf{beendet}},
  future = {werden + \textbf{beenden}},
  pluperfect = {hatten + \textbf{beendet}},
  prepositions = {---},
  separable = {---}
}

% Migrated from verbs/b1/b.tex:8
\VerbTable{
  name = {begegnen},
  english = {to meet / encounter},
  type = {regular},
  auxiliary = {sein},
  style = {default},
  present = {begegne,begegnest,begegnet,begegnen,begegnet,begegnen},
  preterite = {begegnete,begegnetest,begegnete,begegneten,begegnetet,begegneten},
  subjunctive = {},
  imperative = {begegne (du)!,begegnen wir!,begegnet (ihr)!,begegnen Sie!},
  perfect = {sein + \textbf{begegnet}},
  future = {werden + \textbf{begegnen}},
  pluperfect = {waren + \textbf{begegnet}},
  prepositions = {---},
  separable = {---}
}

% Migrated from verbs/b.tex:5
\VerbTable{
  name = {beginnen},
  english = {to begin},
  type = {irregular, + Akkusativ},
  auxiliary = {haben},
  style = {acc},
  present = {beginne,beginnst,beginnt,beginnen,beginnt,beginnen},
  preterite = {begann,begannst,begann,begannen,begannt,begannen},
  subjunctive = {},
  imperative = {beginn(e) (du)!,beginnen wir!,beginnt (ihr)!,beginnen Sie!},
  perfect = {haben + \textbf{begonnen}},
  future = {werden + \textbf{beginnen}},
  pluperfect = {hatten + \textbf{begonnen}},
  prepositions = {---},
  separable = {---}
}

% Migrated from verbs/b1/b.tex:9
\VerbTable{
  name = {begleiten},
  english = {to accompany},
  type = {regular},
  auxiliary = {haben},
  style = {default},
  present = {begleite,begleitest,begleitet,begleiten,begleitet,begleiten},
  preterite = {begleitete,begleitetest,begleitete,begleiteten,begleitetet,begleiteten},
  subjunctive = {},
  imperative = {begleite (du)!,begleiten wir!,begleitet (ihr)!,begleiten Sie!},
  perfect = {haben + \textbf{begleitet}},
  future = {werden + \textbf{begleiten}},
  pluperfect = {hatten + \textbf{begleitet}},
  prepositions = {---},
  separable = {---}
}

% Migrated from verbs/b1/b.tex:10
\VerbTable{
  name = {begründen},
  english = {to justify},
  type = {regular},
  auxiliary = {haben},
  style = {default},
  present = {begründe,begründest,begründet,begründen,begründet,begründen},
  preterite = {begründete,begründetest,begründete,begründeten,begründetet,begründeten},
  subjunctive = {},
  imperative = {begründe (du)!,begründen wir!,begründet (ihr)!,begründen Sie!},
  perfect = {haben + \textbf{begründet}},
  future = {werden + \textbf{begründen}},
  pluperfect = {hatten + \textbf{begründet}},
  prepositions = {---},
  separable = {---}
}

% Migrated from verbs/b.tex:6
\VerbTable{
  name = {begrüßen},
  english = {to greet},
  type = {regular, + Akkusativ},
  auxiliary = {haben},
  style = {acc},
  present = {begrüße,begrüßt,begrüßt,begrüßen,begrüßt,begrüßen},
  preterite = {begrüßte,begrüßtest,begrüßte,begrüßten,begrüßtet,begrüßten},
  subjunctive = {},
  imperative = {begrüße (du)!,begrüßen wir!,begrüßt (ihr)!,begrüßen Sie!},
  perfect = {haben + \textbf{begrüßt}},
  future = {werden + \textbf{begrüßen}},
  pluperfect = {hatten + \textbf{begrüßt}},
  prepositions = {---},
  separable = {---}
}

% Migrated from verbs/b1/b.tex:11
\VerbTable{
  name = {behalten},
  english = {to keep},
  type = {strong / irregular},
  auxiliary = {haben},
  style = {default},
  present = {behalte,behältst,behält,behalten,behaltet,behalten},
  preterite = {behielt,behieltest,behielt,behielten,behieltet,behielten},
  subjunctive = {},
  imperative = {behalte (du)!,behalten wir!,behaltet (ihr)!,behalten Sie!},
  perfect = {haben + \textbf{behalten}},
  future = {werden + \textbf{behalten}},
  pluperfect = {hatten + \textbf{behalten}},
  prepositions = {---},
  separable = {---}
}

% Migrated from verbs/b.tex:7
\VerbTable{
  name = {behandeln},
  english = {to treat},
  type = {regular},
  auxiliary = {haben},
  style = {default},
  present = {behandle,behandelst,behandelt,behandeln,behandelt,behandeln},
  preterite = {behandelte,behandeltest,behandelte,behandelten,behandeltet,behandelten},
  subjunctive = {},
  imperative = {behandle (du)!,behandeln wir!,behandelt (ihr)!,behandeln Sie!},
  perfect = {haben + \textbf{behandelt}},
  future = {werde + \textbf{behandeln}},
  pluperfect = {hatten + \textbf{behandelt}},
  prepositions = {
    & \textbf{behandeln mit} + Dat & (to treat with)\\
    & \textbf{behandeln wegen} + Gen & (to treat because of / for)\\
  },
  separable = {---}
}

% Migrated from verbs/b1/b.tex:12
\VerbTable{
  name = {behaupten},
  english = {to claim},
  type = {regular},
  auxiliary = {haben},
  style = {default},
  present = {behaupte,behauptest,behauptet,behaupten,behauptet,behaupten},
  preterite = {behauptete,behauptetest,behauptete,behaupteten,behauptetet,behaupteten},
  subjunctive = {},
  imperative = {behaupte (du)!,behaupten wir!,behauptet (ihr)!,behaupten Sie!},
  perfect = {haben + \textbf{behauptet}},
  future = {werden + \textbf{behaupten}},
  pluperfect = {hatten + \textbf{behauptet}},
  prepositions = {---},
  separable = {---}
}

% Migrated from verbs/b1/b.tex:13
\VerbTable{
  name = {behindern},
  english = {to obstruct},
  type = {regular},
  auxiliary = {haben},
  style = {default},
  present = {behindre,behinderst,behindert,behindern,behindert,behindern},
  preterite = {behinderte,behindertest,behinderte,behinderten,behindertet,behinderten},
  subjunctive = {},
  imperative = {behindre (du)!,behindern wir!,behindert (ihr)!,behindern Sie!},
  perfect = {haben + \textbf{behindert}},
  future = {werden + \textbf{behindern}},
  pluperfect = {hatten + \textbf{behindert}},
  prepositions = {---},
  separable = {---}
}

% Migrated from verbs/b.tex:8
\VerbTable{
  name = {beißen},
  english = {to bite},
  type = {irregular, strong},
  auxiliary = {haben},
  style = {default},
  present = {beiße,beißt,beißt,beißen,beißt,beißen},
  preterite = {biss,bissest,biss,bissen,bisst,bissen},
  subjunctive = {},
  imperative = {beiß(e) (du)!,beißen wir!,beißt (ihr)!,beißen Sie!},
  perfect = {haben + \textbf{gebissen}},
  future = {werde + \textbf{beißen}},
  pluperfect = {hatten + \textbf{gebissen}},
  prepositions = {
    & \textbf{beißen in} + Akk & (to bite into)\\
    & \textbf{beißen nach} + Dat & (to snap at / bite at)\\
  },
  separable = {
    & \textbf{ab$|$beißen} & (to bite off)\\
    & \textbf{zu$|$beißen} & (to bite down / snap)\\
  }
}

% Migrated from verbs/b1/b.tex:14
\VerbTable{
  name = {bekanntgeben},
  english = {to announce},
  type = {separable, strong / irregular},
  auxiliary = {haben},
  style = {default},
  present = {gebe bekannt,gibst bekannt,gibt bekannt,geben bekannt,gebt bekannt,geben bekannt},
  preterite = {gab bekannt,gabst bekannt,gab bekannt,gaben bekannt,gabt bekannt,gaben bekannt},
  subjunctive = {},
  imperative = {gib (du) bekannt!,geben wir bekannt!,gebt (ihr) bekannt!,geben Sie bekannt!},
  perfect = {haben + \textbf{bekannt}},
  future = {werden + \textbf{bekanntgeben}},
  pluperfect = {hatten + \textbf{bekannt}},
  prepositions = {---},
  separable = {---}
}

% Migrated from verbs/b.tex:9
\VerbTable{
  name = {bekommen},
  english = {to receive},
  type = {irregular, + Akkusativ},
  auxiliary = {haben},
  style = {acc},
  present = {bekomme,bekommst,bekommt,bekommen,bekommt,bekommen},
  preterite = {bekam,bekamst,bekam,bekamen,bekamt,bekamen},
  subjunctive = {},
  imperative = {bekomm(e) (du)!,bekommen wir!,bekommt (ihr)!,bekommen Sie!},
  perfect = {haben + \textbf{bekommen}},
  future = {werden + \textbf{bekommen}},
  pluperfect = {hatten + \textbf{bekommen}},
  prepositions = {---},
  separable = {---}
}

% Migrated from verbs/b.tex:10
\VerbTable{
  name = {belasten},
  english = {to burden / to strain},
  type = {regular},
  auxiliary = {haben},
  style = {default},
  present = {belaste,belastest,belastet,belasten,belastet,belasten},
  preterite = {belastete,belastetest,belastete,belasteten,belastetet,belasteten},
  subjunctive = {},
  imperative = {belaste (du)!,belasten wir!,belastet (ihr)!,belasten Sie!},
  perfect = {haben + \textbf{belastet}},
  future = {werde + \textbf{belasten}},
  pluperfect = {hatten + \textbf{belastet}},
  prepositions = {
    & \textbf{belasten mit} + Dat & (to burden with)\\
    & \textbf{belasten durch} + Akk & (to strain/burden through)\\
  },
  separable = {---}
}

% Migrated from verbs/b.tex:11
\VerbTable{
  name = {belegen},
  english = {to occupy / register / cover},
  type = {regular, + Akkusativ},
  auxiliary = {haben},
  style = {acc},
  present = {belege,belegst,belegt,belegen,belegt,belegen},
  preterite = {belegte,belegtest,belegte,belegten,belegtet,belegten},
  subjunctive = {},
  imperative = {belege (du)!,belegen wir!,belegt (ihr)!,belegen Sie!},
  perfect = {haben + \textbf{belegt}},
  future = {werden + \textbf{belegen}},
  pluperfect = {hatten + \textbf{belegt}},
  prepositions = {& \textbf{belegen mit} + Dat & (to cover with)\\},
  separable = {---}
}

% Migrated from verbs/b1/b.tex:15
\VerbTable{
  name = {beleidigen},
  english = {to insult},
  type = {regular},
  auxiliary = {haben},
  style = {default},
  present = {beleidige,beleidigst,beleidigt,beleidigen,beleidigt,beleidigen},
  preterite = {beleidigte,beleidigtest,beleidigte,beleidigten,beleidigtet,beleidigten},
  subjunctive = {},
  imperative = {beleidig (du)!,beleidigen wir!,beleidigt (ihr)!,beleidigen Sie!},
  perfect = {haben + \textbf{beleidigt}},
  future = {werden + \textbf{beleidigen}},
  pluperfect = {hatten + \textbf{beleidigt}},
  prepositions = {---},
  separable = {---}
}

% Migrated from verbs/b.tex:12
\VerbTable{
  name = {bemerken},
  english = {to notice},
  type = {regular, + Akkusativ},
  auxiliary = {haben},
  style = {acc},
  present = {bemerke,bemerkst,bemerkt,bemerken,bemerkt,bemerken},
  preterite = {bemerkte,bemerktest,bemerkte,bemerkten,bemerktet,bemerkten},
  subjunctive = {},
  imperative = {bemerke (du)!,bemerken wir!,bemerkt (ihr)!,bemerken Sie!},
  perfect = {haben + \textbf{bemerkt}},
  future = {werden + \textbf{bemerken}},
  pluperfect = {hatten + \textbf{bemerkt}},
  prepositions = {---},
  separable = {---}
}

% Migrated from verbs/b1/b.tex:16
\VerbTable{
  name = {benötigen},
  english = {to need},
  type = {regular},
  auxiliary = {haben},
  style = {default},
  present = {benötige,benötigst,benötigt,benötigen,benötigt,benötigen},
  preterite = {benötigte,benötigtest,benötigte,benötigten,benötigtet,benötigten},
  subjunctive = {},
  imperative = {benötig (du)!,benötigen wir!,benötigt (ihr)!,benötigen Sie!},
  perfect = {haben + \textbf{benötigt}},
  future = {werden + \textbf{benötigen}},
  pluperfect = {hatten + \textbf{benötigt}},
  prepositions = {---},
  separable = {---}
}

% Migrated from verbs/b.tex:13
\VerbTable{
  name = {benutzen},
  english = {to use},
  type = {regular, + Akkusativ},
  auxiliary = {haben},
  style = {acc},
  present = {benutze,benutzt,benutz,benutzen,benutzt,benutzen},
  preterite = {benutzte,benutztest,benutzte,benutzten,benutztet,benutzten},
  subjunctive = {},
  imperative = {benutz(e) (du)!,benutzen wir!,benutzt (ihr)!,benutzen Sie!},
  perfect = {haben + \textbf{benutzt}},
  future = {werden + \textbf{benutzen}},
  pluperfect = {hatten + \textbf{benutzt}},
  prepositions = {---},
  separable = {---}
}

% Migrated from verbs/b.tex:14
\VerbTable{
  name = {beobachten},
  english = {to observe},
  type = {regular, + Akkusativ},
  auxiliary = {haben},
  style = {acc},
  present = {beobachte,beobachtest,beobachtet,beobachten,beobachtet,beobachten},
  preterite = {beobachtete,beobachtetest,beobachtete,beobachteten,beobachtetet,beobachteten},
  subjunctive = {},
  imperative = {beobachte (du)!,beobachten wir!,beobachtet (ihr)!,beobachten Sie!},
  perfect = {haben + \textbf{beobachtet}},
  future = {werden + \textbf{beobachten}},
  pluperfect = {hatten + \textbf{beobachtet}},
  prepositions = {& \textbf{beobachten Akk} & (to observe something)\\},
  separable = {---}
}

% Migrated from verbs/b.tex:15
\VerbTable{
  name = {beraten},
  english = {to advise},
  type = {irregular, + Akkusativ},
  auxiliary = {haben},
  style = {acc},
  present = {berate,berätst,berät,beraten,beratet,beraten},
  preterite = {beriet,berietst,beriet,berieten,berietet,berieten},
  subjunctive = {},
  imperative = {berate (du)!,beraten wir!,beratet (ihr)!,beraten Sie!},
  perfect = {haben + \textbf{beraten}},
  future = {werden + \textbf{beraten}},
  pluperfect = {hatten + \textbf{beraten}},
  prepositions = {
    & \textbf{beraten über} + Akk & (to advise about)\\
    & \textbf{beraten bei} + Dat & (to advise during / assist with)\\
  },
  separable = {---}
}

% Migrated from verbs/b1/b.tex:17
\VerbTable{
  name = {berechnen},
  english = {to calculate},
  type = {regular},
  auxiliary = {haben},
  style = {default},
  present = {berechne,berechnest,berechnet,berechnen,berechnet,berechnen},
  preterite = {berechnete,berechnetest,berechnete,berechneten,berechnetet,berechneten},
  subjunctive = {},
  imperative = {berechne (du)!,berechnen wir!,berechnet (ihr)!,berechnen Sie!},
  perfect = {haben + \textbf{berechnet}},
  future = {werden + \textbf{berechnen}},
  pluperfect = {hatten + \textbf{berechnet}},
  prepositions = {---},
  separable = {---}
}

% Migrated from verbs/b.tex:16
\VerbTable{
  name = {bereiten},
  english = {to prepare / cause},
  type = {regular, + Akkusativ},
  auxiliary = {haben},
  style = {acc},
  present = {bereite,bereitest,bereitet,bereiten,bereitet,bereiten},
  preterite = {bereitete,bereitetest,bereitete,bereiteten,bereitetet,bereiteten},
  subjunctive = {},
  imperative = {bereite (du)!,bereiten wir!,bereitet (ihr)!,bereiten Sie!},
  perfect = {haben + \textbf{bereitet}},
  future = {werden + \textbf{bereiten}},
  pluperfect = {hatten + \textbf{bereitet}},
  prepositions = {& \textbf{bereiten auf} + Akk & (to prepare for)\\},
  separable = {& \textbf{vor$|$bereiten} & (to prepare)\\}
}

% Migrated from verbs/b1/b.tex:18
\VerbTable{
  name = {berichten},
  english = {to report},
  type = {regular},
  auxiliary = {haben},
  style = {default},
  present = {berichte,berichtest,berichtet,berichten,berichtet,berichten},
  preterite = {berichtete,berichtetest,berichtete,berichteten,berichtetet,berichteten},
  subjunctive = {},
  imperative = {berichte (du)!,berichten wir!,berichtet (ihr)!,berichten Sie!},
  perfect = {haben + \textbf{berichtet}},
  future = {werden + \textbf{berichten}},
  pluperfect = {hatten + \textbf{berichtet}},
  prepositions = {---},
  separable = {---}
}

% Migrated from verbs/b1/b.tex:19
\VerbTable{
  name = {beruhigen},
  english = {to reassure},
  type = {regular},
  auxiliary = {haben},
  style = {default},
  present = {beruhige,beruhigst,beruhigt,beruhigen,beruhigt,beruhigen},
  preterite = {beruhigte,beruhigtest,beruhigte,beruhigten,beruhigtet,beruhigten},
  subjunctive = {},
  imperative = {beruhig (du)!,beruhigen wir!,beruhigt (ihr)!,beruhigen Sie!},
  perfect = {haben + \textbf{beruhigt}},
  future = {werden + \textbf{beruhigen}},
  pluperfect = {hatten + \textbf{beruhigt}},
  prepositions = {---},
  separable = {---}
}

% Migrated from verbs/b1/b.tex:20
\VerbTable{
  name = {beschädigen},
  english = {to damage},
  type = {regular},
  auxiliary = {haben},
  style = {default},
  present = {beschädige,beschädigst,beschädigt,beschädigen,beschädigt,beschädigen},
  preterite = {beschädigte,beschädigtest,beschädigte,beschädigten,beschädigtet,beschädigten},
  subjunctive = {},
  imperative = {beschädig (du)!,beschädigen wir!,beschädigt (ihr)!,beschädigen Sie!},
  perfect = {haben + \textbf{beschädigt}},
  future = {werden + \textbf{beschädigen}},
  pluperfect = {hatten + \textbf{beschädigt}},
  prepositions = {---},
  separable = {---}
}

% Migrated from verbs/b.tex:18
\VerbTable{
  name = {beschließen},
  english = {to decide},
  type = {irregular, + Akkusativ},
  auxiliary = {haben},
  style = {acc},
  present = {beschließe,beschließt,beschließt,beschließen,beschließt,beschließen},
  preterite = {beschloss,beschlossest,beschloss,beschlossen,beschlosst,beschlossen},
  subjunctive = {},
  imperative = {beschließ(e) (du)!,beschließen wir!,beschließt (ihr)!,beschließen Sie!},
  perfect = {haben + \textbf{beschlossen}},
  future = {werden + \textbf{beschließen}},
  pluperfect = {hatten + \textbf{beschlossen}},
  prepositions = {& \textbf{beschließen Akk} & (to decide something)\\},
  separable = {---}
}

% Migrated from verbs/b.tex:19
\VerbTable{
  name = {beschreiben},
  english = {to describe},
  type = {irregular, + Akkusativ},
  auxiliary = {haben},
  style = {acc},
  present = {beschreibe,beschreibst,beschreibt,beschreiben,beschreibt,beschreiben},
  preterite = {beschrieb,beschriebst,beschrieb,beschrieben,beschriebt,beschrieben},
  subjunctive = {},
  imperative = {beschreib(e) (du)!,beschreiben wir!,beschreibt (ihr)!,beschreiben Sie!},
  perfect = {haben + \textbf{beschrieben}},
  future = {werden + \textbf{beschreiben}},
  pluperfect = {hatten + \textbf{beschrieben}},
  prepositions = {---},
  separable = {---}
}

% Migrated from verbs/b.tex:17
\VerbTable{
  name = {besetzen},
  english = {to occupy / fill a position},
  type = {regular, + Akkusativ},
  auxiliary = {haben},
  style = {acc},
  present = {besetze,besetzt,besetzt,besetzen,besetzt,besetzen},
  preterite = {besetzte,besetztest,besetzte,besetzten,besetztet,besetzten},
  subjunctive = {},
  imperative = {besetze (du)!,besetzen wir!,besetzt (ihr)!,besetzen Sie!},
  perfect = {haben + \textbf{besetzt}},
  future = {werden + \textbf{besetzen}},
  pluperfect = {hatten + \textbf{besetzt}},
  prepositions = {& \textbf{besetzen mit} + Dat & (to fill with / occupy with)\\},
  separable = {---}
}

% Migrated from verbs/b1/b.tex:21
\VerbTable{
  name = {besichtigen},
  english = {to inspect / tour},
  type = {regular},
  auxiliary = {haben},
  style = {default},
  present = {besichtige,besichtigst,besichtigt,besichtigen,besichtigt,besichtigen},
  preterite = {besichtigte,besichtigtest,besichtigte,besichtigten,besichtigtet,besichtigten},
  subjunctive = {},
  imperative = {besichtig (du)!,besichtigen wir!,besichtigt (ihr)!,besichtigen Sie!},
  perfect = {haben + \textbf{besessen}},
  future = {werden + \textbf{besichtigen}},
  pluperfect = {hatten + \textbf{besessen}},
  prepositions = {---},
  separable = {---}
}

% Migrated from verbs/b1/b.tex:22
\VerbTable{
  name = {besitzen},
  english = {to own},
  type = {strong / irregular},
  auxiliary = {haben},
  style = {default},
  present = {besitze,besitzt,besitzt,besitzen,besitzt,besitzen},
  preterite = {besaß,besaßt,besaß,besaßen,besaßt,besaßen},
  subjunctive = {},
  imperative = {besitz (du)!,besitzen wir!,besitzt (ihr)!,besitzen Sie!},
  perfect = {haben + \textbf{besessen}},
  future = {werden + \textbf{besitzen}},
  pluperfect = {hatten + \textbf{besessen}},
  prepositions = {---},
  separable = {---}
}

% Migrated from verbs/b1/b.tex:23
\VerbTable{
  name = {besorgen},
  english = {to obtain / take care of},
  type = {regular},
  auxiliary = {haben},
  style = {default},
  present = {besorge,besorgst,besorgt,besorgen,besorgt,besorgen},
  preterite = {besorgte,besorgtest,besorgte,besorgten,besorgtet,besorgten},
  subjunctive = {},
  imperative = {besorg (du)!,besorgen wir!,besorgt (ihr)!,besorgen Sie!},
  perfect = {haben + \textbf{besorgt}},
  future = {werden + \textbf{besorgen}},
  pluperfect = {hatten + \textbf{besorgt}},
  prepositions = {---},
  separable = {---}
}

% Migrated from verbs/b1/b.tex:24
\VerbTable{
  name = {besprechen},
  english = {to discuss},
  type = {strong / irregular},
  auxiliary = {haben},
  style = {default},
  present = {bespreche,besprichst,bespricht,besprechen,besprecht,besprechen},
  preterite = {besprach,besprachst,besprach,besprachen,bespracht,besprachen},
  subjunctive = {},
  imperative = {besprich (du)!,besprechen wir!,besprecht (ihr)!,besprechen Sie!},
  perfect = {haben + \textbf{besprochen}},
  future = {werden + \textbf{besprechen}},
  pluperfect = {hatten + \textbf{besprochen}},
  prepositions = {---},
  separable = {---}
}

% Migrated from verbs/b.tex:20
\VerbTable{
  name = {bestätigen},
  english = {to confirm},
  type = {regular, + Akkusativ},
  auxiliary = {haben},
  style = {acc},
  present = {bestätige,bestätigst,bestätigt,bestätigen,bestätigt,bestätigen},
  preterite = {bestätigte,bestätigtest,bestätigte,bestätigten,bestätigtet,bestätigten},
  subjunctive = {},
  imperative = {bestätige (du)!,bestätigen wir!,bestätigt (ihr)!,bestätigen Sie!},
  perfect = {haben + \textbf{bestätigt}},
  future = {werden + \textbf{bestätigen}},
  pluperfect = {hatten + \textbf{bestätigt}},
  prepositions = {& \textbf{bestätigen gegenüber} + Dat & (to confirm to someone)\\},
  separable = {---}
}

% Migrated from verbs/b.tex:21
\VerbTable{
  name = {bestehen},
  english = {to pass},
  type = {irregular, + Akkusativ},
  auxiliary = {haben},
  style = {acc},
  present = {bestehe,bestehst,besteht,bestehen,besteht,bestehen},
  preterite = {bestand,bestandest,bestand,bestanden,bestandet,bestanden},
  subjunctive = {},
  imperative = {besteh(e) (du)!,bestehen wir!,besteht (ihr)!,bestehen Sie!},
  perfect = {haben + \textbf{bestanden}},
  future = {werden + \textbf{bestehen}},
  pluperfect = {hatten + \textbf{bestanden}},
  prepositions = {---},
  separable = {---}
}

% Migrated from verbs/b.tex:22
\VerbTable{
  name = {bestellen},
  english = {to order},
  type = {regular, + Akkusativ},
  auxiliary = {haben},
  style = {acc},
  present = {bestelle,bestellst,bestellt,bestellen,bestellt,bestellen},
  preterite = {bestellte,bestelltest,bestellte,bestellten,bestelltet,bestellten},
  subjunctive = {},
  imperative = {bestell(e) (du)!,bestellen wir!,bestellt (ihr)!,bestellen Sie!},
  perfect = {haben + \textbf{bestellt}},
  future = {werden + \textbf{bestellen}},
  pluperfect = {hatten + \textbf{bestellt}},
  prepositions = {---},
  separable = {---}
}

% Migrated from verbs/b1/b.tex:25
\VerbTable{
  name = {bestrafen},
  english = {to punish},
  type = {regular},
  auxiliary = {haben},
  style = {default},
  present = {bestrafe,bestrafst,bestraft,bestrafen,bestraft,bestrafen},
  preterite = {bestrafte,bestraftest,bestrafte,bestraften,bestraftet,bestraften},
  subjunctive = {},
  imperative = {bestraf (du)!,bestrafen wir!,bestraft (ihr)!,bestrafen Sie!},
  perfect = {haben + \textbf{bestraft}},
  future = {werden + \textbf{bestrafen}},
  pluperfect = {hatten + \textbf{bestraft}},
  prepositions = {---},
  separable = {---}
}

% Migrated from verbs/b1/b.tex:26
\VerbTable{
  name = {besuchen},
  english = {to visit},
  type = {regular},
  auxiliary = {haben},
  style = {default},
  present = {besuche,besuchst,besucht,besuchen,besucht,besuchen},
  preterite = {besuchte,besuchtest,besuchte,besuchten,besuchtet,besuchten},
  subjunctive = {},
  imperative = {besuch (du)!,besuchen wir!,besucht (ihr)!,besuchen Sie!},
  perfect = {haben + \textbf{besucht}},
  future = {werden + \textbf{besuchen}},
  pluperfect = {hatten + \textbf{besucht}},
  prepositions = {---},
  separable = {---}
}

% Migrated from verbs/b.tex:23
\VerbTable{
  name = {betragen},
  english = {to amount to / to be},
  type = {irregular, + Akkusativ},
  auxiliary = {haben},
  style = {acc},
  present = {betrage,beträgst,beträgt,betragen,betragt,betragen},
  preterite = {betrug,betrugst,betrug,betrugen,betrugt,betrugen},
  subjunctive = {},
  imperative = {---},
  perfect = {haben + \textbf{betragen}},
  future = {werden + \textbf{betragen}},
  pluperfect = {hatten + \textbf{betragen}},
  prepositions = {---},
  separable = {---}
}

% Migrated from verbs/b.tex:24
\VerbTable{
  name = {betreuen},
  english = {to supervise / take care of},
  type = {regular, + Akkusativ},
  auxiliary = {haben},
  style = {acc},
  present = {betreue,betreust,betreut,betreuen,betreut,betreuen},
  preterite = {betreute,betreutest,betreute,betreuten,betreutet,betreuten},
  subjunctive = {},
  imperative = {betreue (du)!,betreuen wir!,betreut (ihr)!,betreuen Sie!},
  perfect = {haben + \textbf{betreut}},
  future = {werden + \textbf{betreuen}},
  pluperfect = {hatten + \textbf{betreut}},
  prepositions = {
    & \textbf{betreuen bei} + Dat & (to assist during)\\
    & \textbf{betreuen in} + Dat & (to supervise in)\\
  },
  separable = {---}
}

% Migrated from verbs/b1/b.tex:28
\VerbTable{
  name = {betrügen},
  english = {to cheat},
  type = {strong / irregular},
  auxiliary = {haben},
  style = {default},
  present = {betrüge,betrügst,betrügt,betrügen,betrügt,betrügen},
  preterite = {betrog,betrogst,betrog,betrogen,betrogt,betrogen},
  subjunctive = {},
  imperative = {betrüg (du)!,betrügen wir!,betrügt (ihr)!,betrügen Sie!},
  perfect = {haben + \textbf{betrogen}},
  future = {werden + \textbf{betrügen}},
  pluperfect = {hatten + \textbf{betrogen}},
  prepositions = {---},
  separable = {---}
}

% Migrated from verbs/b1/b.tex:29
\VerbTable{
  name = {beweisen},
  english = {to prove},
  type = {strong / irregular},
  auxiliary = {haben},
  style = {default},
  present = {beweise,beweist,beweist,beweisen,beweist,beweisen},
  preterite = {bewies,bewiest,bewies,bewiesen,bewiest,bewiesen},
  subjunctive = {},
  imperative = {beweis (du)!,beweisen wir!,beweist (ihr)!,beweisen Sie!},
  perfect = {haben + \textbf{bewiesen}},
  future = {werden + \textbf{beweisen}},
  pluperfect = {hatten + \textbf{bewiesen}},
  prepositions = {---},
  separable = {---}
}

% Migrated from verbs/b.tex:25
\VerbTable{
  name = {bewerben},
  english = {to advertise / promote},
  type = {irregular, + Akkusativ},
  auxiliary = {haben},
  style = {acc},
  present = {bewerbe,bewirbst,bewirbt,bewerben,bewerbt,bewerben},
  preterite = {bewarb,bewarbst,bewarb,bewarben,bewarbt,bewarben},
  subjunctive = {},
  imperative = {bewirb (du)!,bewerben wir!,bewerbt (ihr)!,bewerben Sie!},
  perfect = {haben + \textbf{beworben}},
  future = {werden + \textbf{bewerben}},
  pluperfect = {hatten + \textbf{beworben}},
  prepositions = {---},
  separable = {---}
}

% Migrated from verbs/b.tex:26
\VerbTable{
  name = {bezahlen},
  english = {to pay},
  type = {regular, + Akkusativ},
  auxiliary = {haben},
  style = {acc},
  present = {bezahle,bezahlst,bezahlt,bezahlen,bezahlt,bezahlen},
  preterite = {bezahlte,bezahltest,bezahlte,bezahlten,bezahltet,bezahlten},
  subjunctive = {},
  imperative = {bezahl(e) (du)!,bezahlen wir!,bezahlt (ihr)!,bezahlen Sie!},
  perfect = {haben + \textbf{bezahlt}},
  future = {werden + \textbf{bezahlen}},
  pluperfect = {hatten + \textbf{bezahlt}},
  prepositions = {---},
  separable = {---}
}

% Migrated from verbs/b.tex:27
\VerbTable{
  name = {beziehen},
  english = {to obtain / refer to},
  type = {irregular, + Akkusativ},
  auxiliary = {haben},
  style = {acc},
  present = {beziehe,beziehst,bezieht,beziehen,bezieht,beziehen},
  preterite = {bezog,bezogst,bezog,bezogen,bezogt,bezogen},
  subjunctive = {},
  imperative = {bezieh(e) (du)!,beziehen wir!,bezieht (ihr)!,beziehen Sie!},
  perfect = {haben + \textbf{bezogen}},
  future = {werden + \textbf{beziehen}},
  pluperfect = {hatten + \textbf{bezogen}},
  prepositions = {
    & \textbf{beziehen Akk} & (to obtain something)\\
    & \textbf{beziehen auf} + Akk & (to refer to)\\
  },
  separable = {
    & \textbf{ein$|$beziehen} & (to involve / include)\\
    & \textbf{bevor$|$ziehen} & (to prefer)\\
  }
}

% Migrated from verbs/b.tex:28
\VerbTable{
  name = {bieten},
  english = {to offer},
  type = {irregular},
  auxiliary = {haben},
  style = {default},
  present = {biete,bietest,bietet,bieten,bietet,bieten},
  preterite = {bot,botest,bot,boten,botet,boten},
  subjunctive = {},
  imperative = {biete (du)!,bieten wir!,bietet (ihr)!,bieten Sie!},
  perfect = {haben + \textbf{geboten}},
  future = {werden + \textbf{bieten}},
  pluperfect = {hatten + \textbf{geboten}},
  prepositions = {
    & \textbf{bieten} + Dat + Akk & (to offer someone something)\\
    & \textbf{bieten für} + Akk & (to bid for)\\
  },
  separable = {& \textbf{an$|$bieten} & (to offer)\\}
}

% Migrated from verbs/b.tex:29
\VerbTable{
  name = {bitten},
  english = {to ask / request},
  type = {irregular},
  auxiliary = {haben},
  style = {default},
  present = {bitte,bittest,bittet,bitten,bittet,bitten},
  preterite = {bat,batst,bat,baten,batet,baten},
  subjunctive = {},
  imperative = {bitte (du)!,bitten wir!,bittet (ihr)!,bitten Sie!},
  perfect = {haben + \textbf{gebeten}},
  future = {werden + \textbf{bitten}},
  pluperfect = {hatten + \textbf{gebeten}},
  prepositions = {& \textbf{bitten um} + Akk & (to ask for / request)\\},
  separable = {---}
}

% Migrated from verbs/b.tex:30
\VerbTable{
  name = {bleiben},
  english = {to stay},
  type = {irregular},
  auxiliary = {sein},
  style = {default},
  present = {bleibe,bleibst,bleibt,bleiben,bleibt,bleiben},
  preterite = {blieb,bliebst,blieb,blieben,bliebt,blieben},
  subjunctive = {},
  imperative = {bleib(e) (du)!,bleiben wir!,bleibt (ihr)!,bleiben Sie!},
  perfect = {sein + \textbf{geblieben}},
  future = {werden + \textbf{bleiben}},
  pluperfect = {waren + \textbf{geblieben}},
  prepositions = {---},
  separable = {---}
}

% Migrated from verbs/b1/b.tex:30
\VerbTable{
  name = {blitzen},
  english = {to flash},
  type = {regular},
  auxiliary = {haben},
  style = {default},
  present = {blitze,blitzt,blitzt,blitzen,blitzt,blitzen},
  preterite = {blitzte,blitztest,blitzte,blitzten,blitztet,blitzten},
  subjunctive = {},
  imperative = {blitz (du)!,blitzen wir!,blitzt (ihr)!,blitzen Sie!},
  perfect = {haben + \textbf{geblitzt}},
  future = {werden + \textbf{blitzen}},
  pluperfect = {hatten + \textbf{geblitzt}},
  prepositions = {---},
  separable = {---}
}

% Migrated from verbs/b.tex:31
\VerbTable{
  name = {blühen},
  english = {to bloom / flourish},
  type = {regular},
  auxiliary = {haben},
  style = {default},
  present = {blühe,blühst,blüht,blühen,blüht,blühen},
  preterite = {blühte,blühtest,blühte,blühten,blühtet,blühten},
  subjunctive = {},
  imperative = {blühe (du)!,blühen wir!,blüht (ihr)!,blühen Sie!},
  perfect = {haben + \textbf{geblüht}},
  future = {werden + \textbf{blühen}},
  pluperfect = {hatten + \textbf{geblüht}},
  prepositions = {---},
  separable = {& \textbf{auf$|$blühen} & (to blossom / flourish)\\}
}

% Migrated from verbs/b1/b.tex:31
\VerbTable{
  name = {bluten},
  english = {to bleed},
  type = {regular},
  auxiliary = {haben},
  style = {default},
  present = {blute,blutest,blutet,bluten,blutet,bluten},
  preterite = {blutete,blutetest,blutete,bluteten,blutetet,bluteten},
  subjunctive = {},
  imperative = {blute (du)!,bluten wir!,blutet (ihr)!,bluten Sie!},
  perfect = {haben + \textbf{geblutet}},
  future = {werden + \textbf{bluten}},
  pluperfect = {hatten + \textbf{geblutet}},
  prepositions = {---},
  separable = {---}
}

% Migrated from verbs/b.tex:32
\VerbTable{
  name = {bohren},
  english = {to drill},
  type = {regular},
  auxiliary = {haben},
  style = {default},
  present = {bohre,bohrst,bohrt,bohren,bohrt,bohren},
  preterite = {bohrte,bohrtest,bohrte,bohrten,bohrtet,bohrten},
  subjunctive = {},
  imperative = {bohre (du)!,bohren wir!,bohrt (ihr)!,bohren Sie!},
  perfect = {haben + \textbf{gebohrt}},
  future = {werden + \textbf{bohren}},
  pluperfect = {hatten + \textbf{gebohrt}},
  prepositions = {
    & \textbf{bohren in} + Dat & (to drill in / into)\\
    & \textbf{bohren durch} + Akk & (to drill through)\\
  },
  separable = {& \textbf{auf$|$bohren} & (to drill on / onto surface)\\}
}

% Migrated from verbs/b1/b.tex:32
\VerbTable{
  name = {braten},
  english = {to fry / roast},
  type = {strong / irregular},
  auxiliary = {haben},
  style = {default},
  present = {brate,brätst,brät,braten,bratet,braten},
  preterite = {briet,brietest,briet,brieten,brietet,brieten},
  subjunctive = {},
  imperative = {brate (du)!,braten wir!,bratet (ihr)!,braten Sie!},
  perfect = {haben + \textbf{gebraten}},
  future = {werden + \textbf{braten}},
  pluperfect = {hatten + \textbf{gebraten}},
  prepositions = {---},
  separable = {---}
}

% Migrated from verbs/b.tex:33
\VerbTable{
  name = {brauchen},
  english = {to need},
  type = {regular, + Akkusativ},
  auxiliary = {haben},
  style = {acc},
  present = {brauche,brauchst,braucht,brauchen,braucht,brauchen},
  preterite = {brauchte,brauchtest,brauchte,brauchten,brauchtet,brauchten},
  subjunctive = {},
  imperative = {brauch(e) (du)!,brauchen wir!,braucht (ihr)!,brauchen Sie!},
  perfect = {haben + \textbf{gebraucht}},
  future = {werden + \textbf{brauchen}},
  pluperfect = {hatten + \textbf{gebraucht}},
  prepositions = {---},
  separable = {---}
}

% Migrated from verbs/b1/b.tex:33
\VerbTable{
  name = {brechen},
  english = {to break},
  type = {strong / irregular},
  auxiliary = {haben},
  style = {default},
  present = {breche,brichst,bricht,brechen,brecht,brechen},
  preterite = {brach,brachst,brach,brachen,bracht,brachen},
  subjunctive = {},
  imperative = {brich (du)!,brechen wir!,brecht (ihr)!,brechen Sie!},
  perfect = {haben + \textbf{gebrochen}},
  future = {werden + \textbf{brechen}},
  pluperfect = {hatten + \textbf{gebrochen}},
  prepositions = {---},
  separable = {---}
}

% Migrated from verbs/b1/b.tex:34
\VerbTable{
  name = {bremsen},
  english = {to brake},
  type = {regular},
  auxiliary = {haben},
  style = {default},
  present = {bremse,bremst,bremst,bremsen,bremst,bremsen},
  preterite = {bremste,bremstest,bremste,bremsten,bremstet,bremsten},
  subjunctive = {},
  imperative = {brems (du)!,bremsen wir!,bremst (ihr)!,bremsen Sie!},
  perfect = {haben + \textbf{gebremst}},
  future = {werden + \textbf{bremsen}},
  pluperfect = {hatten + \textbf{gebremst}},
  prepositions = {---},
  separable = {---}
}

% Migrated from verbs/b1/b.tex:35
\VerbTable{
  name = {brennen},
  english = {to burn},
  type = {mixed},
  auxiliary = {haben},
  style = {default},
  present = {brenne,brennst,brennt,brennen,brennt,brennen},
  preterite = {brannte,branntest,brannte,brannten,branntet,brannten},
  subjunctive = {},
  imperative = {brenn (du)!,brennen wir!,brennt (ihr)!,brennen Sie!},
  perfect = {haben + \textbf{gebrannt}},
  future = {werden + \textbf{brennen}},
  pluperfect = {hatten + \textbf{gebrannt}},
  prepositions = {---},
  separable = {---}
}

% Migrated from verbs/b.tex:34
\VerbTable{
  name = {bringen},
  english = {to bring},
  type = {irregular, + Akkusativ},
  auxiliary = {haben},
  style = {acc},
  present = {bringe,bringst,bringt,bringen,bringt,bringen},
  preterite = {brachte,brachtest,brachte,brachten,brachtet,brachten},
  subjunctive = {},
  imperative = {bring(e) (du)!,bringen wir!,bringt (ihr)!,bringen Sie!},
  perfect = {haben + \textbf{gebracht}},
  future = {werden + \textbf{bringen}},
  pluperfect = {hatten + \textbf{gebracht}},
  prepositions = {
    & \textbf{bringen zu} + Dat & (to bring to someone / place)\\
    &\textbf{bringen auf} + Akk & (to raise / summon)\\
    &\textbf{bringen für} + Akk & (to bring for)\\
    &\textbf{bringen von} + Dat & (to bring from)\\
    &\textbf{bringen in} + Akk & (to move into)\\
  },
  separable = {
    & \textbf{um$|$bringen} & (to kill)\\
    & \textbf{mit$|$bringen} & (to bring along)\\
    & \textbf{nach$|$bringen} & (to bring later)\\
    & \textbf{hin$|$bringen} & (to bring there)\\
    & \textbf{her$|$bringen} & (to bring here)\\
    & \textbf{auf$|$bringen} & (to raise money / summon)\\
    & \textbf{ein$|$bringen} & (to yield / introduce)\\
    & \textbf{bei$|$bringen} & (to teach / instill)\\
    & \textbf{unter$|$bringen} & (to accomodate)\\
    & \textbf{weg$|$bringen} & (to take away)\\
    & \textbf{zurück$|$bringen} & (to bring back)\\
    & \textbf{vorbei$|$bringen} & (to bring by / drop off)\\
  }
}

% Migrated from verbs/b.tex:35
\VerbTable{
  name = {buchen},
  english = {to book / reserve},
  type = {regular, + Akkusativ},
  auxiliary = {haben},
  style = {acc},
  present = {buche,buchst,bucht,buchen,bucht,buchen},
  preterite = {buchte,buchtest,buchte,buchten,buchtet,buchten},
  subjunctive = {},
  imperative = {buche (du)!,buchen wir!,bucht (ihr)!,buchen Sie!},
  perfect = {haben + \textbf{gebucht}},
  future = {werden + \textbf{buchen}},
  pluperfect = {hatten + \textbf{gebucht}},
  prepositions = {
    & \textbf{buchen für} + Akk & (to book for)\\
    & \textbf{buchen bei} + Dat & (to book at / with)\\
  },
  separable = {---}
}

% Migrated from verbs/b1/b.tex:36
\VerbTable{
  name = {buchstabieren},
  english = {to spell},
  type = {regular},
  auxiliary = {haben},
  style = {default},
  present = {buchstabiere,buchstabierst,buchstabiert,buchstabieren,buchstabiert,buchstabieren},
  preterite = {buchstabierte,buchstabiertest,buchstabierte,buchstabierten,buchstabiertet,buchstabierten},
  subjunctive = {},
  imperative = {buchstabier (du)!,buchstabieren wir!,buchstabiert (ihr)!,buchstabieren Sie!},
  perfect = {haben + \textbf{buchstabiert}},
  future = {werden + \textbf{buchstabieren}},
  pluperfect = {hatten + \textbf{buchstabiert}},
  prepositions = {---},
  separable = {---}
}

\section{C}

% Migrated from verbs/c.tex:1
\VerbTable{
  name = {checken},
  english = {to check / verify},
  type = {regular, colloquial},
  auxiliary = {haben},
  style = {default},
  present = {checke,checkst,checkt,checken,checkt,checken},
  preterite = {checkte,checktest,checkte,checkten,checktet,checkten},
  subjunctive = {},
  imperative = {checke (du)!,checken wir!,checkt (ihr)!,checken Sie!},
  perfect = {haben + \textbf{gecheckt}},
  future = {werden + \textbf{checken}},
  pluperfect = {hatten + \textbf{gecheckt}},
  prepositions = {---},
  separable = {---}
}

\section{D}

% Migrated from verbs/b1/d.tex:1
\VerbTable{
  name = {danken},
  english = {to thank},
  type = {regular},
  auxiliary = {haben},
  style = {default},
  present = {danke,dankst,dankt,danken,dankt,danken},
  preterite = {dankte,danktest,dankte,dankten,danktet,dankten},
  subjunctive = {},
  imperative = {dank (du)!,danken wir!,dankt (ihr)!,danken Sie!},
  perfect = {haben + \textbf{gedankt}},
  future = {werden + \textbf{danken}},
  pluperfect = {hatten + \textbf{gedankt}},
  prepositions = {---},
  separable = {---}
}

% Migrated from verbs/b1/d.tex:2
\VerbTable{
  name = {darstellen},
  english = {to represent},
  type = {separable, regular},
  auxiliary = {haben},
  style = {default},
  present = {stelle dar,stellst dar,stellt dar,stellen dar,stellt dar,stellen dar},
  preterite = {stellte dar,stelltest dar,stellte dar,stellten dar,stelltet dar,stellten dar},
  subjunctive = {},
  imperative = {stell (du) dar!,stellen wir dar!,stellt (ihr) dar!,stellen Sie dar!},
  perfect = {haben + \textbf{dargestellt}},
  future = {werden + \textbf{darstellen}},
  pluperfect = {hatten + \textbf{dargestellt}},
  prepositions = {---},
  separable = {---}
}

% Migrated from verbs/d.tex:1
\VerbTable{
  name = {dauern},
  english = {to last / take time},
  type = {regular},
  auxiliary = {haben},
  style = {default},
  present = {dauere,dauerst,dauert,dauern,dauert,dauern},
  preterite = {dauerte,dauertest,dauerte,dauerten,dauertet,dauerten},
  subjunctive = {},
  imperative = {dauere (du)!,dauern wir!,dauert (ihr)!,dauern Sie!},
  perfect = {haben + \textbf{gedauert}},
  future = {werden + \textbf{dauern}},
  pluperfect = {hatten + \textbf{gedauert}},
  prepositions = {
    & \textbf{dauern bis} + Akk & (to last until)\\
    & \textbf{dauern über} + Akk & (to last over / longer than)\\
  },
  separable = {& \textbf{an$|$dauern} & (to continue / persist)\\}
}

% Migrated from verbs/b1/d.tex:3
\VerbTable{
  name = {dekorieren},
  english = {to decorate},
  type = {regular},
  auxiliary = {haben},
  style = {default},
  present = {dekoriere,dekorierst,dekoriert,dekorieren,dekoriert,dekorieren},
  preterite = {dekorierte,dekoriertest,dekorierte,dekorierten,dekoriertet,dekorierten},
  subjunctive = {},
  imperative = {dekorier (du)!,dekorieren wir!,dekoriert (ihr)!,dekorieren Sie!},
  perfect = {haben + \textbf{dekoriert}},
  future = {werden + \textbf{dekorieren}},
  pluperfect = {hatten + \textbf{dekoriert}},
  prepositions = {---},
  separable = {---}
}

% Migrated from verbs/d.tex:2
\VerbTable{
  name = {demonstrieren},
  english = {to demonstrate / protest},
  type = {regular},
  auxiliary = {haben},
  style = {default},
  present = {demonstriere,demonstrierst,demonstriert,demonstrieren,demonstriert,demonstrieren},
  preterite = {demonstrierte,demonstriertest,demonstrierte,demonstrierten,demonstriertet,demonstrierten},
  subjunctive = {},
  imperative = {demonstriere (du)!,demonstrieren wir!,demonstriert (ihr)!,demonstrieren Sie!},
  perfect = {haben + \textbf{demonstriert}},
  future = {werden + \textbf{demonstrieren}},
  pluperfect = {hatten + \textbf{demonstriert}},
  prepositions = {
    & \textbf{demonstrieren gegen} + Akk & (to demonstrate against)\\
    & \textbf{demonstrieren für} + Akk & (to demonstrate for)\\
  },
  separable = {---}
}

% Migrated from verbs/d.tex:3
\VerbTable{
  name = {denken},
  english = {to think},
  type = {mixed, + Akkusativ},
  auxiliary = {haben},
  style = {acc},
  present = {denke,denkst,denkt,denken,denkt,denken},
  preterite = {dachte,dachtest,dachte,dachten,dachtet,dachten},
  subjunctive = {},
  imperative = {denk(e) (du)!,denken wir!,denkt (ihr)!,denken Sie!},
  perfect = {haben + \textbf{gedacht}},
  future = {werden + \textbf{denken}},
  pluperfect = {hatten + \textbf{gedacht}},
  prepositions = {---},
  separable = {---}
}

% Migrated from verbs/d.tex:4
\VerbTable{
  name = {diskutieren},
  english = {to discuss},
  type = {irregular, + Akkusativ},
  auxiliary = {haben},
  style = {acc},
  present = {diskutiere,diskutierst,diskutiert,diskutieren,diskutiert,diskutieren},
  preterite = {diskutierte,diskutiertest,diskutierte,diskutierten,diskutiertet,diskutierten},
  subjunctive = {},
  imperative = {diskutier(e) (du)!,diskutieren wir!,diskutiert (ihr)!,diskutieren Sie!},
  perfect = {haben + \textbf{diskutiert}},
  future = {werden + \textbf{diskutieren}},
  pluperfect = {hatten + \textbf{diskutiert}},
  prepositions = {---},
  separable = {---}
}

% Migrated from verbs/b1/d.tex:4
\VerbTable{
  name = {donnern},
  english = {to thunder},
  type = {regular},
  auxiliary = {haben},
  style = {default},
  present = {donnre,donnerst,donnert,donnern,donnert,donnern},
  preterite = {donnerte,donnertest,donnerte,donnerten,donnertet,donnerten},
  subjunctive = {},
  imperative = {donnre (du)!,donnern wir!,donnert (ihr)!,donnern Sie!},
  perfect = {haben + \textbf{gedonnert}},
  future = {werden + \textbf{donnern}},
  pluperfect = {hatten + \textbf{gedonnert}},
  prepositions = {---},
  separable = {---}
}

% Migrated from verbs/d.tex:5
\VerbTable{
  name = {drehen},
  english = {to turn / rotate},
  type = {regular},
  auxiliary = {haben},
  style = {default},
  present = {drehe,drehst,dreht,drehen,dreht,drehen},
  preterite = {drehte,drehtest,drehte,drehten,drehtet,drehten},
  subjunctive = {},
  imperative = {dreh(e) (du)!,drehen wir!,dreht (ihr)!,drehen Sie!},
  perfect = {haben + \textbf{gedreht}},
  future = {werden + \textbf{drehen}},
  pluperfect = {hatten + \textbf{gedreht}},
  prepositions = {
    & \textbf{drehen an} + Dat & (to turn / adjust)\\
    & \textbf{drehen um} + Akk & (to turn around)\\
    & \textbf{drehen nach} + Dat & (to turn toward)\\
  },
  separable = {
    & \textbf{ab$|$drehen} & (to turn off)\\
    & \textbf{auf$|$drehen} & (to turn up)\\
    & \textbf{um$|$drehen} & (to turn around)\\
    & \textbf{zu$|$drehen} & (to close / screw shut)\\
  }
}

% Migrated from verbs/b1/d.tex:5
\VerbTable{
  name = {drucken},
  english = {to print},
  type = {regular},
  auxiliary = {haben},
  style = {default},
  present = {drucke,druckst,druckt,drucken,druckt,drucken},
  preterite = {druckte,drucktest,druckte,druckten,drucktet,druckten},
  subjunctive = {},
  imperative = {druck (du)!,drucken wir!,druckt (ihr)!,drucken Sie!},
  perfect = {haben + \textbf{gedruckt}},
  future = {werden + \textbf{drucken}},
  pluperfect = {hatten + \textbf{gedruckt}},
  prepositions = {---},
  separable = {---}
}

% Migrated from verbs/b1/d.tex:6
\VerbTable{
  name = {drücken},
  english = {to press},
  type = {regular},
  auxiliary = {haben},
  style = {default},
  present = {drücke,drückst,drückt,drücken,drückt,drücken},
  preterite = {drückte,drücktest,drückte,drückten,drücktet,drückten},
  subjunctive = {},
  imperative = {drück (du)!,drücken wir!,drückt (ihr)!,drücken Sie!},
  perfect = {haben + \textbf{gedrückt}},
  future = {werden + \textbf{drücken}},
  pluperfect = {hatten + \textbf{gedrückt}},
  prepositions = {---},
  separable = {---}
}

% Migrated from verbs/b1/d.tex:7
\VerbTable{
  name = {duschen},
  english = {to shower},
  type = {regular},
  auxiliary = {haben},
  style = {default},
  present = {dusche,duschst,duscht,duschen,duscht,duschen},
  preterite = {duschte,duschtest,duschte,duschten,duschtet,duschten},
  subjunctive = {},
  imperative = {dusch (du)!,duschen wir!,duscht (ihr)!,duschen Sie!},
  perfect = {haben + \textbf{geduscht}},
  future = {werden + \textbf{duschen}},
  pluperfect = {hatten + \textbf{geduscht}},
  prepositions = {---},
  separable = {---}
}

% Migrated from verbs/d.tex:6
\VerbTable{
  name = {duzen},
  english = {to address someone with \enquote{du}},
  type = {regular},
  auxiliary = {haben},
  style = {default},
  present = {duze,duzt,duzt,duzen,duzt,duzen},
  preterite = {duzte,duztest,duzte,duzten,duztet,duzten},
  subjunctive = {},
  imperative = {duze (du)!,duzen wir!,duzt (ihr)!,duzen Sie!},
  perfect = {haben + \textbf{geduzt}},
  future = {werden + \textbf{duzen}},
  pluperfect = {hatten + \textbf{geduzt}},
  prepositions = {& \textbf{duzen mit} + Dat & (to use "du" with someone)\\},
  separable = {---}
}

\section{E}

% Migrated from verbs/b1/e.tex:1
\VerbTable{
  name = {eilen},
  english = {to hurry},
  type = {regular},
  auxiliary = {haben / sein},
  style = {default},
  present = {eile,eilst,eilt,eilen,eilt,eilen},
  preterite = {eilte,eiltest,eilte,eilten,eiltet,eilten},
  subjunctive = {},
  imperative = {eil (du)!,eilen wir!,eilt (ihr)!,eilen Sie!},
  perfect = {haben / sein + \textbf{geeilt}},
  future = {werden + \textbf{eilen}},
  pluperfect = {hatten / waren + \textbf{geeilt}},
  prepositions = {---},
  separable = {---}
}

% Migrated from verbs/b1/e.tex:2
\VerbTable{
  name = {einbrechen},
  english = {to break in},
  type = {separable, strong / irregular},
  auxiliary = {sein},
  style = {default},
  present = {breche ein,brichst ein,bricht ein,brechen ein,brecht ein,brechen ein},
  preterite = {brach ein,brachst ein,brach ein,brachen ein,bracht ein,brachen ein},
  subjunctive = {},
  imperative = {brich (du) ein!,brechen wir ein!,brecht (ihr) ein!,brechen Sie ein!},
  perfect = {sein + \textbf{eingebrochen}},
  future = {werden + \textbf{einbrechen}},
  pluperfect = {waren + \textbf{eingebrochen}},
  prepositions = {---},
  separable = {---}
}

% Migrated from verbs/b1/e.tex:3
\VerbTable{
  name = {einfallen},
  english = {to occur to},
  type = {separable, strong / irregular},
  auxiliary = {sein},
  style = {default},
  present = {falle ein,fällst ein,fällt ein,fallen ein,fallt ein,fallen ein},
  preterite = {fiel ein,fielst ein,fiel ein,fielen ein,fielt ein,fielen ein},
  subjunctive = {},
  imperative = {fall (du) ein!,fallen wir ein!,fallt (ihr) ein!,fallen Sie ein!},
  perfect = {sein + \textbf{eingefallen}},
  future = {werden + \textbf{einfallen}},
  pluperfect = {waren + \textbf{eingefallen}},
  prepositions = {---},
  separable = {---}
}

% Migrated from verbs/b1/e.tex:4
\VerbTable{
  name = {einfügen},
  english = {to insert},
  type = {separable, regular},
  auxiliary = {haben},
  style = {default},
  present = {füge ein,fügst ein,fügt ein,fügen ein,fügt ein,fügen ein},
  preterite = {fügte ein,fügtest ein,fügte ein,fügten ein,fügtet ein,fügten ein},
  subjunctive = {},
  imperative = {füg (du) ein!,fügen wir ein!,fügt (ihr) ein!,fügen Sie ein!},
  perfect = {haben + \textbf{eingefügt}},
  future = {werden + \textbf{einfügen}},
  pluperfect = {hatten + \textbf{eingefügt}},
  prepositions = {---},
  separable = {---}
}

% Migrated from verbs/b1/e.tex:5
\VerbTable{
  name = {einführen},
  english = {to introduce},
  type = {separable, regular},
  auxiliary = {haben},
  style = {default},
  present = {führe ein,führst ein,führt ein,führen ein,führt ein,führen ein},
  preterite = {führte ein,führtest ein,führte ein,führten ein,führtet ein,führten ein},
  subjunctive = {},
  imperative = {führ (du) ein!,führen wir ein!,führt (ihr) ein!,führen Sie ein!},
  perfect = {haben + \textbf{eingeführt}},
  future = {werden + \textbf{einführen}},
  pluperfect = {hatten + \textbf{eingeführt}},
  prepositions = {---},
  separable = {---}
}

% Migrated from verbs/b1/e.tex:7
\VerbTable{
  name = {einkaufen},
  english = {to shop},
  type = {separable, regular},
  auxiliary = {haben},
  style = {default},
  present = {kaufe ein,kaufst ein,kauft ein,kaufen ein,kauft ein,kaufen ein},
  preterite = {kaufte ein,kauftest ein,kaufte ein,kauften ein,kauftet ein,kauften ein},
  subjunctive = {},
  imperative = {kauf (du) ein!,kaufen wir ein!,kauft (ihr) ein!,kaufen Sie ein!},
  perfect = {haben + \textbf{eingekauft}},
  future = {werden + \textbf{einkaufen}},
  pluperfect = {hatten + \textbf{eingekauft}},
  prepositions = {---},
  separable = {---}
}

% Migrated from verbs/b1/e.tex:8
\VerbTable{
  name = {einladen},
  english = {to invite},
  type = {separable, strong / irregular},
  auxiliary = {haben},
  style = {default},
  present = {lade ein,lädst ein,lädt ein,laden ein,ladet ein,laden ein},
  preterite = {lud ein,ludest ein,lud ein,luden ein,ludet ein,luden ein},
  subjunctive = {},
  imperative = {lad (du) ein!,laden wir ein!,ladet (ihr) ein!,laden Sie ein!},
  perfect = {haben + \textbf{eingeladen}},
  future = {werden + \textbf{einladen}},
  pluperfect = {hatten + \textbf{eingeladen}},
  prepositions = {---},
  separable = {---}
}

% Migrated from verbs/b1/e.tex:9
\VerbTable{
  name = {einnehmen},
  english = {to take / collect},
  type = {separable, strong / irregular},
  auxiliary = {haben},
  style = {default},
  present = {nehme ein,nimmst ein,nimmt ein,nehmen ein,nehmt ein,nehmen ein},
  preterite = {nahm ein,nahmst ein,nahm ein,nahmen ein,nahmt ein,nahmen ein},
  subjunctive = {},
  imperative = {nimm (du) ein!,nehmen wir ein!,nehmt (ihr) ein!,nehmen Sie ein!},
  perfect = {haben + \textbf{eingenommen}},
  future = {werden + \textbf{einnehmen}},
  pluperfect = {hatten + \textbf{eingenommen}},
  prepositions = {---},
  separable = {---}
}

% Migrated from verbs/b1/e.tex:10
\VerbTable{
  name = {einpacken},
  english = {to pack / wrap},
  type = {separable, regular},
  auxiliary = {haben},
  style = {default},
  present = {packe ein,packst ein,packt ein,packen ein,packt ein,packen ein},
  preterite = {packte ein,packtest ein,packte ein,packten ein,packtet ein,packten ein},
  subjunctive = {},
  imperative = {pack (du) ein!,packen wir ein!,packt (ihr) ein!,packen Sie ein!},
  perfect = {haben + \textbf{eingepackt}},
  future = {werden + \textbf{einpacken}},
  pluperfect = {hatten + \textbf{eingepackt}},
  prepositions = {---},
  separable = {---}
}

% Migrated from verbs/b1/e.tex:11
\VerbTable{
  name = {einrichten},
  english = {to furnish / set up},
  type = {separable, regular},
  auxiliary = {haben},
  style = {default},
  present = {richte ein,richtest ein,richtet ein,richten ein,richtet ein,richten ein},
  preterite = {richtete ein,richtetest ein,richtete ein,richteten ein,richtetet ein,richteten ein},
  subjunctive = {},
  imperative = {richte (du) ein!,richten wir ein!,richtet (ihr) ein!,richten Sie ein!},
  perfect = {haben + \textbf{eingerichtet}},
  future = {werden + \textbf{einrichten}},
  pluperfect = {hatten + \textbf{eingerichtet}},
  prepositions = {---},
  separable = {---}
}

% Migrated from verbs/b1/e.tex:12
\VerbTable{
  name = {einschalten},
  english = {to switch on},
  type = {separable, regular},
  auxiliary = {haben},
  style = {default},
  present = {schalte ein,schaltest ein,schaltet ein,schalten ein,schaltet ein,schalten ein},
  preterite = {schaltete ein,schaltetest ein,schaltete ein,schalteten ein,schaltetet ein,schalteten ein},
  subjunctive = {},
  imperative = {schalte (du) ein!,schalten wir ein!,schaltet (ihr) ein!,schalten Sie ein!},
  perfect = {haben + \textbf{eingeschaltet}},
  future = {werden + \textbf{einschalten}},
  pluperfect = {hatten + \textbf{eingeschaltet}},
  prepositions = {---},
  separable = {---}
}

% Migrated from verbs/b1/e.tex:13
\VerbTable{
  name = {einsetzen},
  english = {to insert / deploy},
  type = {separable, regular},
  auxiliary = {haben},
  style = {default},
  present = {setze ein,setzt ein,setzt ein,setzen ein,setzt ein,setzen ein},
  preterite = {setzte ein,setztest ein,setzte ein,setzten ein,setztet ein,setzten ein},
  subjunctive = {},
  imperative = {setz (du) ein!,setzen wir ein!,setzt (ihr) ein!,setzen Sie ein!},
  perfect = {haben + \textbf{eingesetzt}},
  future = {werden + \textbf{einsetzen}},
  pluperfect = {hatten + \textbf{eingesetzt}},
  prepositions = {---},
  separable = {---}
}

% Migrated from verbs/b1/e.tex:14
\VerbTable{
  name = {einsteigen},
  english = {to get in},
  type = {separable, strong / irregular},
  auxiliary = {sein},
  style = {default},
  present = {steige ein,steigst ein,steigt ein,steigen ein,steigt ein,steigen ein},
  preterite = {stieg ein,stiegst ein,stieg ein,stiegen ein,stiegt ein,stiegen ein},
  subjunctive = {},
  imperative = {steig (du) ein!,steigen wir ein!,steigt (ihr) ein!,steigen Sie ein!},
  perfect = {sein + \textbf{eingestiegen}},
  future = {werden + \textbf{einsteigen}},
  pluperfect = {waren + \textbf{eingestiegen}},
  prepositions = {---},
  separable = {---}
}

% Migrated from verbs/b1/e.tex:15
\VerbTable{
  name = {einstellen},
  english = {to hire / set / discontinue},
  type = {separable, regular},
  auxiliary = {haben},
  style = {default},
  present = {stelle ein,stellst ein,stellt ein,stellen ein,stellt ein,stellen ein},
  preterite = {stellte ein,stelltest ein,stellte ein,stellten ein,stelltet ein,stellten ein},
  subjunctive = {},
  imperative = {stell (du) ein!,stellen wir ein!,stellt (ihr) ein!,stellen Sie ein!},
  perfect = {haben + \textbf{eingestellt}},
  future = {werden + \textbf{einstellen}},
  pluperfect = {hatten + \textbf{eingestellt}},
  prepositions = {---},
  separable = {---}
}

% Migrated from verbs/b1/e.tex:16
\VerbTable{
  name = {eintragen},
  english = {to enter},
  type = {separable, strong / irregular},
  auxiliary = {haben},
  style = {default},
  present = {trage ein,trägst ein,trägt ein,tragen ein,tragt ein,tragen ein},
  preterite = {trug ein,trugst ein,trug ein,trugen ein,trugt ein,trugen ein},
  subjunctive = {},
  imperative = {trag (du) ein!,tragen wir ein!,tragt (ihr) ein!,tragen Sie ein!},
  perfect = {haben + \textbf{eingetragen}},
  future = {werden + \textbf{eintragen}},
  pluperfect = {hatten + \textbf{eingetragen}},
  prepositions = {---},
  separable = {---}
}

% Migrated from verbs/b1/e.tex:17
\VerbTable{
  name = {eintreten},
  english = {to enter},
  type = {separable, strong / irregular},
  auxiliary = {sein},
  style = {default},
  present = {trete ein,trittst ein,tritt ein,treten ein,tretet ein,treten ein},
  preterite = {trat ein,tratest ein,trat ein,traten ein,tratet ein,traten ein},
  subjunctive = {},
  imperative = {tritt (du) ein!,treten wir ein!,tretet (ihr) ein!,treten Sie ein!},
  perfect = {sein + \textbf{eingetreten}},
  future = {werden + \textbf{eintreten}},
  pluperfect = {waren + \textbf{eingetreten}},
  prepositions = {---},
  separable = {---}
}

% Migrated from verbs/b1/e.tex:18
\VerbTable{
  name = {einzahlen},
  english = {to deposit},
  type = {separable, regular},
  auxiliary = {haben},
  style = {default},
  present = {zahle ein,zahlst ein,zahlt ein,zahlen ein,zahlt ein,zahlen ein},
  preterite = {zahlte ein,zahltest ein,zahlte ein,zahlten ein,zahltet ein,zahlten ein},
  subjunctive = {},
  imperative = {zahl (du) ein!,zahlen wir ein!,zahlt (ihr) ein!,zahlen Sie ein!},
  perfect = {haben + \textbf{eingezahlt}},
  future = {werden + \textbf{einzahlen}},
  pluperfect = {hatten + \textbf{eingezahlt}},
  prepositions = {---},
  separable = {---}
}

% Migrated from verbs/b1/e.tex:19
\VerbTable{
  name = {einziehen},
  english = {to move in},
  type = {separable, strong / irregular},
  auxiliary = {sein},
  style = {default},
  present = {ziehe ein,ziehst ein,zieht ein,ziehen ein,zieht ein,ziehen ein},
  preterite = {zog ein,zogst ein,zog ein,zogen ein,zogt ein,zogen ein},
  subjunctive = {},
  imperative = {zieh (du) ein!,ziehen wir ein!,zieht (ihr) ein!,ziehen Sie ein!},
  perfect = {sein + \textbf{eingezogen}},
  future = {werden + \textbf{einziehen}},
  pluperfect = {waren + \textbf{eingezogen}},
  prepositions = {---},
  separable = {---}
}

% Migrated from verbs/b1/e.tex:20
\VerbTable{
  name = {empfangen},
  english = {to receive},
  type = {strong / irregular},
  auxiliary = {haben},
  style = {default},
  present = {empfange,empfängst,empfängt,empfangen,empfangt,empfangen},
  preterite = {empfing,empfingst,empfing,empfingen,empfingt,empfingen},
  subjunctive = {},
  imperative = {empfang (du)!,empfangen wir!,empfangt (ihr)!,empfangen Sie!},
  perfect = {haben + \textbf{empfangen}},
  future = {werden + \textbf{empfangen}},
  pluperfect = {hatten + \textbf{empfangen}},
  prepositions = {---},
  separable = {---}
}

% Migrated from verbs/e.tex:1
\VerbTable{
  name = {empfehlen},
  english = {to recommend},
  type = {irregular},
  auxiliary = {haben},
  style = {default},
  present = {empfehle,empfiehlst,empfiehlt,empfehlen,empfehlt,empfehlen},
  preterite = {empfahl,empfahlst,empfahl,empfahlen,empfahlt,empfahlen},
  subjunctive = {},
  imperative = {empfiehl (du)!,empfehlen wir!,empfehlt (ihr)!,empfehlen Sie!},
  perfect = {haben + \textbf{empfohlen}},
  future = {werden + \textbf{empfehlen}},
  pluperfect = {hatten + \textbf{empfohlen}},
  prepositions = {
    & \textbf{empfehlen für} + Akk & (to recommend for)\\
    & \textbf{empfehlen an} + Akk & (to recommend to)\\
  },
  separable = {---}
}

% Migrated from verbs/b1/e.tex:21
\VerbTable{
  name = {enden},
  english = {to end},
  type = {regular},
  auxiliary = {haben},
  style = {default},
  present = {ende,endest,endet,enden,endet,enden},
  preterite = {endete,endetest,endete,endeten,endetet,endeten},
  subjunctive = {},
  imperative = {ende (du)!,enden wir!,endet (ihr)!,enden Sie!},
  perfect = {haben + \textbf{geendet}},
  future = {werden + \textbf{enden}},
  pluperfect = {hatten + \textbf{geendet}},
  prepositions = {---},
  separable = {---}
}

% Migrated from verbs/e.tex:2
\VerbTable{
  name = {engagieren},
  english = {to hire / engage},
  type = {regular, + Akkusativ},
  auxiliary = {haben},
  style = {acc},
  present = {engagiere,engagierst,engagiert,engagieren,engagiert,engagieren},
  preterite = {engagierte,engagiertest,engagierte,engagierten,engagiertet,engagierten},
  subjunctive = {},
  imperative = {engagiere (du)!,engagieren wir!,engagiert (ihr)!,engagieren Sie!},
  perfect = {haben + \textbf{engagiert}},
  future = {werden + \textbf{engagieren}},
  pluperfect = {hatten + \textbf{engagiert}},
  prepositions = {---},
  separable = {---}
}

% Migrated from verbs/e.tex:3
\VerbTable{
  name = {entdecken},
  english = {to discover},
  type = {regular, + Akkusativ},
  auxiliary = {haben},
  style = {acc},
  present = {entdecke,entdeckst,entdeckt,entdecken,entdeckt,entdecken},
  preterite = {entdeckte,entdecktest,entdeckte,entdeckten,entdecktet,entdeckten},
  subjunctive = {},
  imperative = {entdecke (du)!,entdecken wir!,entdeckt (ihr)!,entdecken Sie!},
  perfect = {haben + \textbf{entdeckt}},
  future = {werden + \textbf{entdecken}},
  pluperfect = {hatten + \textbf{entdeckt}},
  prepositions = {---},
  separable = {---}
}

% Migrated from verbs/e.tex:4
\VerbTable{
  name = {entfernen},
  english = {to remove},
  type = {regular, + Akkusativ},
  auxiliary = {haben},
  style = {acc},
  present = {entferne,entfernst,entfernt,entfernen,entfernt,entfernen},
  preterite = {entfernte,entferntest,entfernte,entfernten,entferntet,entfernten},
  subjunctive = {},
  imperative = {entferne (du)!,entfernen wir!,entfernt (ihr)!,entfernen Sie!},
  perfect = {haben + \textbf{entfernt}},
  future = {werden + \textbf{entfernen}},
  pluperfect = {hatten + \textbf{entfernt}},
  prepositions = {& \textbf{entfernen aus} + Dat & (to remove from)\\},
  separable = {---}
}

% Migrated from verbs/b1/e.tex:22
\VerbTable{
  name = {enthalten},
  english = {to contain},
  type = {strong / irregular},
  auxiliary = {haben},
  style = {default},
  present = {enthalte,enthältst,enthält,enthalten,enthaltet,enthalten},
  preterite = {enthielt,enthieltest,enthielt,enthielten,enthieltet,enthielten},
  subjunctive = {},
  imperative = {enthalte (du)!,enthalten wir!,enthaltet (ihr)!,enthalten Sie!},
  perfect = {haben + \textbf{enthalten}},
  future = {werden + \textbf{enthalten}},
  pluperfect = {hatten + \textbf{enthalten}},
  prepositions = {---},
  separable = {---}
}

% Migrated from verbs/e.tex:5
\VerbTable{
  name = {entlassen},
  english = {to dismiss / to release},
  type = {irregular},
  auxiliary = {haben},
  style = {default},
  present = {entlasse,entlässt,entlässt,entlassen,entlasst,entlassen},
  preterite = {entließ,entließt,entließ,entließen,entließt,entließen},
  subjunctive = {},
  imperative = {entlass(e) (du)!,entlassen wir!,entlasst (ihr)!,entlassen Sie!},
  perfect = {haben + \textbf{entlassen}},
  future = {werden + \textbf{entlassen}},
  pluperfect = {hatten + \textbf{entlassen}},
  prepositions = {
    & \textbf{entlassen aus} + Dat & (to release from / dismiss from)\\
    & \textbf{entlassen wegen} + Gen & (to dismiss because of)\\
  },
  separable = {---}
}

% Migrated from verbs/e.tex:6
\VerbTable{
  name = {entscheiden},
  english = {to decide},
  type = {irregular, + Akkusativ},
  auxiliary = {haben},
  style = {acc},
  present = {entscheide,entscheidest,entscheidet,entscheiden,entscheidet,entscheiden},
  preterite = {entschied,entschiedst,entschied,entschieden,entschiedet,entschieden},
  subjunctive = {},
  imperative = {entscheid(e) (du)!,entscheiden wir!,entscheidet (ihr)!,entscheiden Sie!},
  perfect = {haben + \textbf{entschieden}},
  future = {werden + \textbf{entscheiden}},
  pluperfect = {hatten + \textbf{entschieden}},
  prepositions = {---},
  separable = {---}
}

% Migrated from verbs/e.tex:7
\VerbTable{
  name = {entsorgen},
  english = {to dispose of},
  type = {regular},
  auxiliary = {haben},
  style = {default},
  present = {entsorge,entsorgst,entsorgt,entsorgen,entsorgt,entsorgen},
  preterite = {entsorgte,entsorgtest,entsorgte,entsorgten,entsorgtet,entsorgten},
  subjunctive = {},
  imperative = {entsorg(e) (du)!,entsorgen wir!,entsorgt (ihr)!,entsorgen Sie!},
  perfect = {haben + \textbf{entsorgt}},
  future = {werde + \textbf{entsorgen}},
  pluperfect = {hatten + \textbf{entsorgt}},
  prepositions = {
    & \textbf{entsorgen über} + Akk & (to dispose of via)\\
    & \textbf{entsorgen in} + Dat & (to dispose of in)\\
  },
  separable = {---}
}

% Migrated from verbs/e.tex:8
\VerbTable{
  name = {entstehen},
  english = {to arise},
  type = {irregular},
  auxiliary = {sein},
  style = {default},
  present = {entstehe,entstehst,entsteht,entstehen,entsteht,entstehen},
  preterite = {entstand,entstandest,entstand,entstanden,entstandet,entstanden},
  subjunctive = {},
  imperative = {entsteh(e) (du)!,entstehen wir!,entsteht (ihr)!,entstehen Sie!},
  perfect = {sein + \textbf{entstanden}},
  future = {werden + \textbf{entstehen}},
  pluperfect = {waren + \textbf{entstanden}},
  prepositions = {---},
  separable = {---}
}

% Migrated from verbs/b1/e.tex:24
\VerbTable{
  name = {enttäuschen},
  english = {to disappoint},
  type = {regular},
  auxiliary = {haben},
  style = {default},
  present = {enttäusche,enttäuschst,enttäuscht,enttäuschen,enttäuscht,enttäuschen},
  preterite = {enttäuschte,enttäuschtest,enttäuschte,enttäuschten,enttäuschtet,enttäuschten},
  subjunctive = {},
  imperative = {enttäusch (du)!,enttäuschen wir!,enttäuscht (ihr)!,enttäuschen Sie!},
  perfect = {haben + \textbf{enttäuscht}},
  future = {werden + \textbf{enttäuschen}},
  pluperfect = {hatten + \textbf{enttäuscht}},
  prepositions = {---},
  separable = {---}
}

% Migrated from verbs/e.tex:9
\VerbTable{
  name = {entzünden},
  english = {to ignite / inflame},
  type = {regular},
  auxiliary = {haben},
  style = {default},
  present = {entzünde,entzündest,entzündet,entzünden,entzündet,entzünden},
  preterite = {entzündete,entzündetest,entzündete,entzündeten,entzündetet,entzündeten},
  subjunctive = {},
  imperative = {entzünde (du)!,entzünden wir!,entzündet (ihr)!,entzünden Sie!},
  perfect = {haben + \textbf{entzündet}},
  future = {werden + \textbf{entzünden}},
  pluperfect = {hatten + \textbf{entzündet}},
  prepositions = {& \textbf{entzünden an} + Dat & (to ignite on / by)\\},
  separable = {---}
}

% Migrated from verbs/e.tex:10
\VerbTable{
  name = {erfahren},
  english = {to experience},
  type = {irregular, + Akkusativ},
  auxiliary = {haben},
  style = {acc},
  present = {erfahre,erf\"ahrst,erf\"ahrt,erfahren,erfahrt,erfahren},
  preterite = {erfuhr,erfuhrst,erfuhr,erfuhren,erfuhrt,erfuhren},
  subjunctive = {},
  imperative = {erfahr(e) (du)!,erfahren wir!,erfahrt (ihr)!,erfahren Sie!},
  perfect = {haben + \textbf{erfahren}},
  future = {werden + \textbf{erfahren}},
  pluperfect = {hatten + \textbf{erfahren}},
  prepositions = {---},
  separable = {---}
}

% Migrated from verbs/b1/e.tex:26
\VerbTable{
  name = {erfinden},
  english = {to invent},
  type = {strong / irregular},
  auxiliary = {haben},
  style = {default},
  present = {erfinde,erfindest,erfindet,erfinden,erfindet,erfinden},
  preterite = {erfand,erfandest,erfand,erfanden,erfandet,erfanden},
  subjunctive = {},
  imperative = {erfinde (du)!,erfinden wir!,erfindet (ihr)!,erfinden Sie!},
  perfect = {haben + \textbf{erfunden}},
  future = {werden + \textbf{erfinden}},
  pluperfect = {hatten + \textbf{erfunden}},
  prepositions = {---},
  separable = {---}
}

% Migrated from verbs/b1/e.tex:27
\VerbTable{
  name = {erfordern},
  english = {to require},
  type = {regular},
  auxiliary = {haben},
  style = {default},
  present = {erfordre,erforderst,erfordert,erfordern,erfordert,erfordern},
  preterite = {erforderte,erfordertest,erforderte,erforderten,erfordertet,erforderten},
  subjunctive = {},
  imperative = {erfordre (du)!,erfordern wir!,erfordert (ihr)!,erfordern Sie!},
  perfect = {haben + \textbf{erfordert}},
  future = {werden + \textbf{erfordern}},
  pluperfect = {hatten + \textbf{erfordert}},
  prepositions = {---},
  separable = {---}
}

% Migrated from verbs/b1/e.tex:28
\VerbTable{
  name = {erfüllen},
  english = {to fulfill},
  type = {regular},
  auxiliary = {haben},
  style = {default},
  present = {erfülle,erfüllst,erfüllt,erfüllen,erfüllt,erfüllen},
  preterite = {erfüllte,erfülltest,erfüllte,erfüllten,erfülltet,erfüllten},
  subjunctive = {},
  imperative = {erfüll (du)!,erfüllen wir!,erfüllt (ihr)!,erfüllen Sie!},
  perfect = {haben + \textbf{erfüllt}},
  future = {werden + \textbf{erfüllen}},
  pluperfect = {hatten + \textbf{erfüllt}},
  prepositions = {---},
  separable = {---}
}

% Migrated from verbs/e.tex:11
\VerbTable{
  name = {ergänzen},
  english = {to complete / supplement},
  type = {regular},
  auxiliary = {haben},
  style = {default},
  present = {ergänze,ergänzt,ergänzt,ergänzen,ergänzt,ergänzen},
  preterite = {ergänzte,ergänztest,ergänzte,ergänzten,ergänztet,ergänzten},
  subjunctive = {},
  imperative = {ergänze (du)!,ergänzen wir!,ergänzt (ihr)!,ergänzen Sie!},
  perfect = {haben + \textbf{ergänzt}},
  future = {werden + \textbf{ergänzen}},
  pluperfect = {hatten + \textbf{ergänzt}},
  prepositions = {& \textbf{ergänzen} + Akk & (to complete something)\\},
  separable = {---}
}

% Migrated from verbs/e.tex:12
\VerbTable{
  name = {erhalten},
  english = {to receive / to preserve},
  type = {irregular},
  auxiliary = {haben},
  style = {default},
  present = {erhalte,erhältst,erhält,erhalten,erhaltet,erhalten},
  preterite = {erhielt,erhieltest,erhielt,erhielten,erhieltet,erhielten},
  subjunctive = {},
  imperative = {erhalte (du)!,erhalten wir!,erhaltet (ihr)!,erhalten Sie!},
  perfect = {haben + \textbf{erhalten}},
  future = {werde + \textbf{erhalten}},
  pluperfect = {hatten + \textbf{erhalten}},
  prepositions = {
    & \textbf{erhalten von} + Dat & (to receive from)\\
    & \textbf{erhalten für} + Akk & (to receive for)\\
  },
  separable = {---}
}

% Migrated from verbs/b1/e.tex:29
\VerbTable{
  name = {erhöhen},
  english = {to increase},
  type = {regular},
  auxiliary = {haben},
  style = {default},
  present = {erhöhe,erhöhst,erhöht,erhöhen,erhöht,erhöhen},
  preterite = {erhöhte,erhöhtest,erhöhte,erhöhten,erhöhtet,erhöhten},
  subjunctive = {},
  imperative = {erhöh (du)!,erhöhen wir!,erhöht (ihr)!,erhöhen Sie!},
  perfect = {haben + \textbf{erhöht}},
  future = {werden + \textbf{erhöhen}},
  pluperfect = {hatten + \textbf{erhöht}},
  prepositions = {---},
  separable = {---}
}

% Migrated from verbs/e.tex:13
\VerbTable{
  name = {erinnern},
  english = {to remember},
  type = {irregular, + Akkusativ},
  auxiliary = {haben},
  style = {acc},
  present = {erinnere,erinnerst,erinnert,erinnern,erinnert,erinnern},
  preterite = {erinnerte,erinnertest,erinnerte,erinnerten,erinnertet,erinnerten},
  subjunctive = {},
  imperative = {erinner(e) (du)!,erinnern wir!,erinnert (ihr)!,erinnern Sie!},
  perfect = {haben + \textbf{erinnert}},
  future = {werden + \textbf{erinnern}},
  pluperfect = {hatten + \textbf{erinnert}},
  prepositions = {---},
  separable = {---}
}

% Migrated from verbs/b1/e.tex:31
\VerbTable{
  name = {erkennen},
  english = {to recognize},
  type = {regular},
  auxiliary = {haben},
  style = {default},
  present = {erkenne,erkennst,erkennt,erkennen,erkennt,erkennen},
  preterite = {erkannte,erkanntest,erkannte,erkannten,erkanntet,erkannten},
  subjunctive = {},
  imperative = {erkenn (du)!,erkennen wir!,erkennt (ihr)!,erkennen Sie!},
  perfect = {haben + \textbf{erkannt}},
  future = {werden + \textbf{erkennen}},
  pluperfect = {hatten + \textbf{erkannt}},
  prepositions = {---},
  separable = {---}
}

% Migrated from verbs/e.tex:14
\VerbTable{
  name = {erklären},
  english = {to explain},
  type = {irregular, + Dativ},
  auxiliary = {haben},
  style = {dat},
  present = {erkläre,erklärst,erklärt,erklären,erklärt,erklären},
  preterite = {erklärte,erklärtest,erklärte,erklärten,erklärtet,erklärten},
  subjunctive = {},
  imperative = {erklär(e) (du)!,erklären wir!,erklärt (ihr)!,erklären Sie!},
  perfect = {haben + \textbf{erklärt}},
  future = {werden + \textbf{erklären}},
  pluperfect = {hatten + \textbf{erklärt}},
  prepositions = {---},
  separable = {---}
}

% Migrated from verbs/e.tex:15
\VerbTable{
  name = {erkunden},
  english = {to explore},
  type = {regular, + Akkusativ},
  auxiliary = {haben},
  style = {acc},
  present = {erkunde,erkundest,erkundet,erkunden,erkundet,erkunden},
  preterite = {erkundete,erkundetest,erkundete,erkundeten,erkundetet,erkundeten},
  subjunctive = {},
  imperative = {erkunde (du)!,erkunden wir!,erkundet (ihr)!,erkunden Sie!},
  perfect = {haben + \textbf{erkundet}},
  future = {werden + \textbf{erkunden}},
  pluperfect = {hatten + \textbf{erkundet}},
  prepositions = {---},
  separable = {---}
}

% Migrated from verbs/e.tex:16
\VerbTable{
  name = {erlauben},
  english = {to allow},
  type = {regular, + Akkusativ},
  auxiliary = {haben},
  style = {acc},
  present = {erlaube,erlaubst,erlaubt,erlauben,erlaubt,erlauben},
  preterite = {erlaubte,erlaubtest,erlaubte,erlaubten,erlaubtet,erlaubten},
  subjunctive = {},
  imperative = {erlaub(e) (du)!,erlauben wir!,erlaubt (ihr)!,erlauben Sie!},
  perfect = {haben + \textbf{erlaubt}},
  future = {werden + \textbf{erlauben}},
  pluperfect = {hatten + \textbf{erlaubt}},
  prepositions = {---},
  separable = {---}
}

% Migrated from verbs/e.tex:17
\VerbTable{
  name = {erleben},
  english = {to experience},
  type = {regular, + Akkusativ},
  auxiliary = {haben},
  style = {acc},
  present = {erlebe,erlebst,erlebt,erleben,erlebt,erleben},
  preterite = {erlebte,erlebtest,erlebte,erlebten,erlebtet,erlebten},
  subjunctive = {},
  imperative = {erleb(e) (du)!,erleben wir!,erlebt (ihr)!,erleben Sie!},
  perfect = {haben + \textbf{erlebt}},
  future = {werden + \textbf{erleben}},
  pluperfect = {hatten + \textbf{erlebt}},
  prepositions = {---},
  separable = {---}
}

% Migrated from verbs/e.tex:18
\VerbTable{
  name = {erledigen},
  english = {to complete / take care of},
  type = {regular},
  auxiliary = {haben},
  style = {default},
  present = {erledige,erledigst,erledigt,erledigen,erledigt,erledigen},
  preterite = {erledigte,erledigtest,erledigte,erledigten,erledigtet,erledigten},
  subjunctive = {},
  imperative = {erledige (du)!,erledigen wir!,erledigt (ihr)!,erledigen Sie!},
  perfect = {haben + \textbf{erledigt}},
  future = {werden + \textbf{erledigen}},
  pluperfect = {hatten + \textbf{erledigt}},
  prepositions = {& \textbf{erledigen} + Akk & (to complete something)\\},
  separable = {---}
}

% Migrated from verbs/b1/e.tex:32
\VerbTable{
  name = {erleichtern},
  english = {to facilitate},
  type = {regular},
  auxiliary = {haben},
  style = {default},
  present = {erleichtre,erleichterst,erleichtert,erleichtern,erleichtert,erleichtern},
  preterite = {erleichterte,erleichtertest,erleichterte,erleichterten,erleichtertet,erleichterten},
  subjunctive = {},
  imperative = {erleichtre (du)!,erleichtern wir!,erleichtert (ihr)!,erleichtern Sie!},
  perfect = {haben + \textbf{erleichtert}},
  future = {werden + \textbf{erleichtern}},
  pluperfect = {hatten + \textbf{erleichtert}},
  prepositions = {---},
  separable = {---}
}

% Migrated from verbs/e.tex:19
\VerbTable{
  name = {erlernen},
  english = {to learn / acquire a skill},
  type = {regular, + Akkusativ},
  auxiliary = {haben},
  style = {acc},
  present = {erlerne,erlernst,erlernt,erlernen,erlernt,erlernen},
  preterite = {erlernte,erlerntest,erlernte,erlernten,erlerntet,erlernten},
  subjunctive = {},
  imperative = {erlern(e) (du)!,erlernen wir!,erlernt (ihr)!,erlernen Sie!},
  perfect = {haben + \textbf{erlernt}},
  future = {werden + \textbf{erlernen}},
  pluperfect = {hatten + \textbf{erlernt}},
  prepositions = {---},
  separable = {---}
}

% Migrated from verbs/e.tex:22
\VerbTable{
  name = {erneuern},
  english = {to renew / renovate},
  type = {regular},
  auxiliary = {haben},
  style = {default},
  present = {erneuere,erneuerst,erneuert,erneuern,erneuert,erneuern},
  preterite = {erneuerte,erneuertest,erneuerte,erneuerten,erneuertet,erneuerten},
  subjunctive = {},
  imperative = {erneuere (du)!,erneuern wir!,erneuert (ihr)!,erneuern Sie!},
  perfect = {haben + \textbf{erneuert}},
  future = {werden + \textbf{erneuern}},
  pluperfect = {hatten + \textbf{erneuert}},
  prepositions = {
    & \textbf{erneuern von} + Dat & (to renew from / of)\\
    & \textbf{erneuern durch} + Akk & (to renew through / by)\\
    & \textbf{erneuern zu} + Dat & (to renew into)\\
  },
  separable = {& \textbf{ein$|$erneuern} & (to renew in / reinstall)\\}
}

% Migrated from verbs/e.tex:20
\VerbTable{
  name = {eröffnen},
  english = {to open / inaugurate},
  type = {regular},
  auxiliary = {haben},
  style = {default},
  present = {eröffne,eröffnest,eröffnet,eröffnen,eröffnet,eröffnen},
  preterite = {eröffnete,eröffnetest,eröffnete,eröffneten,eröffnetet,eröffneten},
  subjunctive = {},
  imperative = {eröffne (du)!,eröffnen wir!,eröffnet (ihr)!,eröffnen Sie!},
  perfect = {haben + \textbf{eröffnet}},
  future = {werden + \textbf{eröffnen}},
  pluperfect = {hatten + \textbf{eröffnet}},
  prepositions = {
    & \textbf{eröffnen für} + Akk & (to open for)\\
    & \textbf{eröffnen zu} + Dat & (to inaugurate at / on)\\
  },
  separable = {---}
}

% Migrated from verbs/e.tex:21
\VerbTable{
  name = {erreichen},
  english = {to reach, achieve},
  type = {regular, + Akkusativ},
  auxiliary = {haben},
  style = {acc},
  present = {erreiche,erreichst,erreicht,erreichen,erreicht,erreichen},
  preterite = {erreichte,erreichtest,erreichte,erreichten,erreichtet,erreichten},
  subjunctive = {},
  imperative = {erreich(e) (du)!,erreichen wir!,erreicht (ihr)!,erreichen Sie!},
  perfect = {haben + \textbf{erreicht}},
  future = {werden + \textbf{erreichen}},
  pluperfect = {hatten + \textbf{erreicht}},
  prepositions = {---},
  separable = {---}
}

% Migrated from verbs/e.tex:23
\VerbTable{
  name = {erscheinen},
  english = {to appear},
  type = {irregular, inseparable},
  auxiliary = {sein},
  style = {dat},
  present = {erscheine,erscheinst,erscheint,erscheinen,erscheint,erscheinen},
  preterite = {erschien,erschienst,erschien,erschienen,erschient,erschienen},
  subjunctive = {},
  imperative = {erschein(e) (du)!,erscheinen wir!,erscheint (ihr)!,erscheinen Sie!},
  perfect = {sein + \textbf{erschienen}},
  future = {werden + \textbf{erscheinen}},
  pluperfect = {waren + \textbf{erschienen}},
  prepositions = {
    & \textbf{erscheinen in} + Dat & (to appear in)\\
    & \textbf{erscheinen vor} + Dat & (to appear before)\\
  },
  separable = {& -- & --\\}
}

% Migrated from verbs/b1/e.tex:33
\VerbTable{
  name = {erschrecken (intransitive)},
  english = {to be startled},
  type = {strong / irregular},
  auxiliary = {sein},
  style = {default},
  present = {erschrecke,erschrickst,erschrickt,erschrecken,erschreckt,erschrecken},
  preterite = {erschrak,erschrakst,erschrak,erschraken,erschrakt,erschraken},
  subjunctive = {},
  imperative = {erschrick (du)!,erschrecken wir!,erschreckt (ihr)!,erschrecken Sie!},
  perfect = {sein + \textbf{erschrocken}},
  future = {werden + \textbf{erschrecken}},
  pluperfect = {waren + \textbf{erschrocken}},
  prepositions = {---},
  separable = {---}
}

% Migrated from verbs/b1/e.tex:34
\VerbTable{
  name = {erschrecken (transitive)},
  english = {to frighten / scare},
  type = {regular},
  auxiliary = {haben},
  style = {default},
  present = {erschrecke,erschreckst,erschreckt,erschrecken,erschreckt,erschrecken},
  preterite = {erschreckte,erschrecktest,erschreckte,erschreckten,erschrecktet,erschreckten},
  subjunctive = {},
  imperative = {erschreck (du)!,erschrecken wir!,erschreckt (ihr)!,erschrecken Sie!},
  perfect = {haben + \textbf{erschreckt}},
  future = {werden + \textbf{erschrecken}},
  pluperfect = {hatten + \textbf{erschreckt}},
  prepositions = {---},
  separable = {---}
}

% Migrated from verbs/e.tex:24
\VerbTable{
  name = {ersetzen},
  english = {to replace},
  type = {regular, inseparable},
  auxiliary = {haben},
  style = {default},
  present = {ersetze,ersetzt,ersetzt,ersetzen,ersetzt,ersetzen},
  preterite = {ersetzte,ersetztest,ersetzte,ersetzten,ersetztet,ersetzten},
  subjunctive = {},
  imperative = {ersetz(e) (du)!,ersetzen wir!,ersetzt (ihr)!,ersetzen Sie!},
  perfect = {haben + \textbf{ersetzt}},
  future = {werden + \textbf{ersetzen}},
  pluperfect = {hatten + \textbf{ersetzt}},
  prepositions = {
    & \textbf{ersetzen durch} + Akk & (to replace with / by)\\
    & \textbf{ersetzen mit} + Dat & (to replace with)\\
  },
  separable = {---}
}

% Migrated from verbs/e.tex:25
\VerbTable{
  name = {erstellen},
  english = {to create / to prepare},
  type = {regular, inseparable},
  auxiliary = {haben},
  style = {default},
  present = {erstelle,erstellst,erstellt,erstellen,erstellt,erstellen},
  preterite = {erstellte,erstelltest,erstellte,erstellten,erstelltet,erstellten},
  subjunctive = {},
  imperative = {erstelle (du)!,erstellen wir!,erstellt (ihr)!,erstellen Sie!},
  perfect = {haben + \textbf{erstellt}},
  future = {werden + \textbf{erstellen}},
  pluperfect = {hatten + \textbf{erstellt}},
  prepositions = {
    & \textbf{erstellen für} + Akk & (to create for)\\
    & \textbf{erstellen mit} + Dat & (to create with)\\
  },
  separable = {---}
}

% Migrated from verbs/e.tex:26
\VerbTable{
  name = {erwarten},
  english = {to expect},
  type = {regular},
  auxiliary = {haben},
  style = {default},
  present = {erwarte,erwartest,erwartet,erwarten,erwartet,erwarten},
  preterite = {erwartete,erwartetest,erwartete,erwarteten,erwartetet,erwarteten},
  subjunctive = {},
  imperative = {erwarte (du)!,erwarten wir!,erwartet (ihr)!,erwarten Sie!},
  perfect = {haben + \textbf{erwartet}},
  future = {werden + \textbf{erwarten}},
  pluperfect = {hatten + \textbf{erwartet}},
  prepositions = {& \textbf{erwarten von} + Dat & (to expect from)\\},
  separable = {---}
}

% Migrated from verbs/e.tex:27
\VerbTable{
  name = {erweitern},
  english = {to expand},
  type = {regular, + Akkusativ},
  auxiliary = {haben},
  style = {acc},
  present = {erweitere,erweiterst,erweitert,erweitern,erweitert,erweitern},
  preterite = {erweiterte,erweitertest,erweiterte,erweiterten,erweitertet,erweiterten},
  subjunctive = {},
  imperative = {erweiter(e) (du)!,erweitern wir!,erweitert (ihr)!,erweitern Sie!},
  perfect = {haben + \textbf{erweitert}},
  future = {werden + \textbf{erweitern}},
  pluperfect = {hatten + \textbf{erweitert}},
  prepositions = {& \textbf{erweitern um} + Akk & (to expand by)\\},
  separable = {---}
}

% Migrated from verbs/e.tex:28
\VerbTable{
  name = {erzählen},
  english = {to tell},
  type = {regular, + Dativ},
  auxiliary = {haben},
  style = {dat},
  present = {erzähle,erzählst,erzählt,erzählen,erzählt,erzählen},
  preterite = {erzählte,erzähltest,erzählte,erzählten,erzähltet,erzählten},
  subjunctive = {},
  imperative = {erzähl(e) (du)!,erzählen wir!,erzählt (ihr)!,erzählen Sie!},
  perfect = {haben + \textbf{erzählt}},
  future = {werden + \textbf{erzählen}},
  pluperfect = {hatten + \textbf{erzählen}},
  prepositions = {---},
  separable = {---}
}

% Migrated from verbs/e.tex:29
\VerbTable{
  name = {erziehen},
  english = {to raise / to educate},
  type = {irregular},
  auxiliary = {haben},
  style = {default},
  present = {erziehe,erziehst,erzieht,erziehen,erzieht,erziehen},
  preterite = {erzog,erzogst,erzog,erzogen,erzogt,erzogen},
  subjunctive = {},
  imperative = {erzieh(e) (du)!,erziehen wir!,erzieht (ihr)!,erziehen Sie!},
  perfect = {haben + \textbf{erzogen}},
  future = {werde + \textbf{erziehen}},
  pluperfect = {hatten + \textbf{erzogen}},
  prepositions = {
    & \textbf{erziehen zu} + Dat & (to raise/educate to be)\\
    & \textbf{erziehen mit} + Dat & (to raise/educate with)\\
  },
  separable = {---}
}

% Migrated from verbs/e.tex:30
\VerbTable{
  name = {essen},
  english = {to eat},
  type = {irregular, + Akkusativ},
  auxiliary = {haben},
  style = {acc},
  present = {esse,isst,isst,essen,esst,essen},
  preterite = {aß,aßest,aß,aßen,aßt,aßen},
  subjunctive = {},
  imperative = {iss (du)!,essen wir!,esst (ihr)!,essen Sie!},
  perfect = {haben + \textbf{gegessen}},
  future = {werden + \textbf{essen}},
  pluperfect = {hatten + \textbf{gegessen}},
  prepositions = {---},
  separable = {---}
}

\section{F}

% Migrated from verbs/f.tex:1
\VerbTable{
  name = {fahren},
  english = {to drive},
  type = {irregular, + Akkusativ},
  auxiliary = {haben/sein},
  style = {acc},
  present = {fahre,f\"ahrst,f\"ahrt,fahren,fahrt,fahren},
  preterite = {fuhr,fuhrst,fuhr,fuhren,fuhrt,fuhren},
  subjunctive = {},
  imperative = {fahr(e) (du)!,fahren wir!,fahrt (ihr)!,fahren Sie!},
  perfect = {sein/haben + \textbf{gefahren}},
  future = {werden + \textbf{fahren}},
  pluperfect = {sein/hatten + \textbf{gefahren}},
  prepositions = {
    & \textbf{fahren nach} + Ort & (to go to city / country)\\
    &\textbf{fahren in} + Akk & (to go into)\\
    &\textbf{fahren zu} + Dat & (to go to place / person)\\
    &\textbf{fahren mit} + Dat & (to go with)\\
    &\textbf{fahren über} + Akk & (to go via)\\
    &\textbf{fahren von} + Akk  \textbf{nach} + Akk & (to go from ... to)\\
    &\textbf{fahren durch} + Akk & (to go through)\\
  },
  separable = {
    & \textbf{ab$|$fahren} & (to depart)\\
    & \textbf{an$|$fahren} & (to start driving)\\
    & \textbf{aus$|$fahren} & (to drive out)\\
    & \textbf{ein$|$fahren} & (to enter / pull in)\\
    & \textbf{los$|$fahren} & (to set off)\\
    & \textbf{mit$|$fahren} & (to ride along)\\
    & \textbf{vor$|$fahren} & (to pull up)\\
    & \textbf{weg$|$fahren} & (to drive away)\\
    & \textbf{zurück$|$fahren} & (to bring back)\\
  }
}

% Migrated from verbs/f.tex:2
\VerbTable{
  name = {fallen},
  english = {to fall},
  type = {irregular},
  auxiliary = {sein},
  style = {default},
  present = {falle,fällst,fällt,fallen,fallt,fallen},
  preterite = {fiel,fielst,fiel,fielen,fielt,fielen},
  subjunctive = {},
  imperative = {fall(e) (du)!,fallen wir!,fallt (ihr)!,fallen Sie!},
  perfect = {sein + \textbf{gefallen}},
  future = {werden + \textbf{fallen}},
  pluperfect = {waren + \textbf{gefallen}},
  prepositions = {
    & \textbf{fallen in} + Akk & (to fall into)\\
    & \textbf{fallen auf} + Akk & (to fall onto)\\
    &\textbf{fallen von} + Dat & (to fall from)\\
    &\textbf{fallen über} + Akk & (to fall over)\\
    &\textbf{fallen an} + Akk & (to accrue / arise costs or interest)\\
  },
  separable = {
    & \textbf{ab$|$fallen} & (to fall / drop off)\\
    & \textbf{auf$|$fallen} & (to stand out)\\
    & \textbf{aus$|$fallen} & (to be cancelled / to fail)\\
    & \textbf{ein$|$fallen} & (to occur to someone)\\
    & \textbf{um$|$fallen} & (to fall over)\\
    & \textbf{zurück$|$fallen} & (to fall behind)\\
  }
}

% Migrated from verbs/f.tex:3
\VerbTable{
  name = {fangen},
  english = {to catch / capture},
  type = {irregular, + Akkusativ},
  auxiliary = {haben},
  style = {acc},
  present = {fange,fängst,fängt,fangen,fangt,fangen},
  preterite = {fing,fingst,fing,fingen,fingt,fingen},
  subjunctive = {},
  imperative = {fang(e) (du)!,fangen wir!,fangt (ihr)!,fangen Sie!},
  perfect = {haben + \textbf{gefangen}},
  future = {werden + \textbf{fangen}},
  pluperfect = {hatten + \textbf{gefangen}},
  prepositions = {
    & \textbf{fangen nach} + Dat & (to reach for / try to catch)\\
    &\textbf{anfangen mit} + Dat & (to begin with)\\
    &\textbf{anfangen bei} + Dat & (to begin at (point / person))\\
    & \textbf{einfangen in} + Akk & (to capture in)\\
    & \textbf{abfangen an} + Dat & (to intercept at)
  },
  separable = {
    & \textbf{an$|$fangen} & (to begin)\\
    & \textbf{auf$|$fangen} & (to catch / absorb / cushion)\\
    & \textbf{ein$|$fangen} & (to catch / capture)\\
    & \textbf{ab$|$fangen} & (to intercept)\\
    & \textbf{weg$|$fangen} & (to snatch away)\\
    & \textbf{über$|$fangen} & (to cover / overtake)
  }
}

% Migrated from verbs/b1/f.tex:1
\VerbTable{
  name = {fassen},
  english = {to grasp / hold},
  type = {regular},
  auxiliary = {haben},
  style = {default},
  present = {fasse,fasst,fasst,fassen,fasst,fassen},
  preterite = {fasste,fasstest,fasste,fassten,fasstet,fassten},
  subjunctive = {},
  imperative = {fass (du)!,fassen wir!,fasst (ihr)!,fassen Sie!},
  perfect = {haben + \textbf{gefasst}},
  future = {werden + \textbf{fassen}},
  pluperfect = {hatten + \textbf{gefasst}},
  prepositions = {---},
  separable = {---}
}

% Migrated from verbs/f.tex:4
\VerbTable{
  name = {faszinieren},
  english = {to fascinate},
  type = {regular, + Akkusativ},
  auxiliary = {haben},
  style = {acc},
  present = {fasziniere,faszinierst,fasziniert,faszinieren,fasziniert,faszinieren},
  preterite = {faszinierte,fasziniertest,faszinierte,faszinierten,fasziniertet,faszinierten},
  subjunctive = {},
  imperative = {fasziniere (du)!,faszinieren wir!,fasziniert (ihr)!,faszinieren Sie!},
  perfect = {haben + \textbf{fasziniert}},
  future = {werden + \textbf{faszinieren}},
  pluperfect = {hatten + \textbf{fasziniert}},
  prepositions = {---},
  separable = {---}
}

% Migrated from verbs/b1/f.tex:2
\VerbTable{
  name = {faulenzen},
  english = {to laze},
  type = {regular},
  auxiliary = {haben},
  style = {default},
  present = {faulenze,faulenzt,faulenzt,faulenzen,faulenzt,faulenzen},
  preterite = {faulenzte,faulenztest,faulenzte,faulenzten,faulenztet,faulenzten},
  subjunctive = {},
  imperative = {faulenz (du)!,faulenzen wir!,faulenzt (ihr)!,faulenzen Sie!},
  perfect = {haben + \textbf{gefaulenzt}},
  future = {werden + \textbf{faulenzen}},
  pluperfect = {hatten + \textbf{gefaulenzt}},
  prepositions = {---},
  separable = {---}
}

% Migrated from verbs/f.tex:5
\VerbTable{
  name = {fehlen},
  english = {to be missing},
  type = {regular},
  auxiliary = {haben},
  style = {default},
  present = {fehle,fehlst,fehlt,fehlen,fehlt,fehlen},
  preterite = {fehlte,fehltest,fehlte,fehlten,fehltet,fehlten},
  subjunctive = {},
  imperative = {fehl(e) (du)!,fehlen wir!,fehlt (ihr)!,fehlen Sie!},
  perfect = {haben + \textbf{gefehlt}},
  future = {werden + \textbf{fehlen}},
  pluperfect = {hatten + \textbf{gefehlt}},
  prepositions = {---},
  separable = {---}
}

% Migrated from verbs/f.tex:6
\VerbTable{
  name = {feiern},
  english = {to celebrate},
  type = {regular, + Akkusativ},
  auxiliary = {haben},
  style = {acc},
  present = {feiere,feierst,feiert,feiern,feiert,feiern},
  preterite = {feierte,feiertest,feierte,feierten,feiertet,feierten},
  subjunctive = {},
  imperative = {feier(e) (du)!,feiern wir!,feiert (ihr)!,feiern Sie!},
  perfect = {haben + \textbf{gefeiert}},
  future = {werden + \textbf{feiern}},
  pluperfect = {hatten + \textbf{gefeiert}},
  prepositions = {---},
  separable = {---}
}

% Migrated from verbs/b1/f.tex:3
\VerbTable{
  name = {fernsehen},
  english = {to watch television},
  type = {separable, strong / irregular},
  auxiliary = {haben},
  style = {default},
  present = {sehe fern,siehst fern,sieht fern,sehen fern,seht fern,sehen fern},
  preterite = {sah fern,sahst fern,sah fern,sahen fern,saht fern,sahen fern},
  subjunctive = {},
  imperative = {sieh (du) fern!,sehen wir fern!,seht (ihr) fern!,sehen Sie fern!},
  perfect = {haben + \textbf{ferngesehen}},
  future = {werden + \textbf{fernsehen}},
  pluperfect = {hatten + \textbf{ferngesehen}},
  prepositions = {---},
  separable = {---}
}

% Migrated from verbs/b1/f.tex:5
\VerbTable{
  name = {festhalten},
  english = {to hold on / hold tight},
  type = {separable, strong / irregular},
  auxiliary = {haben},
  style = {default},
  present = {halte fest,hältst fest,hält fest,halten fest,haltet fest,halten fest},
  preterite = {hielt fest,hieltest fest,hielt fest,hielten fest,hieltet fest,hielten fest},
  subjunctive = {},
  imperative = {halt (du) fest!,halten wir fest!,haltet (ihr) fest!,halten Sie fest!},
  perfect = {haben + \textbf{festgehalten}},
  future = {werden + \textbf{festhalten}},
  pluperfect = {hatten + \textbf{festgehalten}},
  prepositions = {---},
  separable = {---}
}

% Migrated from verbs/b1/f.tex:4
\VerbTable{
  name = {festlegen},
  english = {to set / determine},
  type = {separable, regular},
  auxiliary = {haben},
  style = {default},
  present = {lege fest,legst fest,legt fest,legen fest,legt fest,legen fest},
  preterite = {legte fest,legtest fest,legte fest,legten fest,legtet fest,legten fest},
  subjunctive = {},
  imperative = {leg (du) fest!,legen wir fest!,legt (ihr) fest!,legen Sie fest!},
  perfect = {haben + \textbf{festgelegt}},
  future = {werden + \textbf{festlegen}},
  pluperfect = {hatten + \textbf{festgelegt}},
  prepositions = {---},
  separable = {---}
}

% Migrated from verbs/b1/f.tex:6
\VerbTable{
  name = {festsetzen},
  english = {to set / fix},
  type = {separable, regular},
  auxiliary = {haben},
  style = {default},
  present = {setze fest,setzt fest,setzt fest,setzen fest,setzt fest,setzen fest},
  preterite = {setzte fest,setztest fest,setzte fest,setzten fest,setztet fest,setzten fest},
  subjunctive = {},
  imperative = {setz (du) fest!,setzen wir fest!,setzt (ihr) fest!,setzen Sie fest!},
  perfect = {haben + \textbf{festgesetzt}},
  future = {werden + \textbf{festsetzen}},
  pluperfect = {hatten + \textbf{festgesetzt}},
  prepositions = {---},
  separable = {---}
}

% Migrated from verbs/b1/f.tex:7
\VerbTable{
  name = {feststehen},
  english = {to be certain / be fixed},
  type = {separable, strong / irregular},
  auxiliary = {sein},
  style = {default},
  present = {stehe fest,stehst fest,steht fest,stehen fest,steht fest,stehen fest},
  preterite = {stand fest,standest fest,stand fest,standen fest,standet fest,standen fest},
  subjunctive = {},
  imperative = {steh (du) fest!,stehen wir fest!,steht (ihr) fest!,stehen Sie fest!},
  perfect = {sein + \textbf{festgestanden}},
  future = {werden + \textbf{feststehen}},
  pluperfect = {waren + \textbf{festgestanden}},
  prepositions = {---},
  separable = {---}
}

% Migrated from verbs/b1/f.tex:8
\VerbTable{
  name = {feststellen},
  english = {to determine},
  type = {separable, regular},
  auxiliary = {haben},
  style = {default},
  present = {stelle fest,stellst fest,stellt fest,stellen fest,stellt fest,stellen fest},
  preterite = {stellte fest,stelltest fest,stellte fest,stellten fest,stelltet fest,stellten fest},
  subjunctive = {},
  imperative = {stell (du) fest!,stellen wir fest!,stellt (ihr) fest!,stellen Sie fest!},
  perfect = {haben + \textbf{festgestellt}},
  future = {werden + \textbf{feststellen}},
  pluperfect = {hatten + \textbf{festgestellt}},
  prepositions = {---},
  separable = {---}
}

% Migrated from verbs/f.tex:7
\VerbTable{
  name = {finanzieren},
  english = {to finance / fund},
  type = {regular},
  auxiliary = {haben},
  style = {default},
  present = {finanziere,finanzierst,finanziert,finanzieren,finanziert,finanzieren},
  preterite = {finanzierte,finanziertest,finanzierte,finanzierten,finanziertet,finanzierten},
  subjunctive = {},
  imperative = {finanziere (du)!,finanzieren wir!,finanziert (ihr)!,finanzieren Sie!},
  perfect = {haben + \textbf{finanziert}},
  future = {werden + \textbf{finanzieren}},
  pluperfect = {hatten + \textbf{finanziert}},
  prepositions = {
    & \textbf{finanzieren mit} + Dat & (to finance with)\\
    & \textbf{finanzieren durch} + Akk & (to finance through / by)\\
  },
  separable = {---}
}

% Migrated from verbs/f.tex:8
\VerbTable{
  name = {finden},
  english = {to find},
  type = {irregular, + Akkusativ},
  auxiliary = {haben},
  style = {acc},
  present = {finde,findest,findet,finden,findet,finden},
  preterite = {fand,fandest,fand,fanden,fandet,fanden},
  subjunctive = {},
  imperative = {find(e) (du)!,finden wir!,findet (ihr)!,finden Sie!},
  perfect = {haben + \textbf{gefunden}},
  future = {werden + \textbf{finden}},
  pluperfect = {hatten + \textbf{gefunden}},
  prepositions = {
    & \textbf{finden für} + Akk & (to find suitable for)\\
    & \textbf{finden zu} + Dat & (to find to / tend to)\\
    & \textbf{finden in} + Dat & (to find in / inside)\\
  },
  separable = {
    & \textbf{wieder$|$finden} & (to recover / find again)\\
    & \textbf{statt$|$finden} & (to take place / happen)\\
  }
}

% Migrated from verbs/f.tex:9
\VerbTable{
  name = {fliegen},
  english = {to fly},
  type = {irregular},
  auxiliary = {sein},
  style = {default},
  present = {fliege,fliegst,fliegt,fliegen,fliegt,fliegen},
  preterite = {flog,flogst,flog,flogen,flogt,flogen},
  subjunctive = {},
  imperative = {flieg(e) (du)!,fliegen wir!,fliegt (ihr)!,fliegen Sie!},
  perfect = {sein + \textbf{geflogen}},
  future = {werden + \textbf{fliegen}},
  pluperfect = {waren + \textbf{geflogen}},
  prepositions = {
    & \textbf{fliegen nach} + Dat & (to fly to – cities/countries)\\
    & \textbf{fliegen in} + Akk & (to fly into)\\
    & \textbf{fliegen über} + Akk & (to fly over)\\
    & \textbf{fliegen von} + Dat & (to fly from)\\
  },
  separable = {
    & \textbf{ab$|$fliegen} & (to depart by plane)\\
    & \textbf{an$|$fliegen} & (to approach by flying)\\
    & \textbf{ein$|$fliegen} & (to fly in)\\
    & \textbf{weg$|$fliegen} & (to fly away)\\
    & \textbf{über$|$fliegen} & (to fly over)\\
  }
}

% Migrated from verbs/f.tex:10
\VerbTable{
  name = {fliehen},
  english = {to flee},
  type = {irregular},
  auxiliary = {sein},
  style = {default},
  present = {fliehe,fliehst,flieht,fliehen,flieht,fliehen},
  preterite = {floh,flohst,floh,flohen,floht,flohen},
  subjunctive = {},
  imperative = {flieh(e) (du)!,fliehen wir!,flieht (ihr)!,fliehen Sie!},
  perfect = {sein + \textbf{geflohen}},
  future = {werden + \textbf{fliehen}},
  pluperfect = {waren + \textbf{geflohen}},
  prepositions = {
    & \textbf{fliehen vor} + Dat & (to flee from)\\
    & \textbf{fliehen aus} + Dat & (to flee out of / from)\\
  },
  separable = {& \textbf{ent$|$fliehen} & (to escape from)\\}
}

% Migrated from verbs/b1/f.tex:9
\VerbTable{
  name = {fließen},
  english = {to flow},
  type = {strong / irregular},
  auxiliary = {sein},
  style = {default},
  present = {fließe,fließt,fließt,fließen,fließt,fließen},
  preterite = {floss,flosst,floss,flossen,flosst,flossen},
  subjunctive = {},
  imperative = {fließ (du)!,fließen wir!,fließt (ihr)!,fließen Sie!},
  perfect = {sein + \textbf{geflossen}},
  future = {werden + \textbf{fließen}},
  pluperfect = {waren + \textbf{geflossen}},
  prepositions = {---},
  separable = {---}
}

% Migrated from verbs/b1/f.tex:11
\VerbTable{
  name = {fördern},
  english = {to promote},
  type = {regular},
  auxiliary = {haben},
  style = {default},
  present = {fördre,förderst,fördert,fördern,fördert,fördern},
  preterite = {förderte,fördertest,förderte,förderten,fördertet,förderten},
  subjunctive = {},
  imperative = {fördre (du)!,fördern wir!,fördert (ihr)!,fördern Sie!},
  perfect = {haben + \textbf{gefördert}},
  future = {werden + \textbf{fördern}},
  pluperfect = {hatten + \textbf{gefördert}},
  prepositions = {---},
  separable = {---}
}

% Migrated from verbs/f.tex:11
\VerbTable{
  name = {folgen},
  english = {to follow},
  type = {regular},
  auxiliary = {sein},
  style = {dat},
  present = {folge,folgst,folgt,folgen,folgt,folgen},
  preterite = {folgte,folgtest,folgte,folgten,folgtet,folgten},
  subjunctive = {},
  imperative = {folg(e) (du)!,folgen wir!,folgt (ihr)!,folgen Sie!},
  perfect = {sein + \textbf{gefolgt}},
  future = {werden + \textbf{folgen}},
  pluperfect = {waren + \textbf{gefolgt}},
  prepositions = {& \textbf{folgen auf} + Akk & (to follow after)\\},
  separable = {
    & \textbf{nach$|$folgen} & (to succeed / follow after)\\
    & \textbf{ver$|$folgen} & (not separable: to pursue)\\
  }
}

% Migrated from verbs/b1/f.tex:10
\VerbTable{
  name = {fordern},
  english = {to demand},
  type = {regular},
  auxiliary = {haben},
  style = {default},
  present = {fordre,forderst,fordert,fordern,fordert,fordern},
  preterite = {forderte,fordertest,forderte,forderten,fordertet,forderten},
  subjunctive = {},
  imperative = {fordre (du)!,fordern wir!,fordert (ihr)!,fordern Sie!},
  perfect = {haben + \textbf{gefordert}},
  future = {werden + \textbf{fordern}},
  pluperfect = {hatten + \textbf{gefordert}},
  prepositions = {---},
  separable = {---}
}

% Migrated from verbs/b1/f.tex:12
\VerbTable{
  name = {fortsetzen},
  english = {to continue},
  type = {separable, regular},
  auxiliary = {haben},
  style = {default},
  present = {setze fort,setzt fort,setzt fort,setzen fort,setzt fort,setzen fort},
  preterite = {setzte fort,setztest fort,setzte fort,setzten fort,setztet fort,setzten fort},
  subjunctive = {},
  imperative = {setz (du) fort!,setzen wir fort!,setzt (ihr) fort!,setzen Sie fort!},
  perfect = {haben + \textbf{fortgesetzt}},
  future = {werden + \textbf{fortsetzen}},
  pluperfect = {hatten + \textbf{fortgesetzt}},
  prepositions = {---},
  separable = {---}
}

% Migrated from verbs/b1/f.tex:13
\VerbTable{
  name = {fotografieren},
  english = {to photograph},
  type = {regular},
  auxiliary = {haben},
  style = {default},
  present = {fotografiere,fotografierst,fotografiert,fotografieren,fotografiert,fotografieren},
  preterite = {fotografierte,fotografiertest,fotografierte,fotografierten,fotografiertet,fotografierten},
  subjunctive = {},
  imperative = {fotografier (du)!,fotografieren wir!,fotografiert (ihr)!,fotografieren Sie!},
  perfect = {haben + \textbf{fotografiert}},
  future = {werden + \textbf{fotografieren}},
  pluperfect = {hatten + \textbf{fotografiert}},
  prepositions = {---},
  separable = {---}
}

% Migrated from verbs/f.tex:12
\VerbTable{
  name = {fragen},
  english = {to ask},
  type = {regular, + Akkusativ},
  auxiliary = {haben},
  style = {acc},
  present = {frage,fragst,fragt,fragen,fragt,fragen},
  preterite = {fragte,fragtest,fragte,fragten,fragtet,fragten},
  subjunctive = {},
  imperative = {frag(e) (du)!,fragen wir!,fragt (ihr)!,fragen Sie!},
  perfect = {haben + \textbf{gefragt}},
  future = {werden + \textbf{fragen}},
  pluperfect = {hatten + \textbf{gefragt}},
  prepositions = {---},
  separable = {---}
}

% Migrated from verbs/f.tex:13
\VerbTable{
  name = {fressen},
  english = {to eat — animals},
  type = {irregular},
  auxiliary = {haben},
  style = {default},
  present = {fresse,frisst,frisst,fressen,fresst,fressen},
  preterite = {fraß,fraßest,fraß,fraßen,fraßt,fraßen},
  subjunctive = {},
  imperative = {friss (du)!,fressen wir!,fresst (ihr)!,fressen Sie!},
  perfect = {haben + \textbf{gefressen}},
  future = {werde + \textbf{fressen}},
  pluperfect = {hatten + \textbf{gefressen}},
  prepositions = {
    & \textbf{fressen von} + Dat & (to eat from / of)\\
    & \textbf{fressen aus} + Dat & (to eat from)\\
  },
  separable = {
    & \textbf{auf$|$fressen} & (to eat up)\\
    & \textbf{an$|$fressen} & (to nibble at / eat into)\\
  }
}

% Migrated from verbs/b1/f.tex:14
\VerbTable{
  name = {frieren},
  english = {to freeze / be cold},
  type = {strong / irregular},
  auxiliary = {haben},
  style = {default},
  present = {friere,frierst,friert,frieren,friert,frieren},
  preterite = {fror,frorst,fror,froren,frort,froren},
  subjunctive = {},
  imperative = {frier (du)!,frieren wir!,friert (ihr)!,frieren Sie!},
  perfect = {haben + \textbf{gefroren}},
  future = {werden + \textbf{frieren}},
  pluperfect = {hatten + \textbf{gefroren}},
  prepositions = {---},
  separable = {---}
}

% Migrated from verbs/b1/f.tex:15
\VerbTable{
  name = {frühstücken},
  english = {to have breakfast},
  type = {regular},
  auxiliary = {haben},
  style = {default},
  present = {frühstücke,frühstückst,frühstückt,frühstücken,frühstückt,frühstücken},
  preterite = {frühstückte,frühstücktest,frühstückte,frühstückten,frühstücktet,frühstückten},
  subjunctive = {},
  imperative = {frühstück (du)!,frühstücken wir!,frühstückt (ihr)!,frühstücken Sie!},
  perfect = {haben + \textbf{gefrühstückt}},
  future = {werden + \textbf{frühstücken}},
  pluperfect = {hatten + \textbf{gefrühstückt}},
  prepositions = {---},
  separable = {---}
}

% Migrated from verbs/f.tex:14
\VerbTable{
  name = {fügen},
  english = {to add / insert},
  type = {regular},
  auxiliary = {haben},
  style = {default},
  present = {füge,fügst,fügt,fügen,fügt,fügen},
  preterite = {fügte,fügtest,fügte,fügten,fügtet,fügten},
  subjunctive = {},
  imperative = {füg(e) (du)!,fügen wir!,fügt (ihr)!,fügen Sie!},
  perfect = {haben + \textbf{gefügt}},
  future = {werden + \textbf{fügen}},
  pluperfect = {hatten + \textbf{gefügt}},
  prepositions = {& \textbf{sich fügen in} + Akk & (to submit to)\\},
  separable = {
    & \textbf{ein$|$fügen} & (to insert)\\
    & \textbf{hinzu$|$fügen} & (to add)\\
  }
}

% Migrated from verbs/f.tex:15
\VerbTable{
  name = {führen},
  english = {to lead / guide / drive},
  type = {regular},
  auxiliary = {haben},
  style = {default},
  present = {führe,führst,führt,führen,führt,führen},
  preterite = {führte,führtest,führte,führten,führtet,führten},
  subjunctive = {},
  imperative = {führe (du)!,führen wir!,führt (ihr)!,führen Sie!},
  perfect = {haben + \textbf{geführt}},
  future = {werden + \textbf{führen}},
  pluperfect = {hatten + \textbf{geführt}},
  prepositions = {
    & \textbf{führen zu} + Dat & (to lead to / result in)\\
    & \textbf{führen durch} + Akk & (to guide through)\\
    & \textbf{führen über} + Akk & (to lead across / over)\\
  },
  separable = {
    & \textbf{an$|$führen} & (to lead to / bring to)\\
    & \textbf{durch$|$führen} & (to guide through)\\
  }
}

% Migrated from verbs/f.tex:16
\VerbTable{
  name = {füllen},
  english = {to fill},
  type = {regular},
  auxiliary = {haben},
  style = {default},
  present = {fülle,füllst,füllt,füllen,füllt,füllen},
  preterite = {füllte,fülltest,füllte,füllten,fülltet,füllten},
  subjunctive = {},
  imperative = {füll(e) (du)!,füllen wir!,füllt (ihr)!,füllen Sie!},
  perfect = {haben + \textbf{gefüllt}},
  future = {werde + \textbf{füllen}},
  pluperfect = {hatten + \textbf{gefüllt}},
  prepositions = {
    & \textbf{füllen mit} + Dat & (to fill with)\\
    & \textbf{füllen in} + Akk & (to fill into)\\
  },
  separable = {
    & \textbf{aus$|$füllen} & (to fill out)\\
    & \textbf{auf$|$füllen} & (to refill / top up)\\
  }
}

% Migrated from verbs/f.tex:18
\VerbTable{
  name = {füttern},
  english = {to feed},
  type = {regular},
  auxiliary = {haben},
  style = {default},
  present = {füttere,fütterst,füttert,füttern,füttert,füttern},
  preterite = {fütterte,füttertest,fütterte,fütterten,füttertet,fütterten},
  subjunctive = {},
  imperative = {füttere (du)!,füttern wir!,füttert (ihr)!,füttern Sie!},
  perfect = {haben + \textbf{gefüttert}},
  future = {werde + \textbf{füttern}},
  pluperfect = {hatten + \textbf{gefüttert}},
  prepositions = {
    & \textbf{füttern mit} + Dat & (to feed with)\\
    & \textbf{füttern an} + Akk & (to feed to)\\
  },
  separable = {
    & \textbf{an$|$füttern} & (to start feeding / bait)\\
    & \textbf{zu$|$füttern} & (to give supplementary food)\\
  }
}

% Migrated from verbs/f.tex:17
\VerbTable{
  name = {funktionieren},
  english = {to function / work},
  type = {regular},
  auxiliary = {haben},
  style = {default},
  present = {funktioniere,funktionierst,funktioniert,funktionieren,funktioniert,funktionieren},
  preterite = {funktionierte,funktioniertest,funktionierte,funktionierten,funktioniertet,funktionierten},
  subjunctive = {},
  imperative = {funktioniere (du)!,funktionieren wir!,funktioniert (ihr)!,funktionieren Sie!},
  perfect = {haben + \textbf{funktioniert}},
  future = {werden + \textbf{funktionieren}},
  pluperfect = {hatten + \textbf{funktioniert}},
  prepositions = {& \textbf{funktionieren mit} + Dat & (to function with)\\},
  separable = {---}
}

\section{G}

% Migrated from verbs/b1/g.tex:1
\VerbTable{
  name = {garantieren},
  english = {to guarantee},
  type = {regular},
  auxiliary = {haben},
  style = {default},
  present = {garantiere,garantierst,garantiert,garantieren,garantiert,garantieren},
  preterite = {garantierte,garantiertest,garantierte,garantierten,garantiertet,garantierten},
  subjunctive = {},
  imperative = {garantier (du)!,garantieren wir!,garantiert (ihr)!,garantieren Sie!},
  perfect = {haben + \textbf{garantiert}},
  future = {werden + \textbf{garantieren}},
  pluperfect = {hatten + \textbf{garantiert}},
  prepositions = {---},
  separable = {---}
}

% Migrated from verbs/g.tex:1
\VerbTable{
  name = {geben},
  english = {to give},
  type = {irregular, + Dativ},
  auxiliary = {haben},
  style = {dat},
  present = {gebe,gibst,gibt,geben,gebt,geben},
  preterite = {gab,gabst,gab,gaben,gabt,gaben},
  subjunctive = {},
  imperative = {gib (du)!,geben wir!,gebt (ihr)!,geben Sie!},
  perfect = {haben + \textbf{gegeben}},
  future = {werden + \textbf{geben}},
  pluperfect = {hatten + \textbf{gegeben}},
  prepositions = {---},
  separable = {---}
}

% Migrated from verbs/b1/g.tex:2
\VerbTable{
  name = {gebrauchen},
  english = {to use},
  type = {regular},
  auxiliary = {haben},
  style = {default},
  present = {gebrauche,gebrauchst,gebraucht,gebrauchen,gebraucht,gebrauchen},
  preterite = {gebrauchte,gebrauchtest,gebrauchte,gebrauchten,gebrauchtet,gebrauchten},
  subjunctive = {},
  imperative = {gebrauch (du)!,gebrauchen wir!,gebraucht (ihr)!,gebrauchen Sie!},
  perfect = {haben + \textbf{gebraucht}},
  future = {werden + \textbf{gebrauchen}},
  pluperfect = {hatten + \textbf{gebraucht}},
  prepositions = {---},
  separable = {---}
}

% Migrated from verbs/g.tex:2
\VerbTable{
  name = {gefallen},
  english = {to like},
  type = {irregular, + Dativ},
  auxiliary = {haben},
  style = {dat},
  present = {gefalle,gefällst,gefällt,gefallen,gefallt,gefallen},
  preterite = {gefiel,gefielst,gefiel,gefielen,gefielt,gefielen},
  subjunctive = {},
  imperative = {gefall(e) (du)!,gefallen wir!,gefallt (ihr)!,gefallen Sie!},
  perfect = {haben + \textbf{gefallen}},
  future = {werden + \textbf{gefallen}},
  pluperfect = {hatten + \textbf{gefallen}},
  prepositions = {---},
  separable = {---}
}

% Migrated from verbs/g.tex:4
\VerbTable{
  name = {gehen},
  english = {to go},
  type = {irregular},
  auxiliary = {sein},
  style = {acc},
  present = {gehe,gehst,geht,gehen,geht,gehen},
  preterite = {ging,gingst,ging,gingen,gingt,gingen},
  subjunctive = {},
  imperative = {geh(e) (du)!,gehen wir!,geht (ihr)!,gehen Sie!},
  perfect = {sein + \textbf{gegangen}},
  future = {werden + \textbf{gehen}},
  pluperfect = {waren + \textbf{gegangen}},
  prepositions = {---},
  separable = {---}
}

% Migrated from verbs/g.tex:3
\VerbTable{
  name = {gehören},
  english = {to belong},
  type = {regular, + Dativ},
  auxiliary = {haben},
  style = {dat},
  present = {gehöre,gehörst,gehört,gehören,gehört,gehören},
  preterite = {gehörte,gehörtest,gehörte,gehörten,gehörtet,gehörten},
  subjunctive = {},
  imperative = {gehör(e) (du)!,gehören wir!,gehört (ihr)!,gehören Sie!},
  perfect = {haben + \textbf{gehört}},
  future = {werden + \textbf{gehören}},
  pluperfect = {hatten + \textbf{gehört}},
  prepositions = {---},
  separable = {---}
}

% Migrated from verbs/b1/g.tex:3
\VerbTable{
  name = {gelingen},
  english = {to succeed},
  type = {strong / irregular},
  auxiliary = {sein},
  style = {default},
  present = {gelinge,gelingst,gelingt,gelingen,gelingt,gelingen},
  preterite = {gelang,gelangst,gelang,gelangen,gelangt,gelangen},
  subjunctive = {},
  imperative = {geling (du)!,gelingen wir!,gelingt (ihr)!,gelingen Sie!},
  perfect = {sein + \textbf{gelungen}},
  future = {werden + \textbf{gelingen}},
  pluperfect = {waren + \textbf{gelungen}},
  prepositions = {---},
  separable = {---}
}

% Migrated from verbs/g.tex:5
\VerbTable{
  name = {gelten},
  english = {to be valid / to apply},
  type = {irregular},
  auxiliary = {haben},
  style = {default},
  present = {gelte,giltst,gilt,gelten,geltet,gelten},
  preterite = {galt,galtest,galt,galten,galtet,galten},
  subjunctive = {},
  imperative = {gilt (du)!,gelten wir!,geltet (ihr)!,gelten Sie!},
  perfect = {haben + \textbf{gegolten}},
  future = {werden + \textbf{gelten}},
  pluperfect = {hatten + \textbf{gegolten}},
  prepositions = {
    & \textbf{gelten für} + Akk & (to apply to / be valid for)\\
    & \textbf{gelten als} + Nom & (to be considered as)\\
  },
  separable = {---}
}

% Migrated from verbs/b1/g.tex:4
\VerbTable{
  name = {genehmigen},
  english = {to approve},
  type = {regular},
  auxiliary = {haben},
  style = {default},
  present = {genehmige,genehmigst,genehmigt,genehmigen,genehmigt,genehmigen},
  preterite = {genehmigte,genehmigtest,genehmigte,genehmigten,genehmigtet,genehmigten},
  subjunctive = {},
  imperative = {genehmig (du)!,genehmigen wir!,genehmigt (ihr)!,genehmigen Sie!},
  perfect = {haben + \textbf{genehmigt}},
  future = {werden + \textbf{genehmigen}},
  pluperfect = {hatten + \textbf{genehmigt}},
  prepositions = {---},
  separable = {---}
}

% Migrated from verbs/g.tex:7
\VerbTable{
  name = {genießen},
  english = {to enjoy},
  type = {irregular, + Akkusativ},
  auxiliary = {haben},
  style = {acc},
  present = {genieße,genießt,genießt,genießen,genießt,genießen},
  preterite = {genoss,genossest,genoss,genossen,genosst,genossen},
  subjunctive = {},
  imperative = {genieß(e) (du)!,genießen wir!,genießt (ihr)!,genießen Sie!},
  perfect = {haben + \textbf{genossen}},
  future = {werden + \textbf{genießen}},
  pluperfect = {hatten + \textbf{genossen}},
  prepositions = {---},
  separable = {---}
}

% Migrated from verbs/b1/g.tex:5
\VerbTable{
  name = {genügen},
  english = {to suffice},
  type = {regular},
  auxiliary = {haben},
  style = {default},
  present = {genüge,genügst,genügt,genügen,genügt,genügen},
  preterite = {genügte,genügtest,genügte,genügten,genügtet,genügten},
  subjunctive = {},
  imperative = {genüg (du)!,genügen wir!,genügt (ihr)!,genügen Sie!},
  perfect = {haben + \textbf{genügt}},
  future = {werden + \textbf{genügen}},
  pluperfect = {hatten + \textbf{genügt}},
  prepositions = {---},
  separable = {---}
}

% Migrated from verbs/b1/g.tex:6
\VerbTable{
  name = {geschehen},
  english = {to happen},
  type = {strong / irregular},
  auxiliary = {sein},
  style = {default},
  present = {geschehe,geschiehst,geschieht,geschehen,gescheht,geschehen},
  preterite = {geschah,geschahst,geschah,geschahen,geschaht,geschahen},
  subjunctive = {},
  imperative = {geschieh (du)!,geschehen wir!,gescheht (ihr)!,geschehen Sie!},
  perfect = {sein + \textbf{geschehen}},
  future = {werden + \textbf{geschehen}},
  pluperfect = {waren + \textbf{geschehen}},
  prepositions = {---},
  separable = {---}
}

% Migrated from verbs/g.tex:8
\VerbTable{
  name = {gestikulieren},
  english = {to gesticulate},
  type = {regular},
  auxiliary = {haben},
  style = {default},
  present = {gestikuliere,gestikulierst,gestikuliert,gestikulieren,gestikuliert,gestikulieren},
  preterite = {gestikulierte,gestikuliertest,gestikulierte,gestikulierten,gestikuliertet,gestikulierten},
  subjunctive = {},
  imperative = {gestikuliere (du)!,gestikulieren wir!,gestikuliert (ihr)!,gestikulieren Sie!},
  perfect = {haben + \textbf{gestikuliert}},
  future = {werden + \textbf{gestikulieren}},
  pluperfect = {hatten + \textbf{gestikuliert}},
  prepositions = {& \textbf{gestikulieren mit} + Dat & (to gesticulate with)\\},
  separable = {---}
}

% Migrated from verbs/g.tex:9
\VerbTable{
  name = {gewinnen},
  english = {to win},
  type = {irregular, + Akkusativ},
  auxiliary = {haben},
  style = {acc},
  present = {gewinne,gewinnst,gewinnt,gewinnen,gewinnt,gewinnen},
  preterite = {gewann,gewannst,gewann,gewannen,gewannt,gewannen},
  subjunctive = {},
  imperative = {gewinn(e) (du)!,gewinnen wir!,gewinnt (ihr)!,gewinnen Sie!},
  perfect = {haben + \textbf{gewonnen}},
  future = {werden + \textbf{gewinnen}},
  pluperfect = {hatten + \textbf{gewonnen}},
  prepositions = {---},
  separable = {---}
}

% Migrated from verbs/g.tex:6
\VerbTable{
  name = {gießen},
  english = {to pour / to water},
  type = {irregular},
  auxiliary = {haben},
  style = {default},
  present = {gieße,gießt,gießt,gießen,gießt,gießen},
  preterite = {goss,gossest,goss,gossen,gosst,gossen},
  subjunctive = {},
  imperative = {gieß(e) (du)!,gießen wir!,gießt (ihr)!,gießen Sie!},
  perfect = {haben + \textbf{gegossen}},
  future = {werden + \textbf{gießen}},
  pluperfect = {hatten + \textbf{gegossen}},
  prepositions = {
    & \textbf{gießen in} + Akk & (to pour into)\\
    & \textbf{gießen über} + Akk & (to pour over)\\
  },
  separable = {
    & \textbf{ein$|$gießen} & (to pour in)\\
    & \textbf{aus$|$gießen} & (to pour out)\\
  }
}

% Migrated from verbs/g.tex:10
\VerbTable{
  name = {glauben},
  english = {to believe},
  type = {regular, + Dativ},
  auxiliary = {haben},
  style = {dat},
  present = {glaube,glaubst,glaubt,glauben,glaubt,glauben},
  preterite = {glaubte,glaubtest,glaubte,glaubten,glaubtet,glaubten},
  subjunctive = {},
  imperative = {glaub(e) (du)!,glauben wir!,glaubt (ihr)!,glauben Sie!},
  perfect = {haben + \textbf{geglaubt}},
  future = {werden + \textbf{glauben}},
  pluperfect = {hatten + \textbf{geglaubt}},
  prepositions = {& \textbf{glauben an} + Akk & (to believe in)\\},
  separable = {---}
}

% Migrated from verbs/b1/g.tex:7
\VerbTable{
  name = {gratulieren},
  english = {to congratulate},
  type = {regular},
  auxiliary = {haben},
  style = {default},
  present = {gratuliere,gratulierst,gratuliert,gratulieren,gratuliert,gratulieren},
  preterite = {gratulierte,gratuliertest,gratulierte,gratulierten,gratuliertet,gratulierten},
  subjunctive = {},
  imperative = {gratulier (du)!,gratulieren wir!,gratuliert (ihr)!,gratulieren Sie!},
  perfect = {haben + \textbf{gratuliert}},
  future = {werden + \textbf{gratulieren}},
  pluperfect = {hatten + \textbf{gratuliert}},
  prepositions = {---},
  separable = {---}
}

% Migrated from verbs/b1/g.tex:8
\VerbTable{
  name = {greifen},
  english = {to grab},
  type = {strong / irregular},
  auxiliary = {haben},
  style = {default},
  present = {greife,greifst,greift,greifen,greift,greifen},
  preterite = {griff,griffst,griff,griffen,grifft,griffen},
  subjunctive = {},
  imperative = {greif (du)!,greifen wir!,greift (ihr)!,greifen Sie!},
  perfect = {haben + \textbf{gegriffen}},
  future = {werden + \textbf{greifen}},
  pluperfect = {hatten + \textbf{gegriffen}},
  prepositions = {---},
  separable = {---}
}

% Migrated from verbs/g.tex:11
\VerbTable{
  name = {grillen},
  english = {to grill / barbecue},
  type = {regular, + Akkusativ},
  auxiliary = {haben},
  style = {acc},
  present = {grille,grillst,grillt,grillen,grillt,grillen},
  preterite = {grillte,grilltest,grillte,grillten,grilltet,grillten},
  subjunctive = {},
  imperative = {grille (du)!,grillen wir!,grillt (ihr)!,grillen Sie!},
  perfect = {haben + \textbf{gegrillt}},
  future = {werden + \textbf{grillen}},
  pluperfect = {hatten + \textbf{gegrillt}},
  prepositions = {& \textbf{grillen auf} + Dat & (to grill on)\\},
  separable = {& an$|$grillen & (to lightly grill / start grilling)\\}
}

% Migrated from verbs/b1/g.tex:9
\VerbTable{
  name = {grillieren},
  english = {to grill},
  type = {regular},
  auxiliary = {haben},
  style = {default},
  present = {grilliere,grillierst,grilliert,grillieren,grilliert,grillieren},
  preterite = {grillierte,grilliertest,grillierte,grillierten,grilliertet,grillierten},
  subjunctive = {},
  imperative = {grillier (du)!,grillieren wir!,grilliert (ihr)!,grillieren Sie!},
  perfect = {haben + \textbf{grilliert}},
  future = {werden + \textbf{grillieren}},
  pluperfect = {hatten + \textbf{grilliert}},
  prepositions = {---},
  separable = {---}
}

% Migrated from verbs/b1/g.tex:10
\VerbTable{
  name = {gründen},
  english = {to establish},
  type = {regular},
  auxiliary = {haben},
  style = {default},
  present = {gründe,gründest,gründet,gründen,gründet,gründen},
  preterite = {gründete,gründetest,gründete,gründeten,gründetet,gründeten},
  subjunctive = {},
  imperative = {gründe (du)!,gründen wir!,gründet (ihr)!,gründen Sie!},
  perfect = {haben + \textbf{gegründet}},
  future = {werden + \textbf{gründen}},
  pluperfect = {hatten + \textbf{gegründet}},
  prepositions = {---},
  separable = {---}
}

% Migrated from verbs/b1/g.tex:11
\VerbTable{
  name = {grüßen},
  english = {to greet},
  type = {regular},
  auxiliary = {haben},
  style = {default},
  present = {grüße,grüßt,grüßt,grüßen,grüßt,grüßen},
  preterite = {grüßte,grüßtest,grüßte,grüßten,grüßtet,grüßten},
  subjunctive = {},
  imperative = {grüß (du)!,grüßen wir!,grüßt (ihr)!,grüßen Sie!},
  perfect = {haben + \textbf{gegrüßt}},
  future = {werden + \textbf{grüßen}},
  pluperfect = {hatten + \textbf{gegrüßt}},
  prepositions = {---},
  separable = {---}
}

% Migrated from verbs/b1/g.tex:12
\VerbTable{
  name = {gucken},
  english = {to look},
  type = {regular},
  auxiliary = {haben},
  style = {default},
  present = {gucke,guckst,guckt,gucken,guckt,gucken},
  preterite = {guckte,gucktest,guckte,guckten,gucktet,guckten},
  subjunctive = {},
  imperative = {guck (du)!,gucken wir!,guckt (ihr)!,gucken Sie!},
  perfect = {haben + \textbf{geguckt}},
  future = {werden + \textbf{gucken}},
  pluperfect = {hatten + \textbf{geguckt}},
  prepositions = {---},
  separable = {---}
}

% Migrated from verbs/g.tex:12
\VerbTable{
  name = {gurgeln},
  english = {to gargle},
  type = {regular},
  auxiliary = {haben},
  style = {default},
  present = {gurgle,gurgelst,gurgelt,gurgeln,gurgelt,gurgeln},
  preterite = {gurgelte,gurgeltest,gurgelte,gurgelten,gurgeltet,gurgelten},
  subjunctive = {},
  imperative = {gurgle (du)!,gurgeln wir!,gurgelt (ihr)!,gurgeln Sie!},
  perfect = {haben + \textbf{gegurgelt}},
  future = {werden + \textbf{gurgeln}},
  pluperfect = {hatten + \textbf{gegurgelt}},
  prepositions = {& \textbf{gurgeln mit} + Dat & (to gargle with)\\},
  separable = {---}
}

\section{H}

% Migrated from verbs/b1/h.tex:3
\VerbTable{
  name = {hängen},
  english = {to hang},
  type = {strong / regular (meaning-dependent)},
  auxiliary = {haben},
  style = {default},
  present = {hänge,hängst,hängt,hängen,hängt,hängen},
  preterite = {hing / hängte,hingst / hängtest,hing / hängte,hingen / hängten,hingt / hängtet,hingen / hängten},
  subjunctive = {},
  imperative = {häng (du)!,hängen wir!,hängt (ihr)!,hängen Sie!},
  perfect = {haben + \textbf{gehangen / gehängt}},
  future = {werden + \textbf{hängen}},
  pluperfect = {hatten + \textbf{gehangen / gehängt}},
  prepositions = {---},
  separable = {---}
}

% Migrated from verbs/b1/h.tex:1
\VerbTable{
  name = {hageln},
  english = {to hail},
  type = {regular},
  auxiliary = {haben},
  style = {default},
  present = {hagle,hagelst,hagelt,hageln,hagelt,hageln},
  preterite = {hagelte,hageltest,hagelte,hagelten,hageltet,hagelten},
  subjunctive = {},
  imperative = {---},
  perfect = {haben + \textbf{gehagelt}},
  future = {werden + \textbf{hageln}},
  pluperfect = {hatten + \textbf{gehagelt}},
  prepositions = {---},
  separable = {---}
}

% Migrated from verbs/h.tex:1
\VerbTable{
  name = {halten},
  english = {to hold, stop},
  type = {irregular},
  auxiliary = {haben},
  style = {default},
  present = {halte,hältst,hält,halten,haltet,halten},
  preterite = {hielt,hieltst,hielt,hielten,hieltet,hielten},
  subjunctive = {},
  imperative = {halt(e) (du)!,halten wir!,haltet (ihr)!,halten Sie!},
  perfect = {haben + \textbf{gehalten}},
  future = {werden + \textbf{halten}},
  pluperfect = {hatten + \textbf{gehalten}},
  prepositions = {---},
  separable = {---}
}

% Migrated from verbs/b1/h.tex:2
\VerbTable{
  name = {handeln},
  english = {to act / deal},
  type = {regular},
  auxiliary = {haben},
  style = {default},
  present = {handle,handelst,handelt,handeln,handelt,handeln},
  preterite = {handelte,handeltest,handelte,handelten,handeltet,handelten},
  subjunctive = {},
  imperative = {handle (du)!,handeln wir!,handelt (ihr)!,handeln Sie!},
  perfect = {haben + \textbf{gehandelt}},
  future = {werden + \textbf{handeln}},
  pluperfect = {hatten + \textbf{gehandelt}},
  prepositions = {---},
  separable = {---}
}

% Migrated from verbs/b1/h.tex:4
\VerbTable{
  name = {hassen},
  english = {to hate},
  type = {regular},
  auxiliary = {haben},
  style = {default},
  present = {hasse,hasst,hasst,hassen,hasst,hassen},
  preterite = {hasste,hasstest,hasste,hassten,hasstet,hassten},
  subjunctive = {},
  imperative = {hass (du)!,hassen wir!,hasst (ihr)!,hassen Sie!},
  perfect = {haben + \textbf{gehasst}},
  future = {werden + \textbf{hassen}},
  pluperfect = {hatten + \textbf{gehasst}},
  prepositions = {---},
  separable = {---}
}

% Migrated from verbs/h.tex:2
\VerbTable{
  name = {heben},
  english = {to lift / raise},
  type = {irregular},
  auxiliary = {haben},
  style = {default},
  present = {hebe,hebst,hebt,heben,hebt,heben},
  preterite = {hob,hobst,hob,hoben,hobt,hoben},
  subjunctive = {},
  imperative = {heb(e) (du)!,heben wir!,hebt (ihr)!,heben Sie!},
  perfect = {haben + \textbf{gehoben}},
  future = {werden + \textbf{heben}},
  pluperfect = {hatten + \textbf{gehoben}},
  prepositions = {
    & \textbf{heben auf} + Akk & (to pick up / keep / store)\\
    & \textbf{heben von} + Dat & (to lift from)\\
    & \textbf{heben aus} + Dat & (to lift out of)\\
  },
  separable = {
    & \textbf{ab$|$heben} & (to lift off / withdraw money)\\
    & \textbf{auf$|$heben} & (to pick up / cancel / keep)\\
    & \textbf{an$|$heben} & (to begin / raise slightly)\\
    & \textbf{heraus$|$heben} & (to lift out)\\
  }
}

% Migrated from verbs/h.tex:3
\VerbTable{
  name = {heiraten},
  english = {to marry},
  type = {regular},
  auxiliary = {haben},
  style = {default},
  present = {heirate,heiratest,heiratet,heiraten,heiratet,heiraten},
  preterite = {heiratete,heiratetest,heiratete,heirateten,heiratetet,heirateten},
  subjunctive = {},
  imperative = {heirate (du)!,heiraten wir!,heiratet (ihr)!,heiraten Sie!},
  perfect = {haben + \textbf{geheiratet}},
  future = {werden + \textbf{heiraten}},
  pluperfect = {hatten + \textbf{geheiratet}},
  prepositions = {
    & \textbf{heiraten} + Akk & (to marry someone)\\
    & \textbf{heiraten mit} + Dat & (to marry, formal)\\
  },
  separable = {---}
}

% Migrated from verbs/h.tex:4
\VerbTable{
  name = {heißen},
  english = {to be called},
  type = {irregular},
  auxiliary = {haben},
  style = {default},
  present = {heiße,heißt,heißt,heißen,heißt,heißen},
  preterite = {hieß,hießest,hieß,hießen,hießt,hießen},
  subjunctive = {},
  imperative = {heiß(e) (du)!,heißen wir!,heißt (ihr)!,heißen Sie!},
  perfect = {haben + \textbf{geheißen}},
  future = {werden + \textbf{heißen}},
  pluperfect = {hatten + \textbf{geheißen}},
  prepositions = {---},
  separable = {---}
}

% Migrated from verbs/b1/h.tex:5
\VerbTable{
  name = {heizen},
  english = {to heat},
  type = {regular},
  auxiliary = {haben},
  style = {default},
  present = {heize,heizt,heizt,heizen,heizt,heizen},
  preterite = {heizte,heiztest,heizte,heizten,heiztet,heizten},
  subjunctive = {},
  imperative = {heiz (du)!,heizen wir!,heizt (ihr)!,heizen Sie!},
  perfect = {haben + \textbf{geheizt}},
  future = {werden + \textbf{heizen}},
  pluperfect = {hatten + \textbf{geheizt}},
  prepositions = {---},
  separable = {---}
}

% Migrated from verbs/h.tex:5
\VerbTable{
  name = {helfen},
  english = {to help},
  type = {irregular, + Dativ},
  auxiliary = {haben},
  style = {dat},
  present = {helfe,hilfst,hilft,helfen,helft,helfen},
  preterite = {half,halfst,half,halfen,halft,halfen},
  subjunctive = {},
  imperative = {hilf (du)!,helfen wir!,helft (ihr)!,helfen Sie!},
  perfect = {haben + \textbf{geholfen}},
  future = {werden + \textbf{helfen}},
  pluperfect = {hatten + \textbf{geholfen}},
  prepositions = {---},
  separable = {---}
}

% Migrated from verbs/b1/h.tex:6
\VerbTable{
  name = {herausfinden},
  english = {to find out},
  type = {separable, strong / irregular},
  auxiliary = {haben},
  style = {default},
  present = {finde heraus,findest heraus,findet heraus,finden heraus,findet heraus,finden heraus},
  preterite = {fand heraus,fandest heraus,fand heraus,fanden heraus,fandet heraus,fanden heraus},
  subjunctive = {},
  imperative = {finde (du) heraus!,finden wir heraus!,findet (ihr) heraus!,finden Sie heraus!},
  perfect = {haben + \textbf{herausgefunden}},
  future = {werden + \textbf{herausfinden}},
  pluperfect = {hatten + \textbf{herausgefunden}},
  prepositions = {---},
  separable = {---}
}

% Migrated from verbs/b1/h.tex:7
\VerbTable{
  name = {herstellen},
  english = {to produce / manufacture},
  type = {separable, regular},
  auxiliary = {haben},
  style = {default},
  present = {stelle her,stellst her,stellt her,stellen her,stellt her,stellen her},
  preterite = {stellte her,stelltest her,stellte her,stellten her,stelltet her,stellten her},
  subjunctive = {},
  imperative = {stell (du) her!,stellen wir her!,stellt (ihr) her!,stellen Sie her!},
  perfect = {haben + \textbf{hergestellt}},
  future = {werden + \textbf{herstellen}},
  pluperfect = {hatten + \textbf{hergestellt}},
  prepositions = {---},
  separable = {---}
}

% Migrated from verbs/b1/h.tex:8
\VerbTable{
  name = {herunterfahren},
  english = {to shut down},
  type = {separable, strong / irregular},
  auxiliary = {haben},
  style = {default},
  present = {fahre herunter,fährst herunter,fährt herunter,fahren herunter,fahrt herunter,fahren herunter},
  preterite = {fuhr herunter,fuhrst herunter,fuhr herunter,fuhren herunter,fuhrt herunter,fuhren herunter},
  subjunctive = {},
  imperative = {fahr (du) herunter!,fahren wir herunter!,fahrt (ihr) herunter!,fahren Sie herunter!},
  perfect = {haben + \textbf{heruntergefahren}},
  future = {werden + \textbf{herunterfahren}},
  pluperfect = {hatten + \textbf{heruntergefahren}},
  prepositions = {---},
  separable = {---}
}

% Migrated from verbs/b1/h.tex:9
\VerbTable{
  name = {herunterladen},
  english = {to download},
  type = {separable, strong / irregular},
  auxiliary = {haben},
  style = {default},
  present = {lade herunter,lädst herunter,lädt herunter,laden herunter,ladet herunter,laden herunter},
  preterite = {lud herunter,ludest herunter,lud herunter,luden herunter,ludet herunter,luden herunter},
  subjunctive = {},
  imperative = {lad (du) herunter!,laden wir herunter!,ladet (ihr) herunter!,laden Sie herunter!},
  perfect = {haben + \textbf{heruntergeladen}},
  future = {werden + \textbf{herunterladen}},
  pluperfect = {hatten + \textbf{heruntergeladen}},
  prepositions = {---},
  separable = {---}
}

% Migrated from verbs/b1/h.tex:10
\VerbTable{
  name = {hinterlassen},
  english = {to leave behind},
  type = {strong / irregular},
  auxiliary = {haben},
  style = {default},
  present = {hinterlasse,hinterlässt,hinterlässt,hinterlassen,hinterlasst,hinterlassen},
  preterite = {hinterließ,hinterließt,hinterließ,hinterließen,hinterließt,hinterließen},
  subjunctive = {},
  imperative = {hinterlass (du)!,hinterlassen wir!,hinterlasst (ihr)!,hinterlassen Sie!},
  perfect = {haben + \textbf{hinterlassen}},
  future = {werden + \textbf{hinterlassen}},
  pluperfect = {hatten + \textbf{hinterlassen}},
  prepositions = {---},
  separable = {---}
}

% Migrated from verbs/b1/h.tex:11
\VerbTable{
  name = {hinweisen},
  english = {to point out},
  type = {separable, strong / irregular},
  auxiliary = {haben},
  style = {default},
  present = {weise hin,weist hin,weist hin,weisen hin,weist hin,weisen hin},
  preterite = {wies hin,wiest hin,wies hin,wiesen hin,wiest hin,wiesen hin},
  subjunctive = {},
  imperative = {weis (du) hin!,weisen wir hin!,weist (ihr) hin!,weisen Sie hin!},
  perfect = {haben + \textbf{hingewiesen}},
  future = {werden + \textbf{hinweisen}},
  pluperfect = {hatten + \textbf{hingewiesen}},
  prepositions = {---},
  separable = {---}
}

% Migrated from verbs/b1/h.tex:12
\VerbTable{
  name = {hochladen},
  english = {to upload},
  type = {separable, strong / irregular},
  auxiliary = {haben},
  style = {default},
  present = {lade hoch,lädst hoch,lädt hoch,laden hoch,ladet hoch,laden hoch},
  preterite = {lud hoch,ludest hoch,lud hoch,luden hoch,ludet hoch,luden hoch},
  subjunctive = {},
  imperative = {lad (du) hoch!,laden wir hoch!,ladet (ihr) hoch!,laden Sie hoch!},
  perfect = {haben + \textbf{hochgeladen}},
  future = {werden + \textbf{hochladen}},
  pluperfect = {hatten + \textbf{hochgeladen}},
  prepositions = {---},
  separable = {---}
}

% Migrated from verbs/h.tex:8
\VerbTable{
  name = {hören},
  english = {to hear},
  type = {regular, + Akkusativ},
  auxiliary = {haben},
  style = {acc},
  present = {höre,hörst,hört,hören,hört,hören},
  preterite = {hörte,hörtest,hörte,hörten,hörtet,hörten},
  subjunctive = {},
  imperative = {hör(e) (du)!,hören wir!,hört (ihr)!,hören Sie!},
  perfect = {haben + \textbf{gehört}},
  future = {werden + \textbf{hören}},
  pluperfect = {hatten + \textbf{gehört}},
  prepositions = {
    & \textbf{hören auf} + Akk & (to listen to / obey / pay attention to)\\
    & \textbf{hören von} + Dat & (to hear about / receive news)\\
  },
  separable = {
    & \textbf{ab$|$hören} & (to listen to / examine)\\
    & \textbf{an$|$hören} & (to hear completely)\\
    & \textbf{auf$|$hören} & (to stop / quit)\\
    & \textbf{mit$|$hören} & (to listen in / eavesdrop)\\
    & \textbf{nach$|$hören} & (to check / listen again)\\
    & \textbf{zu$|$hören} & (to listen carefully)\\
  }
}

% Migrated from verbs/h.tex:7
\VerbTable{
  name = {hoffen},
  english = {to hope},
  type = {regular, + Akkusativ},
  auxiliary = {haben},
  style = {acc},
  present = {hoffe,hoffst,hofft,hoffen,hofft,hoffen},
  preterite = {hoffte,hofftest,hoffte,hofften,hofftet,hofften},
  subjunctive = {},
  imperative = {hoff(e) (du)!,hoffen wir!,hofft (ihr)!,hoffen Sie!},
  perfect = {haben + \textbf{gehofft}},
  future = {werden + \textbf{hoffen}},
  pluperfect = {hatten + \textbf{gehofft}},
  prepositions = {---},
  separable = {---}
}

% Migrated from verbs/h.tex:6
\VerbTable{
  name = {holen},
  english = {to fetch},
  type = {regular, + Akkusativ},
  auxiliary = {haben},
  style = {acc},
  present = {hole,holst,holt,holen,holt,holen},
  preterite = {gholte,holtest,holte,holten,holtet,holten},
  subjunctive = {},
  imperative = {hol(e) (du)!,holen wir!,holt (ihr)!,holen Sie!},
  perfect = {haben + \textbf{geholt}},
  future = {werden + \textbf{holen}},
  pluperfect = {hatten + \textbf{geholt}},
  prepositions = {
    & \textbf{holen für} + Akk & (to fetch something for someone)\\
    & \textbf{holen aus} + Dat & (to take out of / fetch from inside)\\
    &\textbf{holen nach} + Dat & (to fetch to a location)\\
    &\textbf{abholen von} + Akk & (to pick up from)
  },
  separable = {
    & \textbf{ab$|$holen} & (to pick up / collect)\\
    & \textbf{auf$|$holen} & (to catch up)\\
    & \textbf{ein$|$holen} & (to catch up / retrieve)\\
    & \textbf{nach$|$holen} & (to make up / catch up later)
  }
}

% Migrated from verbs/b1/h.tex:13
\VerbTable{
  name = {hupen},
  english = {to honk},
  type = {regular},
  auxiliary = {haben},
  style = {default},
  present = {hupe,hupst,hupt,hupen,hupt,hupen},
  preterite = {hupte,huptest,hupte,hupten,huptet,hupten},
  subjunctive = {},
  imperative = {hup (du)!,hupen wir!,hupt (ihr)!,hupen Sie!},
  perfect = {haben + \textbf{gehupt}},
  future = {werden + \textbf{hupen}},
  pluperfect = {hatten + \textbf{gehupt}},
  prepositions = {---},
  separable = {---}
}

% Migrated from verbs/h.tex:9
\VerbTable{
  name = {husten},
  english = {to cough},
  type = {regular},
  auxiliary = {haben},
  style = {default},
  present = {huste,hustest,hustet,husten,hustet,husten},
  preterite = {hustete,hustetest,hustete,husteten,hustetet,husteten},
  subjunctive = {},
  imperative = {huste (du)!,husten wir!,hustet (ihr)!,husten Sie!},
  perfect = {haben + \textbf{gehustet}},
  future = {werden + \textbf{husten}},
  pluperfect = {hatten + \textbf{gehustet}},
  prepositions = {
    & \textbf{husten in} + Akk & (to cough into)\\
    & \textbf{husten vor} + Dat & (to cough from / because of)\\
  },
  separable = {& \textbf{aus$|$husten} & (to cough out)\\}
}

\section{I}

% Migrated from verbs/i.tex:1
\VerbTable{
  name = {impfen},
  english = {to vaccinate / inoculate},
  type = {regular},
  auxiliary = {haben},
  style = {default},
  present = {impfe,impfst,impft,impfen,impft,impfen},
  preterite = {impfte,impftest,impfte,impften,impftet,impften},
  subjunctive = {},
  imperative = {impfe (du)!,impfen wir!,impft (ihr)!,impfen Sie!},
  perfect = {haben + \textbf{geimpft}},
  future = {werden + \textbf{impfen}},
  pluperfect = {hatten + \textbf{geimpft}},
  prepositions = {
    & \textbf{impfen gegen} + Akk & (to vaccinate against)\\
    & \textbf{impfen für} + Akk & (to vaccinate for / on behalf of)\\
  },
  separable = {---}
}

% Migrated from verbs/i.tex:2
\VerbTable{
  name = {informieren},
  english = {to inform},
  type = {regular},
  auxiliary = {haben},
  style = {default},
  present = {informiere,informierst,informiert,informieren,informiert,informieren},
  preterite = {informierte,informiertest,informierte,informierten,informiertet,informierten},
  subjunctive = {},
  imperative = {informiere (du)!,informieren wir!,informiert (ihr)!,informieren Sie!},
  perfect = {haben + \textbf{informiert}},
  future = {werden + \textbf{informieren}},
  pluperfect = {hatten + \textbf{informiert}},
  prepositions = {
    & \textbf{informieren über} + Akk & (to inform about)\\
    & \textbf{informieren von} + Dat & (to inform of / from)\\
    & \textbf{informieren bei} + Dat & (to inform at / with)\\
  },
  separable = {& \textbf{ein$|$informieren} & (to inform in / register information)\\}
}

% Migrated from verbs/i.tex:3
\VerbTable{
  name = {installieren},
  english = {to install},
  type = {regular},
  auxiliary = {haben},
  style = {default},
  present = {installiere,installierst,installiert,installieren,installiert,installieren},
  preterite = {installierte,installiertest,installierte,installierten,installiertet,installierten},
  subjunctive = {},
  imperative = {installiere (du)!,installieren wir!,installiert (ihr)!,installieren Sie!},
  perfect = {haben + \textbf{installiert}},
  future = {werden + \textbf{installieren}},
  pluperfect = {hatten + \textbf{installiert}},
  prepositions = {
    & \textbf{installieren auf} + Dat & (to install on)\\
    & \textbf{installieren in} + Akk & (to install in)\\
  },
  separable = {& \textbf{ein$|$installieren} & (to install into)\\}
}

% Migrated from verbs/b1/i.tex:1
\VerbTable{
  name = {integrieren},
  english = {to integrate},
  type = {regular},
  auxiliary = {haben},
  style = {default},
  present = {integriere,integrierst,integriert,integrieren,integriert,integrieren},
  preterite = {integrierte,integriertest,integrierte,integrierten,integriertet,integrierten},
  subjunctive = {},
  imperative = {integrier (du)!,integrieren wir!,integriert (ihr)!,integrieren Sie!},
  perfect = {haben + \textbf{integriert}},
  future = {werden + \textbf{integrieren}},
  pluperfect = {hatten + \textbf{integriert}},
  prepositions = {---},
  separable = {---}
}

\section{K}

% Migrated from verbs/b1/k.tex:1
\VerbTable{
  name = {kämpfen},
  english = {to fight},
  type = {regular},
  auxiliary = {haben},
  style = {default},
  present = {kämpfe,kämpfst,kämpft,kämpfen,kämpft,kämpfen},
  preterite = {kämpfte,kämpftest,kämpfte,kämpften,kämpftet,kämpften},
  subjunctive = {},
  imperative = {kämpf (du)!,kämpfen wir!,kämpft (ihr)!,kämpfen Sie!},
  perfect = {haben + \textbf{gekämpft}},
  future = {werden + \textbf{kämpfen}},
  pluperfect = {hatten + \textbf{gekämpft}},
  prepositions = {---},
  separable = {---}
}

% Migrated from verbs/k.tex:1
\VerbTable{
  name = {kaufen},
  english = {to buy},
  type = {regular, + Akkusativ},
  auxiliary = {haben},
  style = {acc},
  present = {kaufe,kaufst,kauft,kaufen,kauft,kaufen},
  preterite = {kaufte,kauftest,kaufte,kauften,kauftet,kauften},
  subjunctive = {},
  imperative = {kauf(e) (du)!,kaufen wir!,kauft (ihr)!,kaufen Sie!},
  perfect = {haben + \textbf{gekauft}},
  future = {werden + \textbf{kaufen}},
  pluperfect = {hatten + \textbf{gekauft}},
  prepositions = {---},
  separable = {---}
}

% Migrated from verbs/k.tex:2
\VerbTable{
  name = {kennen},
  english = {to know (person/place)},
  type = {regular, + Akkusativ},
  auxiliary = {haben},
  style = {acc},
  present = {kenne,kennst,kennt,kennen,kennt,kennen},
  preterite = {kannte,kanntest,kannte,kannten,kanntet,kannten},
  subjunctive = {},
  imperative = {kenn(e) (du)!,kennen wir!,kennt (ihr)!,kennen Sie!},
  perfect = {haben + \textbf{gekannt}},
  future = {werden + \textbf{kennen}},
  pluperfect = {hatten + \textbf{gekannt}},
  prepositions = {---},
  separable = {---}
}

% Migrated from verbs/b1/k.tex:2
\VerbTable{
  name = {kennenlernen},
  english = {to get to know each other},
  type = {separable, regular},
  auxiliary = {haben},
  style = {default},
  present = {lerne kennen,lernst kennen,lernt kennen,lernen kennen,lernt kennen,lernen kennen},
  preterite = {lernte kennen,lerntest kennen,lernte kennen,lernten kennen,lerntet kennen,lernten kennen},
  subjunctive = {},
  imperative = {lern (du) kennen!,lernen wir kennen!,lernt (ihr) kennen!,lernen Sie kennen!},
  perfect = {haben + \textbf{kennengelernt}},
  future = {werden + \textbf{kennenlernen}},
  pluperfect = {hatten + \textbf{kennengelernt}},
  prepositions = {---},
  separable = {---}
}

% Migrated from verbs/k.tex:3
\VerbTable{
  name = {klären},
  english = {to clarify},
  type = {regular, + Akkusativ},
  auxiliary = {haben},
  style = {acc},
  present = {kläre,klärst,klärt,klären,klärt,klären},
  preterite = {klärte,klärtest,klärte,klärten,klärtet,klärten},
  subjunctive = {},
  imperative = {klär(e) (du)!,klären wir!,klärt (ihr)!,klären Sie!},
  perfect = {haben + \textbf{geklärt}},
  future = {werden + \textbf{klären}},
  pluperfect = {hatten + \textbf{geklärt}},
  prepositions = {---},
  separable = {---}
}

% Migrated from verbs/b1/k.tex:3
\VerbTable{
  name = {klagen},
  english = {to complain},
  type = {regular},
  auxiliary = {haben},
  style = {default},
  present = {klage,klagst,klagt,klagen,klagt,klagen},
  preterite = {klagte,klagtest,klagte,klagten,klagtet,klagten},
  subjunctive = {},
  imperative = {klag (du)!,klagen wir!,klagt (ihr)!,klagen Sie!},
  perfect = {haben + \textbf{geklagt}},
  future = {werden + \textbf{klagen}},
  pluperfect = {hatten + \textbf{geklagt}},
  prepositions = {---},
  separable = {---}
}

% Migrated from verbs/b1/k.tex:4
\VerbTable{
  name = {klappen},
  english = {to work / succeed},
  type = {regular},
  auxiliary = {haben},
  style = {default},
  present = {klappe,klappst,klappt,klappen,klappt,klappen},
  preterite = {klappte,klapptest,klappte,klappten,klapptet,klappten},
  subjunctive = {},
  imperative = {klapp (du)!,klappen wir!,klappt (ihr)!,klappen Sie!},
  perfect = {haben + \textbf{geklappt}},
  future = {werden + \textbf{klappen}},
  pluperfect = {hatten + \textbf{geklappt}},
  prepositions = {---},
  separable = {---}
}

% Migrated from verbs/b1/k.tex:5
\VerbTable{
  name = {kleben},
  english = {to stick / glue},
  type = {regular},
  auxiliary = {haben},
  style = {default},
  present = {klebe,klebst,klebt,kleben,klebt,kleben},
  preterite = {klebte,klebtest,klebte,klebten,klebtet,klebten},
  subjunctive = {},
  imperative = {kleb (du)!,kleben wir!,klebt (ihr)!,kleben Sie!},
  perfect = {haben + \textbf{geklebt}},
  future = {werden + \textbf{kleben}},
  pluperfect = {hatten + \textbf{geklebt}},
  prepositions = {---},
  separable = {---}
}

% Migrated from verbs/k.tex:4
\VerbTable{
  name = {kleiden},
  english = {to dress / to clothe},
  type = {regular},
  auxiliary = {haben},
  style = {default},
  present = {kleide,kleidest,kleidet,kleiden,kleidet,kleiden},
  preterite = {kleidete,kleidetest,kleidete,kleideten,kleidetet,kleideten},
  subjunctive = {},
  imperative = {kleide (du)!,kleiden wir!,kleidet (ihr)!,kleiden Sie!},
  perfect = {haben + \textbf{gekleidet}},
  future = {werde + \textbf{kleiden}},
  pluperfect = {hatten + \textbf{gekleidet}},
  prepositions = {
    & \textbf{sich kleiden in} + Akk & (to dress in)\\
    & \textbf{kleiden mit} + Dat & (to clothe with)\\
  },
  separable = {
    & \textbf{an$|$kleiden} & (to dress someone)\\
    & \textbf{ver$|$kleiden} & (to disguise / dress up)\\
  }
}

% Migrated from verbs/k.tex:5
\VerbTable{
  name = {klettern},
  english = {to climb},
  type = {regular},
  auxiliary = {sein / haben},
  style = {default},
  present = {klettere,kletterst,klettert,klettern,klettert,klettern},
  preterite = {kletterte,klettertest,kletterte,kletterten,klettertet,kletterten},
  subjunctive = {},
  imperative = {klettere (du)!,klettern wir!,klettert (ihr)!,klettern Sie!},
  perfect = {sein + \textbf{geklettert}},
  future = {werde + \textbf{klettern}},
  pluperfect = {waren + \textbf{geklettert}},
  prepositions = {
    & \textbf{klettern auf} + Akk & (to climb onto)\\
    & \textbf{klettern über} + Akk & (to climb over)\\
    & \textbf{klettern in} + Akk & (to climb into)\\
  },
  separable = {
    & \textbf{hinauf$|$klettern} & (to climb up)\\
    & \textbf{herunter$|$klettern} & (to climb down)\\
  }
}

% Migrated from verbs/k.tex:6
\VerbTable{
  name = {klicken},
  english = {to click},
  type = {regular},
  auxiliary = {haben},
  style = {default},
  present = {klicke,klickst,klickt,klicken,klickt,klicken},
  preterite = {klickte,klicktest,klickte,klickten,klicktet,klickten},
  subjunctive = {},
  imperative = {klicke (du)!,klicken wir!,klickt (ihr)!,klicken Sie!},
  perfect = {haben + \textbf{geklickt}},
  future = {werden + \textbf{klicken}},
  pluperfect = {hatten + \textbf{geklickt}},
  prepositions = {
    & \textbf{klicken auf} + Akk & (to click on)\\
    & \textbf{klicken in} + Akk & (to click into)\\
  },
  separable = {& an$|$klicken & (to click on / activate)\\}
}

% Migrated from verbs/b1/k.tex:6
\VerbTable{
  name = {klingeln},
  english = {to ring},
  type = {regular},
  auxiliary = {haben},
  style = {default},
  present = {klingle,klingelst,klingelt,klingeln,klingelt,klingeln},
  preterite = {klingelte,klingeltest,klingelte,klingelten,klingeltet,klingelten},
  subjunctive = {},
  imperative = {klingle (du)!,klingeln wir!,klingelt (ihr)!,klingeln Sie!},
  perfect = {haben + \textbf{geklingelt}},
  future = {werden + \textbf{klingeln}},
  pluperfect = {hatten + \textbf{geklingelt}},
  prepositions = {---},
  separable = {---}
}

% Migrated from verbs/b1/k.tex:7
\VerbTable{
  name = {klingen},
  english = {to sound},
  type = {strong / irregular},
  auxiliary = {haben},
  style = {default},
  present = {klinge,klingst,klingt,klingen,klingt,klingen},
  preterite = {klang,klangst,klang,klangen,klangt,klangen},
  subjunctive = {},
  imperative = {kling (du)!,klingen wir!,klingt (ihr)!,klingen Sie!},
  perfect = {haben + \textbf{geklungen}},
  future = {werden + \textbf{klingen}},
  pluperfect = {hatten + \textbf{geklungen}},
  prepositions = {---},
  separable = {---}
}

% Migrated from verbs/b1/k.tex:8
\VerbTable{
  name = {klopfen},
  english = {to knock},
  type = {regular},
  auxiliary = {haben},
  style = {default},
  present = {klopfe,klopfst,klopft,klopfen,klopft,klopfen},
  preterite = {klopfte,klopftest,klopfte,klopften,klopftet,klopften},
  subjunctive = {},
  imperative = {klopf (du)!,klopfen wir!,klopft (ihr)!,klopfen Sie!},
  perfect = {haben + \textbf{geklopft}},
  future = {werden + \textbf{klopfen}},
  pluperfect = {hatten + \textbf{geklopft}},
  prepositions = {---},
  separable = {---}
}

% Migrated from verbs/k.tex:7
\VerbTable{
  name = {knien},
  english = {to kneel},
  type = {regular},
  auxiliary = {haben / sein},
  style = {default},
  present = {knie,kniest,kniet,knien,kniet,knien},
  preterite = {kniete,knietest,kniete,knieten,knietet,knieten},
  subjunctive = {},
  imperative = {knie (du)!,knien wir!,kniet (ihr)!,knien Sie!},
  perfect = {haben / sein + \textbf{gekniet}},
  future = {werden + \textbf{knien}},
  pluperfect = {hatten / waren + \textbf{gekniet}},
  prepositions = {---},
  separable = {---}
}

% Migrated from verbs/k.tex:8
\VerbTable{
  name = {knüpfen},
  english = {to tie / establish},
  type = {regular},
  auxiliary = {haben},
  style = {default},
  present = {knüpfe,knüpfst,knüpft,knüpfen,knüpft,knüpfen},
  preterite = {knüpfte,knüpftest,knüpfte,knüpften,knüpftet,knüpften},
  subjunctive = {},
  imperative = {knüpfe (du)!,knüpfen wir!,knüpft (ihr)!,knüpfen Sie!},
  perfect = {haben + \textbf{geknüpft}},
  future = {werden + \textbf{knüpfen}},
  pluperfect = {hatten + \textbf{geknüpft}},
  prepositions = {& \textbf{knüpfen an} + Akk & (to link to)\\},
  separable = {& \textbf{an$|$knüpfen} & (to connect to)\\}
}

% Migrated from verbs/k.tex:9
\VerbTable{
  name = {kochen},
  english = {to cook},
  type = {regular, + Akkusativ},
  auxiliary = {haben},
  style = {acc},
  present = {koche,kochst,kocht,kochen,kocht,kochen},
  preterite = {kochte,kochtest,kochte,kochten,kochtet,kochten},
  subjunctive = {},
  imperative = {koch(e) (du)!,kochen wir!,kocht (ihr)!,kochen Sie!},
  perfect = {haben + \textbf{gekocht}},
  future = {werden + \textbf{kochen}},
  pluperfect = {hatten + \textbf{gekocht}},
  prepositions = {---},
  separable = {---}
}

% Migrated from verbs/k.tex:10
\VerbTable{
  name = {kommen},
  english = {to come},
  type = {irregular},
  auxiliary = {sein},
  style = {default},
  present = {komme,kommst,kommt,kommen,kommt,kommen},
  preterite = {kam,kamst,kam,kamen,kamt,kamen},
  subjunctive = {},
  imperative = {komm(e) (du)!,kommen wir!,kommt (ihr)!,kommen Sie!},
  perfect = {sein + \textbf{gekommen}},
  future = {werden + \textbf{kommen}},
  pluperfect = {waren + \textbf{gekommen}},
  prepositions = {
    & \textbf{kommen aus} + Dat & (to come from)\\
    & \textbf{kommen von} + Dat & (origin / source of)\\
    &\textbf{kommen zu} + Dat & (to approach to)\\
    &\textbf{kommen nach} + Dat & (come to a city / country)\\
    &\textbf{kommen in} + Akk & (to come into)\\
  },
  separable = {
    & \textbf{an$|$kommen} & (to arrive)\\
    & \textbf{mit$|$kommen} & (to come along)\\
    & \textbf{auf$|$kommen} & (to arise / come up)\\
    & \textbf{vor$|$kommen} & (to occur)\\
    & \textbf{nach$|$kommen} & (to follow / comply)\\
    & \textbf{zu$|$kommen} & (to approach)\\
    & \textbf{herein$|$kommen} & (to come in)\\
    & \textbf{hinein$|$kommen} & (to get inside / enter)\\
    & \textbf{dazu$|$kommen} & (to join / be added)\\
    & \textbf{fort$|$kommen} & (to make progress)
  }
}

% Migrated from verbs/k.tex:11
\VerbTable{
  name = {kommunizieren},
  english = {to communicate},
  type = {regular},
  auxiliary = {haben},
  style = {default},
  present = {kommuniziere,kommunizierst,kommuniziert,kommunizieren,kommuniziert,kommunizieren},
  preterite = {kommunizierte,kommuniziertest,kommunizierte,kommunizierten,kommuniziertet,kommunizierten},
  subjunctive = {},
  imperative = {kommuniziere (du)!,kommunizieren wir!,kommuniziert (ihr)!,kommunizieren Sie!},
  perfect = {haben + \textbf{kommuniziert}},
  future = {werden + \textbf{kommunizieren}},
  pluperfect = {hatten + \textbf{kommuniziert}},
  prepositions = {
    & \textbf{kommunizieren mit} + Dat & (to communicate with)\\
    & \textbf{kommunizieren über} + Akk & (to communicate about)\\
  },
  separable = {---}
}

% Migrated from verbs/b1/k.tex:9
\VerbTable{
  name = {konsumieren},
  english = {to consume},
  type = {regular},
  auxiliary = {haben},
  style = {default},
  present = {konsumiere,konsumierst,konsumiert,konsumieren,konsumiert,konsumieren},
  preterite = {konsumierte,konsumiertest,konsumierte,konsumierten,konsumiertet,konsumierten},
  subjunctive = {},
  imperative = {konsumier (du)!,konsumieren wir!,konsumiert (ihr)!,konsumieren Sie!},
  perfect = {haben + \textbf{konsumiert}},
  future = {werden + \textbf{konsumieren}},
  pluperfect = {hatten + \textbf{konsumiert}},
  prepositions = {---},
  separable = {---}
}

% Migrated from verbs/b1/k.tex:10
\VerbTable{
  name = {kontrollieren},
  english = {to control},
  type = {regular},
  auxiliary = {haben},
  style = {default},
  present = {kontrolliere,kontrollierst,kontrolliert,kontrollieren,kontrolliert,kontrollieren},
  preterite = {kontrollierte,kontrolliertest,kontrollierte,kontrollierten,kontrolliertet,kontrollierten},
  subjunctive = {},
  imperative = {kontrollier (du)!,kontrollieren wir!,kontrolliert (ihr)!,kontrollieren Sie!},
  perfect = {haben + \textbf{kontrolliert}},
  future = {werden + \textbf{kontrollieren}},
  pluperfect = {hatten + \textbf{kontrolliert}},
  prepositions = {---},
  separable = {---}
}

% Migrated from verbs/b1/k.tex:11
\VerbTable{
  name = {korrigieren},
  english = {to correct},
  type = {regular},
  auxiliary = {haben},
  style = {default},
  present = {korrigiere,korrigierst,korrigiert,korrigieren,korrigiert,korrigieren},
  preterite = {korrigierte,korrigiertest,korrigierte,korrigierten,korrigiertet,korrigierten},
  subjunctive = {},
  imperative = {korrigier (du)!,korrigieren wir!,korrigiert (ihr)!,korrigieren Sie!},
  perfect = {haben + \textbf{korrigiert}},
  future = {werden + \textbf{korrigieren}},
  pluperfect = {hatten + \textbf{korrigiert}},
  prepositions = {---},
  separable = {---}
}

% Migrated from verbs/k.tex:12
\VerbTable{
  name = {kosten},
  english = {to cost},
  type = {regular, + Akkusativ},
  auxiliary = {haben},
  style = {acc},
  present = {-,-,kostet,-,-,kosten},
  preterite = {-,-,kostete,-,-,kosteten},
  subjunctive = {},
  imperative = {---},
  perfect = {haben + \textbf{gekostet}},
  future = {werden + \textbf{kosten}},
  pluperfect = {hatten + \textbf{gekostet}},
  prepositions = {---},
  separable = {---}
}

% Migrated from verbs/b1/k.tex:12
\VerbTable{
  name = {kriegen},
  english = {to get},
  type = {regular},
  auxiliary = {haben},
  style = {default},
  present = {kriege,kriegst,kriegt,kriegen,kriegt,kriegen},
  preterite = {kriegte,kriegtest,kriegte,kriegten,kriegtet,kriegten},
  subjunctive = {},
  imperative = {krieg (du)!,kriegen wir!,kriegt (ihr)!,kriegen Sie!},
  perfect = {haben + \textbf{gekriegt}},
  future = {werden + \textbf{kriegen}},
  pluperfect = {hatten + \textbf{gekriegt}},
  prepositions = {---},
  separable = {---}
}

% Migrated from verbs/b1/k.tex:13
\VerbTable{
  name = {kritisieren},
  english = {to criticize},
  type = {regular},
  auxiliary = {haben},
  style = {default},
  present = {kritisiere,kritisierst,kritisiert,kritisieren,kritisiert,kritisieren},
  preterite = {kritisierte,kritisiertest,kritisierte,kritisierten,kritisiertet,kritisierten},
  subjunctive = {},
  imperative = {kritisier (du)!,kritisieren wir!,kritisiert (ihr)!,kritisieren Sie!},
  perfect = {haben + \textbf{kritisiert}},
  future = {werden + \textbf{kritisieren}},
  pluperfect = {hatten + \textbf{kritisiert}},
  prepositions = {---},
  separable = {---}
}

% Migrated from verbs/k.tex:13
\VerbTable{
  name = {kucken},
  english = {to look},
  type = {regular},
  auxiliary = {haben},
  style = {default},
  present = {kucke,kuckst,kuckt,kucken,kuckt,kucken},
  preterite = {kuckte,kucktest,kuckte,kuckten,kucktet,kuckten},
  subjunctive = {},
  imperative = {kuck(e) du!,kucken wir!,kuckt (ihr)!,kucken Sie!},
  perfect = {haben + \textbf{gekuckt}},
  future = {werden + \textbf{kucken}},
  pluperfect = {hatten + \textbf{gekuckt}},
  prepositions = {---},
  separable = {---}
}

% Migrated from verbs/k.tex:14
\VerbTable{
  name = {kündigen},
  english = {to quit / to terminate},
  type = {regular},
  auxiliary = {haben},
  style = {default},
  present = {kündige,kündigst,kündigt,kündigen,kündigt,kündigen},
  preterite = {kündigte,kündigtest,kündigte,kündigten,kündigtet,kündigten},
  subjunctive = {},
  imperative = {kündige (du)!,kündigen wir!,kündigt (ihr)!,kündigen Sie!},
  perfect = {haben + \textbf{gekündigt}},
  future = {werde + \textbf{kündigen}},
  pluperfect = {hatten + \textbf{gekündigt}},
  prepositions = {
    & \textbf{kündigen bei} + Dat & (to quit at / resign from)\\
    & \textbf{kündigen wegen} + Gen & (to terminate because of)\\
  },
  separable = {
    & \textbf{an$|$kündigen} & (to announce)\\
    & \textbf{auf$|$kündigen} & (to terminate / cancel)\\
  }
}

% Migrated from verbs/b1/k.tex:14
\VerbTable{
  name = {küssen},
  english = {to kiss},
  type = {regular},
  auxiliary = {haben},
  style = {default},
  present = {küsse,küsst,küsst,küssen,küsst,küssen},
  preterite = {küsste,küsstest,küsste,küssten,küsstet,küssten},
  subjunctive = {},
  imperative = {küss (du)!,küssen wir!,küsst (ihr)!,küssen Sie!},
  perfect = {haben + \textbf{geküsst}},
  future = {werden + \textbf{küssen}},
  pluperfect = {hatten + \textbf{geküsst}},
  prepositions = {---},
  separable = {---}
}

\section{L}

% Migrated from verbs/l.tex:1
\VerbTable{
  name = {lachen},
  english = {to laugh},
  type = {regular},
  auxiliary = {haben},
  style = {default},
  present = {lache,lachst,lacht,lachen,lacht,lachen},
  preterite = {lachte,lachtest,lachte,lachten,lachtet,lachten},
  subjunctive = {},
  imperative = {lach(e) (du)!,lachen wir!,lacht (ihr)!,lachen Sie!},
  perfect = {haben + \textbf{gelacht}},
  future = {werden + \textbf{lachen}},
  pluperfect = {hatten + \textbf{gelacht}},
  prepositions = {& \textbf{lachen über} + Akk & (to laugh about)\\},
  separable = {& aus$|$lachen & (to laugh at / mock)\\}
}

% Migrated from verbs/l.tex:3
\VerbTable{
  name = {laden},
  english = {to load / invite},
  type = {irregular},
  auxiliary = {haben},
  style = {acc},
  present = {lade,lädst,lädt,laden,ladet,laden},
  preterite = {lud,ludst,lud,luden,ludet,luden},
  subjunctive = {},
  imperative = {lad(e) (du)!,laden wir!,ladet (ihr)!,laden Sie!},
  perfect = {haben + \textbf{geladen}},
  future = {werden + \textbf{laden}},
  pluperfect = {hatten + \textbf{geladen}},
  prepositions = {
    & \textbf{laden auf} + Akk & (to load onto)\\
    & \textbf{laden in} + Akk & (to load into)\\
    & \textbf{laden zu} + Dat & (to invite to)\\
  },
  separable = {
    & \textbf{ein$|$laden} & (to invite)\\
    & \textbf{auf$|$laden} & (to load up)\\
  }
}

% Migrated from verbs/l.tex:2
\VerbTable{
  name = {lächeln},
  english = {to smile},
  type = {regular},
  auxiliary = {haben},
  style = {default},
  present = {lächele,lächelst,lächelt,lächeln,lächelt,lächeln},
  preterite = {lächelte,lächeltest,lächelte,lächelten,lächeltet,lächelten},
  subjunctive = {},
  imperative = {lächle (du)!,lächeln wir!,lächelt (ihr)!,lächeln Sie!},
  perfect = {haben + \textbf{gelächelt}},
  future = {werden + \textbf{lächeln}},
  pluperfect = {hatten + \textbf{gelächelt}},
  prepositions = {
    & \textbf{lächeln über} + Akk & (to smile about)\\
    & \textbf{lächeln zu} + Dat & (to smile at / toward)\\
  },
  separable = {& an$|$lächeln & (to smile at)\\}
}

% Migrated from verbs/l.tex:4
\VerbTable{
  name = {landen},
  english = {to land},
  type = {regular},
  auxiliary = {sein},
  style = {default},
  present = {lande,landest,landet,landen,landet,landen},
  preterite = {landete,landetest,landete,landeten,landetet,landeten},
  subjunctive = {},
  imperative = {land(e) (du)!,landen wir!,landet (ihr)!,landen Sie!},
  perfect = {sein + \textbf{gelandet}},
  future = {werden + \textbf{landen}},
  pluperfect = {waren + \textbf{gelandet}},
  prepositions = {
    & \textbf{landen auf} + Dat & (to land on – surface / area)\\
    & \textbf{landen in} + Dat & (to land in – country / city)\\
    & \textbf{landen bei} + Dat & (to land near / at)\\
  },
  separable = {
    & \textbf{an$|$landen} & (to arrive / land at a place)\\
    & \textbf{auf$|$landen} & (to land on – surface)\\
    & \textbf{ein$|$landen} & (to land in / enter)\\
    & \textbf{zurück$|$landen} & (to land back / return)\\
  }
}

% Migrated from verbs/l.tex:5
\VerbTable{
  name = {laufen},
  english = {to run},
  type = {irregular},
  auxiliary = {sein},
  style = {default},
  present = {laufe,läufst,läuft,laufen,lauft,laufen},
  preterite = {lief,liefst,lief,liefen,lieft,liefen},
  subjunctive = {},
  imperative = {lauf(e) (du)!,laufen wir!,lauft (ihr)!,laufen Sie!},
  perfect = {sein + \textbf{gelaufen}},
  future = {werden + \textbf{laufen}},
  pluperfect = {waren + \textbf{gelaufen}},
  prepositions = {---},
  separable = {---}
}

% Migrated from verbs/l.tex:6
\VerbTable{
  name = {leben},
  english = {to live},
  type = {regular},
  auxiliary = {haben},
  style = {default},
  present = {lebe,lebst,lebt,leben,lebt,leben},
  preterite = {lebte,lebtest,lebte,lebten,lebtet,lebten},
  subjunctive = {},
  imperative = {leb(e) (du)!,leben wir!,lebt (ihr)!,leben Sie!},
  perfect = {haben + \textbf{gelebt}},
  future = {werden + \textbf{leben}},
  pluperfect = {hatten + \textbf{gelebt}},
  prepositions = {---},
  separable = {---}
}

% Migrated from verbs/l.tex:7
\VerbTable{
  name = {legen},
  english = {to lay / put},
  type = {regular},
  auxiliary = {haben},
  style = {default},
  present = {lege,legst,legt,legen,legt,legen},
  preterite = {legte,legtest,legte,legten,legtet,legten},
  subjunctive = {},
  imperative = {lege (du)!,legen wir!,legt (ihr)!,legen Sie!},
  perfect = {haben + \textbf{gelegt}},
  future = {werden + \textbf{legen}},
  pluperfect = {hatten + \textbf{gelegt}},
  prepositions = {
    & \textbf{legen auf} + Akk & (to lay on)\\
    & \textbf{legen in} + Akk & (to put into)\\
    & \textbf{legen unter} + Akk & (to lay under)\\
  },
  separable = {
    & \textbf{hin$|$legen} & (to lay down)\\
    & \textbf{auf$|$legen} & (to put on / overlay)\\
  }
}

% Migrated from verbs/l.tex:8
\VerbTable{
  name = {lehnen},
  english = {to lean},
  type = {regular},
  auxiliary = {haben / sein reflexive},
  style = {default},
  present = {lehne,lehnst,lehnt,lehnen,lehnt,lehnen},
  preterite = {lehnte,lehntest,lehnte,lehnten,lehntet,lehnten},
  subjunctive = {},
  imperative = {lehne (du)!,lehnen wir!,lehnt (ihr)!,lehnen Sie!},
  perfect = {haben + \textbf{gelehnt}},
  future = {werden + \textbf{lehnen}},
  pluperfect = {hatten + \textbf{gelehnt}},
  prepositions = {
    & \textbf{lehnen an} + Dat & (to lean against)\\
    & \textbf{sich lehnen gegen} + Akk & (to lean against)\\
  },
  separable = {---}
}

% Migrated from verbs/b1/l.tex:2
\VerbTable{
  name = {leiden},
  english = {to suffer},
  type = {strong / irregular},
  auxiliary = {haben},
  style = {default},
  present = {leide,leidest,leidet,leiden,leidet,leiden},
  preterite = {litt,littest,litt,litten,littet,litten},
  subjunctive = {},
  imperative = {leide (du)!,leiden wir!,leidet (ihr)!,leiden Sie!},
  perfect = {haben + \textbf{gelitten}},
  future = {werden + \textbf{leiden}},
  pluperfect = {hatten + \textbf{gelitten}},
  prepositions = {---},
  separable = {---}
}

% Migrated from verbs/b1/l.tex:3
\VerbTable{
  name = {leidtun},
  english = {to be sorry / hurt},
  type = {separable, strong / irregular},
  auxiliary = {haben},
  style = {default},
  present = {tue leid,tust leid,tut leid,tun leid,tut leid,tun leid},
  preterite = {tat leid,tatst leid,tat leid,taten leid,tatet leid,taten leid},
  subjunctive = {},
  imperative = {tu (du) leid!,tun wir leid!,tut (ihr) leid!,tun Sie leid!},
  perfect = {haben + \textbf{leidgetan}},
  future = {werden + \textbf{leidtun}},
  pluperfect = {hatten + \textbf{leidgetan}},
  prepositions = {---},
  separable = {---}
}

% Migrated from verbs/l.tex:9
\VerbTable{
  name = {leihen},
  english = {to borrow},
  type = {irregular, + Akkusativ},
  auxiliary = {haben},
  style = {acc},
  present = {leihe,leihst,leiht,leihen,leiht,leihen},
  preterite = {lieh,liehst,lieh,liehen,lieht,liehen},
  subjunctive = {},
  imperative = {leih(e) (du)!,leihen wir!,leiht (ihr)!,leihen Sie!},
  perfect = {sein + \textbf{geliehen}},
  future = {werden + \textbf{leihen}},
  pluperfect = {waren + \textbf{geliehen}},
  prepositions = {
    & \textbf{leihen von} + Dat & (to borrow from)\\
    & \textbf{leihen an} + Dat & (to lend to)\\
  },
  separable = {& \textbf{aus$|$leihen} & (to lend out (library / shop))\\}
}

% Migrated from verbs/b1/l.tex:4
\VerbTable{
  name = {leisten},
  english = {to perform / afford},
  type = {regular},
  auxiliary = {haben},
  style = {default},
  present = {leiste,leistest,leistet,leisten,leistet,leisten},
  preterite = {leistete,leistetest,leistete,leisteten,leistetet,leisteten},
  subjunctive = {},
  imperative = {leiste (du)!,leisten wir!,leistet (ihr)!,leisten Sie!},
  perfect = {haben + \textbf{geleistet}},
  future = {werden + \textbf{leisten}},
  pluperfect = {hatten + \textbf{geleistet}},
  prepositions = {---},
  separable = {---}
}

% Migrated from verbs/b1/l.tex:5
\VerbTable{
  name = {leiten},
  english = {to lead / manage},
  type = {regular},
  auxiliary = {haben},
  style = {default},
  present = {leite,leitest,leitet,leiten,leitet,leiten},
  preterite = {leitete,leitetest,leitete,leiteten,leitetet,leiteten},
  subjunctive = {},
  imperative = {leite (du)!,leiten wir!,leitet (ihr)!,leiten Sie!},
  perfect = {haben + \textbf{geleitet}},
  future = {werden + \textbf{leiten}},
  pluperfect = {hatten + \textbf{geleitet}},
  prepositions = {---},
  separable = {---}
}

% Migrated from verbs/l.tex:10
\VerbTable{
  name = {lernen},
  english = {to learn},
  type = {regular, + Akkusativ},
  auxiliary = {haben},
  style = {acc},
  present = {lerne,lernst,lernt,lernen,lernt,lernen},
  preterite = {lernte,lerntest,lernte,lernten,lerntet,lernten},
  subjunctive = {},
  imperative = {lern(e) (du)!,lernen wir!,lernt (ihr)!,lernen Sie!},
  perfect = {haben + \textbf{gelernt}},
  future = {werden + \textbf{lernen}},
  pluperfect = {hatten + \textbf{gelernt}},
  prepositions = {---},
  separable = {---}
}

% Migrated from verbs/l.tex:11
\VerbTable{
  name = {lesen},
  english = {to read},
  type = {irregular, + acusative},
  auxiliary = {haben},
  style = {acc},
  present = {lese,liest,liest,lesen,lest,lesen},
  preterite = {las,lastest,las,lasen,last,lasen},
  subjunctive = {},
  imperative = {lies (du)!,lesen wir!,lest (ihr)!,lesen Sie!},
  perfect = {haben + \textbf{gelesen}},
  future = {werden + \textbf{lesen}},
  pluperfect = {hatten + \textbf{gelesen}},
  prepositions = {---},
  separable = {---}
}

% Migrated from verbs/l.tex:12
\VerbTable{
  name = {lieben},
  english = {to love},
  type = {regular, + Akkusativ},
  auxiliary = {haben},
  style = {acc},
  present = {liebe,liebst,liebt,lieben,liebt,lieben},
  preterite = {liebte,liebtest,liebte,liebten,liebtet,liebten},
  subjunctive = {},
  imperative = {lieb(e) (du)!,lieben wir!,liebt (ihr)!,lieben Sie!},
  perfect = {haben + \textbf{geliebt}},
  future = {werden + \textbf{lieben}},
  pluperfect = {hatten + \textbf{geliebt}},
  prepositions = {---},
  separable = {---}
}

% Migrated from verbs/b1/l.tex:6
\VerbTable{
  name = {liefern},
  english = {to deliver},
  type = {regular},
  auxiliary = {haben},
  style = {default},
  present = {liefre,lieferst,liefert,liefern,liefert,liefern},
  preterite = {lieferte,liefertest,lieferte,lieferten,liefertet,lieferten},
  subjunctive = {},
  imperative = {liefre (du)!,liefern wir!,liefert (ihr)!,liefern Sie!},
  perfect = {haben + \textbf{geliefert}},
  future = {werden + \textbf{liefern}},
  pluperfect = {hatten + \textbf{geliefert}},
  prepositions = {---},
  separable = {---}
}

% Migrated from verbs/b1/l.tex:7
\VerbTable{
  name = {liegen},
  english = {to lie / be located},
  type = {strong / irregular},
  auxiliary = {haben / sein},
  style = {default},
  present = {liege,liegst,liegt,liegen,liegt,liegen},
  preterite = {lag,lagst,lag,lagen,lagt,lagen},
  subjunctive = {},
  imperative = {lieg (du)!,liegen wir!,liegt (ihr)!,liegen Sie!},
  perfect = {haben / sein + \textbf{gelegen}},
  future = {werden + \textbf{liegen}},
  pluperfect = {hatten / waren + \textbf{gelegen}},
  prepositions = {---},
  separable = {---}
}

% Migrated from verbs/l.tex:13
\VerbTable{
  name = {loben},
  english = {to praise},
  type = {regular},
  auxiliary = {haben},
  style = {default},
  present = {lobe,lobst,lobt,loben,lobt,loben},
  preterite = {lobte,lobtest,lobte,lobten,lobtet,lobten},
  subjunctive = {},
  imperative = {lob(e) (du)!,loben wir!,lobt (ihr)!,loben Sie!},
  perfect = {haben + \textbf{gelobt}},
  future = {werden + \textbf{loben}},
  pluperfect = {hatten + \textbf{gelobt}},
  prepositions = {
    & \textbf{loben für} + Akk & (to praise for)\\
    & \textbf{loben wegen} + Gen & (to praise because of)\\
  },
  separable = {---}
}

% Migrated from verbs/b1/l.tex:9
\VerbTable{
  name = {löschen},
  english = {to delete},
  type = {regular},
  auxiliary = {haben},
  style = {default},
  present = {lösche,löschst,löscht,löschen,löscht,löschen},
  preterite = {löschte,löschtest,löschte,löschten,löschtet,löschten},
  subjunctive = {},
  imperative = {lösch (du)!,löschen wir!,löscht (ihr)!,löschen Sie!},
  perfect = {haben + \textbf{gelöscht}},
  future = {werden + \textbf{löschen}},
  pluperfect = {hatten + \textbf{gelöscht}},
  prepositions = {---},
  separable = {---}
}

% Migrated from verbs/b1/l.tex:10
\VerbTable{
  name = {lösen},
  english = {to release},
  type = {regular},
  auxiliary = {haben},
  style = {default},
  present = {löse,löst,löst,lösen,löst,lösen},
  preterite = {löste,löstest,löste,lösten,löstet,lösten},
  subjunctive = {},
  imperative = {lös (du)!,lösen wir!,löst (ihr)!,lösen Sie!},
  perfect = {haben + \textbf{gelöst}},
  future = {werden + \textbf{lösen}},
  pluperfect = {hatten + \textbf{gelöst}},
  prepositions = {---},
  separable = {---}
}

% Migrated from verbs/b1/l.tex:11
\VerbTable{
  name = {losfahren},
  english = {to drive off},
  type = {separable, strong / irregular},
  auxiliary = {sein},
  style = {default},
  present = {fahre los,fährst los,fährt los,fahren los,fahrt los,fahren los},
  preterite = {fuhr los,fuhrst los,fuhr los,fuhren los,fuhrt los,fuhren los},
  subjunctive = {},
  imperative = {fahr (du) los!,fahren wir los!,fahrt (ihr) los!,fahren Sie los!},
  perfect = {sein + \textbf{losgefahren}},
  future = {werden + \textbf{losfahren}},
  pluperfect = {waren + \textbf{losgefahren}},
  prepositions = {---},
  separable = {---}
}

% Migrated from verbs/b1/l.tex:12
\VerbTable{
  name = {lügen},
  english = {to lie},
  type = {strong / irregular},
  auxiliary = {haben},
  style = {default},
  present = {lüge,lügst,lügt,lügen,lügt,lügen},
  preterite = {log,logst,log,logen,logt,logen},
  subjunctive = {},
  imperative = {lüg (du)!,lügen wir!,lügt (ihr)!,lügen Sie!},
  perfect = {haben + \textbf{gelogen}},
  future = {werden + \textbf{lügen}},
  pluperfect = {hatten + \textbf{gelogen}},
  prepositions = {---},
  separable = {---}
}

\section{M}

% Migrated from verbs/m.tex:1
\VerbTable{
  name = {machen},
  english = {to make},
  type = {regular, + Akkusative},
  auxiliary = {haben},
  style = {acc},
  present = {mache,machst,macht,machen,macht,machen},
  preterite = {machte,machtest,machte,machten,machtet,machten},
  subjunctive = {},
  imperative = {mach(e) (du)!,machen wir!,macht (ihr)!,machen Sie!},
  perfect = {haben + \textbf{gemacht}},
  future = {werden + \textbf{machen}},
  pluperfect = {hatten + \textbf{gemacht}},
  prepositions = {
    & \textbf{machen mit} + Dat & (to participate in)\\
    & \textbf{machen für} + Akk & (to do something for someone / purpose)\\
    &\textbf{machen aus} + Dat & (to turn something into something)\\
    &\textbf{machen über} + Akk & (to complain / joke about)\\
    &\textbf{machen von} + Dat & (to make of / think of)\\
    &\textbf{machen bei} + Dat & (to take part in)\\
  },
  separable = {
    & \textbf{an$|$machen} & (to turn on)\\
    & \textbf{aus$|$machen} & (to turn off)\\
    & \textbf{auf$|$machen} & (to open)\\
    & \textbf{zu$|$machen} & (to close)\\
    & \textbf{mit$|$machen} & (to participate)\\
    & \textbf{nach$|$machen} & (to imitate)\\
    & \textbf{weiter$|$machen} & (to continue)\\
    & \textbf{fertig$|$machen} & (to finish /exhaust)\\
    & \textbf{kaputt$|$machen} & (to break)\\
    & \textbf{los$|$machen} & (to loosen / detach)\\
  }
}

% Migrated from verbs/b1/m.tex:1
\VerbTable{
  name = {malen},
  english = {to paint},
  type = {regular},
  auxiliary = {haben},
  style = {default},
  present = {male,malst,malt,malen,malt,malen},
  preterite = {malte,maltest,malte,malten,maltet,malten},
  subjunctive = {},
  imperative = {mal (du)!,malen wir!,malt (ihr)!,malen Sie!},
  perfect = {haben + \textbf{gemalt}},
  future = {werden + \textbf{malen}},
  pluperfect = {hatten + \textbf{gemalt}},
  prepositions = {---},
  separable = {---}
}

% Migrated from verbs/b1/m.tex:2
\VerbTable{
  name = {markieren},
  english = {to mark / highlight},
  type = {regular},
  auxiliary = {haben},
  style = {default},
  present = {markiere,markierst,markiert,markieren,markiert,markieren},
  preterite = {markierte,markiertest,markierte,markierten,markiertet,markierten},
  subjunctive = {},
  imperative = {markier (du)!,markieren wir!,markiert (ihr)!,markieren Sie!},
  perfect = {haben + \textbf{markiert}},
  future = {werden + \textbf{markieren}},
  pluperfect = {hatten + \textbf{markiert}},
  prepositions = {---},
  separable = {---}
}

% Migrated from verbs/m.tex:2
\VerbTable{
  name = {mauern},
  english = {to build walls / brick},
  type = {regular},
  auxiliary = {haben},
  style = {default},
  present = {mauere,mauerst,mauert,mauern,mauert,mauern},
  preterite = {mauerte,mauertest,mauerte,mauerten,mauertet,mauerten},
  subjunctive = {},
  imperative = {mauere (du)!,mauern wir!,mauert (ihr)!,mauern Sie!},
  perfect = {haben + \textbf{gemauert}},
  future = {werden + \textbf{mauern}},
  pluperfect = {hatten + \textbf{gemauert}},
  prepositions = {
    & \textbf{mauern mit} + Dat & (to build with / using)\\
    & \textbf{mauern an} + Dat & (to build onto / on)\\
  },
  separable = {& \textbf{ein$|$mauern} & (to wall in / enclose)\\}
}

% Migrated from verbs/m.tex:3
\VerbTable{
  name = {meinen},
  english = {to mean something},
  type = {regular, + Akkusativ},
  auxiliary = {haben},
  style = {acc},
  present = {meine,meinst,meint,meinen,meint,meinen},
  preterite = {meinte,meintest,meinte,meinten,meintet,meinten},
  subjunctive = {},
  imperative = {mein(e) (du)!,meinen wir!,meint (ihr)!,meinen Sie!},
  perfect = {haben + \textbf{gemeint}},
  future = {werden + \textbf{meinen}},
  pluperfect = {hatten + \textbf{gemeint}},
  prepositions = {---},
  separable = {---}
}

% Migrated from verbs/m.tex:4
\VerbTable{
  name = {melken},
  english = {to milk},
  type = {irregular},
  auxiliary = {haben},
  style = {default},
  present = {melke,melkst,melkt,melken,melkt,melken},
  preterite = {molk,molkst,molk,molken,molkt,molken},
  subjunctive = {},
  imperative = {melk(e) (du)!,melken wir!,melkt (ihr)!,melken Sie!},
  perfect = {haben + \textbf{gemolken}},
  future = {werden + \textbf{melken}},
  pluperfect = {hatten + \textbf{gemolken}},
  prepositions = {
    & \textbf{melken von} + Dat & (to extract from)\\
    & \textbf{melken aus} + Dat & (to extract out of)\\
  },
  separable = {
    & \textbf{ab$|$melken} & (to siphon off / exploit)\\
    & \textbf{aus$|$melken} & (to milk out / exploit fully)\\
  }
}

% Migrated from verbs/m.tex:5
\VerbTable{
  name = {merken},
  english = {to notice / remember},
  type = {regular},
  auxiliary = {haben},
  style = {default},
  present = {merke,merkst,merkt,merken,merkt,merken},
  preterite = {merkte,merktest,merkte,merkten,merktet,merkten},
  subjunctive = {},
  imperative = {merk(e) (du)!,merken wir!,merkt (ihr)!,merken Sie!},
  perfect = {haben + \textbf{gemerkt}},
  future = {werden + \textbf{merken}},
  pluperfect = {hatten + \textbf{gemerkt}},
  prepositions = {
    & \textbf{merken an} + Dat & (to notice in)\\
    & \textbf{merken sich} + Akk & (to remember)\\
  },
  separable = {
    & \textbf{fest$|$merken} & (to take note of)\\
    & \textbf{zurück$|$merken} & (to remember for later)\\
  }
}

% Migrated from verbs/b1/m.tex:3
\VerbTable{
  name = {messen},
  english = {to measure},
  type = {strong / irregular},
  auxiliary = {haben},
  style = {default},
  present = {messe,misst,misst,messen,messt,messen},
  preterite = {maß,maßt,maß,maßen,maßt,maßen},
  subjunctive = {},
  imperative = {miss (du)!,messen wir!,messt (ihr)!,messen Sie!},
  perfect = {haben + \textbf{gemessen}},
  future = {werden + \textbf{messen}},
  pluperfect = {hatten + \textbf{gemessen}},
  prepositions = {---},
  separable = {---}
}

% Migrated from verbs/b1/m.tex:4
\VerbTable{
  name = {mieten},
  english = {to rent},
  type = {regular},
  auxiliary = {haben},
  style = {default},
  present = {miete,mietest,mietet,mieten,mietet,mieten},
  preterite = {mietete,mietetest,mietete,mieteten,mietetet,mieteten},
  subjunctive = {},
  imperative = {miete (du)!,mieten wir!,mietet (ihr)!,mieten Sie!},
  perfect = {haben + \textbf{gemietet}},
  future = {werden + \textbf{mieten}},
  pluperfect = {hatten + \textbf{gemietet}},
  prepositions = {---},
  separable = {---}
}

% Migrated from verbs/b1/m.tex:5
\VerbTable{
  name = {mischen},
  english = {to mix},
  type = {regular},
  auxiliary = {haben},
  style = {default},
  present = {mische,mischst,mischt,mischen,mischt,mischen},
  preterite = {mischte,mischtest,mischte,mischten,mischtet,mischten},
  subjunctive = {},
  imperative = {misch (du)!,mischen wir!,mischt (ihr)!,mischen Sie!},
  perfect = {haben + \textbf{gemischt}},
  future = {werden + \textbf{mischen}},
  pluperfect = {hatten + \textbf{gemischt}},
  prepositions = {---},
  separable = {---}
}

% Migrated from verbs/m.tex:6
\VerbTable{
  name = {missverstehen},
  english = {to misunderstand},
  type = {irregular},
  auxiliary = {haben},
  style = {default},
  present = {missverstehe,missverstehst,missversteht,missverstehen,missversteht,missverstehen},
  preterite = {missverstand,missverstandest,missverstand,missverstanden,missverstandet,missverstanden},
  subjunctive = {},
  imperative = {missversteh(e) (du)!,missverstehen wir!,missversteht (ihr)!,missverstehen Sie!},
  perfect = {haben + \textbf{missverstanden}},
  future = {werde + \textbf{missverstehen}},
  pluperfect = {hatten + \textbf{missverstanden}},
  prepositions = {
    & \textbf{missverstehen als} + Akk & (to misunderstand as)\\
    & \textbf{missverstehen wegen} + Gen & (to misunderstand because of)\\
  },
  separable = {---}
}

% Migrated from verbs/b1/m.tex:6
\VerbTable{
  name = {mitteilen},
  english = {to inform / tell},
  type = {separable, regular},
  auxiliary = {haben},
  style = {default},
  present = {teile mit,teilst mit,teilt mit,teilen mit,teilt mit,teilen mit},
  preterite = {teilte mit,teiltest mit,teilte mit,teilten mit,teiltet mit,teilten mit},
  subjunctive = {},
  imperative = {teil (du) mit!,teilen wir mit!,teilt (ihr) mit!,teilen Sie mit!},
  perfect = {haben + \textbf{mitgeteilt}},
  future = {werden + \textbf{mitteilen}},
  pluperfect = {hatten + \textbf{mitgeteilt}},
  prepositions = {---},
  separable = {---}
}

\section{N}

% Migrated from verbs/b1/n.tex:1
\VerbTable{
  name = {nachdenken},
  english = {to think / reflect},
  type = {separable, regular},
  auxiliary = {haben},
  style = {default},
  present = {denke nach,denkst nach,denkt nach,denken nach,denkt nach,denken nach},
  preterite = {dachte nach,dachtest nach,dachte nach,dachten nach,dachtet nach,dachten nach},
  subjunctive = {},
  imperative = {denk (du) nach!,denken wir nach!,denkt (ihr) nach!,denken Sie nach!},
  perfect = {haben + \textbf{nachgedacht}},
  future = {werden + \textbf{nachdenken}},
  pluperfect = {hatten + \textbf{nachgedacht}},
  prepositions = {---},
  separable = {---}
}

% Migrated from verbs/b1/n.tex:2
\VerbTable{
  name = {nachschlagen},
  english = {to look up},
  type = {separable, strong / irregular},
  auxiliary = {haben},
  style = {default},
  present = {schlage nach,schlägst nach,schlägt nach,schlagen nach,schlagt nach,schlagen nach},
  preterite = {schlug nach,schlugst nach,schlug nach,schlugen nach,schlugt nach,schlugen nach},
  subjunctive = {},
  imperative = {schlag (du) nach!,schlagen wir nach!,schlagt (ihr) nach!,schlagen Sie nach!},
  perfect = {haben + \textbf{nachgeschlagen}},
  future = {werden + \textbf{nachschlagen}},
  pluperfect = {hatten + \textbf{nachgeschlagen}},
  prepositions = {---},
  separable = {---}
}

% Migrated from verbs/n.tex:1
\VerbTable{
  name = {nähen},
  english = {to sew},
  type = {regular},
  auxiliary = {haben},
  style = {default},
  present = {nähe,nähst,näht,nähen,näht,nähen},
  preterite = {nähte,nähntest,nähte,nähten,nähtet,nähten},
  subjunctive = {},
  imperative = {nähe (du)!,nähen wir!,näht (ihr)!,nähen Sie!},
  perfect = {haben + \textbf{genäht}},
  future = {werden + \textbf{nähen}},
  pluperfect = {hatten + \textbf{genäht}},
  prepositions = {
    & \textbf{nähen an} + Dat & (to sew on / at)\\
    & \textbf{nähen für} + Akk & (to sew for / on behalf of)\\
  },
  separable = {& \textbf{auf$|$nähen} & (to sew onto / attach by sewing)\\}
}

% Migrated from verbs/n.tex:2
\VerbTable{
  name = {nehmen},
  english = {to take},
  type = {irregular, + Akkusative},
  auxiliary = {haben},
  style = {acc},
  present = {nehme,nimmst,nimmt,nehmen,nehmt,nehmen},
  preterite = {nahm,nahmst,nahm,nahmen,nahmt,nahmen},
  subjunctive = {},
  imperative = {nimm (du)!,nehmen wir!,nehmt (ihr)!,nehmen Sie!},
  perfect = {haben + \textbf{genommen}},
  future = {werden + \textbf{nehmen}},
  pluperfect = {hatten + \textbf{genommen}},
  prepositions = {
    & \textbf{nehmen aus} + Dat & (to take from)\\
    & \textbf{nehmen von} + Dat & (to take from)\\
    &\textbf{nehmen für} + Akk & (to take for)\\
    &\textbf{nehmen über} + Akk & (to take over / control of)\\
  },
  separable = {
    & \textbf{ab$|$nehmen} & (to turn off / lose weight)\\
    & \textbf{an$|$nehmen} & (to accept / assume)\\
    & \textbf{auf$|$nehmen} & (to pick up / record)\\
    & \textbf{ein$|$nehmen} & (to take something / to occupy)\\
    & \textbf{mit$|$nehmen} & (to take along)\\
    & \textbf{über$|$nehmen} & (to take over)\\
    & \textbf{weg$|$nehmen} & (to take away)\\
    & \textbf{zurück$|$nehmen} & (to take back)\\
    & \textbf{teil$|$nehmen} & (to participate)\\
    & \textbf{zu$|$nehmen} & (to gain weight)\\
  }
}

% Migrated from verbs/n.tex:3
\VerbTable{
  name = {nennen},
  english = {to name},
  type = {irregular, + Akkusativ},
  auxiliary = {haben},
  style = {acc},
  present = {nenne,nennst,nennt,nennen,nennt,nennen},
  preterite = {nannte,nanntest,nannte,nannten,nanntet,nannten},
  subjunctive = {},
  imperative = {nenn(e) (du)!,nennen wir!,nennt (ihr)!,nennen Sie!},
  perfect = {haben + \textbf{genannt}},
  future = {werden + \textbf{nennen}},
  pluperfect = {hatten + \textbf{genannt}},
  prepositions = {---},
  separable = {---}
}

% Migrated from verbs/n.tex:4
\VerbTable{
  name = {niesen},
  english = {to sneeze},
  type = {regular},
  auxiliary = {haben},
  style = {default},
  present = {niese,niest,niest,niesen,niest,niesen},
  preterite = {nieste,niestest,nieste,niesten,niestet,niesten},
  subjunctive = {},
  imperative = {niese (du)!,niesen wir!,niest (ihr)!,niesen Sie!},
  perfect = {haben + \textbf{geniest}},
  future = {werden + \textbf{niesen}},
  pluperfect = {hatten + \textbf{geniest}},
  prepositions = {& \textbf{niesen in} + Akk & (to sneeze into)\\},
  separable = {---}
}

% Migrated from verbs/b1/n.tex:4
\VerbTable{
  name = {nützen},
  english = {to use},
  type = {regular},
  auxiliary = {haben},
  style = {default},
  present = {nütze,nützt,nützt,nützen,nützt,nützen},
  preterite = {nützte,nütztest,nützte,nützten,nütztet,nützten},
  subjunctive = {},
  imperative = {nütz (du)!,nützen wir!,nützt (ihr)!,nützen Sie!},
  perfect = {haben + \textbf{genützt}},
  future = {werden + \textbf{nützen}},
  pluperfect = {hatten + \textbf{genützt}},
  prepositions = {---},
  separable = {---}
}

% Migrated from verbs/n.tex:5
\VerbTable{
  name = {nutzen},
  english = {to use},
  type = {regular, + Akkusativ},
  auxiliary = {haben},
  style = {acc},
  present = {nutze,nutzt,nutzt,nutzen,nutzt,nutzen},
  preterite = {nutzte,nutztest,nutzte,nutzten,nutztet,nutzten},
  subjunctive = {},
  imperative = {nutz(e) (du)!,ven wir!,nutzt (ihr)!,nutzen Sie!},
  perfect = {haben + \textbf{genutzt}},
  future = {werden + \textbf{nutzen}},
  pluperfect = {hatten + \textbf{genutzt}},
  prepositions = {---},
  separable = {---}
}

\section{O}

% Migrated from verbs/o.tex:1
\VerbTable{
  name = {öffnen},
  english = {to open},
  type = {regular, + Akkusativ},
  auxiliary = {haben},
  style = {acc},
  present = {öffne,öffnest,öffnet,öffnen,öffnet,öffnen},
  preterite = {öffnete,öffnetest,öffnete,öffneten,öffnetet,öffneten},
  subjunctive = {},
  imperative = {öffne (du)!,öffnen wir!,öffnet (ihr)!,öffnen Sie!},
  perfect = {haben + \textbf{geöffnet}},
  future = {werden + \textbf{öffnen}},
  pluperfect = {hatten + \textbf{geöffnet}},
  prepositions = {---},
  separable = {---}
}

% Migrated from verbs/b1/o.tex:1
\VerbTable{
  name = {operieren},
  english = {to operate},
  type = {regular},
  auxiliary = {haben},
  style = {default},
  present = {operiere,operierst,operiert,operieren,operiert,operieren},
  preterite = {operierte,operiertest,operierte,operierten,operiertet,operierten},
  subjunctive = {},
  imperative = {operier (du)!,operieren wir!,operiert (ihr)!,operieren Sie!},
  perfect = {haben + \textbf{operiert}},
  future = {werden + \textbf{operieren}},
  pluperfect = {hatten + \textbf{operiert}},
  prepositions = {---},
  separable = {---}
}

% Migrated from verbs/b1/o.tex:2
\VerbTable{
  name = {ordnen},
  english = {to order},
  type = {regular},
  auxiliary = {haben},
  style = {default},
  present = {ordne,ordnest,ordnet,ordnen,ordnet,ordnen},
  preterite = {ordnete,ordnetest,ordnete,ordneten,ordnetet,ordneten},
  subjunctive = {},
  imperative = {ordne (du)!,ordnen wir!,ordnet (ihr)!,ordnen Sie!},
  perfect = {haben + \textbf{geordnet}},
  future = {werden + \textbf{ordnen}},
  pluperfect = {hatten + \textbf{geordnet}},
  prepositions = {---},
  separable = {---}
}

% Migrated from verbs/o.tex:2
\VerbTable{
  name = {organisieren},
  english = {to organize},
  type = {regular, + Akkusativ},
  auxiliary = {haben},
  style = {acc},
  present = {organisiere,organisierst,organisiert,organisieren,organisiert,organisieren},
  preterite = {organisierte,organisiertest,organisierte,organisierten,organisiertet,organisierten},
  subjunctive = {},
  imperative = {organisiere (du)!,organisieren wir!,organisiert (ihr)!,organisieren Sie!},
  perfect = {haben + \textbf{organisiert}},
  future = {werden + \textbf{organisieren}},
  pluperfect = {hatten + \textbf{organisiert}},
  prepositions = {& \textbf{organisieren für} + Akk & (to organize for)\\},
  separable = {---}
}

\section{P}

% Migrated from verbs/p.tex:1
\VerbTable{
  name = {packen},
  english = {to pack},
  type = {regular},
  auxiliary = {haben},
  style = {default},
  present = {packe,packst,packt,packen,packt,packen},
  preterite = {packte,packtest,packte,packten,packtet,packten},
  subjunctive = {},
  imperative = {pack(e) (du)!,packen wir!,packt (ihr)!,packen Sie!},
  perfect = {haben + \textbf{gepackt}},
  future = {werden + \textbf{packen}},
  pluperfect = {hatten + \textbf{gepackt}},
  prepositions = {
    & \textbf{packen in} + Akk & (to pack into)\\
    & \textbf{packen auf} + Akk & (to load onto)\\
  },
  separable = {
    & \textbf{ein$|$packen} & (to pack up)\\
    & \textbf{aus$|$packen} & (to unpack)\\
    & \textbf{an$|$packen} & (to tackle / get started)\\
    & \textbf{mit$|$packen} & (to help pack)\\
  }
}

% Migrated from verbs/p.tex:2
\VerbTable{
  name = {parken},
  english = {to park},
  type = {regular},
  auxiliary = {haben},
  style = {default},
  present = {parke,parkst,parkt,parken,parkt,parken},
  preterite = {parkte,parktest,parkte,parkten,parktet,parkten},
  subjunctive = {},
  imperative = {park(e) (du)!,parken wir!,parket (ihr)!,parken Sie!},
  perfect = {haben + \textbf{geparkt}},
  future = {werden + \textbf{parken}},
  pluperfect = {hatten + \textbf{geparkt}},
  prepositions = {---},
  separable = {---}
}

% Migrated from verbs/p.tex:3
\VerbTable{
  name = {passen},
  english = {to fit / to suit},
  type = {regular},
  auxiliary = {haben},
  style = {dat},
  present = {passe,passt,passt,passen,passt,passen},
  preterite = {passte,passtest,passte,passten,passtet,passten},
  subjunctive = {},
  imperative = {pass(e) (du)!,passen wir!,passt (ihr)!,passen Sie!},
  perfect = {haben + \textbf{gepasst}},
  future = {werden + \textbf{passen}},
  pluperfect = {hatten + \textbf{gepasst}},
  prepositions = {
    & \textbf{passen zu} + Dat & (to suit / go with)\\
    & \textbf{passen in} + Akk & (to fit into)\\
    & \textbf{passen auf} + Akk & (to watch / take care of)
  },
  separable = {
    & \textbf{auf$|$passen} & (to pay attention / watch)\\
    & \textbf{an$|$passen} & (to adapt / adjust)\\
  }
}

% Migrated from verbs/p.tex:4
\VerbTable{
  name = {passieren},
  english = {to happen},
  type = {irregular},
  auxiliary = {sein},
  style = {default},
  present = {-,-,passiert,-,-,passieren},
  preterite = {-,-,passierte,-,-,passierten},
  subjunctive = {},
  imperative = {---},
  perfect = {sein + \textbf{passiert}},
  future = {werden + \textbf{passieren}},
  pluperfect = {waren + \textbf{passiert}},
  prepositions = {---},
  separable = {---}
}

% Migrated from verbs/b1/p.tex:1
\VerbTable{
  name = {pflanzen},
  english = {to plant},
  type = {regular},
  auxiliary = {haben},
  style = {default},
  present = {pflanze,pflanzt,pflanzt,pflanzen,pflanzt,pflanzen},
  preterite = {pflanzte,pflanztest,pflanzte,pflanzten,pflanztet,pflanzten},
  subjunctive = {},
  imperative = {pflanz (du)!,pflanzen wir!,pflanzt (ihr)!,pflanzen Sie!},
  perfect = {haben + \textbf{gepflanzt}},
  future = {werden + \textbf{pflanzen}},
  pluperfect = {hatten + \textbf{gepflanzt}},
  prepositions = {---},
  separable = {---}
}

% Migrated from verbs/b1/p.tex:2
\VerbTable{
  name = {pflegen},
  english = {to care for / maintain},
  type = {regular},
  auxiliary = {haben},
  style = {default},
  present = {pflege,pflegst,pflegt,pflegen,pflegt,pflegen},
  preterite = {pflegte,pflegtest,pflegte,pflegten,pflegtet,pflegten},
  subjunctive = {},
  imperative = {pfleg (du)!,pflegen wir!,pflegt (ihr)!,pflegen Sie!},
  perfect = {haben + \textbf{gepflegt}},
  future = {werden + \textbf{pflegen}},
  pluperfect = {hatten + \textbf{gepflegt}},
  prepositions = {---},
  separable = {---}
}

% Migrated from verbs/p.tex:5
\VerbTable{
  name = {picknicken},
  english = {to have a picnic},
  type = {regular},
  auxiliary = {haben},
  style = {default},
  present = {picknicke,picknickst,picknickt,picknicken,picknickt,picknicken},
  preterite = {picknickte,picknicktest,picknickte,picknickten,picknicktet,picknickten},
  subjunctive = {},
  imperative = {picknicke (du)!,picknicken wir!,picknickt (ihr)!,picknicken Sie!},
  perfect = {haben + \textbf{gepicknickt}},
  future = {werden + \textbf{picknicken}},
  pluperfect = {hatten + \textbf{gepicknickt}},
  prepositions = {
    & \textbf{picknicken im} + Dat & (to have a picnic in)\\
    & \textbf{picknicken auf} + Dat & (to have a picnic on)\\
  },
  separable = {---}
}

% Migrated from verbs/p.tex:6
\VerbTable{
  name = {planen},
  english = {to plan},
  type = {regular, + Akkusativ},
  auxiliary = {haben},
  style = {acc},
  present = {plane,planst,plant,planen,plant,planen},
  preterite = {plante,plantest,plante,planten,plantet,planten},
  subjunctive = {},
  imperative = {plan(e) (du)!,planen wir!,plant (ihr)!,planen Sie!},
  perfect = {haben + \textbf{geplant}},
  future = {werden + \textbf{planen}},
  pluperfect = {hatten + \textbf{geplant}},
  prepositions = {---},
  separable = {---}
}

% Migrated from verbs/p.tex:7
\VerbTable{
  name = {plaudern},
  english = {to chat},
  type = {regular},
  auxiliary = {haben},
  style = {default},
  present = {plaudere,plauderst,plaudert,plaudern,plaudert,plaudern},
  preterite = {plauderte,plaudertest,plauderte,plauderten,plaudertet,plauderten},
  subjunctive = {},
  imperative = {plaudere (du)!,plaudern wir!,plaudert (ihr)!,plaudern Sie!},
  perfect = {haben + \textbf{geplaudert}},
  future = {werden + \textbf{plaudern}},
  pluperfect = {hatten + \textbf{geplaudert}},
  prepositions = {
    & \textbf{plaudern mit} + Dat & (to chat with)\\
    & \textbf{plaudern über} + Akk & (to chat about)\\
  },
  separable = {& aus$|$plaudern & (to let slip / blurt out)\\}
}

% Migrated from verbs/p.tex:8
\VerbTable{
  name = {posten},
  english = {to post},
  type = {regular},
  auxiliary = {haben},
  style = {default},
  present = {poste,postest,postet,posten,postet,posten},
  preterite = {postete,postetest,postete,posteten,postetet,posteten},
  subjunctive = {},
  imperative = {poste (du)!,posten wir!,postet (ihr)!,posten Sie!},
  perfect = {haben + \textbf{gepostet}},
  future = {werden + \textbf{posten}},
  pluperfect = {hatten + \textbf{gepostet}},
  prepositions = {& \textbf{posten auf} + Dat & (to post on)\\},
  separable = {---}
}

% Migrated from verbs/p.tex:9
\VerbTable{
  name = {probieren},
  english = {to try},
  type = {irregular, + Akkusativ},
  auxiliary = {haben},
  style = {acc},
  present = {probiere,probierst,probiert,probieren,probiert,probieren},
  preterite = {probierte,probiertest,probierte,probierten,probiertet,probierten},
  subjunctive = {},
  imperative = {probier(e) (du)!,probieren wir!,probiert (ihr)!,probieren Sie!},
  perfect = {haben + \textbf{probiert}},
  future = {werden + \textbf{probieren}},
  pluperfect = {hatten + \textbf{probiert}},
  prepositions = {---},
  separable = {---}
}

% Migrated from verbs/b1/p.tex:3
\VerbTable{
  name = {produzieren},
  english = {to produce},
  type = {regular},
  auxiliary = {haben},
  style = {default},
  present = {produziere,produzierst,produziert,produzieren,produziert,produzieren},
  preterite = {produzierte,produziertest,produzierte,produzierten,produziertet,produzierten},
  subjunctive = {},
  imperative = {produzier (du)!,produzieren wir!,produziert (ihr)!,produzieren Sie!},
  perfect = {haben + \textbf{produziert}},
  future = {werden + \textbf{produzieren}},
  pluperfect = {hatten + \textbf{produziert}},
  prepositions = {---},
  separable = {---}
}

% Migrated from verbs/b1/p.tex:4
\VerbTable{
  name = {protestieren},
  english = {to protest},
  type = {regular},
  auxiliary = {haben},
  style = {default},
  present = {protestiere,protestierst,protestiert,protestieren,protestiert,protestieren},
  preterite = {protestierte,protestiertest,protestierte,protestierten,protestiertet,protestierten},
  subjunctive = {},
  imperative = {protestier (du)!,protestieren wir!,protestiert (ihr)!,protestieren Sie!},
  perfect = {haben + \textbf{protestiert}},
  future = {werden + \textbf{protestieren}},
  pluperfect = {hatten + \textbf{protestiert}},
  prepositions = {---},
  separable = {---}
}

% Migrated from verbs/p.tex:10
\VerbTable{
  name = {prüfen},
  english = {to check / examine},
  type = {regular},
  auxiliary = {haben},
  style = {default},
  present = {prüfe,prüfst,prüft,prüfen,prüft,prüfen},
  preterite = {prüfte,prüftest,prüfte,prüften,prüftet,prüften},
  subjunctive = {},
  imperative = {prüfe (du)!,prüfen wir!,prüft (ihr)!,prüfen Sie!},
  perfect = {haben + \textbf{geprüft}},
  future = {werden + \textbf{prüfen}},
  pluperfect = {hatten + \textbf{geprüft}},
  prepositions = {& \textbf{prüfen auf} + Akk & (to check for)\\},
  separable = {---}
}

% Migrated from verbs/p.tex:11
\VerbTable{
  name = {putzen},
  english = {to clean},
  type = {regular, + Akkusativ},
  auxiliary = {haben},
  style = {acc},
  present = {putze,putzt,putzt,putzen,putzt,putzen},
  preterite = {putzte,putztest,putzte,putzten,putztet,putzten},
  subjunctive = {},
  imperative = {putz(e) (du)!,putzen wir!,putzt (ihr)!,putzen Sie!},
  perfect = {haben + \textbf{geputzt}},
  future = {werden + \textbf{putzen}},
  pluperfect = {hatten + \textbf{geputzt}},
  prepositions = {---},
  separable = {---}
}

\section{Q}

% Migrated from verbs/q.tex:1
\VerbTable{
  name = {qualifizieren},
  english = {to qualify},
  type = {regular, + Akkusativ},
  auxiliary = {haben},
  style = {acc},
  present = {qualifiziere,qualifizierst,qualifiziert,qualifizieren,qualifiziert,qualifizieren},
  preterite = {qualifizierte,qualifiziertest,qualifizierte,qualifizierten,qualifiziertet,qualifizierten},
  subjunctive = {},
  imperative = {qualifizier(e) (du)!,qualifizieren wir!,qualifiziert (ihr)!,qualifizieren Sie!},
  perfect = {haben + \textbf{qualifiziert}},
  future = {werden + \textbf{qualifizieren}},
  pluperfect = {hatten + \textbf{qualifiziert}},
  prepositions = {& \textbf{sich qualifizieren für} + Akk & (to qualify for)\\},
  separable = {---}
}

\section{R}

% Migrated from verbs/b1/r.tex:1
\VerbTable{
  name = {rasieren},
  english = {to shave},
  type = {regular},
  auxiliary = {haben},
  style = {default},
  present = {rasiere,rasierst,rasiert,rasieren,rasiert,rasieren},
  preterite = {rasierte,rasiertest,rasierte,rasierten,rasiertet,rasierten},
  subjunctive = {},
  imperative = {rasier (du)!,rasieren wir!,rasiert (ihr)!,rasieren Sie!},
  perfect = {haben + \textbf{rasiert}},
  future = {werden + \textbf{rasieren}},
  pluperfect = {hatten + \textbf{rasiert}},
  prepositions = {---},
  separable = {---}
}

% Migrated from verbs/r.tex:1
\VerbTable{
  name = {raten},
  english = {to advise / to guess},
  type = {irregular, + Dativ},
  auxiliary = {haben},
  style = {dat},
  present = {rate,rätst,rät,raten,ratet,raten},
  preterite = {riet,rietest,riet,rieten,rietet,rieten},
  subjunctive = {},
  imperative = {rat(e) (du)!,raten wir!,ratet (ihr)!,raten Sie!},
  perfect = {haben + \textbf{geraten}},
  future = {werden + \textbf{raten}},
  pluperfect = {hatten + \textbf{geraten}},
  prepositions = {
    & \textbf{raten zu} + Dat & (to advise to do / recommend)\\
    & \textbf{raten von} + Dat & (to advise against)\\
  },
  separable = {
    & \textbf{ab$|$raten} & (to advise against)\\
    & \textbf{zu$|$raten} & (to advise to do)\\
  }
}

% Migrated from verbs/b1/r.tex:2
\VerbTable{
  name = {rauchen},
  english = {to smoke},
  type = {regular},
  auxiliary = {haben},
  style = {default},
  present = {rauche,rauchst,raucht,rauchen,raucht,rauchen},
  preterite = {rauchte,rauchtest,rauchte,rauchten,rauchtet,rauchten},
  subjunctive = {},
  imperative = {rauch (du)!,rauchen wir!,raucht (ihr)!,rauchen Sie!},
  perfect = {haben + \textbf{geraucht}},
  future = {werden + \textbf{rauchen}},
  pluperfect = {hatten + \textbf{geraucht}},
  prepositions = {---},
  separable = {---}
}

% Migrated from verbs/r.tex:2
\VerbTable{
  name = {reagieren},
  english = {to react},
  type = {irregular, + Akkusativ},
  auxiliary = {haben},
  style = {acc},
  present = {reagiere,reagierst,reagiert,reagieren,reagiert,reagieren},
  preterite = {reagierte,reagiertest,reagierte,reagierten,reagiertet,reagierten},
  subjunctive = {},
  imperative = {reagier(e) (du)!,reagieren wir!,reagiert (ihr)!,reagieren Sie!},
  perfect = {haben + \textbf{reagiert}},
  future = {werden + \textbf{reagieren}},
  pluperfect = {hatten + \textbf{reagiert}},
  prepositions = {---},
  separable = {---}
}

% Migrated from verbs/b1/r.tex:3
\VerbTable{
  name = {realisieren},
  english = {to realise},
  type = {regular},
  auxiliary = {haben},
  style = {default},
  present = {realisiere,realisierst,realisiert,realisieren,realisiert,realisieren},
  preterite = {realisierte,realisiertest,realisierte,realisierten,realisiertet,realisierten},
  subjunctive = {},
  imperative = {realisier (du)!,realisieren wir!,realisiert (ihr)!,realisieren Sie!},
  perfect = {haben + \textbf{realisiert}},
  future = {werden + \textbf{realisieren}},
  pluperfect = {hatten + \textbf{realisiert}},
  prepositions = {---},
  separable = {---}
}

% Migrated from verbs/b1/r.tex:4
\VerbTable{
  name = {rechnen},
  english = {to calculate / expect},
  type = {regular},
  auxiliary = {haben},
  style = {default},
  present = {rechne,rechnest,rechnet,rechnen,rechnet,rechnen},
  preterite = {rechnete,rechnetest,rechnete,rechneten,rechnetet,rechneten},
  subjunctive = {},
  imperative = {rechne (du)!,rechnen wir!,rechnet (ihr)!,rechnen Sie!},
  perfect = {haben + \textbf{gerechnet}},
  future = {werden + \textbf{rechnen}},
  pluperfect = {hatten + \textbf{gerechnet}},
  prepositions = {---},
  separable = {---}
}

% Migrated from verbs/r.tex:3
\VerbTable{
  name = {reden},
  english = {to talk},
  type = {regular},
  auxiliary = {haben},
  style = {default},
  present = {rede,redest,redet,reden,redet,reden},
  preterite = {redete,redetest,redete,redeten,redetet,redeten},
  subjunctive = {},
  imperative = {red(e) (du)!,reden wir!,redet (ihr)!,reden Sie!},
  perfect = {haben + \textbf{geredet}},
  future = {werden + \textbf{reden}},
  pluperfect = {hatten + \textbf{geredet}},
  prepositions = {---},
  separable = {---}
}

% Migrated from verbs/r.tex:4
\VerbTable{
  name = {reduzieren},
  english = {to reduce},
  type = {regular, + Akkusativ},
  auxiliary = {haben},
  style = {acc},
  present = {reduziere,reduzierst,reduziert,reduzieren,reduziert,reduzieren},
  preterite = {reduzierte,reduziertest,reduzierte,reduzierten,reduziertet,reduzierten},
  subjunctive = {},
  imperative = {reduziere (du)!,reduzieren wir!,reduziert (ihr)!,reduzieren Sie!},
  perfect = {haben + \textbf{reduziert}},
  future = {werden + \textbf{reduzieren}},
  pluperfect = {hatten + \textbf{reduziert}},
  prepositions = {
    & \textbf{reduzieren} + Akk & (to reduce something)\\
    & \textbf{reduzieren auf} + Akk & (to reduce to)\\
    & \textbf{reduzieren um} + Akk & (to reduce by)\\
  },
  separable = {---}
}

% Migrated from verbs/r.tex:6
\VerbTable{
  name = {regeln},
  english = {to regulate / manage},
  type = {regular},
  auxiliary = {haben},
  style = {default},
  present = {regle,regelst,regelt,regeln,regelt,regeln},
  preterite = {regelte,regeltest,regelte,regelten,regeltet,regelten},
  subjunctive = {},
  imperative = {regle (du)!,regeln wir!,regelt (ihr)!,regeln Sie!},
  perfect = {haben + \textbf{geregelt}},
  future = {werden + \textbf{regeln}},
  pluperfect = {hatten + \textbf{geregelt}},
  prepositions = {& \textbf{regeln} + Akk & (to manage something)\\},
  separable = {---}
}

% Migrated from verbs/r.tex:5
\VerbTable{
  name = {registrieren},
  english = {to register},
  type = {regular, -ieren verb},
  auxiliary = {haben},
  style = {default},
  present = {registriere,registrierst,registriert,registrieren,registriert,registrieren},
  preterite = {registrierte,registriertest,registrierte,registrierten,registriertet,registrierten},
  subjunctive = {},
  imperative = {registriere (du)!,registrieren wir!,registriert (ihr)!,registrieren Sie!},
  perfect = {haben + \textbf{registriert}},
  future = {werden + \textbf{registrieren}},
  pluperfect = {hatten + \textbf{registriert}},
  prepositions = {
    & \textbf{sich registrieren bei} + Dat & (to register with)\\
    & \textbf{sich registrieren für} + Akk & (to register for)\\
  },
  separable = {---}
}

% Migrated from verbs/r.tex:7
\VerbTable{
  name = {regnen},
  english = {to rain},
  type = {regular},
  auxiliary = {haben},
  style = {default},
  present = {-,-,regnet,-,-,-},
  preterite = {-,-,regnete,-,-,-},
  subjunctive = {},
  imperative = {---},
  perfect = {haben + \textbf{geredet}},
  future = {werden + \textbf{regnen}},
  pluperfect = {hatten + \textbf{geregnet}},
  prepositions = {---},
  separable = {---}
}

% Migrated from verbs/r.tex:8
\VerbTable{
  name = {reichen},
  english = {to be enough / to hand},
  type = {regular},
  auxiliary = {haben},
  style = {default},
  present = {reiche,reichst,reicht,reichen,reicht,reichen},
  preterite = {reichte,reichtest,reichte,reichten,reichtet,reichten},
  subjunctive = {},
  imperative = {reiche (du)!,reichen wir!,reicht (ihr)!,reichen Sie!},
  perfect = {haben + \textbf{gereicht}},
  future = {werden + \textbf{reichen}},
  pluperfect = {hatten + \textbf{gereicht}},
  prepositions = {
    & \textbf{reichen für} + Akk & (to be enough for)\\
    & \textbf{reichen an} + Akk & (to hand to)\\
  },
  separable = {
    & \textbf{ein$|$reichen} & (to submit)\\
    & \textbf{nach$|$reichen} & (to hand later)\\
  }
}

% Migrated from verbs/r.tex:9
\VerbTable{
  name = {reinigen},
  english = {to clean},
  type = {regular},
  auxiliary = {haben},
  style = {default},
  present = {reinige,reinigst,reinigt,reinigen,reinigt,reinigen},
  preterite = {reinigte,reinigtest,reinigte,reinigten,reinigtet,reinigten},
  subjunctive = {},
  imperative = {reinige (du)!,reinigen wir!,reinigt (ihr)!,reinigen Sie!},
  perfect = {haben + \textbf{gereinigt}},
  future = {werden + \textbf{reinigen}},
  pluperfect = {hatten + \textbf{gereinigt}},
  prepositions = {
    & \textbf{reinigen von} + Dat & (to clean from)\\
    & \textbf{reinigen mit} + Dat & (to clean with)\\
  },
  separable = {
    & \textbf{ab$|$reinigen} & (to clean off)\\
    & \textbf{aus$|$reinigen} & (to clean out)\\
  }
}

% Migrated from verbs/r.tex:10
\VerbTable{
  name = {reisen},
  english = {to travel},
  type = {regular},
  auxiliary = {sein},
  style = {default},
  present = {reise,reist,reist,reisen,reist,reisen},
  preterite = {reiste,reistest,reiste,reisten,reistet,reisten},
  subjunctive = {},
  imperative = {reis(e) (du)!,reisen wir!,reist (ihr)!,reisen Sie!},
  perfect = {sein + \textbf{gereist}},
  future = {werden + \textbf{reisen}},
  pluperfect = {waren + \textbf{gereist}},
  prepositions = {---},
  separable = {---}
}

% Migrated from verbs/b1/r.tex:5
\VerbTable{
  name = {reiten},
  english = {to ride},
  type = {strong / irregular},
  auxiliary = {sein},
  style = {default},
  present = {reite,reitest,reitet,reiten,reitet,reiten},
  preterite = {ritt,rittest,ritt,ritten,rittet,ritten},
  subjunctive = {},
  imperative = {reite (du)!,reiten wir!,reitet (ihr)!,reiten Sie!},
  perfect = {sein + \textbf{geritten}},
  future = {werden + \textbf{reiten}},
  pluperfect = {waren + \textbf{geritten}},
  prepositions = {---},
  separable = {---}
}

% Migrated from verbs/r.tex:11
\VerbTable{
  name = {reklamieren},
  english = {to complain / claim},
  type = {regular},
  auxiliary = {haben},
  style = {default},
  present = {reklamiere,reklamierst,reklamiert,reklamieren,reklamiert,reklamieren},
  preterite = {reklamierte,reklamiertest,reklamierte,reklamierten,reklamiertet,reklamierten},
  subjunctive = {},
  imperative = {reklamiere (du)!,reklamieren wir!,reklamiert (ihr)!,reklamieren Sie!},
  perfect = {haben + \textbf{reklamiert}},
  future = {werden + \textbf{reklamieren}},
  pluperfect = {hatten + \textbf{reklamiert}},
  prepositions = {
    & \textbf{reklamieren bei} + Dat & (to complain to)\\
    & \textbf{reklamieren für} + Akk & (to claim for)\\
  },
  separable = {---}
}

% Migrated from verbs/b1/r.tex:6
\VerbTable{
  name = {rennen},
  english = {to run},
  type = {mixed},
  auxiliary = {sein},
  style = {default},
  present = {renne,rennst,rennt,rennen,rennt,rennen},
  preterite = {rannte,ranntest,rannte,rannten,ranntet,rannten},
  subjunctive = {},
  imperative = {renn (du)!,rennen wir!,rennt (ihr)!,rennen Sie!},
  perfect = {sein + \textbf{gerannt}},
  future = {werden + \textbf{rennen}},
  pluperfect = {waren + \textbf{gerannt}},
  prepositions = {---},
  separable = {---}
}

% Migrated from verbs/b1/r.tex:7
\VerbTable{
  name = {reparieren},
  english = {to repair},
  type = {regular},
  auxiliary = {haben},
  style = {default},
  present = {repariere,reparierst,repariert,reparieren,repariert,reparieren},
  preterite = {reparierte,repariertest,reparierte,reparierten,repariertet,reparierten},
  subjunctive = {},
  imperative = {reparier (du)!,reparieren wir!,repariert (ihr)!,reparieren Sie!},
  perfect = {haben + \textbf{repariert}},
  future = {werden + \textbf{reparieren}},
  pluperfect = {hatten + \textbf{repariert}},
  prepositions = {---},
  separable = {---}
}

% Migrated from verbs/b1/r.tex:8
\VerbTable{
  name = {reservieren},
  english = {to reserve},
  type = {regular},
  auxiliary = {haben},
  style = {default},
  present = {reserviere,reservierst,reserviert,reservieren,reserviert,reservieren},
  preterite = {reservierte,reserviertest,reservierte,reservierten,reserviertet,reservierten},
  subjunctive = {},
  imperative = {reservier (du)!,reservieren wir!,reserviert (ihr)!,reservieren Sie!},
  perfect = {haben + \textbf{reserviert}},
  future = {werden + \textbf{reservieren}},
  pluperfect = {hatten + \textbf{reserviert}},
  prepositions = {---},
  separable = {---}
}

% Migrated from verbs/b1/r.tex:9
\VerbTable{
  name = {retten},
  english = {to rescue},
  type = {regular},
  auxiliary = {haben},
  style = {default},
  present = {rette,rettest,rettet,retten,rettet,retten},
  preterite = {rettete,rettetest,rettete,retteten,rettetet,retteten},
  subjunctive = {},
  imperative = {rette (du)!,retten wir!,rettet (ihr)!,retten Sie!},
  perfect = {haben + \textbf{gerettet}},
  future = {werden + \textbf{retten}},
  pluperfect = {hatten + \textbf{gerettet}},
  prepositions = {---},
  separable = {---}
}

% Migrated from verbs/b1/r.tex:10
\VerbTable{
  name = {riechen},
  english = {to smell},
  type = {strong / irregular},
  auxiliary = {haben},
  style = {default},
  present = {rieche,riechst,riecht,riechen,riecht,riechen},
  preterite = {roch,rochst,roch,rochen,rocht,rochen},
  subjunctive = {},
  imperative = {riech (du)!,riechen wir!,riecht (ihr)!,riechen Sie!},
  perfect = {haben + \textbf{gerochen}},
  future = {werden + \textbf{riechen}},
  pluperfect = {hatten + \textbf{gerochen}},
  prepositions = {---},
  separable = {---}
}

% Migrated from verbs/r.tex:12
\VerbTable{
  name = {rufen},
  english = {to call / shout},
  type = {irregular, + Akkusativ},
  auxiliary = {haben},
  style = {acc},
  present = {rufe,rufst,ruft,rufen,ruft,rufen},
  preterite = {rief,riefst,rief,riefen,rieft,riefen},
  subjunctive = {},
  imperative = {ruf(e) (du)!,rufen wir!,ruft (ihr)!,rufen Sie!},
  perfect = {haben + \textbf{gerufen}},
  future = {werden + \textbf{rufen}},
  pluperfect = {hatten + \textbf{gerufen}},
  prepositions = {
    & \textbf{rufen nach} + Dat & (to cry out for)\\
    & \textbf{rufen zu} + Dat & (to call to)\\
    &\textbf{rufen um} + Akk & (to plead for)\\
    &\textbf{anrufen bei} + Dat & (to call at (office))\\
    &\textbf{anrufen wegen} + Gen/Dat & (to call about)\\
    &\textbf{zurufen zu} + Dat & (to shout to)\\
    &\textbf{aufrufen zu} + Dat & (to call on to do)\\
    &\textbf{abrufen von} + Dat & (to retrieve from)\\
  },
  separable = {
    & \textbf{an$|$rufen} & (to call (phone))\\
    & \textbf{ab$|$rufen} & (to retrieve / access)\\
    & \textbf{zu$|$rufen} & (to call out)\\
    & \textbf{heraus$|$rufen} & (to call out (from inside))\\
    & \textbf{hinein$|$rufen} & (to call into)\\
    & \textbf{zurück$|$rufen} & (to call back)\\
    & \textbf{auf$|$rufen} & (to call up / apppeal)\\
    & \textbf{herbei$|$rufen} & (to summon)\\
  }
}

% Migrated from verbs/r.tex:13
\VerbTable{
  name = {ruhen},
  english = {to rest},
  type = {regular},
  auxiliary = {haben},
  style = {default},
  present = {ruhe,ruhst,ruht,ruhen,ruht,ruhen},
  preterite = {ruhte,ruhtest,ruhte,ruhten,ruhtet,ruhten},
  subjunctive = {},
  imperative = {ruhe (du)!,ruhen wir!,ruht (ihr)!,ruhen Sie!},
  perfect = {haben + \textbf{geruht}},
  future = {werden + \textbf{ruhen}},
  pluperfect = {hatten + \textbf{geruht}},
  prepositions = {& \textbf{ruhen auf} + Dat & (to rest on)\\},
  separable = {---}
}

% Migrated from verbs/b1/r.tex:11
\VerbTable{
  name = {runterwerfen},
  english = {to throw down},
  type = {separable, strong / irregular},
  auxiliary = {haben},
  style = {default},
  present = {werfe runter,wirfst runter,wirft runter,werfen runter,werft runter,werfen runter},
  preterite = {warf runter,warfst runter,warf runter,warfen runter,warft runter,warfen runter},
  subjunctive = {},
  imperative = {wirf (du) runter!,werfen wir runter!,werft (ihr) runter!,werfen Sie runter!},
  perfect = {haben + \textbf{runtergeworfen}},
  future = {werden + \textbf{runterwerfen}},
  pluperfect = {hatten + \textbf{runtergeworfen}},
  prepositions = {---},
  separable = {---}
}

\section{S}

% Migrated from verbs/s.tex:1
\VerbTable{
  name = {sagen},
  english = {to say},
  type = {regular, + Akkusativ},
  auxiliary = {haben},
  style = {acc},
  present = {sage,sagst,sagt,sagen,sagt,sagen},
  preterite = {sagte,sagtest,sagte,sagten,sagtet,sagten},
  subjunctive = {},
  imperative = {sag(e) (du)!,sagen wir!,sagt (ihr)!,sagen Sie!},
  perfect = {haben + \textbf{gesagt}},
  future = {werden + \textbf{sagen}},
  pluperfect = {hatten + \textbf{gesagt}},
  prepositions = {---},
  separable = {---}
}

% Migrated from verbs/s.tex:2
\VerbTable{
  name = {sammeln},
  english = {to collect / gather},
  type = {regular},
  auxiliary = {haben},
  style = {acc},
  present = {sammle,sammelst,sammelt,sammeln,sammelt,sammeln},
  preterite = {sammelte,sammeltest,sammelte,sammelten,sammeltet,sammelten},
  subjunctive = {},
  imperative = {sammle (du)!,sammeln wir!,sammelt (ihr)!,sammeln Sie!},
  perfect = {haben + \textbf{gesammelt}},
  future = {werden + \textbf{sammeln}},
  pluperfect = {hatten + \textbf{gesammelt}},
  prepositions = {
    & \textbf{sammeln für} + Akk & (to collect for / gather for)\\
    & \textbf{sammeln von} + Dat & (to collect from / gather from)\\
    & \textbf{sammeln in} + Dat & (to collect in / within)\\
  },
  separable = {
    & \textbf{zusammen$|$sammeln} & (to gather together / collect together)\\
    & \textbf{ein$|$sammeln} & (to collect in / pick up)\\
  },
  bottom-layout = {stacked}
}

% Migrated from verbs/s.tex:3
\VerbTable{
  name = {saugen},
  english = {to vacuum / suck},
  type = {regular, + Akkusativ},
  auxiliary = {haben},
  style = {acc},
  present = {sauge,saugst,saugt,saugen,saugt,saugen},
  preterite = {saugte,saugtest,saugte,saugten,saugtet,saugten},
  subjunctive = {},
  imperative = {saug(e) (du)!,saugen wir!,saugt (ihr)!,saugen Sie!},
  perfect = {haben + \textbf{gesaugt}},
  future = {werden + \textbf{saugen}},
  pluperfect = {hatten + \textbf{gesaugt}},
  prepositions = {---},
  separable = {
    & \textbf{auf$|$saugen} & (to absorb)\\
    & \textbf{ein$|$saugen} & (to inhale / absorb)\\
  }
}

% Migrated from verbs/s.tex:4
\VerbTable{
  name = {schaden},
  english = {to harm / damage},
  type = {regular},
  auxiliary = {haben},
  style = {default},
  present = {schade,schadest,schadet,schaden,schadet,schaden},
  preterite = {schadete,schadetest,schadete,schadeten,schadetet,schadeten},
  subjunctive = {},
  imperative = {schade (du)!,schaden wir!,schadet (ihr)!,schaden Sie!},
  perfect = {haben + \textbf{geschadet}},
  future = {werden + \textbf{schaden}},
  pluperfect = {hatten + \textbf{geschadet}},
  prepositions = {& \textbf{schaden durch} + Akk & (to be harmed by)\\},
  separable = {---}
}

% Migrated from verbs/b1/s.tex:1
\VerbTable{
  name = {schätzen},
  english = {to estimate},
  type = {regular},
  auxiliary = {haben},
  style = {default},
  present = {schätze,schätzt,schätzt,schätzen,schätzt,schätzen},
  preterite = {schätzte,schätztest,schätzte,schätzten,schätztet,schätzten},
  subjunctive = {},
  imperative = {schätz (du)!,schätzen wir!,schätzt (ihr)!,schätzen Sie!},
  perfect = {haben + \textbf{geschätzt}},
  future = {werden + \textbf{schätzen}},
  pluperfect = {hatten + \textbf{geschätzt}},
  prepositions = {---},
  separable = {---}
}

% Migrated from verbs/s.tex:5
\VerbTable{
  name = {schaffen},
  english = {to manage, create},
  type = {irregular, + Akkusativ},
  auxiliary = {haben},
  style = {acc},
  present = {schaffe,schaffst,schafft,schaffen,schafft,schaffen},
  preterite = {schuf,schufst,schuf,schufen,schuft,schufen},
  subjunctive = {},
  imperative = {schaff(e) (du)!,schaffen wir!,schafft (ihr)!,schaffen Sie!},
  perfect = {haben + \textbf{geschaffen}},
  future = {werden + \textbf{schaffen}},
  pluperfect = {hatten + \textbf{geschaffen}},
  prepositions = {---},
  separable = {---}
}

% Migrated from verbs/s.tex:6
\VerbTable{
  name = {schalten},
  english = {to switch},
  type = {regular},
  auxiliary = {haben},
  style = {default},
  present = {schalte,schaltest,schaltet,schalten,schaltet,schalten},
  preterite = {schaltete,schaltetest,schaltete,schalteten,schaltetet,schalteten},
  subjunctive = {},
  imperative = {schalte (du)!,schalten wir!,schaltet (ihr)!,schalten Sie!},
  perfect = {haben + \textbf{geschaltet}},
  future = {werden + \textbf{schalten}},
  pluperfect = {hatten + \textbf{geschaltet}},
  prepositions = {& \textbf{schalten auf} + Akk & (to switch to)\\},
  separable = {
    & \textbf{ein$|$schalten} & (to switch on)\\
    & \textbf{aus$|$schalten} & (to switch off)\\
  }
}

% Migrated from verbs/s.tex:7
\VerbTable{
  name = {schauen},
  english = {to look / watch},
  type = {regular},
  auxiliary = {haben},
  style = {default},
  present = {schaue,schaust,schaut,schauen,schaut,schauen},
  preterite = {schaute,schautest,schaute,schauten,schautet,schauten},
  subjunctive = {},
  imperative = {schaue (du)!,schauen wir!,schaut (ihr)!,schauen Sie!},
  perfect = {haben + \textbf{geschaut}},
  future = {werden + \textbf{schauen}},
  pluperfect = {hatten + \textbf{geschaut}},
  prepositions = {& \textbf{schauen auf} + Akk & (to look at)\\},
  separable = {& \textbf{an$|$schauen} & (to look at)\\}
}

% Migrated from verbs/s.tex:8
\VerbTable{
  name = {scheinen},
  english = {to shine},
  type = {irregular},
  auxiliary = {haben},
  style = {default},
  present = {scheine,scheinst,scheint,scheinen,scheint,scheinen},
  preterite = {schien,schienst,schien,schienen,schient,schienen},
  subjunctive = {},
  imperative = {schein(e) (du)!,scheinen wir!,scheint (ihr)!,scheinen Sie!},
  perfect = {haben + \textbf{geschienen}},
  future = {werden + \textbf{scheinen}},
  pluperfect = {hatten + \textbf{geschienen}},
  prepositions = {---},
  separable = {---}
}

% Migrated from verbs/b1/s.tex:2
\VerbTable{
  name = {schenken},
  english = {to give as a gift},
  type = {regular},
  auxiliary = {haben},
  style = {default},
  present = {schenke,schenkst,schenkt,schenken,schenkt,schenken},
  preterite = {schenkte,schenktest,schenkte,schenkten,schenktet,schenkten},
  subjunctive = {},
  imperative = {schenk (du)!,schenken wir!,schenkt (ihr)!,schenken Sie!},
  perfect = {haben + \textbf{geschenkt}},
  future = {werden + \textbf{schenken}},
  pluperfect = {hatten + \textbf{geschenkt}},
  prepositions = {---},
  separable = {---}
}

% Migrated from verbs/s.tex:9
\VerbTable{
  name = {schicken},
  english = {to send},
  type = {regular},
  auxiliary = {haben},
  style = {default},
  present = {schicke,schickst,schickt,schicken,schickt,schicken},
  preterite = {schickte,schicktest,schickte,schickten,schicktet,schickten},
  subjunctive = {},
  imperative = {schick(e) (du)!,schicken wir!,schickt (ihr)!,schicken Sie!},
  perfect = {haben + \textbf{geschickt}},
  future = {werden + \textbf{schicken}},
  pluperfect = {hatten + \textbf{geschickt}},
  prepositions = {
    & \textbf{schicken an} + Akk & (to send to)\\
    & \textbf{schicken mit} + Dat & (to send with)\\
  },
  separable = {
    & \textbf{ab$|$schicken} & (to send off)\\
    & \textbf{ein$|$schicken} & (to submit / send in)\\
    & \textbf{zurück$|$schicken} & (to send back)\\
  }
}

% Migrated from verbs/b1/s.tex:3
\VerbTable{
  name = {schieben},
  english = {to push},
  type = {strong / irregular},
  auxiliary = {haben},
  style = {default},
  present = {schiebe,schiebst,schiebt,schieben,schiebt,schieben},
  preterite = {schob,schobst,schob,schoben,schobt,schoben},
  subjunctive = {},
  imperative = {schieb (du)!,schieben wir!,schiebt (ihr)!,schieben Sie!},
  perfect = {haben + \textbf{geschoben}},
  future = {werden + \textbf{schieben}},
  pluperfect = {hatten + \textbf{geschoben}},
  prepositions = {---},
  separable = {---}
}

% Migrated from verbs/b1/s.tex:4
\VerbTable{
  name = {schießen},
  english = {to shoot},
  type = {strong / irregular},
  auxiliary = {haben},
  style = {default},
  present = {schieße,schießt,schießt,schießen,schießt,schießen},
  preterite = {schoss,schosst,schoss,schossen,schosst,schossen},
  subjunctive = {},
  imperative = {schieß (du)!,schießen wir!,schießt (ihr)!,schießen Sie!},
  perfect = {haben + \textbf{geschossen}},
  future = {werden + \textbf{schießen}},
  pluperfect = {hatten + \textbf{geschossen}},
  prepositions = {---},
  separable = {---}
}

% Migrated from verbs/b1/s.tex:5
\VerbTable{
  name = {schimpfen},
  english = {to scold / complain},
  type = {regular},
  auxiliary = {haben},
  style = {default},
  present = {schimpfe,schimpfst,schimpft,schimpfen,schimpft,schimpfen},
  preterite = {schimpfte,schimpftest,schimpfte,schimpften,schimpftet,schimpften},
  subjunctive = {},
  imperative = {schimpf (du)!,schimpfen wir!,schimpft (ihr)!,schimpfen Sie!},
  perfect = {haben + \textbf{geschimpft}},
  future = {werden + \textbf{schimpfen}},
  pluperfect = {hatten + \textbf{geschimpft}},
  prepositions = {---},
  separable = {---}
}

% Migrated from verbs/s.tex:10
\VerbTable{
  name = {schlafen},
  english = {to sleep},
  type = {irregular},
  auxiliary = {haben},
  style = {default},
  present = {schlafe,schläfst,schläft,schlafen,schlaft,schlafen},
  preterite = {schlief,schliefest,schlief,schliefen,schlieft,schlieften},
  subjunctive = {},
  imperative = {schlaf(e) (du)!,schlafen wir!,schlaft (ihr)!,schlafen Sie!},
  perfect = {haben + \textbf{geschlafen}},
  future = {werden + \textbf{schlafen}},
  pluperfect = {hatten + \textbf{geschlafen}},
  prepositions = {---},
  separable = {---}
}

% Migrated from verbs/s.tex:11
\VerbTable{
  name = {schlagen},
  english = {to hit / beat / strike},
  type = {irregular, + Akkusativ},
  auxiliary = {haben},
  style = {acc},
  present = {schlage,schlägst,schlägt,schlagen,schlagt,schlagen},
  preterite = {schlug,schlugst,schlug,schlugen,schlugt,schlugen},
  subjunctive = {},
  imperative = {schlag (du)!,schlagen wir!,schlagt (ihr)!,schlagen Sie!},
  perfect = {haben + \textbf{geschlagen}},
  future = {werden + \textbf{schlagen}},
  pluperfect = {hatten + \textbf{geschlagen}},
  prepositions = {
    & \textbf{schlagen auf} + Akk & (to hit on / strike)\\
    & \textbf{schlagen gegen} + Akk & (to hit against / collide)\\
    & \textbf{schlagen für} + Akk & (to vote / advocate for)\\
  },
  separable = {
    & \textbf{auf$|$schlagen} & (to hit open / pitch / unfold)\\
    & \textbf{ein$|$schlagen} & (to break / strike / hammer in)\\
    & \textbf{vor$|$schlagen} & (to propose / suggest)\\
    & \textbf{zurück$|$schlagen} & (to hit back / retaliate)\\
    & \textbf{zusammen$|$schlagen} & (to beat / knock together)\\
  }
}

% Migrated from verbs/s.tex:12
\VerbTable{
  name = {schleifen},
  english = {to sand / grind / sharpen},
  type = {regular},
  auxiliary = {haben},
  style = {default},
  present = {schleife,schleifst,schleift,schleifen,schleift,schleifen},
  preterite = {schliff,schliffst,schliff,schliffen,schlifft,schliffen},
  subjunctive = {},
  imperative = {schleife (du)!,schleifen wir!,schleift (ihr)!,schleifen Sie!},
  perfect = {haben + \textbf{geschliffen}},
  future = {werden + \textbf{schleifen}},
  pluperfect = {hatten + \textbf{geschliffen}},
  prepositions = {
    & \textbf{schleifen an} + Dat & (to sharpen on / work on)\\
    & \textbf{schleifen mit} + Dat & (to sand / grind with)\\
  },
  separable = {
    & \textbf{ab$|$schleifen} & (to grind off / remove)\\
    & \textbf{ein$|$schleifen} & (to polish / fit in)\\
  }
}

% Migrated from verbs/s.tex:13
\VerbTable{
  name = {schlendern},
  english = {to stroll},
  type = {regular},
  auxiliary = {sein},
  style = {default},
  present = {schlendere,schlenderst,schlendert,schlendern,schlendert,schlendern},
  preterite = {schlenderte,schlendertest,schlenderte,schlenderten,schlendertet,schlenderten},
  subjunctive = {},
  imperative = {schlendere (du)!,schlendern wir!,schlendert (ihr)!,schlendern Sie!},
  perfect = {sein + \textbf{geschlendert}},
  future = {werden + \textbf{schlendern}},
  pluperfect = {waren + \textbf{geschlendert}},
  prepositions = {
    & \textbf{schlendern durch} + Akk & (to stroll through)\\
    & \textbf{schlendern über} + Akk & (to stroll across)\\
  },
  separable = {---}
}

% Migrated from verbs/s.tex:14
\VerbTable{
  name = {schließen},
  english = {to close},
  type = {irregular, + Akkusativ},
  auxiliary = {haben},
  style = {acc},
  present = {schließe,schließt,schließt,schließen,schließt,schließen},
  preterite = {schloss,schlossest,schloss,schlossen,schlosst,schlossen},
  subjunctive = {},
  imperative = {schließ(e) (du)!,schließen wir!,schließt (ihr)!,schließen Sie!},
  perfect = {haben + \textbf{geschlossen}},
  future = {werden + \textbf{schließen}},
  pluperfect = {hatten + \textbf{geschlossen}},
  prepositions = {
    & \textbf{schließen auf} + Akk & (to infer / conclude)\\
    &\textbf{schließen aus} + Dat & (to infer from)\\
    &\textbf{schließen mit} + Dat & (to close with / end with)\\
    & \textbf{anschließen an} + Akk & (to conect to)\\
    & \textbf{abschließen mit} + Dat & (to finish with)\\
    &\textbf{einschließen in} + Akk & (to include in)\\
    &\textbf{ausschließen von} + Dat & (to exlude from)
  },
  separable = {
    & \textbf{ab$|$schließen} & (to lock)\\
    & \textbf{auf$|$schließen} & (to unlock)\\
    & \textbf{ein$|$schließen} & (to include / enclose)\\
    & \textbf{aus$|$schließen} & (to exclude)\\
    & \textbf{an$|$schließen} & (to connect / join)\\
    & \textbf{zu$|$schließen} & (to lock (shut))
  }
}

% Migrated from verbs/s.tex:15
\VerbTable{
  name = {schmecken},
  english = {to taste},
  type = {regular, + Akkusativ/Dativ},
  auxiliary = {haben},
  style = {dat},
  present = {schmecke,schmeckst,schmeckt,schmecken,schmeckt,schmecken},
  preterite = {schmeckte,schmecktest,schmeckte,schmeckten,schmecktet,schmeckten},
  subjunctive = {},
  imperative = {schmeck(e) (du)!,schmecken wir!,schmeckt (ihr)!,schmecken Sie!},
  perfect = {haben + \textbf{geschmeckt}},
  future = {werden + \textbf{schmecken}},
  pluperfect = {hatten + \textbf{geschmeckt}},
  prepositions = {---},
  separable = {---}
}

% Migrated from verbs/s.tex:16
\VerbTable{
  name = {schmücken},
  english = {to decorate},
  type = {regular},
  auxiliary = {haben},
  style = {default},
  present = {schmücke,schmückst,schmückt,schmücken,schmückt,schmücken},
  preterite = {schmückte,schmücktest,schmückte,schmückten,schmücktet,schmückten},
  subjunctive = {},
  imperative = {schmücke (du)!,schmücken wir!,schmückt (ihr)!,schmücken Sie!},
  perfect = {haben + \textbf{geschmückt}},
  future = {werden + \textbf{schmücken}},
  pluperfect = {hatten + \textbf{geschmückt}},
  prepositions = {
    & \textbf{schmücken mit} + Dat & (to decorate with)\\
    & \textbf{schmücken für} + Akk & (to decorate for)\\
  },
  separable = {& \textbf{aus$|$schmücken} & (to decorate elaborately)\\}
}

% Migrated from verbs/b1/s.tex:6
\VerbTable{
  name = {schneiden},
  english = {to cut},
  type = {strong / irregular},
  auxiliary = {haben},
  style = {default},
  present = {schneide,schneidest,schneidet,schneiden,schneidet,schneiden},
  preterite = {schnitt,schnittest,schnitt,schnitten,schnittet,schnitten},
  subjunctive = {},
  imperative = {schneide (du)!,schneiden wir!,schneidet (ihr)!,schneiden Sie!},
  perfect = {haben + \textbf{geschnitten}},
  future = {werden + \textbf{schneiden}},
  pluperfect = {hatten + \textbf{geschnitten}},
  prepositions = {---},
  separable = {---}
}

% Migrated from verbs/s.tex:17
\VerbTable{
  name = {schneien},
  english = {to snow},
  type = {irregular},
  auxiliary = {haben},
  style = {default},
  present = {-,-,schneit,-,-,-},
  preterite = {-,-,schneite,-,-,-},
  subjunctive = {},
  imperative = {---},
  perfect = {haben + \textbf{geschneit}},
  future = {werden + \textbf{schneien}},
  pluperfect = {hatten + \textbf{geschneit}},
  prepositions = {---},
  separable = {---}
}

% Migrated from verbs/s.tex:18
\VerbTable{
  name = {schreiben},
  english = {to write},
  type = {irregular, + Akkusativ},
  auxiliary = {haben},
  style = {acc},
  present = {schreibe,schreibst,schreibt,schreiben,schreibt,schreiben},
  preterite = {schrieb,schriebst,schrieb,schrieben,schriebt,schrieben},
  subjunctive = {},
  imperative = {schreib(e) (du)!,schreiben wir!,schreibt (ihr)!,schreiben Sie!},
  perfect = {haben + \textbf{geschrieben}},
  future = {werden + \textbf{schreiben}},
  pluperfect = {hatten + \textbf{geschrieben}},
  prepositions = {---},
  separable = {---}
}

% Migrated from verbs/b1/s.tex:7
\VerbTable{
  name = {schreien},
  english = {to shout / scream},
  type = {strong / irregular},
  auxiliary = {haben},
  style = {default},
  present = {schreie,schreist,schreit,schreien,schreit,schreien},
  preterite = {schrie,schriest,schrie,schrieen,schriet,schrieen},
  subjunctive = {},
  imperative = {schrei (du)!,schreien wir!,schreit (ihr)!,schreien Sie!},
  perfect = {haben + \textbf{geschrien}},
  future = {werden + \textbf{schreien}},
  pluperfect = {hatten + \textbf{geschrien}},
  prepositions = {---},
  separable = {---}
}

% Migrated from verbs/b1/s.tex:8
\VerbTable{
  name = {schütteln},
  english = {to shake},
  type = {regular},
  auxiliary = {haben},
  style = {default},
  present = {schüttle,schüttelst,schüttelt,schütteln,schüttelt,schütteln},
  preterite = {schüttelte,schütteltest,schüttelte,schüttelten,schütteltet,schüttelten},
  subjunctive = {},
  imperative = {schüttle (du)!,schütteln wir!,schüttelt (ihr)!,schütteln Sie!},
  perfect = {haben + \textbf{geschüttelt}},
  future = {werden + \textbf{schütteln}},
  pluperfect = {hatten + \textbf{geschüttelt}},
  prepositions = {---},
  separable = {---}
}

% Migrated from verbs/b1/s.tex:9
\VerbTable{
  name = {schweigen},
  english = {to be silent},
  type = {strong / irregular},
  auxiliary = {haben},
  style = {default},
  present = {schweige,schweigst,schweigt,schweigen,schweigt,schweigen},
  preterite = {schwieg,schwiegst,schwieg,schwiegen,schwiegt,schwiegen},
  subjunctive = {},
  imperative = {schweig (du)!,schweigen wir!,schweigt (ihr)!,schweigen Sie!},
  perfect = {haben + \textbf{geschwiegen}},
  future = {werden + \textbf{schweigen}},
  pluperfect = {hatten + \textbf{geschwiegen}},
  prepositions = {---},
  separable = {---}
}

% Migrated from verbs/b1/s.tex:10
\VerbTable{
  name = {schwimmen},
  english = {to swim},
  type = {strong / irregular},
  auxiliary = {sein},
  style = {default},
  present = {schwimme,schwimmst,schwimmt,schwimmen,schwimmt,schwimmen},
  preterite = {schwamm,schwammst,schwamm,schwammen,schwammt,schwammen},
  subjunctive = {},
  imperative = {schwimm (du)!,schwimmen wir!,schwimmt (ihr)!,schwimmen Sie!},
  perfect = {sein + \textbf{geschwommen}},
  future = {werden + \textbf{schwimmen}},
  pluperfect = {waren + \textbf{geschwommen}},
  prepositions = {---},
  separable = {---}
}

% Migrated from verbs/s.tex:19
\VerbTable{
  name = {schwitzen},
  english = {to sweat},
  type = {regular},
  auxiliary = {haben},
  style = {default},
  present = {schwitze,schwitzt,schwitzt,schwitzen,schwitzt,schwitzen},
  preterite = {schwitzte,schwitztest,schwitzte,schwitzten,schwitztet,schwitzten},
  subjunctive = {},
  imperative = {schwitze (du)!,schwitzen wir!,schwitzt (ihr)!,schwitzen Sie!},
  perfect = {haben + \textbf{geschwitzt}},
  future = {werden + \textbf{schwitzen}},
  pluperfect = {hatten + \textbf{geschwitzt}},
  prepositions = {
    & \textbf{schwitzen bei} + Dat & (to sweat during)\\
    & \textbf{schwitzen vor} + Dat & (to sweat from)\\
  },
  separable = {---}
}

% Migrated from verbs/s.tex:20
\VerbTable{
  name = {sehen},
  english = {to see},
  type = {irregular, + Akkusativ},
  auxiliary = {haben},
  style = {acc},
  present = {sehe,siehst,sieht,sehen,seht,sehen},
  preterite = {sah,sahst,sah,sahen,saht,sahen},
  subjunctive = {},
  imperative = {sieh(e) (du)!,sehen wir!,seht (ihr)!,sehen Sie!},
  perfect = {haben + \textbf{gesehen}},
  future = {werden + \textbf{sehen}},
  pluperfect = {hatten + \textbf{gesehen}},
  prepositions = {
    & \textbf{sehen auf} + Akk & (to pay attention to)\\
    & \textbf{sehen in} + Akk & (to look into)\\
    & \textbf{sehen nach} + Dat & (to check / look after)\\
    & \textbf{absehen von} + Dat & (to refrain from)\\
    & \textbf{ansehen als} + Akk & (to regard as)\\
    & \textbf{ansehen von} + Dat & (to look from)\\
    & \textbf{hinsehen auf} + Akk & (to look at)\\
    & \textbf{wegsehen bei} + Dat & (to ignore during)\\
    & \textbf{zusehen bei} + Dat & (to watch while)\\
  },
  separable = {
    & \textbf{an$|$sehen} & (to regard as)\\
    & \textbf{ab$|$sehen} & (to foresee / refrain)\\
    & \textbf{auf$|$sehen} & (to look up)\\
    & \textbf{aus$|$sehen} & (to appear)\\
    & \textbf{durch$|$sehen} & (to look through / examine)\\
    & \textbf{ein$|$sehen} & (to review)\\
    & \textbf{hin$|$sehen} & (to look (in a direction))\\
    & \textbf{um$|$sehen} & (to look around)\\
    & \textbf{weg$|$sehen} & (to look away / ignore)\\
    & \textbf{zu$|$sehen} & (to watch)\\
  }
}

% Migrated from verbs/b1/s.tex:11
\VerbTable{
  name = {senden},
  english = {to send},
  type = {regular / irregular variants},
  auxiliary = {haben},
  style = {default},
  present = {sende,sendest,sendet,senden,sendet,senden},
  preterite = {sendete / sandte,sendetest / sandtest,sendete / sandte,sendeten / sandten,sendetet / sandtet,sendeten / sandten},
  subjunctive = {},
  imperative = {sende (du)!,senden wir!,sendet (ihr)!,senden Sie!},
  perfect = {haben + \textbf{gesendet / gesandt}},
  future = {werden + \textbf{senden}},
  pluperfect = {hatten + \textbf{gesendet / gesandt}},
  prepositions = {---},
  separable = {---}
}

% Migrated from verbs/s.tex:21
\VerbTable{
  name = {setzen},
  english = {to set / place / put},
  type = {regular},
  auxiliary = {haben},
  style = {default},
  present = {setze,setzt,setzt,setzen,setzt,setzen},
  preterite = {setzte,setztest,setzte,setzten,setztet,setzten},
  subjunctive = {},
  imperative = {setze (du)!,setzen wir!,setzt (ihr)!,setzen Sie!},
  perfect = {haben + \textbf{gesetzt}},
  future = {werden + \textbf{setzen}},
  pluperfect = {hatten + \textbf{gesetzt}},
  prepositions = {
    & \textbf{setzen auf} + Akk & (to place on / set on)\\
    & \textbf{setzen in} + Akk & (to place into / put in)\\
    & \textbf{sich setzen auf} + Akk & (to sit down on)\\
    & \textbf{setzen für} + Akk & (to set for / dedicate)\\
  },
  separable = {
    & \textbf{ein$|$setzen} & (to insert / deploy / put in)\\
    & \textbf{auf$|$setzen} & (to place on / put onto)\\
    & \textbf{durch$|$setzen} & (to enforce / push through)\\
  }
}

% Migrated from verbs/b1/s.tex:12
\VerbTable{
  name = {sichern},
  english = {to secure / safeguard},
  type = {regular},
  auxiliary = {haben},
  style = {default},
  present = {sichre,sicherst,sichert,sichern,sichert,sichern},
  preterite = {sicherte,sichertest,sicherte,sicherten,sichertet,sicherten},
  subjunctive = {},
  imperative = {sichre (du)!,sichern wir!,sichert (ihr)!,sichern Sie!},
  perfect = {haben + \textbf{gesichert}},
  future = {werden + \textbf{sichern}},
  pluperfect = {hatten + \textbf{gesichert}},
  prepositions = {---},
  separable = {---}
}

% Migrated from verbs/b1/s.tex:13
\VerbTable{
  name = {siegen},
  english = {to win / be victorious},
  type = {regular},
  auxiliary = {haben},
  style = {default},
  present = {siege,siegst,siegt,siegen,siegt,siegen},
  preterite = {siegte,siegtest,siegte,siegten,siegtet,siegten},
  subjunctive = {},
  imperative = {sieg (du)!,siegen wir!,siegt (ihr)!,siegen Sie!},
  perfect = {haben + \textbf{gesiegt}},
  future = {werden + \textbf{siegen}},
  pluperfect = {hatten + \textbf{gesiegt}},
  prepositions = {---},
  separable = {---}
}

% Migrated from verbs/s.tex:22
\VerbTable{
  name = {siezen},
  english = {to address someone with \enquote{Sie}},
  type = {regular},
  auxiliary = {haben},
  style = {default},
  present = {sieze,siezt,siezt,siezen,siezt,siezen},
  preterite = {siezte,sieztest,siezte,siezten,sieztet,siezten},
  subjunctive = {},
  imperative = {sieze (du)!,siezen wir!,siezt (ihr)!,siezen Sie!},
  perfect = {haben + \textbf{gesiezt}},
  future = {werden + \textbf{siezen}},
  pluperfect = {hatten + \textbf{gesiezt}},
  prepositions = {& \textbf{siezen mit} + Dat & (to use "Sie" with someone)\\},
  separable = {---}
}

% Migrated from verbs/b1/s.tex:14
\VerbTable{
  name = {singen},
  english = {to sing},
  type = {strong / irregular},
  auxiliary = {haben},
  style = {default},
  present = {singe,singst,singt,singen,singt,singen},
  preterite = {sang,sangst,sang,sangen,sangt,sangen},
  subjunctive = {},
  imperative = {sing (du)!,singen wir!,singt (ihr)!,singen Sie!},
  perfect = {haben + \textbf{gesungen}},
  future = {werden + \textbf{singen}},
  pluperfect = {hatten + \textbf{gesungen}},
  prepositions = {---},
  separable = {---}
}

% Migrated from verbs/b1/s.tex:15
\VerbTable{
  name = {sinken},
  english = {to sink},
  type = {strong / irregular},
  auxiliary = {sein},
  style = {default},
  present = {sinke,sinkst,sinkt,sinken,sinkt,sinken},
  preterite = {sank,sankst,sank,sanken,sankt,sanken},
  subjunctive = {},
  imperative = {sink (du)!,sinken wir!,sinkt (ihr)!,sinken Sie!},
  perfect = {sein + \textbf{gesunken}},
  future = {werden + \textbf{sinken}},
  pluperfect = {waren + \textbf{gesunken}},
  prepositions = {---},
  separable = {---}
}

% Migrated from verbs/s.tex:23
\VerbTable{
  name = {sitzen},
  english = {to sit},
  type = {irregular},
  auxiliary = {haben / sein context-dependent},
  style = {default},
  present = {sitze,sitzt,sitzt,sitzen,sitzt,sitzen},
  preterite = {saß,saßt,saß,saßen,saßt,saßen},
  subjunctive = {},
  imperative = {sitze (du)!,sitzen wir!,sitzt (ihr)!,sitzen Sie!},
  perfect = {haben + \textbf{gesessen}},
  future = {werden + \textbf{sitzen}},
  pluperfect = {hatten + \textbf{gesessen}},
  prepositions = {
    & \textbf{sitzen auf} + Dat & (to sit on)\\
    & \textbf{sitzen in} + Dat & (to sit in / inside)\\
    & \textbf{sitzen an} + Dat & (to sit at / by)\\
  },
  separable = {
    & \textbf{abs$|$sitzen} & (to be seated / sit down)\\
    & \textbf{zusammen$|$sitzen} & (to sit together / be grouped)\\
    & \textbf{nach$|$sitzen} & (to stay after school / serve detention)\\
  }
}

% Migrated from verbs/s.tex:24
\VerbTable{
  name = {sparen},
  english = {to save},
  type = {regular, + Akkusativ},
  auxiliary = {haben},
  style = {acc},
  present = {spare,sparst,spart,sparen,spart,sparen},
  preterite = {sparte,spartest,sparte,sparten,spartet,sparten},
  subjunctive = {},
  imperative = {spar(e) (du)!,sparen wir!,sparet (ihr)!,sparen Sie!},
  perfect = {haben + \textbf{gespart}},
  future = {werden + \textbf{sparen}},
  pluperfect = {hatten + \textbf{gesparet}},
  prepositions = {---},
  separable = {---}
}

% Migrated from verbs/b1/s.tex:16
\VerbTable{
  name = {spazierengehen},
  english = {to go for a walk},
  type = {separable, strong / irregular},
  auxiliary = {sein},
  style = {default},
  present = {gehe spazieren,gehst spazieren,geht spazieren,gehen spazieren,geht spazieren,gehen spazieren},
  preterite = {ging spazieren,gingst spazieren,ging spazieren,gingen spazieren,gingt spazieren,gingen spazieren},
  subjunctive = {},
  imperative = {geh (du) spazieren!,gehen wir spazieren!,geht (ihr) spazieren!,gehen Sie spazieren!},
  perfect = {sein + \textbf{spazieren gegangen}},
  future = {werden + \textbf{spazierengehen}},
  pluperfect = {waren + \textbf{spazieren gegangen}},
  prepositions = {---},
  separable = {---}
}

% Migrated from verbs/b1/s.tex:17
\VerbTable{
  name = {speichern},
  english = {to save},
  type = {regular},
  auxiliary = {haben},
  style = {default},
  present = {speichre,speicherst,speichert,speichern,speichert,speichern},
  preterite = {speicherte,speichertest,speicherte,speicherten,speichertet,speicherten},
  subjunctive = {},
  imperative = {speichre (du)!,speichern wir!,speichert (ihr)!,speichern Sie!},
  perfect = {haben + \textbf{gespeichert}},
  future = {werden + \textbf{speichern}},
  pluperfect = {hatten + \textbf{gespeichert}},
  prepositions = {---},
  separable = {---}
}

% Migrated from verbs/s.tex:25
\VerbTable{
  name = {spielen},
  english = {to play},
  type = {regular, + Akkusativ},
  auxiliary = {haben},
  style = {acc},
  present = {spiele,spielst,spielt,spielen,spielt,spielen},
  preterite = {spielte,spieltest,spielte,spielten,spieltet,spielten},
  subjunctive = {},
  imperative = {spiel(e) (du)!,spielen wir!,spielt (ihr)!,spielen Sie!},
  perfect = {haben + \textbf{gespielt}},
  future = {werden + \textbf{spielen}},
  pluperfect = {hatten + \textbf{gespielt}},
  prepositions = {---},
  separable = {---}
}

% Migrated from verbs/s.tex:26
\VerbTable{
  name = {sprechen},
  english = {to speak},
  type = {irregular},
  auxiliary = {haben},
  style = {default},
  present = {spreche,sprichst,spricht,sprechen,sprecht,sprechen},
  preterite = {sprach,sprachst,sprach,sprachen,spracht,sprachen},
  subjunctive = {},
  imperative = {sprich (du)!,sprechen wir!,sprecht (ihr)!,sprechen Sie!},
  perfect = {haben + \textbf{gesprochen}},
  future = {werden + \textbf{sprechen}},
  pluperfect = {hatten + \textbf{gesprochen}},
  prepositions = {
    &\textbf{sprechen für} + Akk & (to speak for)\\
    & \textbf{sprechen über} + Akk & (to speak about)\\
    & \textbf{sprechen mit} + Dat & (to speak with)\\
    &\textbf{sprechen von} + Dat & (to speak of)
  },
  separable = {
    & \textbf{aus$|$sprechen} & (to pronounce)\\
    &\textbf{ab$|$sprechen} & (to agree / cancel)\\
    &\textbf{nach$|$sprechen} & (to repeat)
  }
}

% Migrated from verbs/b1/s.tex:18
\VerbTable{
  name = {springen},
  english = {to jump},
  type = {strong / irregular},
  auxiliary = {sein},
  style = {default},
  present = {springe,springst,springt,springen,springt,springen},
  preterite = {sprang,sprangst,sprang,sprangen,sprangt,sprangen},
  subjunctive = {},
  imperative = {spring (du)!,springen wir!,springt (ihr)!,springen Sie!},
  perfect = {sein + \textbf{gesprungen}},
  future = {werden + \textbf{springen}},
  pluperfect = {waren + \textbf{gesprungen}},
  prepositions = {---},
  separable = {---}
}

% Migrated from verbs/b1/s.tex:19
\VerbTable{
  name = {spülen},
  english = {to rinse / wash up},
  type = {regular},
  auxiliary = {haben},
  style = {default},
  present = {spüle,spülst,spült,spülen,spült,spülen},
  preterite = {spülte,spültest,spülte,spülten,spültet,spülten},
  subjunctive = {},
  imperative = {spül (du)!,spülen wir!,spült (ihr)!,spülen Sie!},
  perfect = {haben + \textbf{gespült}},
  future = {werden + \textbf{spülen}},
  pluperfect = {hatten + \textbf{gespült}},
  prepositions = {---},
  separable = {---}
}

% Migrated from verbs/b1/s.tex:20
\VerbTable{
  name = {spüren},
  english = {to sense / feel},
  type = {regular},
  auxiliary = {haben},
  style = {default},
  present = {spüre,spürst,spürt,spüren,spürt,spüren},
  preterite = {spürte,spürtest,spürte,spürten,spürtet,spürten},
  subjunctive = {},
  imperative = {spür (du)!,spüren wir!,spürt (ihr)!,spüren Sie!},
  perfect = {haben + \textbf{gespürt}},
  future = {werden + \textbf{spüren}},
  pluperfect = {hatten + \textbf{gespürt}},
  prepositions = {---},
  separable = {---}
}

% Migrated from verbs/b1/s.tex:21
\VerbTable{
  name = {stammen},
  english = {to come from},
  type = {regular},
  auxiliary = {haben},
  style = {default},
  present = {stamme,stammst,stammt,stammen,stammt,stammen},
  preterite = {stammte,stammtest,stammte,stammten,stammtet,stammten},
  subjunctive = {},
  imperative = {stamm (du)!,stammen wir!,stammt (ihr)!,stammen Sie!},
  perfect = {haben + \textbf{gestammt}},
  future = {werden + \textbf{stammen}},
  pluperfect = {hatten + \textbf{gestammt}},
  prepositions = {---},
  separable = {---}
}

% Migrated from verbs/s.tex:27
\VerbTable{
  name = {stapeln},
  english = {to stack / pile up},
  type = {regular},
  auxiliary = {haben},
  style = {default},
  present = {staple,stapelst,stapelt,stapeln,stapelt,stapeln},
  preterite = {stapelte,stapeltest,stapelte,stapelten,stapeltet,stapelten},
  subjunctive = {},
  imperative = {staple (du)!,stapeln wir!,stapelt (ihr)!,stapeln Sie!},
  perfect = {haben + \textbf{gestapelt}},
  future = {werden + \textbf{stapeln}},
  pluperfect = {hatten + \textbf{gestapelt}},
  prepositions = {
    & \textbf{stapeln auf} + Akk & (to stack onto)\\
    & \textbf{stapeln in} + Akk & (to stack into)\\
  },
  separable = {& \textbf{auf$|$stapeln} & (to stack up)\\}
}

% Migrated from verbs/s.tex:28
\VerbTable{
  name = {starten},
  english = {to start / launch},
  type = {regular},
  auxiliary = {haben},
  style = {default},
  present = {starte,startest,startet,starten,startet,starten},
  preterite = {startete,startetest,startete,starteten,startetet,starteten},
  subjunctive = {},
  imperative = {starte (du)!,starten wir!,startet (ihr)!,starten Sie!},
  perfect = {haben + \textbf{gestartet}},
  future = {werden + \textbf{starten}},
  pluperfect = {hatten + \textbf{gestartet}},
  prepositions = {
    & \textbf{starten zu} + Dat & (to start to / head towards)\\
    & \textbf{starten mit} + Dat & (to start with)\\
    & \textbf{starten von} + Dat & (to start from)\\
  },
  separable = {
    & \textbf{ab$|$starten} & (to take off / depart)\\
    & \textbf{auf$|$starten} & (to launch / initiate)\\
  }
}

% Migrated from verbs/b1/s.tex:22
\VerbTable{
  name = {stattfinden},
  english = {to take place},
  type = {separable, strong / irregular},
  auxiliary = {haben},
  style = {default},
  present = {finde statt,findest statt,findet statt,finden statt,findet statt,finden statt},
  preterite = {fand statt,fandest statt,fand statt,fanden statt,fandet statt,fanden statt},
  subjunctive = {},
  imperative = {finde (du) statt!,finden wir statt!,findet (ihr) statt!,finden Sie statt!},
  perfect = {haben + \textbf{stattgefunden}},
  future = {werden + \textbf{stattfinden}},
  pluperfect = {hatten + \textbf{stattgefunden}},
  prepositions = {---},
  separable = {---}
}

% Migrated from verbs/s.tex:29
\VerbTable{
  name = {staubsaugen},
  english = {to vacuum},
  type = {mixed / separable},
  auxiliary = {haben},
  style = {default},
  present = {sauge Staub,saugst Staub,saugt Staub,saugen Staub,saugt Staub,saugen Staub},
  preterite = {saugte Staub,saugtest Staub,saugte Staub,saugten Staub,saugtet Staub,saugten Staub},
  subjunctive = {},
  imperative = {saug Staub (du)!,saugen wir Staub!,saugt Staub (ihr)!,saugen Sie Staub!},
  perfect = {haben + \textbf{staubgesaugt}},
  future = {werden + \textbf{staubsaugen}},
  pluperfect = {hatten + \textbf{staubgesaugt}},
  prepositions = {& \textbf{staubsaugen} + Akk & (to vacuum something)\\},
  separable = {---}
}

% Migrated from verbs/s.tex:30
\VerbTable{
  name = {stechen},
  english = {to sting / to stab},
  type = {irregular},
  auxiliary = {haben},
  style = {default},
  present = {steche,stichst,sticht,stechen,stecht,stechen},
  preterite = {stach,stachst,stach,stachen,stacht,stachen},
  subjunctive = {},
  imperative = {stich (du)!,stechen wir!,stecht (ihr)!,stechen Sie!},
  perfect = {haben + \textbf{gestochen}},
  future = {werde + \textbf{stechen}},
  pluperfect = {hatten + \textbf{gestochen}},
  prepositions = {
    & \textbf{stechen in} + Akk & (to sting/stab into)\\
    & \textbf{stechen mit} + Dat & (to sting/stab with)\\
  },
  separable = {
    & \textbf{zu$|$stechen} & (to stab)\\
    & \textbf{ein$|$stechen} & (to stab into / pierce)\\
  }
}

% Migrated from verbs/b1/s.tex:23
\VerbTable{
  name = {stecken},
  english = {to stick / be located},
  type = {regular},
  auxiliary = {haben},
  style = {default},
  present = {stecke,steckst,steckt,stecken,steckt,stecken},
  preterite = {steckte,stecktest,steckte,steckten,stecktet,steckten},
  subjunctive = {},
  imperative = {steck (du)!,stecken wir!,steckt (ihr)!,stecken Sie!},
  perfect = {haben + \textbf{gesteckt}},
  future = {werden + \textbf{stecken}},
  pluperfect = {hatten + \textbf{gesteckt}},
  prepositions = {---},
  separable = {---}
}

% Migrated from verbs/s.tex:31
\VerbTable{
  name = {stehen},
  english = {to stand},
  type = {irregular},
  auxiliary = {haben/sein},
  style = {default},
  present = {stehe,stehst,steht,stehen,steht,stehen},
  preterite = {stand,standst,stand,standen,standet,standen},
  subjunctive = {},
  imperative = {steh(e) (du)!,stehen wir!,steht (ihr)!,stehen Sie!},
  perfect = {haben / sein + \textbf{gestanden}},
  future = {werden + \textbf{stehen}},
  pluperfect = {hatten / waren + \textbf{gestanden}},
  prepositions = {
    & \textbf{stehen für} + Akk & (to represent)\\
    & \textbf{stehen zu} + Dat & (to stand by)\\
    &\textbf{stehen vor} + Dat & (to stand facing / in front of)\\
    &\textbf{stehen in} + Dat & (to stand in)\\
    &\textbf{stehen auf} + Dat & (to stand on)\\
    &\textbf{stehen an} + Dat & (to stand at)\\
    &\textbf{stehen neben} + Dat & (to stand next to)\\
    &\textbf{stehen mit} + Dat & (to be connected with)\\
  },
  separable = {
    & \textbf{auf$|$stehen} & (to stand up, get up)\\
    & \textbf{an$|$stehen} & (to queue)\\
    & \textbf{fest$|$stehen} & (to be certain)\\
    & \textbf{zu$|$stehen} & (to be entitled to)\\
  }
}

% Migrated from verbs/b1/s.tex:24
\VerbTable{
  name = {stehlen},
  english = {to steal},
  type = {strong / irregular},
  auxiliary = {haben},
  style = {default},
  present = {stehle,stiehlst,stiehlt,stehlen,stehlt,stehlen},
  preterite = {stahl,stahlst,stahl,stahlen,stahlt,stahlen},
  subjunctive = {},
  imperative = {stiehl (du)!,stehlen wir!,stehlt (ihr)!,stehlen Sie!},
  perfect = {haben + \textbf{gestohlen}},
  future = {werden + \textbf{stehlen}},
  pluperfect = {hatten + \textbf{gestohlen}},
  prepositions = {---},
  separable = {---}
}

% Migrated from verbs/s.tex:32
\VerbTable{
  name = {steigen},
  english = {to climb / rise},
  type = {irregular},
  auxiliary = {sein},
  style = {default},
  present = {steige,steigst,steigt,steigen,steigt,steigen},
  preterite = {stieg,stiegst,stieg,stiegen,stiegt,stiegen},
  subjunctive = {},
  imperative = {steige (du)!,steigen wir!,steigt (ihr)!,steigen Sie!},
  perfect = {sein + \textbf{gestiegen}},
  future = {werden + \textbf{steigen}},
  pluperfect = {waren + \textbf{gestiegen}},
  prepositions = {
    & \textbf{steigen auf} + Akk & (to climb onto / rise to)\\
    & \textbf{steigen in} + Akk & (to climb into / rise into)\\
    & \textbf{steigen aus} + Dat & (to climb out of)\\
  },
  separable = {
    & \textbf{auf$|$steigen} & (to take off / rise up)\\
    & \textbf{ein$|$steigen} & (to get on / board)\\
    & \textbf{aus$|$steigen} & (to get off / disembark)\\
    & \textbf{ab$|$steigen} & (to descend / go down)\\
    & \textbf{um$|$steigen} & (to transfer)\\
  }
}

% Migrated from verbs/s.tex:33
\VerbTable{
  name = {stellen},
  english = {to place},
  type = {regular, + Akkusativ},
  auxiliary = {haben},
  style = {acc},
  present = {stelle,stellst,stellt,stellen,stellt,stellen},
  preterite = {stellte,stelltest,stellte,stellten,stelltet,stellten},
  subjunctive = {},
  imperative = {stell(e) (du)!,stellen wir!,stellt (ihr)!,stellen Sie!},
  perfect = {haben + \textbf{gestellt}},
  future = {werden + \textbf{stellen}},
  pluperfect = {hatten + \textbf{gestellt}},
  prepositions = {
    &\textbf{sich stellen auf} + Akk & (to prepare for / brace for)\\
    & \textbf{sich stellen gegan} + Akk & (to oppose)\\
    & \textbf{sich stellen vor} + Akk & (to side with)\\
    &\textbf{sich stellen zu} + Dat & (to stand in front of / defend)\\
  },
  separable = {
    & \textbf{auf$|$stellen} & (to set up / put up)\\
    &\textbf{hin$|$stellen} & (to put (something) there)\\
    &\textbf{um$|$stellen} & (to rearrange)\\
    &\textbf{vor$|$stellen} & (to introduce)\\
    &\textbf{ver$|$stellen} & (to block)\\
    &\textbf{ein$|$stellen} & (to adjust, hire)\\
    &\textbf{her$|$stellen} & (to produce)\\
    &\textbf{dar$|$stellen} & (to portray)\\
    &\textbf{fest$|$stellen} & (to determine, establish)\\
    &\textbf{zurück$|$stellen} & (to postpone)
  }
}

% Migrated from verbs/b1/s.tex:25
\VerbTable{
  name = {sterben},
  english = {to die},
  type = {strong / irregular},
  auxiliary = {sein},
  style = {default},
  present = {sterbe,stirbst,stirbt,sterben,sterbt,sterben},
  preterite = {starb,starbst,starb,starben,starbt,starben},
  subjunctive = {},
  imperative = {stirb (du)!,sterben wir!,sterbt (ihr)!,sterben Sie!},
  perfect = {sein + \textbf{gestorben}},
  future = {werden + \textbf{sterben}},
  pluperfect = {waren + \textbf{gestorben}},
  prepositions = {---},
  separable = {---}
}

% Migrated from verbs/s.tex:34
\VerbTable{
  name = {stimmen},
  english = {to be correct / to vote},
  type = {regular},
  auxiliary = {haben},
  style = {dat},
  present = {stimme,stimmst,stimmt,stimmen,stimmt,stimmen},
  preterite = {stimmte,stimmtest,stimmte,stimmten,stimmtet,stimmten},
  subjunctive = {},
  imperative = {stimm(e) (du)!,stimmen wir!,stimmt (ihr)!,stimmen Sie!},
  perfect = {haben + \textbf{gestimmt}},
  future = {werden + \textbf{stimmen}},
  pluperfect = {hatten + \textbf{gestimmt}},
  prepositions = {
    & \textbf{stimmen für} + Akk & (to vote for)\\
    & \textbf{stimmen gegen} + Akk & (to vote against)\\
  },
  separable = {
    & \textbf{ab$|$stimmen} & (to vote / coordinate)\\
    & \textbf{überein$|$stimmen} & (to agree / correspond)\\
  }
}

% Migrated from verbs/b1/s.tex:26
\VerbTable{
  name = {stinken},
  english = {to stink},
  type = {strong / irregular},
  auxiliary = {haben},
  style = {default},
  present = {stinke,stinkst,stinkt,stinken,stinkt,stinken},
  preterite = {stank,stankst,stank,stanken,stankt,stanken},
  subjunctive = {},
  imperative = {stink (du)!,stinken wir!,stinkt (ihr)!,stinken Sie!},
  perfect = {haben + \textbf{gestunken}},
  future = {werden + \textbf{stinken}},
  pluperfect = {hatten + \textbf{gestunken}},
  prepositions = {---},
  separable = {---}
}

% Migrated from verbs/b1/s.tex:27
\VerbTable{
  name = {stören},
  english = {to disturb},
  type = {regular},
  auxiliary = {haben},
  style = {default},
  present = {störe,störst,stört,stören,stört,stören},
  preterite = {störte,störtest,störte,störten,störtet,störten},
  subjunctive = {},
  imperative = {stör (du)!,stören wir!,stört (ihr)!,stören Sie!},
  perfect = {haben + \textbf{gestört}},
  future = {werden + \textbf{stören}},
  pluperfect = {hatten + \textbf{gestört}},
  prepositions = {---},
  separable = {---}
}

% Migrated from verbs/s.tex:35
\VerbTable{
  name = {stoppen},
  english = {to stop},
  type = {regular},
  auxiliary = {haben},
  style = {default},
  present = {stoppe,stoppst,stoppt,stoppen,stoppt,stoppen},
  preterite = {stoppte,stopptest,stoppte,stoppten,stopptet,stoppten},
  subjunctive = {},
  imperative = {stopp(e) (du)!,stoppen wir!,stoppt (ihr)!,stoppen Sie!},
  perfect = {haben + \textbf{gestoppt}},
  future = {werden + \textbf{stoppen}},
  pluperfect = {hatten + \textbf{gestoppt}},
  prepositions = {
    & \textbf{stoppen vor} + Dat & (to stop in front of)\\
    & \textbf{stoppen bei} + Dat & (to stop at)\\
  },
  separable = {& \textbf{an$|$stoppen} & (to stop / halt briefly)\\}
}

% Migrated from verbs/b1/s.tex:28
\VerbTable{
  name = {stoßen},
  english = {to bump},
  type = {strong / irregular},
  auxiliary = {haben},
  style = {default},
  present = {stoße,stößt,stößt,stoßen,stoßt,stoßen},
  preterite = {stieß,stießt,stieß,stießen,stießt,stießen},
  subjunctive = {},
  imperative = {stoß (du)!,stoßen wir!,stoßt (ihr)!,stoßen Sie!},
  perfect = {haben + \textbf{gestoßen}},
  future = {werden + \textbf{stoßen}},
  pluperfect = {hatten + \textbf{gestoßen}},
  prepositions = {---},
  separable = {---}
}

% Migrated from verbs/s.tex:36
\VerbTable{
  name = {streichen},
  english = {to paint / stroke},
  type = {irregular},
  auxiliary = {haben},
  style = {default},
  present = {streiche,streichst,streicht,streichen,streicht,streichen},
  preterite = {strich,strichst,strich,strichen,stricht,strichen},
  subjunctive = {},
  imperative = {streiche (du)!,streichen wir!,streicht (ihr)!,streichen Sie!},
  perfect = {haben + \textbf{gestrichen}},
  future = {werden + \textbf{streichen}},
  pluperfect = {hatten + \textbf{gestrichen}},
  prepositions = {
    & \textbf{streichen über} + Akk & (to stroke over / to paint over)\\
    & \textbf{streichen mit} + Dat & (to paint / smear with)\\
  },
  separable = {
    & \textbf{ab$|$streichen} & (to cancel / remove / coat off)\\
    & \textbf{ein$|$streichen} & (to apply / cover in)\\
    & \textbf{auf$|$streichen} & (to spread / coat onto)\\
  }
}

% Migrated from verbs/b1/s.tex:29
\VerbTable{
  name = {streiken},
  english = {to strike},
  type = {regular},
  auxiliary = {haben},
  style = {default},
  present = {streike,streikst,streikt,streiken,streikt,streiken},
  preterite = {streikte,streiktest,streikte,streikten,streiktet,streikten},
  subjunctive = {},
  imperative = {streik (du)!,streiken wir!,streikt (ihr)!,streiken Sie!},
  perfect = {haben + \textbf{gestreikt}},
  future = {werden + \textbf{streiken}},
  pluperfect = {hatten + \textbf{gestreikt}},
  prepositions = {---},
  separable = {---}
}

% Migrated from verbs/b1/s.tex:30
\VerbTable{
  name = {streiten},
  english = {to argue},
  type = {strong / irregular},
  auxiliary = {haben},
  style = {default},
  present = {streite,streitest,streitet,streiten,streitet,streiten},
  preterite = {stritt,strittest,stritt,stritten,strittet,stritten},
  subjunctive = {},
  imperative = {streite (du)!,streiten wir!,streitet (ihr)!,streiten Sie!},
  perfect = {haben + \textbf{gestritten}},
  future = {werden + \textbf{streiten}},
  pluperfect = {hatten + \textbf{gestritten}},
  prepositions = {---},
  separable = {---}
}

% Migrated from verbs/s.tex:37
\VerbTable{
  name = {studieren},
  english = {to study},
  type = {regular},
  auxiliary = {haben},
  style = {default},
  present = {studiere,studierst,studiert,studieren,studiert,studieren},
  preterite = {studierte,studiertest,studierte,studierten,studiertet,studierten},
  subjunctive = {},
  imperative = {studier(e) (du)!,studieren wir!,studiert (ihr)!,studieren Sie!},
  perfect = {haben + \textbf{studiert}},
  future = {werden + \textbf{studieren}},
  pluperfect = {hatten + \textbf{studiert}},
  prepositions = {---},
  separable = {---}
}

% Migrated from verbs/b1/s.tex:31
\VerbTable{
  name = {stürzen},
  english = {to fall / crash},
  type = {regular},
  auxiliary = {sein},
  style = {default},
  present = {stürze,stürzt,stürzt,stürzen,stürzt,stürzen},
  preterite = {stürzte,stürztest,stürzte,stürzten,stürztet,stürzten},
  subjunctive = {},
  imperative = {stürz (du)!,stürzen wir!,stürzt (ihr)!,stürzen Sie!},
  perfect = {sein + \textbf{gestürzt}},
  future = {werden + \textbf{stürzen}},
  pluperfect = {waren + \textbf{gestürzt}},
  prepositions = {---},
  separable = {---}
}

% Migrated from verbs/s.tex:38
\VerbTable{
  name = {suchen},
  english = {to search},
  type = {regular, + Akkusativ},
  auxiliary = {haben},
  style = {acc},
  present = {suche,suchst,sucht,suchen,sucht,suchen},
  preterite = {suchte,suchtest,suchte,suchten,suchtet,suchten},
  subjunctive = {},
  imperative = {such(e) (du)!,suchen wir!,sucht (ihr)!,suchen Sie!},
  perfect = {haben + \textbf{gesucht}},
  future = {werden + \textbf{suchen}},
  pluperfect = {hatten + \textbf{gesucht}},
  prepositions = {---},
  separable = {---}
}

\section{T}

% Migrated from verbs/b1/t.tex:1
\VerbTable{
  name = {tanken},
  english = {to refuel},
  type = {regular},
  auxiliary = {haben},
  style = {default},
  present = {tanke,tankst,tankt,tanken,tankt,tanken},
  preterite = {tankte,tanktest,tankte,tankten,tanktet,tankten},
  subjunctive = {},
  imperative = {tank (du)!,tanken wir!,tankt (ihr)!,tanken Sie!},
  perfect = {haben + \textbf{getankt}},
  future = {werden + \textbf{tanken}},
  pluperfect = {hatten + \textbf{getankt}},
  prepositions = {---},
  separable = {---}
}

% Migrated from verbs/b1/t.tex:2
\VerbTable{
  name = {tanzen},
  english = {to dance},
  type = {regular},
  auxiliary = {haben},
  style = {default},
  present = {tanze,tanzt,tanzt,tanzen,tanzt,tanzen},
  preterite = {tanzte,tanztest,tanzte,tanzten,tanztet,tanzten},
  subjunctive = {},
  imperative = {tanz (du)!,tanzen wir!,tanzt (ihr)!,tanzen Sie!},
  perfect = {haben + \textbf{getanzt}},
  future = {werden + \textbf{tanzen}},
  pluperfect = {hatten + \textbf{getanzt}},
  prepositions = {---},
  separable = {---}
}

% Migrated from verbs/b1/t.tex:3
\VerbTable{
  name = {tauchen},
  english = {to dive},
  type = {regular},
  auxiliary = {haben / sein},
  style = {default},
  present = {tauche,tauchst,taucht,tauchen,taucht,tauchen},
  preterite = {tauchte,tauchtest,tauchte,tauchten,tauchtet,tauchten},
  subjunctive = {},
  imperative = {tauch (du)!,tauchen wir!,taucht (ihr)!,tauchen Sie!},
  perfect = {haben / sein + \textbf{getaucht}},
  future = {werden + \textbf{tauchen}},
  pluperfect = {hatten / waren + \textbf{getaucht}},
  prepositions = {---},
  separable = {---}
}

% Migrated from verbs/t.tex:1
\VerbTable{
  name = {tauschen},
  english = {to exchange / swap},
  type = {regular},
  auxiliary = {haben},
  style = {default},
  present = {tausche,tauschst,tauscht,tauschen,tauscht,tauschen},
  preterite = {tauschte,tauschtest,tauschte,tauschten,tauschtet,tauschten},
  subjunctive = {},
  imperative = {tausche (du)!,tauschen wir!,tauscht (ihr)!,tauschen Sie!},
  perfect = {haben + \textbf{getauscht}},
  future = {werden + \textbf{tauschen}},
  pluperfect = {hatten + \textbf{getauscht}},
  prepositions = {
    & \textbf{tauschen gegen} + Akk & (to exchange for)\\
    & \textbf{tauschen mit} + Dat & (to exchange with)\\
  },
  separable = {
    & \textbf{aus$|$tauschen} & (to exchange / replace)\\
    & \textbf{um$|$tauschen} & (to return (a purchased item))\\
  }
}

% Migrated from verbs/t.tex:2
\VerbTable{
  name = {teilen},
  english = {to share / divide},
  type = {regular},
  auxiliary = {haben},
  style = {default},
  present = {teile,teilst,teilt,teilen,teilt,teilen},
  preterite = {teilte,teiltest,teilte,teilten,teiltet,teilten},
  subjunctive = {},
  imperative = {teil(e) (du)!,teilen wir!,teilt (ihr)!,teilen Sie!},
  perfect = {haben + \textbf{geteilt}},
  future = {werden + \textbf{teilen}},
  pluperfect = {hatten + \textbf{geteilt}},
  prepositions = {
    & \textbf{teilen mit} + Dat & (to share with)\\
    & \textbf{teilen in} + Akk & (to divide into)\\
  },
  separable = {
    & \textbf{auf$|$teilen} & (to split up / divide)\\
    & \textbf{ab$|$teilen} & (to allocate / section off)\\
  }
}

% Migrated from verbs/b1/t.tex:4
\VerbTable{
  name = {teilnehmen},
  english = {to take part},
  type = {separable, strong / irregular},
  auxiliary = {haben},
  style = {default},
  present = {nehme teil,nimmst teil,nimmt teil,nehmen teil,nehmt teil,nehmen teil},
  preterite = {nahm teil,nahmst teil,nahm teil,nahmen teil,nahmt teil,nahmen teil},
  subjunctive = {},
  imperative = {nimm (du) teil!,nehmen wir teil!,nehmt (ihr) teil!,nehmen Sie teil!},
  perfect = {haben + \textbf{teilgenommen}},
  future = {werden + \textbf{teilnehmen}},
  pluperfect = {hatten + \textbf{teilgenommen}},
  prepositions = {---},
  separable = {---}
}

% Migrated from verbs/b1/t.tex:5
\VerbTable{
  name = {telefonieren},
  english = {to make a telephone call},
  type = {regular},
  auxiliary = {haben},
  style = {default},
  present = {telefoniere,telefonierst,telefoniert,telefonieren,telefoniert,telefonieren},
  preterite = {telefonierte,telefoniertest,telefonierte,telefonierten,telefoniertet,telefonierten},
  subjunctive = {},
  imperative = {telefonier (du)!,telefonieren wir!,telefoniert (ihr)!,telefonieren Sie!},
  perfect = {haben + \textbf{telefoniert}},
  future = {werden + \textbf{telefonieren}},
  pluperfect = {hatten + \textbf{telefoniert}},
  prepositions = {---},
  separable = {---}
}

% Migrated from verbs/b1/t.tex:6
\VerbTable{
  name = {testen},
  english = {to test},
  type = {regular},
  auxiliary = {haben},
  style = {default},
  present = {teste,testest,testet,testen,testet,testen},
  preterite = {testete,testetest,testete,testeten,testetet,testeten},
  subjunctive = {},
  imperative = {teste (du)!,testen wir!,testet (ihr)!,testen Sie!},
  perfect = {haben + \textbf{getestet}},
  future = {werden + \textbf{testen}},
  pluperfect = {hatten + \textbf{getestet}},
  prepositions = {---},
  separable = {---}
}

% Migrated from verbs/b1/t.tex:7
\VerbTable{
  name = {tippen},
  english = {to type / tap},
  type = {regular},
  auxiliary = {haben},
  style = {default},
  present = {tippe,tippst,tippt,tippen,tippt,tippen},
  preterite = {tippte,tipptest,tippte,tippten,tipptet,tippten},
  subjunctive = {},
  imperative = {tipp (du)!,tippen wir!,tippt (ihr)!,tippen Sie!},
  perfect = {haben + \textbf{getippt}},
  future = {werden + \textbf{tippen}},
  pluperfect = {hatten + \textbf{getippt}},
  prepositions = {---},
  separable = {---}
}

% Migrated from verbs/t.tex:5
\VerbTable{
  name = {träumen},
  english = {to dream},
  type = {regular},
  auxiliary = {haben},
  style = {default},
  present = {träume,träumst,träumt,träumen,träumt,träumen},
  preterite = {träumte,träumtest,träumte,träumten,träumtet,träumten},
  subjunctive = {},
  imperative = {träum(e) (du)!,träumen wir!,träumt (ihr)!,träumen Sie!},
  perfect = {haben + \textbf{geträumt}},
  future = {werden + \textbf{träumen}},
  pluperfect = {hatten + \textbf{geträumt}},
  prepositions = {---},
  separable = {---}
}

% Migrated from verbs/t.tex:3
\VerbTable{
  name = {tragen},
  english = {to wear, carry},
  type = {irregular, + Akkusativ},
  auxiliary = {haben},
  style = {acc},
  present = {trage,trägst,trägt,tragen,tragt,tragen},
  preterite = {trug,trugst,trug,trugen,trugt,trugen},
  subjunctive = {},
  imperative = {trag(e) (du)!,tragen wir!,tragt (ihr)!,tragen Sie!},
  perfect = {haben + \textbf{getragen}},
  future = {werden + \textbf{tragen}},
  pluperfect = {hatten + \textbf{getragen}},
  prepositions = {---},
  separable = {---}
}

% Migrated from verbs/t.tex:4
\VerbTable{
  name = {trainieren},
  english = {to train / practice},
  type = {regular},
  auxiliary = {haben},
  style = {default},
  present = {trainiere,trainierst,trainiert,trainieren,trainiert,trainieren},
  preterite = {trainierte,trainiertest,trainierte,trainierten,trainiertet,trainierten},
  subjunctive = {},
  imperative = {trainier(e) (du)!,trainieren wir!,trainiert (ihr)!,trainieren Sie!},
  perfect = {haben + \textbf{trainiert}},
  future = {werden + \textbf{trainieren}},
  pluperfect = {hatten + \textbf{trainiert}},
  prepositions = {
    & \textbf{trainieren für} + Akk & (to train for)\\
    & \textbf{trainieren mit} + Dat & (to train with)\\
    & \textbf{trainieren gegen} + Akk & (to train against / play against)\\
  },
  separable = {& \textbf{ab$|$trainieren} & (to train off / reduce through training)\\},
  bottom-layout = {stacked}
}

% Migrated from verbs/b1/t.tex:8
\VerbTable{
  name = {transportieren},
  english = {to transport},
  type = {regular},
  auxiliary = {haben},
  style = {default},
  present = {transportiere,transportierst,transportiert,transportieren,transportiert,transportieren},
  preterite = {transportierte,transportiertest,transportierte,transportierten,transportiertet,transportierten},
  subjunctive = {},
  imperative = {transportier (du)!,transportieren wir!,transportiert (ihr)!,transportieren Sie!},
  perfect = {haben + \textbf{transportiert}},
  future = {werden + \textbf{transportieren}},
  pluperfect = {hatten + \textbf{transportiert}},
  prepositions = {---},
  separable = {---}
}

% Migrated from verbs/t.tex:8
\VerbTable{
  name = {treffen},
  english = {to meet},
  type = {irregular, + Akkusativ},
  auxiliary = {haben},
  style = {acc},
  present = {treffe,triffst,trifft,treffen,trefft,treffen},
  preterite = {traf,trafst,traf,trafen,traft,trafen},
  subjunctive = {},
  imperative = {triff (du)!,treffen wir!,trefft (ihr)!,treffen Sie!},
  perfect = {haben + \textbf{getroffen}},
  future = {werden + \textbf{treffen}},
  pluperfect = {hatten + \textbf{getroffen}},
  prepositions = {---},
  separable = {---}
}

% Migrated from verbs/b1/t.tex:9
\VerbTable{
  name = {treiben},
  english = {to drive / do},
  type = {strong / irregular},
  auxiliary = {haben},
  style = {default},
  present = {treibe,treibst,treibt,treiben,treibt,treiben},
  preterite = {trieb,triebst,trieb,trieben,triebt,trieben},
  subjunctive = {},
  imperative = {treib (du)!,treiben wir!,treibt (ihr)!,treiben Sie!},
  perfect = {haben + \textbf{getrieben}},
  future = {werden + \textbf{treiben}},
  pluperfect = {hatten + \textbf{getrieben}},
  prepositions = {---},
  separable = {---}
}

% Migrated from verbs/t.tex:6
\VerbTable{
  name = {treten},
  english = {to step},
  type = {irregular},
  auxiliary = {sein},
  style = {default},
  present = {trete,trittst,tritt,treten,tretet,treten},
  preterite = {trat,tratst,trat,traten,tratet,traten},
  subjunctive = {},
  imperative = {tritt (du)!,treten wir!,tretet (ihr)!,treten Sie!},
  perfect = {sein + \textbf{getreten}},
  future = {werden + \textbf{treten}},
  pluperfect = {waren + \textbf{getreten}},
  prepositions = {---},
  separable = {---}
}

% Migrated from verbs/t.tex:7
\VerbTable{
  name = {trinken},
  english = {to drink},
  type = {irregular, + Akkusativ},
  auxiliary = {haben},
  style = {acc},
  present = {trinke,trinkst,trinkt,trinken,trinkt,trinken},
  preterite = {trank,trankest,trank,tranken,trankt,tranken},
  subjunctive = {},
  imperative = {trink(e) (du)!,trinken wir!,trinkt (ihr)!,trinken Sie!},
  perfect = {haben + \textbf{getrunken}},
  future = {werden + \textbf{trinken}},
  pluperfect = {hatten + \textbf{getrunken}},
  prepositions = {---},
  separable = {---}
}

% Migrated from verbs/b1/t.tex:10
\VerbTable{
  name = {trocknen},
  english = {to dry},
  type = {regular},
  auxiliary = {haben / sein},
  style = {default},
  present = {trockne,trocknest,trocknet,trocknen,trocknet,trocknen},
  preterite = {trocknete,trocknetest,trocknete,trockneten,trocknetet,trockneten},
  subjunctive = {},
  imperative = {trockne (du)!,trocknen wir!,trocknet (ihr)!,trocknen Sie!},
  perfect = {haben / sein + \textbf{getrocknet}},
  future = {werden + \textbf{trocknen}},
  pluperfect = {hatten / waren + \textbf{getrocknet}},
  prepositions = {---},
  separable = {---}
}

% Migrated from verbs/t.tex:9
\VerbTable{
  name = {trösten},
  english = {to comfort},
  type = {regular},
  auxiliary = {haben},
  style = {default},
  present = {tröste,tröstest,tröstet,trösten,tröstet,trösten},
  preterite = {tröstete,tröstetest,tröstete,trösteten,tröstetet,trösteten},
  subjunctive = {},
  imperative = {tröste (du)!,trösten wir!,tröstet (ihr)!,trösten Sie!},
  perfect = {haben + \textbf{getröstet}},
  future = {werden + \textbf{trösten}},
  pluperfect = {hatten + \textbf{getröstet}},
  prepositions = {
    & \textbf{trösten mit} + Dat & (to comfort with)\\
    & \textbf{trösten über} + Akk & (to console about)\\
  },
  separable = {---}
}

% Migrated from verbs/t.tex:10
\VerbTable{
  name = {tun},
  english = {to do},
  type = {irregular, + Akkusativ},
  auxiliary = {haben},
  style = {acc},
  present = {tue,tust,tut,tun,tut,tun},
  preterite = {tat,tatest,tat,taten,tatet,taten},
  subjunctive = {},
  imperative = {tu(e) (du)!,tun wir!,tut (ihr)!,tun Sie!},
  perfect = {haben + \textbf{getan}},
  future = {werden + \textbf{tun}},
  pluperfect = {hatten + \textbf{getan}},
  prepositions = {
    & \textbf{tun für} + Akk & (to do for)\\
    & \textbf{tun gegen} + Akk & (to do against / about)\\
  },
  separable = {
    & \textbf{mit$|$tun} & (to participate / be involved)\\
    & \textbf{ab$|$tun} & (to dismiss / remove)\\
    & \textbf{weg$|$tun} & (to put away / remove)\\
    & \textbf{hinzu$|$tun} & (to add)\\
    & \textbf{weh$|$tun} & (to hurt / cause pain)\\
  }
}

\section{U}

% Migrated from verbs/u.tex:1
\VerbTable{
  name = {üben},
  english = {to practice},
  type = {regular},
  auxiliary = {haben},
  style = {default},
  present = {übe,übst,übt,üben,übt,üben},
  preterite = {übte,übtest,übte,übten,übtet,übten},
  subjunctive = {},
  imperative = {üb(e) (du)!,üben wir!,übt (ihr)!,üben Sie!},
  perfect = {haben + \textbf{geübt}},
  future = {werde + \textbf{üben}},
  pluperfect = {hatten + \textbf{geübt}},
  prepositions = {
    & \textbf{üben für} + Akk & (to practice for)\\
    & \textbf{üben mit} + Dat & (to practice with)\\
  },
  separable = {
    & \textbf{ein$|$üben} & (to rehearse / practice)\\
    & \textbf{aus$|$üben} & (to exercise / carry out)\\
  }
}

% Migrated from verbs/b1/u.tex:1
\VerbTable{
  name = {überfahren},
  english = {to run over},
  type = {strong / irregular},
  auxiliary = {haben},
  style = {default},
  present = {überfahre,überfährst,überfährt,überfahren,überfahrt,überfahren},
  preterite = {überfuhr,überfuhrst,überfuhr,überfuhren,überfuhrt,überfuhren},
  subjunctive = {},
  imperative = {überfahr (du)!,überfahren wir!,überfahrt (ihr)!,überfahren Sie!},
  perfect = {haben + \textbf{überfahren}},
  future = {werden + \textbf{überfahren}},
  pluperfect = {hatten + \textbf{überfahren}},
  prepositions = {---},
  separable = {---}
}

% Migrated from verbs/b1/u.tex:2
\VerbTable{
  name = {überholen},
  english = {to overtake},
  type = {regular},
  auxiliary = {haben},
  style = {default},
  present = {überhole,überholst,überholt,überholen,überholt,überholen},
  preterite = {überholte,überholtest,überholte,überholten,überholtet,überholten},
  subjunctive = {},
  imperative = {überhol (du)!,überholen wir!,überholt (ihr)!,überholen Sie!},
  perfect = {haben + \textbf{überholt}},
  future = {werden + \textbf{überholen}},
  pluperfect = {hatten + \textbf{überholt}},
  prepositions = {---},
  separable = {---}
}

% Migrated from verbs/u.tex:2
\VerbTable{
  name = {überlegen},
  english = {to think about},
  type = {irregular, + Akkusativ},
  auxiliary = {haben},
  style = {acc},
  present = {überlege,überlegst,überlegt,überlegen,überlegt,überlegen},
  preterite = {überlegte,überlegtest,überlegte,überlegten,überlegtet,überlegten},
  subjunctive = {},
  imperative = {überleg(e) (du)!,überlegen wir!,überlegt (ihr)!,überlegen Sie!},
  perfect = {haben + \textbf{überlegt}},
  future = {werden + \textbf{überlegen}},
  pluperfect = {hatten + \textbf{überlegt}},
  prepositions = {---},
  separable = {---}
}

% Migrated from verbs/b1/u.tex:3
\VerbTable{
  name = {übernachten},
  english = {to stay overnight},
  type = {regular},
  auxiliary = {haben},
  style = {default},
  present = {übernachte,übernachtest,übernachtet,übernachten,übernachtet,übernachten},
  preterite = {übernachtete,übernachtetest,übernachtete,übernachteten,übernachtetet,übernachteten},
  subjunctive = {},
  imperative = {übernachte (du)!,übernachten wir!,übernachtet (ihr)!,übernachten Sie!},
  perfect = {haben + \textbf{übernachtet}},
  future = {werden + \textbf{übernachten}},
  pluperfect = {hatten + \textbf{übernachtet}},
  prepositions = {---},
  separable = {---}
}

% Migrated from verbs/u.tex:3
\VerbTable{
  name = {übernehmen},
  english = {to take over / to assume},
  type = {irregular},
  auxiliary = {haben},
  style = {default},
  present = {übernehme,übernimmst,übernimmt,übernehmen,übernehmt,übernehmen},
  preterite = {übernahm,übernahmst,übernahm,übernahmen,übernahmt,übernahmen},
  subjunctive = {},
  imperative = {übernimm (du)!,übernehmen wir!,übernehmt (ihr)!,übernehmen Sie!},
  perfect = {haben + \textbf{übernommen}},
  future = {werden + \textbf{übernehmen}},
  pluperfect = {hatten + \textbf{übernommen}},
  prepositions = {
    & \textbf{übernehmen von} + Dat & (to take over from)\\
    & \textbf{übernehmen für} + Akk & (to assume responsibility for)\\
  },
  separable = {---}
}

% Migrated from verbs/b1/u.tex:4
\VerbTable{
  name = {überprüfen},
  english = {to check},
  type = {regular},
  auxiliary = {haben},
  style = {default},
  present = {überprüfe,überprüfst,überprüft,überprüfen,überprüft,überprüfen},
  preterite = {überprüfte,überprüftest,überprüfte,überprüften,überprüftet,überprüften},
  subjunctive = {},
  imperative = {überprüf (du)!,überprüfen wir!,überprüft (ihr)!,überprüfen Sie!},
  perfect = {haben + \textbf{überprüft}},
  future = {werden + \textbf{überprüfen}},
  pluperfect = {hatten + \textbf{überprüft}},
  prepositions = {---},
  separable = {---}
}

% Migrated from verbs/b1/u.tex:5
\VerbTable{
  name = {überqueren},
  english = {to cross},
  type = {regular},
  auxiliary = {haben},
  style = {default},
  present = {überquere,überquerst,überquert,überqueren,überquert,überqueren},
  preterite = {überquerte,überquertest,überquerte,überquerten,überquertet,überquerten},
  subjunctive = {},
  imperative = {überquer (du)!,überqueren wir!,überquert (ihr)!,überqueren Sie!},
  perfect = {haben + \textbf{überquert}},
  future = {werden + \textbf{überqueren}},
  pluperfect = {hatten + \textbf{überquert}},
  prepositions = {---},
  separable = {---}
}

% Migrated from verbs/b1/u.tex:6
\VerbTable{
  name = {überraschen},
  english = {to surprise},
  type = {regular},
  auxiliary = {haben},
  style = {default},
  present = {überrasche,überraschst,überrascht,überraschen,überrascht,überraschen},
  preterite = {überraschte,überraschtest,überraschte,überraschten,überraschtet,überraschten},
  subjunctive = {},
  imperative = {überrasch (du)!,überraschen wir!,überrascht (ihr)!,überraschen Sie!},
  perfect = {haben + \textbf{überrascht}},
  future = {werden + \textbf{überraschen}},
  pluperfect = {hatten + \textbf{überrascht}},
  prepositions = {---},
  separable = {---}
}

% Migrated from verbs/b1/u.tex:7
\VerbTable{
  name = {überreden},
  english = {to persuade},
  type = {regular},
  auxiliary = {haben},
  style = {default},
  present = {überrede,überredest,überredet,überreden,überredet,überreden},
  preterite = {überredete,überredetest,überredete,überredeten,überredetet,überredeten},
  subjunctive = {},
  imperative = {überrede (du)!,überreden wir!,überredet (ihr)!,überreden Sie!},
  perfect = {haben + \textbf{überredet}},
  future = {werden + \textbf{überreden}},
  pluperfect = {hatten + \textbf{überredet}},
  prepositions = {---},
  separable = {---}
}

% Migrated from verbs/b1/u.tex:8
\VerbTable{
  name = {übersetzen},
  english = {to translate},
  type = {regular},
  auxiliary = {haben},
  style = {default},
  present = {übersetze,übersetzt,übersetzt,übersetzen,übersetzt,übersetzen},
  preterite = {übersetzte,übersetztest,übersetzte,übersetzten,übersetztet,übersetzten},
  subjunctive = {},
  imperative = {übersetz (du)!,übersetzen wir!,übersetzt (ihr)!,übersetzen Sie!},
  perfect = {haben + \textbf{übersetzt}},
  future = {werden + \textbf{übersetzen}},
  pluperfect = {hatten + \textbf{übersetzt}},
  prepositions = {---},
  separable = {---}
}

% Migrated from verbs/b1/u.tex:9
\VerbTable{
  name = {übertreiben},
  english = {to exaggerate},
  type = {strong / irregular},
  auxiliary = {haben},
  style = {default},
  present = {übertreibe,übertreibst,übertreibt,übertreiben,übertreibt,übertreiben},
  preterite = {übertrieb,übertriebst,übertrieb,übertrieben,übertriebt,übertrieben},
  subjunctive = {},
  imperative = {übertreib (du)!,übertreiben wir!,übertreibt (ihr)!,übertreiben Sie!},
  perfect = {haben + \textbf{übertrieben}},
  future = {werden + \textbf{übertreiben}},
  pluperfect = {hatten + \textbf{übertrieben}},
  prepositions = {---},
  separable = {---}
}

% Migrated from verbs/u.tex:4
\VerbTable{
  name = {überweisen},
  english = {to transfer money / refer},
  type = {regular},
  auxiliary = {haben},
  style = {default},
  present = {überweise,überweist,überweist,überweisen,überweist,überweisen},
  preterite = {überwies,überwiesest,überwies,überwiesen,überwieset,überwiesen},
  subjunctive = {},
  imperative = {überweise (du)!,überweisen wir!,überweist (ihr)!,überweisen Sie!},
  perfect = {haben + \textbf{überwiesen}},
  future = {werden + \textbf{überweisen}},
  pluperfect = {hatten + \textbf{überwiesen}},
  prepositions = {& \textbf{überweisen an} + Akk & (to transfer to / refer to)\\},
  separable = {---}
}

% Migrated from verbs/b1/u.tex:10
\VerbTable{
  name = {überzeugen},
  english = {to convince},
  type = {regular},
  auxiliary = {haben},
  style = {default},
  present = {überzeuge,überzeugst,überzeugt,überzeugen,überzeugt,überzeugen},
  preterite = {überzeugte,überzeugtest,überzeugte,überzeugten,überzeugtet,überzeugten},
  subjunctive = {},
  imperative = {überzeug (du)!,überzeugen wir!,überzeugt (ihr)!,überzeugen Sie!},
  perfect = {haben + \textbf{überzeugt}},
  future = {werden + \textbf{überzeugen}},
  pluperfect = {hatten + \textbf{überzeugt}},
  prepositions = {---},
  separable = {---}
}

% Migrated from verbs/u.tex:5
\VerbTable{
  name = {umarmen},
  english = {to hug},
  type = {regular, + Akkusativ},
  auxiliary = {haben},
  style = {acc},
  present = {umarme,umarmst,umarmt,umarmen,umarmt,umarmen},
  preterite = {umarmte,umarmtest,umarmte,umarmten,umarmtet,umarmten},
  subjunctive = {},
  imperative = {umarme (du)!,umarmen wir!,umarmt (ihr)!,umarmen Sie!},
  perfect = {haben + \textbf{umarmt}},
  future = {werden + \textbf{umarmen}},
  pluperfect = {hatten + \textbf{umarmt}},
  prepositions = {---},
  separable = {---}
}

% Migrated from verbs/b1/u.tex:11
\VerbTable{
  name = {umdrehen},
  english = {to turn around},
  type = {separable, regular},
  auxiliary = {haben},
  style = {default},
  present = {drehe um,drehst um,dreht um,drehen um,dreht um,drehen um},
  preterite = {drehte um,drehtest um,drehte um,drehten um,drehtet um,drehten um},
  subjunctive = {},
  imperative = {dreh (du) um!,drehen wir um!,dreht (ihr) um!,drehen Sie um!},
  perfect = {haben + \textbf{umgedreht}},
  future = {werden + \textbf{umdrehen}},
  pluperfect = {hatten + \textbf{umgedreht}},
  prepositions = {---},
  separable = {---}
}

% Migrated from verbs/b1/u.tex:12
\VerbTable{
  name = {umgehen},
  english = {to deal with / go around},
  type = {separable, strong / irregular},
  auxiliary = {sein},
  style = {default},
  present = {gehe um,gehst um,geht um,gehen um,geht um,gehen um},
  preterite = {ging um,gingst um,ging um,gingen um,gingt um,gingen um},
  subjunctive = {},
  imperative = {geh (du) um!,gehen wir um!,geht (ihr) um!,gehen Sie um!},
  perfect = {sein + \textbf{umgegangen}},
  future = {werden + \textbf{umgehen}},
  pluperfect = {waren + \textbf{umgegangen}},
  prepositions = {---},
  separable = {---}
}

% Migrated from verbs/b1/u.tex:13
\VerbTable{
  name = {umsteigen},
  english = {to change trains / transfer},
  type = {separable, strong / irregular},
  auxiliary = {sein},
  style = {default},
  present = {steige um,steigst um,steigt um,steigen um,steigt um,steigen um},
  preterite = {stieg um,stiegst um,stieg um,stiegen um,stiegt um,stiegen um},
  subjunctive = {},
  imperative = {steig (du) um!,steigen wir um!,steigt (ihr) um!,steigen Sie um!},
  perfect = {sein + \textbf{umgestiegen}},
  future = {werden + \textbf{umsteigen}},
  pluperfect = {waren + \textbf{umgestiegen}},
  prepositions = {---},
  separable = {---}
}

% Migrated from verbs/b1/u.tex:14
\VerbTable{
  name = {umtauschen},
  english = {to exchange},
  type = {separable, regular},
  auxiliary = {haben},
  style = {default},
  present = {tausche um,tauschst um,tauscht um,tauschen um,tauscht um,tauschen um},
  preterite = {tauschte um,tauschtest um,tauschte um,tauschten um,tauschtet um,tauschten um},
  subjunctive = {},
  imperative = {tausch (du) um!,tauschen wir um!,tauscht (ihr) um!,tauschen Sie um!},
  perfect = {haben + \textbf{umgetauscht}},
  future = {werden + \textbf{umtauschen}},
  pluperfect = {hatten + \textbf{umgetauscht}},
  prepositions = {---},
  separable = {---}
}

% Migrated from verbs/b1/u.tex:15
\VerbTable{
  name = {umziehen},
  english = {to move house},
  type = {separable, strong / irregular},
  auxiliary = {sein},
  style = {default},
  present = {ziehe um,ziehst um,zieht um,ziehen um,zieht um,ziehen um},
  preterite = {zog um,zogst um,zog um,zogen um,zogt um,zogen um},
  subjunctive = {},
  imperative = {zieh (du) um!,ziehen wir um!,zieht (ihr) um!,ziehen Sie um!},
  perfect = {sein + \textbf{umgezogen}},
  future = {werden + \textbf{umziehen}},
  pluperfect = {waren + \textbf{umgezogen}},
  prepositions = {---},
  separable = {---}
}

% Migrated from verbs/b1/u.tex:17
\VerbTable{
  name = {unterbrechen},
  english = {to interrupt},
  type = {strong / irregular},
  auxiliary = {haben},
  style = {default},
  present = {unterbreche,unterbrichst,unterbricht,unterbrechen,unterbrecht,unterbrechen},
  preterite = {unterbrach,unterbrachst,unterbrach,unterbrachen,unterbracht,unterbrachen},
  subjunctive = {},
  imperative = {unterbrich (du)!,unterbrechen wir!,unterbrecht (ihr)!,unterbrechen Sie!},
  perfect = {haben + \textbf{unterbrochen}},
  future = {werden + \textbf{unterbrechen}},
  pluperfect = {hatten + \textbf{unterbrochen}},
  prepositions = {---},
  separable = {---}
}

% Migrated from verbs/u.tex:6
\VerbTable{
  name = {unterhalten},
  english = {to entertain},
  type = {irregular, + Akkusativ},
  auxiliary = {haben},
  style = {acc},
  present = {unterhalte,unterhältst,unterhält,unterhalten,unterhaltet,unterhalten},
  preterite = {unterhielt,unterhieltest,unterhielt,unterhielten,unterhieltet,unterhielten},
  subjunctive = {},
  imperative = {unterhalte (du)!,unterhalten wir!,unterhaltet (ihr)!,unterhalten Sie!},
  perfect = {haben + \textbf{unterhalten}},
  future = {werden + \textbf{unterhalten}},
  pluperfect = {hatten + \textbf{unterhalten}},
  prepositions = {---},
  separable = {---}
}

% Migrated from verbs/b1/u.tex:18
\VerbTable{
  name = {unterlassen},
  english = {to refrain from},
  type = {strong / irregular},
  auxiliary = {haben},
  style = {default},
  present = {unterlasse,unterlässt,unterlässt,unterlassen,unterlasst,unterlassen},
  preterite = {unterließ,unterließt,unterließ,unterließen,unterließt,unterließen},
  subjunctive = {},
  imperative = {unterlass (du)!,unterlassen wir!,unterlasst (ihr)!,unterlassen Sie!},
  perfect = {haben + \textbf{unterlassen}},
  future = {werden + \textbf{unterlassen}},
  pluperfect = {hatten + \textbf{unterlassen}},
  prepositions = {---},
  separable = {---}
}

% Migrated from verbs/u.tex:7
\VerbTable{
  name = {unternehmen},
  english = {to undertake / to do},
  type = {irregular},
  auxiliary = {haben},
  style = {default},
  present = {unternehme,unternimmst,unternimmt,unternehmen,unternehmt,unternehmen},
  preterite = {unternahm,unternahmst,unternahm,unternahmen,unternahmt,unternahmen},
  subjunctive = {},
  imperative = {unternimm (du)!,unternehmen wir!,unternehmt (ihr)!,unternehmen Sie!},
  perfect = {haben + \textbf{unternommen}},
  future = {werden + \textbf{unternehmen}},
  pluperfect = {hatten + \textbf{unternommen}},
  prepositions = {
    & \textbf{unternehmen gegen} + Akk & (to undertake against)\\
    & \textbf{unternehmen mit} + Dat & (to do with)\\
  },
  separable = {---}
}

% Migrated from verbs/b1/u.tex:19
\VerbTable{
  name = {unterrichten},
  english = {to teach},
  type = {regular},
  auxiliary = {haben},
  style = {default},
  present = {unterrichte,unterrichtest,unterrichtet,unterrichten,unterrichtet,unterrichten},
  preterite = {unterrichtete,unterrichtetest,unterrichtete,unterrichteten,unterrichtetet,unterrichteten},
  subjunctive = {},
  imperative = {unterrichte (du)!,unterrichten wir!,unterrichtet (ihr)!,unterrichten Sie!},
  perfect = {haben + \textbf{unterrichtet}},
  future = {werden + \textbf{unterrichten}},
  pluperfect = {hatten + \textbf{unterrichtet}},
  prepositions = {---},
  separable = {---}
}

% Migrated from verbs/u.tex:8
\VerbTable{
  name = {unterscheiden},
  english = {to distinguish},
  type = {irregular},
  auxiliary = {haben},
  style = {default},
  present = {unterscheide,unterscheidest,unterscheidet,unterscheiden,unterscheidet,unterscheiden},
  preterite = {unterschied,unterschiedst,unterschied,unterschieden,unterschiedet,unterschieden},
  subjunctive = {},
  imperative = {unterscheide (du)!,unterscheiden wir!,unterscheidet (ihr)!,unterscheiden Sie!},
  perfect = {haben + \textbf{unterschieden}},
  future = {werden + \textbf{unterscheiden}},
  pluperfect = {hatten + \textbf{unterschieden}},
  prepositions = {& \textbf{unterscheiden zwischen} + Dat & (to distinguish between)\\},
  separable = {---}
}

% Migrated from verbs/u.tex:9
\VerbTable{
  name = {unterschreiben},
  english = {to sign something},
  type = {irregular, + Akkusativ},
  auxiliary = {haben},
  style = {acc},
  present = {unterschreibe,unterschreibst,unterschreibt,unterschreiben,unterschreibt,unterschreiben},
  preterite = {unterschrieb,unterschriebst,unterschreib,unterschrieben,unterschriebt,unterschrieben},
  subjunctive = {},
  imperative = {unterschreib(e) (du)!,unterschreiben wir!,unterschreibt (ihr)!,unterschreiben Sie!},
  perfect = {haben + \textbf{unterschrieben}},
  future = {werden + \textbf{unterschreiben}},
  pluperfect = {hatten + \textbf{unterschrieben}},
  prepositions = {---},
  separable = {---}
}

% Migrated from verbs/b1/u.tex:20
\VerbTable{
  name = {unterstreichen},
  english = {to underline},
  type = {strong / irregular},
  auxiliary = {haben},
  style = {default},
  present = {unterstreiche,unterstreichst,unterstreicht,unterstreichen,unterstreicht,unterstreichen},
  preterite = {unterstrich,unterstrichst,unterstrich,unterstrichen,unterstricht,unterstrichen},
  subjunctive = {},
  imperative = {unterstreich (du)!,unterstreichen wir!,unterstreicht (ihr)!,unterstreichen Sie!},
  perfect = {haben + \textbf{unterstrichen}},
  future = {werden + \textbf{unterstreichen}},
  pluperfect = {hatten + \textbf{unterstrichen}},
  prepositions = {---},
  separable = {---}
}

% Migrated from verbs/u.tex:10
\VerbTable{
  name = {unterstützen},
  english = {to support},
  type = {regular},
  auxiliary = {haben},
  style = {default},
  present = {unterstütze,unterstützt,unterstützt,unterstützen,unterstützt,unterstützen},
  preterite = {unterstützte,unterstütztest,unterstützte,unterstützten,unterstütztet,unterstützten},
  subjunctive = {},
  imperative = {unterstütz(e) (du)!,unterstützen wir!,unterstützt (ihr)!,unterstützen Sie!},
  perfect = {haben + \textbf{unterstützt}},
  future = {werden + \textbf{unterstützen}},
  pluperfect = {hatten + \textbf{unterstützt}},
  prepositions = {
    & \textbf{unterstützen bei} + Dat & (to support with / help with)\\
    & \textbf{unterstützen mit} + Dat & (to support with / by means of)\\
  },
  separable = {---}
}

% Migrated from verbs/u.tex:11
\VerbTable{
  name = {untersuchen},
  english = {to examine},
  type = {irregular},
  auxiliary = {haben},
  style = {default},
  present = {untersuche,untersuchst,untersucht,untersuchen,untersucht,untersuchen},
  preterite = {untersuchte,untersuchtest,untersuchte,untersuchten,untersuchtet,untersuchten},
  subjunctive = {},
  imperative = {untersuch(e) (du)!,untersuchen wir!,untersucht (ihr)!,untersuchen Sie!},
  perfect = {haben + \textbf{untersucht}},
  future = {werden + \textbf{untersuchen}},
  pluperfect = {hatten + \textbf{untersucht}},
  prepositions = {---},
  separable = {---}
}

\section{V}

% Migrated from verbs/b1/v.tex:1
\VerbTable{
  name = {verabschieden},
  english = {to say goodbye},
  type = {regular},
  auxiliary = {haben},
  style = {default},
  present = {verabschiede,verabschiedest,verabschiedet,verabschieden,verabschiedet,verabschieden},
  preterite = {verabschiedete,verabschiedetest,verabschiedete,verabschiedeten,verabschiedetet,verabschiedeten},
  subjunctive = {},
  imperative = {verabschiede (du)!,verabschieden wir!,verabschiedet (ihr)!,verabschieden Sie!},
  perfect = {haben + \textbf{verabschiedet}},
  future = {werden + \textbf{verabschieden}},
  pluperfect = {hatten + \textbf{verabschiedet}},
  prepositions = {---},
  separable = {---}
}

% Migrated from verbs/v.tex:1
\VerbTable{
  name = {verändern},
  english = {to change},
  type = {irregular, + Akkusativ},
  auxiliary = {haben},
  style = {acc},
  present = {verändere,veränderst,verändert,verändern,verändert,verändern},
  preterite = {veränderte,verändertest,veränderte,veränderten,verändertet,veränderten},
  subjunctive = {},
  imperative = {veränder(e) (du)!,verändern wir!,verändert (ihr)!,verändern Sie!},
  perfect = {haben + \textbf{verändert}},
  future = {werden + \textbf{verändern}},
  pluperfect = {hatten + \textbf{verändert}},
  prepositions = {---},
  separable = {---}
}

% Migrated from verbs/b1/v.tex:2
\VerbTable{
  name = {verbessern},
  english = {to improve},
  type = {regular},
  auxiliary = {haben},
  style = {default},
  present = {verbessre,verbesserst,verbessert,verbessern,verbessert,verbessern},
  preterite = {verbesserte,verbessertest,verbesserte,verbesserten,verbessertet,verbesserten},
  subjunctive = {},
  imperative = {verbessre (du)!,verbessern wir!,verbessert (ihr)!,verbessern Sie!},
  perfect = {haben + \textbf{verbessert}},
  future = {werden + \textbf{verbessern}},
  pluperfect = {hatten + \textbf{verbessert}},
  prepositions = {---},
  separable = {---}
}

% Migrated from verbs/b1/v.tex:3
\VerbTable{
  name = {verbieten},
  english = {to prohibit},
  type = {strong / irregular},
  auxiliary = {haben},
  style = {default},
  present = {verbiete,verbietest,verbietet,verbieten,verbietet,verbieten},
  preterite = {verbot,verbotest,verbot,verboten,verbotet,verboten},
  subjunctive = {},
  imperative = {verbiete (du)!,verbieten wir!,verbietet (ihr)!,verbieten Sie!},
  perfect = {haben + \textbf{verboten}},
  future = {werden + \textbf{verbieten}},
  pluperfect = {hatten + \textbf{verboten}},
  prepositions = {---},
  separable = {---}
}

% Migrated from verbs/v.tex:2
\VerbTable{
  name = {verbinden},
  english = {to connect / link},
  type = {regular},
  auxiliary = {haben},
  style = {default},
  present = {verbinde,verbindest,verbindet,verbinden,verbindet,verbinden},
  preterite = {verband,verbandst,verband,verbanden,verbandet,verbanden},
  subjunctive = {},
  imperative = {verbinde (du)!,verbinden wir!,verbindet (ihr)!,verbinden Sie!},
  perfect = {haben + \textbf{verbunden}},
  future = {werden + \textbf{verbinden}},
  pluperfect = {hatten + \textbf{verbunden}},
  prepositions = {
    & \textbf{verbinden mit} + Dat & (to connect with)\\
    & \textbf{verbinden zu} + Dat & (to link to / form into)\\
  },
  separable = {---}
}

% Migrated from verbs/b1/v.tex:4
\VerbTable{
  name = {verbrauchen},
  english = {to consume / use up},
  type = {regular},
  auxiliary = {haben},
  style = {default},
  present = {verbrauche,verbrauchst,verbraucht,verbrauchen,verbraucht,verbrauchen},
  preterite = {verbrauchte,verbrauchtest,verbrauchte,verbrauchten,verbrauchtet,verbrauchten},
  subjunctive = {},
  imperative = {verbrauch (du)!,verbrauchen wir!,verbraucht (ihr)!,verbrauchen Sie!},
  perfect = {haben + \textbf{verbraucht}},
  future = {werden + \textbf{verbrauchen}},
  pluperfect = {hatten + \textbf{verbraucht}},
  prepositions = {---},
  separable = {---}
}

% Migrated from verbs/v.tex:3
\VerbTable{
  name = {verbreiten},
  english = {to spread / distribute},
  type = {regular},
  auxiliary = {haben},
  style = {default},
  present = {verbreite,verbreitest,verbreitet,verbreiten,verbreitet,verbreiten},
  preterite = {verbreitete,verbreitetest,verbreitete,verbreiteten,verbreitetet,verbreiteten},
  subjunctive = {},
  imperative = {verbreite (du)!,verbreiten wir!,verbreitet (ihr)!,verbreiten Sie!},
  perfect = {haben + \textbf{verbreitet}},
  future = {werden + \textbf{verbreiten}},
  pluperfect = {hatten + \textbf{verbreitet}},
  prepositions = {
    & \textbf{verbreiten unter} + Dat & (to spread among)\\
    & \textbf{verbreiten über} + Akk & (to spread about)\\
    & \textbf{verbreiten durch} + Akk & (to spread through)\\
  },
  separable = {---}
}

% Migrated from verbs/b1/v.tex:5
\VerbTable{
  name = {verbrennen},
  english = {to burn},
  type = {regular},
  auxiliary = {haben},
  style = {default},
  present = {verbrenne,verbrennst,verbrennt,verbrennen,verbrennt,verbrennen},
  preterite = {verbrannte,verbranntest,verbrannte,verbrannten,verbranntet,verbrannten},
  subjunctive = {},
  imperative = {verbrenn (du)!,verbrennen wir!,verbrennt (ihr)!,verbrennen Sie!},
  perfect = {haben + \textbf{verbrannt}},
  future = {werden + \textbf{verbrennen}},
  pluperfect = {hatten + \textbf{verbrannt}},
  prepositions = {---},
  separable = {---}
}

% Migrated from verbs/v.tex:4
\VerbTable{
  name = {verbringen},
  english = {to spend time},
  type = {irregular, + Akkusativ},
  auxiliary = {haben},
  style = {acc},
  present = {verbringe,verbringst,verbringt,verbringen,verbringt,verbringen},
  preterite = {verbrachte,verbrachtest,verbrachte,verbrachten,verbrachtet,verbrachten},
  subjunctive = {},
  imperative = {verbring(e) (du)!,verbringen wir!,verbringt (ihr)!,verbringen Sie!},
  perfect = {haben + \textbf{verbracht}},
  future = {werden + \textbf{verbringen}},
  pluperfect = {hatten + \textbf{verbracht}},
  prepositions = {---},
  separable = {---}
}

% Migrated from verbs/b1/v.tex:6
\VerbTable{
  name = {verdienen},
  english = {to earn},
  type = {regular},
  auxiliary = {haben},
  style = {default},
  present = {verdiene,verdienst,verdient,verdienen,verdient,verdienen},
  preterite = {verdiente,verdientest,verdiente,verdienten,verdientet,verdienten},
  subjunctive = {},
  imperative = {verdien (du)!,verdienen wir!,verdient (ihr)!,verdienen Sie!},
  perfect = {haben + \textbf{verdient}},
  future = {werden + \textbf{verdienen}},
  pluperfect = {hatten + \textbf{verdient}},
  prepositions = {---},
  separable = {---}
}

% Migrated from verbs/v.tex:5
\VerbTable{
  name = {vereinbaren},
  english = {to arrange / agree on},
  type = {regular},
  auxiliary = {haben},
  style = {default},
  present = {vereinbare,vereinbarst,vereinbart,vereinbaren,vereinbart,vereinbaren},
  preterite = {vereinbarte,vereinbartest,vereinbarte,vereinbarten,vereinbartet,vereinbarten},
  subjunctive = {},
  imperative = {vereinbare (du)!,vereinbaren wir!,vereinbart (ihr)!,vereinbaren Sie!},
  perfect = {haben + \textbf{vereinbart}},
  future = {werden + \textbf{vereinbaren}},
  pluperfect = {hatten + \textbf{vereinbart}},
  prepositions = {
    & \textbf{vereinbaren mit} + Dat & (to arrange with / agree with someone)\\
    & \textbf{vereinbaren über} + Akk & (to agree on / about something)\\
  },
  separable = {---}
}

% Migrated from verbs/v.tex:7
\VerbTable{
  name = {vergessen},
  english = {to forget},
  type = {irregular, + Akkusativ},
  auxiliary = {haben},
  style = {acc},
  present = {vergesse,vergisst,vergisst,vergessen,vergesst,vergessen},
  preterite = {vergaß,vergaßest,vergaß,vergaßen,vergaßt,vergaßen},
  subjunctive = {},
  imperative = {vergiss (du)!,vergessen wir!,vergesst (ihr)!,vergessen Sie!},
  perfect = {haben + \textbf{vergessen}},
  future = {werden + \textbf{vergessen}},
  pluperfect = {hatten + \textbf{vergessen}},
  prepositions = {---},
  separable = {---}
}

% Migrated from verbs/v.tex:6
\VerbTable{
  name = {vergleichen},
  english = {to compare},
  type = {irregular, + Accuative + Dativ},
  auxiliary = {haben},
  style = {dat},
  present = {vergleiche,vergleichst,vergleicht,vergleichen,vergleicht,vergleichen},
  preterite = {verglich,verglichst,verglich,verglichen,verglicht,verglichen},
  subjunctive = {},
  imperative = {vergleich(e) (du)!,vergleichen wir!,vergleicht (ihr)!,vergleichen Sie!},
  perfect = {haben + \textbf{verglichen}},
  future = {werden + \textbf{vergleichen}},
  pluperfect = {hatten + \textbf{verglichen}},
  prepositions = {---},
  separable = {---}
}

% Migrated from verbs/b1/v.tex:8
\VerbTable{
  name = {vergrößern},
  english = {to enlarge},
  type = {regular},
  auxiliary = {haben},
  style = {default},
  present = {vergrößre,vergrößerst,vergrößert,vergrößern,vergrößert,vergrößern},
  preterite = {vergrößerte,vergrößertest,vergrößerte,vergrößerten,vergrößertet,vergrößerten},
  subjunctive = {},
  imperative = {vergrößre (du)!,vergrößern wir!,vergrößert (ihr)!,vergrößern Sie!},
  perfect = {haben + \textbf{vergrößert}},
  future = {werden + \textbf{vergrößern}},
  pluperfect = {hatten + \textbf{vergrößert}},
  prepositions = {---},
  separable = {---}
}

% Migrated from verbs/b1/v.tex:9
\VerbTable{
  name = {verhaften},
  english = {to arrest},
  type = {regular},
  auxiliary = {haben},
  style = {default},
  present = {verhafte,verhaftest,verhaftet,verhaften,verhaftet,verhaften},
  preterite = {verhaftete,verhaftetest,verhaftete,verhafteten,verhaftetet,verhafteten},
  subjunctive = {},
  imperative = {verhafte (du)!,verhaften wir!,verhaftet (ihr)!,verhaften Sie!},
  perfect = {haben + \textbf{verhaftet}},
  future = {werden + \textbf{verhaften}},
  pluperfect = {hatten + \textbf{verhaftet}},
  prepositions = {---},
  separable = {---}
}

% Migrated from verbs/v.tex:8
\VerbTable{
  name = {verhindern},
  english = {to prevent},
  type = {regular, inseparable},
  auxiliary = {haben},
  style = {default},
  present = {verhindere,verhinderst,verhindert,verhindern,verhindert,verhindern},
  preterite = {verhinderte,verhindertest,verhinderte,verhinderten,verhindertet,verhinderten},
  subjunctive = {},
  imperative = {verhindere (du)!,verhindern wir!,verhindert (ihr)!,verhindern Sie!},
  perfect = {haben + \textbf{verhindert}},
  future = {werden + \textbf{verhindern}},
  pluperfect = {hatten + \textbf{verhindert}},
  prepositions = {
    & \textbf{verhindern an} + Dat & (to prevent from)\\
    & \textbf{verhindern durch} + Akk & (to prevent by means of)\\
  },
  separable = {---}
}

% Migrated from verbs/b1/v.tex:11
\VerbTable{
  name = {verkaufen},
  english = {to sell},
  type = {regular},
  auxiliary = {haben},
  style = {default},
  present = {verkaufe,verkaufst,verkauft,verkaufen,verkauft,verkaufen},
  preterite = {verkaufte,verkauftest,verkaufte,verkauften,verkauftet,verkauften},
  subjunctive = {},
  imperative = {verkauf (du)!,verkaufen wir!,verkauft (ihr)!,verkaufen Sie!},
  perfect = {haben + \textbf{verkauft}},
  future = {werden + \textbf{verkaufen}},
  pluperfect = {hatten + \textbf{verkauft}},
  prepositions = {---},
  separable = {---}
}

% Migrated from verbs/b1/v.tex:13
\VerbTable{
  name = {verlängern},
  english = {to extend},
  type = {regular},
  auxiliary = {haben},
  style = {default},
  present = {verlängre,verlängerst,verlängert,verlängern,verlängert,verlängern},
  preterite = {verlängerte,verlängertest,verlängerte,verlängerten,verlängertet,verlängerten},
  subjunctive = {},
  imperative = {verlängre (du)!,verlängern wir!,verlängert (ihr)!,verlängern Sie!},
  perfect = {haben + \textbf{verlängert}},
  future = {werden + \textbf{verlängern}},
  pluperfect = {hatten + \textbf{verlängert}},
  prepositions = {---},
  separable = {---}
}

% Migrated from verbs/b1/v.tex:12
\VerbTable{
  name = {verlangen},
  english = {to demand},
  type = {regular},
  auxiliary = {haben},
  style = {default},
  present = {verlange,verlangst,verlangt,verlangen,verlangt,verlangen},
  preterite = {verlangte,verlangtest,verlangte,verlangten,verlangtet,verlangten},
  subjunctive = {},
  imperative = {verlang (du)!,verlangen wir!,verlangt (ihr)!,verlangen Sie!},
  perfect = {haben + \textbf{verlangt}},
  future = {werden + \textbf{verlangen}},
  pluperfect = {hatten + \textbf{verlangt}},
  prepositions = {---},
  separable = {---}
}

% Migrated from verbs/v.tex:9
\VerbTable{
  name = {verleihen},
  english = {to lend / award},
  type = {regular},
  auxiliary = {haben},
  style = {default},
  present = {verleihe,verleihst,verleiht,verleihen,verleiht,verleihen},
  preterite = {verlieh,verliehst,verlieh,verliehen,verlieht,verliehen},
  subjunctive = {},
  imperative = {verleihe (du)!,verleihen wir!,verleiht (ihr)!,verleihen Sie!},
  perfect = {haben + \textbf{verliehen}},
  future = {werden + \textbf{verleihen}},
  pluperfect = {hatten + \textbf{verliehen}},
  prepositions = {
    & \textbf{verleihen an} + Akk & (to lend to)\\
    & \textbf{verleihen für} + Akk & (to award for)\\
  },
  separable = {& \textbf{ver$|$leihen} & (to lend out / award)\\}
}

% Migrated from verbs/v.tex:10
\VerbTable{
  name = {verletzen},
  english = {to injure / to hurt},
  type = {regular},
  auxiliary = {haben},
  style = {default},
  present = {verletze,verletzt,verletzt,verletzen,verletzt,verletzen},
  preterite = {verletzte,verletztest,verletzte,verletzten,verletztet,verletzten},
  subjunctive = {},
  imperative = {verletz(e) (du)!,verletzen wir!,verletzt (ihr)!,verletzen Sie!},
  perfect = {haben + \textbf{verletzt}},
  future = {werde + \textbf{verletzen}},
  pluperfect = {hatten + \textbf{verletzt}},
  prepositions = {
    & \textbf{sich verletzen an} + Dat & (to injure oneself on)\\
    & \textbf{verletzen mit} + Dat & (to injure with)\\
  },
  separable = {---}
}

% Migrated from verbs/v.tex:11
\VerbTable{
  name = {verlieren},
  english = {to lose},
  type = {irregular},
  auxiliary = {haben},
  style = {default},
  present = {verliere,verlierst,verliert,verlieren,verliert,verlieren},
  preterite = {verlor,verlorst,verlor,verloren,verlort,verloren},
  subjunctive = {},
  imperative = {verlier(e) (du)!,verlieren wir!,verliert (ihr)!,verlieren Sie!},
  perfect = {haben + \textbf{verloren}},
  future = {werde + \textbf{verlieren}},
  pluperfect = {hatten + \textbf{verloren}},
  prepositions = {
    & \textbf{verlieren an} + Akk & (to lose to)\\
    & \textbf{verlieren durch} + Akk & (to lose through / because of)\\
  },
  separable = {---}
}

% Migrated from verbs/v.tex:12
\VerbTable{
  name = {vermeiden},
  english = {to avoid},
  type = {irregular},
  auxiliary = {haben},
  style = {default},
  present = {vermeide,vermeidest,vermeidet,vermeiden,vermeidet,vermeiden},
  preterite = {vermied,vermiedest,vermied,vermieden,vermiedet,vermieden},
  subjunctive = {},
  imperative = {vermeide (du)!,vermeiden wir!,vermeidet (ihr)!,vermeiden Sie!},
  perfect = {haben + \textbf{vermieden}},
  future = {werden + \textbf{vermeiden}},
  pluperfect = {hatten + \textbf{vermieden}},
  prepositions = {---},
  separable = {---}
}

% Migrated from verbs/b1/v.tex:15
\VerbTable{
  name = {vermieten},
  english = {to rent out},
  type = {regular},
  auxiliary = {haben},
  style = {default},
  present = {vermiete,vermietest,vermietet,vermieten,vermietet,vermieten},
  preterite = {vermietete,vermietetest,vermietete,vermieteten,vermietetet,vermieteten},
  subjunctive = {},
  imperative = {vermiete (du)!,vermieten wir!,vermietet (ihr)!,vermieten Sie!},
  perfect = {haben + \textbf{vermietet}},
  future = {werden + \textbf{vermieten}},
  pluperfect = {hatten + \textbf{vermietet}},
  prepositions = {---},
  separable = {---}
}

% Migrated from verbs/v.tex:13
\VerbTable{
  name = {vermissen},
  english = {to miss someone},
  type = {regular},
  auxiliary = {haben},
  style = {default},
  present = {vermisse,vermisst,vermisst,vermissen,vermisst,vermissen},
  preterite = {vermisste,vermisstest,vermisste,vermissten,vermisstet,vermissten},
  subjunctive = {},
  imperative = {vermisse (du)!,vermissen wir!,vermisst (ihr)!,vermissen Sie!},
  perfect = {haben + \textbf{vermisst}},
  future = {werden + \textbf{vermissen}},
  pluperfect = {hatten + \textbf{vermisst}},
  prepositions = {& \textbf{vermissen} + Akk & (to miss someone/something)\\},
  separable = {---}
}

% Migrated from verbs/b1/v.tex:16
\VerbTable{
  name = {vermuten},
  english = {to suspect},
  type = {regular},
  auxiliary = {haben},
  style = {default},
  present = {vermute,vermutest,vermutet,vermuten,vermutet,vermuten},
  preterite = {vermutete,vermutetest,vermutete,vermuteten,vermutetet,vermuteten},
  subjunctive = {},
  imperative = {vermute (du)!,vermuten wir!,vermutet (ihr)!,vermuten Sie!},
  perfect = {haben + \textbf{vermutet}},
  future = {werden + \textbf{vermuten}},
  pluperfect = {hatten + \textbf{vermutet}},
  prepositions = {---},
  separable = {---}
}

% Migrated from verbs/b1/v.tex:17
\VerbTable{
  name = {veröffentlichen},
  english = {to publish},
  type = {regular},
  auxiliary = {haben},
  style = {default},
  present = {veröffentliche,veröffentlichst,veröffentlicht,veröffentlichen,veröffentlicht,veröffentlichen},
  preterite = {veröffentlichte,veröffentlichtest,veröffentlichte,veröffentlichten,veröffentlichtet,veröffentlichten},
  subjunctive = {},
  imperative = {veröffentlich (du)!,veröffentlichen wir!,veröffentlicht (ihr)!,veröffentlichen Sie!},
  perfect = {haben + \textbf{veröffentlicht}},
  future = {werden + \textbf{veröffentlichen}},
  pluperfect = {hatten + \textbf{veröffentlicht}},
  prepositions = {---},
  separable = {---}
}

% Migrated from verbs/b1/v.tex:18
\VerbTable{
  name = {verpacken},
  english = {to pack / package},
  type = {regular},
  auxiliary = {haben},
  style = {default},
  present = {verpacke,verpackst,verpackt,verpacken,verpackt,verpacken},
  preterite = {verpackte,verpacktest,verpackte,verpackten,verpacktet,verpackten},
  subjunctive = {},
  imperative = {verpack (du)!,verpacken wir!,verpackt (ihr)!,verpacken Sie!},
  perfect = {haben + \textbf{verpackt}},
  future = {werden + \textbf{verpacken}},
  pluperfect = {hatten + \textbf{verpackt}},
  prepositions = {---},
  separable = {---}
}

% Migrated from verbs/b1/v.tex:19
\VerbTable{
  name = {verpassen},
  english = {to miss},
  type = {regular},
  auxiliary = {haben},
  style = {default},
  present = {verpasse,verpasst,verpasst,verpassen,verpasst,verpassen},
  preterite = {verpasste,verpasstest,verpasste,verpassten,verpasstet,verpassten},
  subjunctive = {},
  imperative = {verpass (du)!,verpassen wir!,verpasst (ihr)!,verpassen Sie!},
  perfect = {haben + \textbf{verpasst}},
  future = {werden + \textbf{verpassen}},
  pluperfect = {hatten + \textbf{verpasst}},
  prepositions = {---},
  separable = {---}
}

% Migrated from verbs/b1/v.tex:20
\VerbTable{
  name = {verpflegen},
  english = {to provide food / care for},
  type = {regular},
  auxiliary = {haben},
  style = {default},
  present = {verpflege,verpflegst,verpflegt,verpflegen,verpflegt,verpflegen},
  preterite = {verpflegte,verpflegtest,verpflegte,verpflegten,verpflegtet,verpflegten},
  subjunctive = {},
  imperative = {verpfleg (du)!,verpflegen wir!,verpflegt (ihr)!,verpflegen Sie!},
  perfect = {haben + \textbf{verpflegt}},
  future = {werden + \textbf{verpflegen}},
  pluperfect = {hatten + \textbf{verpflegt}},
  prepositions = {---},
  separable = {---}
}

% Migrated from verbs/v.tex:14
\VerbTable{
  name = {verputzen},
  english = {to plaster / render},
  type = {regular},
  auxiliary = {haben},
  style = {default},
  present = {verputze,verputzt,verputzt,verputzen,verputzt,verputzen},
  preterite = {verputzte,verputztest,verputzte,verputzten,verputztet,verputzten},
  subjunctive = {},
  imperative = {verputze (du)!,verputzen wir!,verputzt (ihr)!,verputzen Sie!},
  perfect = {haben + \textbf{verputzt}},
  future = {werden + \textbf{verputzen}},
  pluperfect = {hatten + \textbf{verputzt}},
  prepositions = {& \textbf{verputzen mit} + Dat & (to plaster with)\\},
  separable = {& \textbf{ein$|$verputzen} & (to plaster in / apply)\\}
}

% Migrated from verbs/b1/v.tex:21
\VerbTable{
  name = {verraten},
  english = {to betray / reveal},
  type = {strong / irregular},
  auxiliary = {haben},
  style = {default},
  present = {verrate,verrätst,verrät,verraten,verratet,verraten},
  preterite = {verriet,verrietest,verriet,verrieten,verrietet,verrieten},
  subjunctive = {},
  imperative = {verrate (du)!,verraten wir!,verratet (ihr)!,verraten Sie!},
  perfect = {haben + \textbf{verraten}},
  future = {werden + \textbf{verraten}},
  pluperfect = {hatten + \textbf{verraten}},
  prepositions = {---},
  separable = {---}
}

% Migrated from verbs/b1/v.tex:22
\VerbTable{
  name = {verreisen},
  english = {to go away / travel},
  type = {regular},
  auxiliary = {sein},
  style = {default},
  present = {verreise,verreist,verreist,verreisen,verreist,verreisen},
  preterite = {verreiste,verreistest,verreiste,verreisten,verreistet,verreisten},
  subjunctive = {},
  imperative = {verreis (du)!,verreisen wir!,verreist (ihr)!,verreisen Sie!},
  perfect = {sein + \textbf{verreist}},
  future = {werden + \textbf{verreisen}},
  pluperfect = {waren + \textbf{verreist}},
  prepositions = {---},
  separable = {---}
}

% Migrated from verbs/b1/v.tex:23
\VerbTable{
  name = {versäumen},
  english = {to miss / fail to do},
  type = {regular},
  auxiliary = {haben},
  style = {default},
  present = {versäume,versäumst,versäumt,versäumen,versäumt,versäumen},
  preterite = {versäumte,versäumtest,versäumte,versäumten,versäumtet,versäumten},
  subjunctive = {},
  imperative = {versäum (du)!,versäumen wir!,versäumt (ihr)!,versäumen Sie!},
  perfect = {haben + \textbf{versäumt}},
  future = {werden + \textbf{versäumen}},
  pluperfect = {hatten + \textbf{versäumt}},
  prepositions = {---},
  separable = {---}
}

% Migrated from verbs/b1/v.tex:24
\VerbTable{
  name = {verschieben},
  english = {to postpone / move},
  type = {strong / irregular},
  auxiliary = {haben},
  style = {default},
  present = {verschiebe,verschiebst,verschiebt,verschieben,verschiebt,verschieben},
  preterite = {verschob,verschobst,verschob,verschoben,verschobt,verschoben},
  subjunctive = {},
  imperative = {verschieb (du)!,verschieben wir!,verschiebt (ihr)!,verschieben Sie!},
  perfect = {haben + \textbf{verschoben}},
  future = {werden + \textbf{verschieben}},
  pluperfect = {hatten + \textbf{verschoben}},
  prepositions = {---},
  separable = {---}
}

% Migrated from verbs/b1/v.tex:25
\VerbTable{
  name = {verschmutzen},
  english = {to contaminate},
  type = {regular},
  auxiliary = {haben},
  style = {default},
  present = {verschmutze,verschmutzt,verschmutzt,verschmutzen,verschmutzt,verschmutzen},
  preterite = {verschmutzte,verschmutztest,verschmutzte,verschmutzten,verschmutztet,verschmutzten},
  subjunctive = {},
  imperative = {verschmutz (du)!,verschmutzen wir!,verschmutzt (ihr)!,verschmutzen Sie!},
  perfect = {haben + \textbf{verschmutzt}},
  future = {werden + \textbf{verschmutzen}},
  pluperfect = {hatten + \textbf{verschmutzt}},
  prepositions = {---},
  separable = {---}
}

% Migrated from verbs/v.tex:15
\VerbTable{
  name = {verschreiben},
  english = {to prescribe},
  type = {irregular},
  auxiliary = {haben / sich + haben},
  style = {default},
  present = {verschreibe,verschreibst,verschreibt,verschreiben,verschreibt,verschreiben},
  preterite = {verschrieb,verschriebst,verschrieb,verschrieben,verschriebt,verschrieben},
  subjunctive = {},
  imperative = {verschreibe (du)!,verschreiben wir!,verschreibt (ihr)!,verschreiben Sie!},
  perfect = {haben + \textbf{verschrieben}},
  future = {werden + \textbf{verschreiben}},
  pluperfect = {hatten + \textbf{verschrieben}},
  prepositions = {---},
  separable = {---}
}

% Migrated from verbs/v.tex:16
\VerbTable{
  name = {verschwenden},
  english = {to waste},
  type = {regular},
  auxiliary = {haben},
  style = {default},
  present = {verschwende,verschwendest,verschwendet,verschwenden,verschwendet,verschwenden},
  preterite = {verschwendete,verschwendetest,verschwendete,verschwendeten,verschwendetet,verschwendeten},
  subjunctive = {},
  imperative = {verschwende (du)!,verschwenden wir!,verschwendet (ihr)!,verschwenden Sie!},
  perfect = {haben + \textbf{verschwendet}},
  future = {werde + \textbf{verschwenden}},
  pluperfect = {hatten + \textbf{verschwendet}},
  prepositions = {
    & \textbf{verschwenden an} + Akk & (to waste on)\\
    & \textbf{verschwenden für} + Akk & (to waste for)\\
  },
  separable = {---}
}

% Migrated from verbs/v.tex:17
\VerbTable{
  name = {verschwinden},
  english = {to disappear},
  type = {irregular},
  auxiliary = {sein},
  style = {default},
  present = {verschwinde,verschwindest,verschwindet,verschwinden,verschwindet,verschwinden},
  preterite = {verschwand,verschwandest,verschwand,verschwanden,verschwandet,verschwanden},
  subjunctive = {},
  imperative = {verschwind(e) (du)!,verschwinden wir!,verschwindet (ihr)!,verschwinden Sie!},
  perfect = {sein + \textbf{verschwunden}},
  future = {werde + \textbf{verschwinden}},
  pluperfect = {waren + \textbf{verschwunden}},
  prepositions = {
    & \textbf{verschwinden aus} + Dat & (to disappear from)\\
    & \textbf{verschwinden in} + Dat & (to disappear in / into)\\
  },
  separable = {---}
}

% Migrated from verbs/b1/v.tex:26
\VerbTable{
  name = {versichern},
  english = {to insure},
  type = {regular},
  auxiliary = {haben},
  style = {default},
  present = {versichre,versicherst,versichert,versichern,versichert,versichern},
  preterite = {versicherte,versichertest,versicherte,versicherten,versichertet,versicherten},
  subjunctive = {},
  imperative = {versichre (du)!,versichern wir!,versichert (ihr)!,versichern Sie!},
  perfect = {haben + \textbf{versichert}},
  future = {werden + \textbf{versichern}},
  pluperfect = {hatten + \textbf{versichert}},
  prepositions = {---},
  separable = {---}
}

% Migrated from verbs/b1/v.tex:27
\VerbTable{
  name = {versprechen},
  english = {to promise},
  type = {strong / irregular},
  auxiliary = {haben},
  style = {default},
  present = {verspreche,versprichst,verspricht,versprechen,versprecht,versprechen},
  preterite = {versprach,versprachst,versprach,versprachen,verspracht,versprachen},
  subjunctive = {},
  imperative = {versprich (du)!,versprechen wir!,versprecht (ihr)!,versprechen Sie!},
  perfect = {haben + \textbf{versprochen}},
  future = {werden + \textbf{versprechen}},
  pluperfect = {hatten + \textbf{versprochen}},
  prepositions = {---},
  separable = {---}
}

% Migrated from verbs/b1/v.tex:28
\VerbTable{
  name = {verstecken},
  english = {to hide},
  type = {regular},
  auxiliary = {haben},
  style = {default},
  present = {verstecke,versteckst,versteckt,verstecken,versteckt,verstecken},
  preterite = {versteckte,verstecktest,versteckte,versteckten,verstecktet,versteckten},
  subjunctive = {},
  imperative = {versteck (du)!,verstecken wir!,versteckt (ihr)!,verstecken Sie!},
  perfect = {haben + \textbf{versteckt}},
  future = {werden + \textbf{verstecken}},
  pluperfect = {hatten + \textbf{versteckt}},
  prepositions = {---},
  separable = {---}
}

% Migrated from verbs/v.tex:18
\VerbTable{
  name = {verstehen},
  english = {to understand},
  type = {irregular, + Akkusativ},
  auxiliary = {haben},
  style = {acc},
  present = {verstehe,verstehst,versteht,verstehen,versteht,verstehen},
  preterite = {verstand,verstandest,verstand,verstanden,verstandet,verstanden},
  subjunctive = {},
  imperative = {versteh(e) (du)!,verstehen wir!,versteht (ihr)!,verstehen Sie!},
  perfect = {haben + \textbf{verstanden}},
  future = {werden + \textbf{verstehen}},
  pluperfect = {hatten + \textbf{verstanden}},
  prepositions = {---},
  separable = {---}
}

% Migrated from verbs/v.tex:19
\VerbTable{
  name = {versuchen},
  english = {to try},
  type = {irregular, + Akkusativ},
  auxiliary = {haben},
  style = {acc},
  present = {versuche,versuchst,versucht,versuchen,versucht,versuchen},
  preterite = {versuchte,versuchtest,versuchte,versuchten,versuchtet,versuchten},
  subjunctive = {},
  imperative = {versuch(e) (du)!,versuchen wir!,versucht (ihr)!,versuchen Sie!},
  perfect = {haben + \textbf{versucht}},
  future = {werden + \textbf{versuchen}},
  pluperfect = {hatten + \textbf{versucht}},
  prepositions = {---},
  separable = {---}
}

% Migrated from verbs/b1/v.tex:29
\VerbTable{
  name = {verteilen},
  english = {to distribute},
  type = {regular},
  auxiliary = {haben},
  style = {default},
  present = {verteile,verteilst,verteilt,verteilen,verteilt,verteilen},
  preterite = {verteilte,verteiltest,verteilte,verteilten,verteiltet,verteilten},
  subjunctive = {},
  imperative = {verteil (du)!,verteilen wir!,verteilt (ihr)!,verteilen Sie!},
  perfect = {haben + \textbf{verteilt}},
  future = {werden + \textbf{verteilen}},
  pluperfect = {hatten + \textbf{verteilt}},
  prepositions = {---},
  separable = {---}
}

% Migrated from verbs/v.tex:20
\VerbTable{
  name = {vertiefen},
  english = {to deepen},
  type = {regular, + Akkusativ},
  auxiliary = {haben},
  style = {acc},
  present = {vertiefe,vertiefst,vertieft,vertiefen,vertieft,vertiefen},
  preterite = {vertiefte,vertieftest,vertiefte,vertieften,vertieftet,vertieften},
  subjunctive = {},
  imperative = {vertiefe (du)!,vertiefen wir!,vertieft (ihr)!,vertiefen Sie!},
  perfect = {haben + \textbf{vertieft}},
  future = {werden + \textbf{vertiefen}},
  pluperfect = {hatten + \textbf{vertieft}},
  prepositions = {& \textbf{vertiefen in} + Akk & (to deepen into / immerse in)\\},
  separable = {---}
}

% Migrated from verbs/v.tex:21
\VerbTable{
  name = {vertrauen},
  english = {to trust},
  type = {regular},
  auxiliary = {haben},
  style = {default},
  present = {vertraue,vertraust,vertraut,vertrauen,vertraut,vertrauen},
  preterite = {vertraute,vertrautest,vertraute,vertrauten,vertrautet,vertrauten},
  subjunctive = {},
  imperative = {vertraue (du)!,vertrauen wir!,vertraut (ihr)!,vertrauen Sie!},
  perfect = {haben + \textbf{vertraut}},
  future = {werden + \textbf{vertrauen}},
  pluperfect = {hatten + \textbf{vertraut}},
  prepositions = {
    & \textbf{vertrauen auf} + Akk & (to rely on)\\
    & \textbf{vertrauen in} + Akk & (to trust in)\\
    & \textbf{vertrauen} + Dat & (to trust someone)\\
  },
  separable = {---}
}

% Migrated from verbs/b1/v.tex:30
\VerbTable{
  name = {vertreten},
  english = {to represent / stand in for},
  type = {strong / irregular},
  auxiliary = {haben},
  style = {default},
  present = {vertrete,vertrittst,vertritt,vertreten,vertretet,vertreten},
  preterite = {vertrat,vertratest,vertrat,vertraten,vertratet,vertraten},
  subjunctive = {},
  imperative = {vertritt (du)!,vertreten wir!,vertretet (ihr)!,vertreten Sie!},
  perfect = {haben + \textbf{vertreten}},
  future = {werden + \textbf{vertreten}},
  pluperfect = {hatten + \textbf{vertreten}},
  prepositions = {---},
  separable = {---}
}

% Migrated from verbs/b1/v.tex:31
\VerbTable{
  name = {verursachen},
  english = {to cause},
  type = {regular},
  auxiliary = {haben},
  style = {default},
  present = {verursache,verursachst,verursacht,verursachen,verursacht,verursachen},
  preterite = {verursachte,verursachtest,verursachte,verursachten,verursachtet,verursachten},
  subjunctive = {},
  imperative = {verursach (du)!,verursachen wir!,verursacht (ihr)!,verursachen Sie!},
  perfect = {haben + \textbf{verursacht}},
  future = {werden + \textbf{verursachen}},
  pluperfect = {hatten + \textbf{verursacht}},
  prepositions = {---},
  separable = {---}
}

% Migrated from verbs/b1/v.tex:32
\VerbTable{
  name = {verurteilen},
  english = {to condemn},
  type = {regular},
  auxiliary = {haben},
  style = {default},
  present = {verurteile,verurteilst,verurteilt,verurteilen,verurteilt,verurteilen},
  preterite = {verurteilte,verurteiltest,verurteilte,verurteilten,verurteiltet,verurteilten},
  subjunctive = {},
  imperative = {verurteil (du)!,verurteilen wir!,verurteilt (ihr)!,verurteilen Sie!},
  perfect = {haben + \textbf{verurteilt}},
  future = {werden + \textbf{verurteilen}},
  pluperfect = {hatten + \textbf{verurteilt}},
  prepositions = {---},
  separable = {---}
}

% Migrated from verbs/b1/v.tex:33
\VerbTable{
  name = {verwechseln},
  english = {to confuse},
  type = {regular},
  auxiliary = {haben},
  style = {default},
  present = {verwechsle,verwechselst,verwechselt,verwechseln,verwechselt,verwechseln},
  preterite = {verwechselte,verwechseltest,verwechselte,verwechselten,verwechseltet,verwechselten},
  subjunctive = {},
  imperative = {verwechsle (du)!,verwechseln wir!,verwechselt (ihr)!,verwechseln Sie!},
  perfect = {haben + \textbf{verwechselt}},
  future = {werden + \textbf{verwechseln}},
  pluperfect = {hatten + \textbf{verwechselt}},
  prepositions = {---},
  separable = {---}
}

% Migrated from verbs/b1/v.tex:34
\VerbTable{
  name = {verwenden},
  english = {to use},
  type = {regular},
  auxiliary = {haben},
  style = {default},
  present = {verwende,verwendest,verwendet,verwenden,verwendet,verwenden},
  preterite = {verwendete,verwendetest,verwendete,verwendeten,verwendetet,verwendeten},
  subjunctive = {},
  imperative = {verwende (du)!,verwenden wir!,verwendet (ihr)!,verwenden Sie!},
  perfect = {haben + \textbf{verwendet}},
  future = {werden + \textbf{verwenden}},
  pluperfect = {hatten + \textbf{verwendet}},
  prepositions = {---},
  separable = {---}
}

% Migrated from verbs/v.tex:22
\VerbTable{
  name = {verzeihen},
  english = {to forgive},
  type = {irregular, + Dativ},
  auxiliary = {haben},
  style = {dat},
  present = {verzeihe,verzeihst,verzeiht,verzeihen,verzeiht,verzeihen},
  preterite = {verzieh,verziehst,verzieh,verziehen,verzieht,verziehen},
  subjunctive = {},
  imperative = {verzeih(e) (du)!,verzeihen wir!,verzeiht (ihr)!,verzeihen Sie!},
  perfect = {haben + \textbf{verziehen}},
  future = {werden + \textbf{verzeihen}},
  pluperfect = {hatten + \textbf{verziehen}},
  prepositions = {& -- & --\\},
  separable = {& -- & --\\}
}

% Migrated from verbs/v.tex:23
\VerbTable{
  name = {verzichten},
  english = {to do without},
  type = {regular},
  auxiliary = {haben},
  style = {default},
  present = {verzichte,verzichtest,verzichtet,verzichten,verzichtet,verzichten},
  preterite = {verzichtete,verzichtetest,verzichtete,verzichteten,verzichtetet,verzichteten},
  subjunctive = {},
  imperative = {verzichte (du)!,verzichten wir!,verzichtet (ihr)!,verzichten Sie!},
  perfect = {haben + \textbf{verzichtet}},
  future = {werden + \textbf{verzichten}},
  pluperfect = {hatten + \textbf{verzichtet}},
  prepositions = {& \textbf{verzichten auf} + Akk & (to give up / do without)\\},
  separable = {---}
}

% Migrated from verbs/v.tex:24
\VerbTable{
  name = {verzieren},
  english = {to decorate / adorn},
  type = {regular, + Akkusativ},
  auxiliary = {haben},
  style = {acc},
  present = {verziere,verzierst,verziert,verzieren,verziert,verzieren},
  preterite = {verzierte,verziertest,verzierte,verzierten,verziertet,verzierten},
  subjunctive = {},
  imperative = {verziere (du)!,verzieren wir!,verziert (ihr)!,verzieren Sie!},
  perfect = {haben + \textbf{verziert}},
  future = {werden + \textbf{verzieren}},
  pluperfect = {hatten + \textbf{verziert}},
  prepositions = {& \textbf{verzieren mit} + Dat & (to decorate with)\\},
  separable = {---}
}

% Migrated from verbs/b1/v.tex:35
\VerbTable{
  name = {vorbereiten},
  english = {to prepare},
  type = {separable, regular},
  auxiliary = {haben},
  style = {default},
  present = {bereite vor,bereitest vor,bereitet vor,bereiten vor,bereitet vor,bereiten vor},
  preterite = {bereitete vor,bereitetest vor,bereitete vor,bereiteten vor,bereitetet vor,bereiteten vor},
  subjunctive = {},
  imperative = {bereite (du) vor!,bereiten wir vor!,bereitet (ihr) vor!,bereiten Sie vor!},
  perfect = {haben + \textbf{vorbereitet}},
  future = {werden + \textbf{vorbereiten}},
  pluperfect = {hatten + \textbf{vorbereitet}},
  prepositions = {---},
  separable = {---}
}

% Migrated from verbs/b1/v.tex:36
\VerbTable{
  name = {vorhaben},
  english = {to intend / plan},
  type = {separable, regular},
  auxiliary = {haben},
  style = {default},
  present = {habe vor,hast vor,hat vor,haben vor,habt vor,haben vor},
  preterite = {hatte vor,hattest vor,hatte vor,hatten vor,hattet vor,hatten vor},
  subjunctive = {},
  imperative = {hab (du) vor!,haben wir vor!,habt (ihr) vor!,haben Sie vor!},
  perfect = {haben + \textbf{vorgehabt}},
  future = {werden + \textbf{vorhaben}},
  pluperfect = {hatten + \textbf{vorgehabt}},
  prepositions = {---},
  separable = {---}
}

% Migrated from verbs/b1/v.tex:37
\VerbTable{
  name = {vorkommen},
  english = {to occur / seem},
  type = {separable, strong / irregular},
  auxiliary = {sein},
  style = {default},
  present = {komme vor,kommst vor,kommt vor,kommen vor,kommt vor,kommen vor},
  preterite = {kam vor,kamst vor,kam vor,kamen vor,kamt vor,kamen vor},
  subjunctive = {},
  imperative = {komm (du) vor!,kommen wir vor!,kommt (ihr) vor!,kommen Sie vor!},
  perfect = {sein + \textbf{vorgekommen}},
  future = {werden + \textbf{vorkommen}},
  pluperfect = {waren + \textbf{vorgekommen}},
  prepositions = {---},
  separable = {---}
}

% Migrated from verbs/b1/v.tex:38
\VerbTable{
  name = {vorlesen},
  english = {to read aloud},
  type = {separable, strong / irregular},
  auxiliary = {haben},
  style = {default},
  present = {lese vor,liest vor,liest vor,lesen vor,lest vor,lesen vor},
  preterite = {las vor,last vor,las vor,lasen vor,last vor,lasen vor},
  subjunctive = {},
  imperative = {lies (du) vor!,lesen wir vor!,lest (ihr) vor!,lesen Sie vor!},
  perfect = {haben + \textbf{vorgelesen}},
  future = {werden + \textbf{vorlesen}},
  pluperfect = {hatten + \textbf{vorgelesen}},
  prepositions = {---},
  separable = {---}
}

% Migrated from verbs/b1/v.tex:39
\VerbTable{
  name = {vorschlagen},
  english = {to suggest},
  type = {separable, strong / irregular},
  auxiliary = {haben},
  style = {default},
  present = {schlage vor,schlägst vor,schlägt vor,schlagen vor,schlagt vor,schlagen vor},
  preterite = {schlug vor,schlugst vor,schlug vor,schlugen vor,schlugt vor,schlugen vor},
  subjunctive = {},
  imperative = {schlag (du) vor!,schlagen wir vor!,schlagt (ihr) vor!,schlagen Sie vor!},
  perfect = {haben + \textbf{vorgeschlagen}},
  future = {werden + \textbf{vorschlagen}},
  pluperfect = {hatten + \textbf{vorgeschlagen}},
  prepositions = {---},
  separable = {---}
}

% Migrated from verbs/b1/v.tex:40
\VerbTable{
  name = {vorstellen},
  english = {to introduce},
  type = {separable, regular},
  auxiliary = {haben},
  style = {default},
  present = {stelle vor,stellst vor,stellt vor,stellen vor,stellt vor,stellen vor},
  preterite = {stellte vor,stelltest vor,stellte vor,stellten vor,stelltet vor,stellten vor},
  subjunctive = {},
  imperative = {stell (du) vor!,stellen wir vor!,stellt (ihr) vor!,stellen Sie vor!},
  perfect = {haben + \textbf{vorgestellt}},
  future = {werden + \textbf{vorstellen}},
  pluperfect = {hatten + \textbf{vorgestellt}},
  prepositions = {---},
  separable = {---}
}

\section{W}

% Migrated from verbs/b1/w.tex:1
\VerbTable{
  name = {wachsen},
  english = {to grow},
  type = {strong / irregular},
  auxiliary = {sein},
  style = {default},
  present = {wachse,wächst,wächst,wachsen,wachst,wachsen},
  preterite = {wuchs,wuchst,wuchs,wuchsen,wuchst,wuchsen},
  subjunctive = {},
  imperative = {wachs (du)!,wachsen wir!,wachst (ihr)!,wachsen Sie!},
  perfect = {sein + \textbf{gewachsen}},
  future = {werden + \textbf{wachsen}},
  pluperfect = {waren + \textbf{gewachsen}},
  prepositions = {---},
  separable = {---}
}

% Migrated from verbs/w.tex:1
\VerbTable{
  name = {wählen},
  english = {to choose / vote},
  type = {regular},
  auxiliary = {haben},
  style = {acc},
  present = {wähle,wählst,wählt,wählen,wählt,wählen},
  preterite = {wählte,wähltest,wählte,wählten,wähltet,wählten},
  subjunctive = {},
  imperative = {wähle (du)!,wählen wir!,wählt (ihr)!,wählen Sie!},
  perfect = {haben + \textbf{gewählt}},
  future = {werden + \textbf{wählen}},
  pluperfect = {hatten + \textbf{gewählt}},
  prepositions = {
    & \textbf{wählen für} + Akk & (to vote for / choose for)\\
    & \textbf{wählen zwischen} + Dat & (to choose between)\\
    & \textbf{wählen zu} + Dat & (to elect as / choose as)\\
  },
  separable = {& \textbf{aus$|$wählen} & (to select / pick out)\\}
}

% Migrated from verbs/b1/w.tex:2
\VerbTable{
  name = {wandern},
  english = {to hike},
  type = {regular},
  auxiliary = {sein},
  style = {default},
  present = {wandre,wanderst,wandert,wandern,wandert,wandern},
  preterite = {wanderte,wandertest,wanderte,wanderten,wandertet,wanderten},
  subjunctive = {},
  imperative = {wandre (du)!,wandern wir!,wandert (ihr)!,wandern Sie!},
  perfect = {sein + \textbf{gewandert}},
  future = {werden + \textbf{wandern}},
  pluperfect = {waren + \textbf{gewandert}},
  prepositions = {---},
  separable = {---}
}

% Migrated from verbs/b1/w.tex:3
\VerbTable{
  name = {warnen},
  english = {to warn},
  type = {regular},
  auxiliary = {haben},
  style = {default},
  present = {warne,warnst,warnt,warnen,warnt,warnen},
  preterite = {warnte,warntest,warnte,warnten,warntet,warnten},
  subjunctive = {},
  imperative = {warn (du)!,warnen wir!,warnt (ihr)!,warnen Sie!},
  perfect = {haben + \textbf{gewarnt}},
  future = {werden + \textbf{warnen}},
  pluperfect = {hatten + \textbf{gewarnt}},
  prepositions = {---},
  separable = {---}
}

% Migrated from verbs/w.tex:2
\VerbTable{
  name = {warten},
  english = {to wait},
  type = {regular},
  auxiliary = {haben},
  style = {default},
  present = {warte,wartest,wartet,warten,wartet,warten},
  preterite = {wartete,wartetest,wartete,warteten,wartetet,warteten},
  subjunctive = {},
  imperative = {wart(e) (du)!,warten wir!,wartet (ihr)!,warten Sie!},
  perfect = {haben + \textbf{gewartet}},
  future = {werden + \textbf{warten}},
  pluperfect = {hatten + \textbf{gewartet}},
  prepositions = {---},
  separable = {---}
}

% Migrated from verbs/w.tex:3
\VerbTable{
  name = {waschen},
  english = {to wash},
  type = {regular, + Akkusativ},
  auxiliary = {haben},
  style = {acc},
  present = {wasche,wäschst,wäscht,waschen,wascht,waschen},
  preterite = {wusch,wuschst,wusch,wuschen,wuscht,wuschen},
  subjunctive = {},
  imperative = {wasch(e) (du)!,waschen wir!,wascht (ihr)!,waschen Sie!},
  perfect = {haben + \textbf{gewaschen}},
  future = {werden + \textbf{waschen}},
  pluperfect = {hatten + \textbf{gewaschen}},
  prepositions = {---},
  separable = {---}
}

% Migrated from verbs/w.tex:4
\VerbTable{
  name = {wechseln},
  english = {to change / switch},
  type = {regular},
  auxiliary = {haben / sein (rarely)},
  style = {default},
  present = {wechsle,wechselst,wechselt,wechseln,wechselt,wechseln},
  preterite = {wechselte,wechseltest,wechselte,wechselten,wechseltet,wechselten},
  subjunctive = {},
  imperative = {wechsle (du)!,wechseln wir!,wechselt (ihr)!,wechseln Sie!},
  perfect = {haben + \textbf{gewechselt}},
  future = {werden + \textbf{wechseln}},
  pluperfect = {hatten + \textbf{gewechselt}},
  prepositions = {
    & \textbf{wechseln zu} + Dat & (to switch to)\\
    & \textbf{wechseln mit} + Dat & (to exchange with)\\
  },
  separable = {& \textbf{um$|$wechseln} & (to change / swap around)\\}
}

% Migrated from verbs/b1/w.tex:4
\VerbTable{
  name = {wecken},
  english = {to wake up},
  type = {regular},
  auxiliary = {haben},
  style = {default},
  present = {wecke,weckst,weckt,wecken,weckt,wecken},
  preterite = {weckte,wecktest,weckte,weckten,wecktet,weckten},
  subjunctive = {},
  imperative = {weck (du)!,wecken wir!,weckt (ihr)!,wecken Sie!},
  perfect = {haben + \textbf{geweckt}},
  future = {werden + \textbf{wecken}},
  pluperfect = {hatten + \textbf{geweckt}},
  prepositions = {---},
  separable = {---}
}

% Migrated from verbs/b1/w.tex:5
\VerbTable{
  name = {wehtun},
  english = {to hurt},
  type = {separable, strong / irregular},
  auxiliary = {haben},
  style = {default},
  present = {tue weh,tust weh,tut weh,tun weh,tut weh,tun weh},
  preterite = {tat weh,tatst weh,tat weh,taten weh,tatet weh,taten weh},
  subjunctive = {},
  imperative = {tu (du) weh!,tun wir weh!,tut (ihr) weh!,tun Sie weh!},
  perfect = {haben + \textbf{wehgetan}},
  future = {werden + \textbf{wehtun}},
  pluperfect = {hatten + \textbf{wehgetan}},
  prepositions = {---},
  separable = {---}
}

% Migrated from verbs/w.tex:5
\VerbTable{
  name = {weinen},
  english = {to cry},
  type = {regular},
  auxiliary = {haben},
  style = {default},
  present = {weine,weinst,weint,weinen,weint,weinen},
  preterite = {weinte,weintest,weinte,weinten,weintet,weinten},
  subjunctive = {},
  imperative = {wein(e) (du)!,weinen wir!,weint (ihr)!,weinen Sie!},
  perfect = {haben + \textbf{geweint}},
  future = {werden + \textbf{weinen}},
  pluperfect = {hatten + \textbf{geweint}},
  prepositions = {& \textbf{weinen über} + Akk & (to cry about)\\},
  separable = {---}
}

% Migrated from verbs/b1/w.tex:7
\VerbTable{
  name = {wenden},
  english = {to turn},
  type = {mixed},
  auxiliary = {haben},
  style = {default},
  present = {wende,wendest,wendet,wenden,wendet,wenden},
  preterite = {wendete,wendetest,wendete,wendeten,wendetet,wendeten},
  subjunctive = {},
  imperative = {wende (du)!,wenden wir!,wendet (ihr)!,wenden Sie!},
  perfect = {haben + \textbf{gewendet}},
  future = {werden + \textbf{wenden}},
  pluperfect = {hatten + \textbf{gewendet}},
  prepositions = {---},
  separable = {---}
}

% Migrated from verbs/w.tex:6
\VerbTable{
  name = {werfen},
  english = {to throw},
  type = {irregular},
  auxiliary = {haben},
  style = {default},
  present = {werfe,wirfst,wirft,werfen,werft,werfen},
  preterite = {warf,warfst,warf,warfen,warft,warfen},
  subjunctive = {},
  imperative = {wirf (du)!,werfen wir!,werft (ihr)!,werfen Sie!},
  perfect = {haben + \textbf{geworfen}},
  future = {werde + \textbf{werfen}},
  pluperfect = {hatten + \textbf{geworfen}},
  prepositions = {
    & \textbf{werfen auf} + Akk & (to throw onto / at)\\
    & \textbf{werfen in} + Akk & (to throw into)\\
  },
  separable = {
    & \textbf{weg$|$werfen} & (to throw away)\\
    & \textbf{ein$|$werfen} & (to throw in / insert)\\
  }
}

% Migrated from verbs/b1/w.tex:8
\VerbTable{
  name = {wetten},
  english = {to bet},
  type = {regular},
  auxiliary = {haben},
  style = {default},
  present = {wette,wettest,wettet,wetten,wettet,wetten},
  preterite = {wettete,wettetest,wettete,wetteten,wettetet,wetteten},
  subjunctive = {},
  imperative = {wette (du)!,wetten wir!,wettet (ihr)!,wetten Sie!},
  perfect = {haben + \textbf{gewettet}},
  future = {werden + \textbf{wetten}},
  pluperfect = {hatten + \textbf{gewettet}},
  prepositions = {---},
  separable = {---}
}

% Migrated from verbs/w.tex:7
\VerbTable{
  name = {widersprechen},
  english = {to contradict},
  type = {irregular, + Dativ},
  auxiliary = {haben},
  style = {dat},
  present = {widerspreche,widersprichst,widerspricht,widersprechen,widersprecht,widersprechen},
  preterite = {widersprach,widersprachst,widersprach,widersprachen,widerspracht,widersprachen},
  subjunctive = {},
  imperative = {widersprich (du)!,widersprechen wir!,widersprecht (ihr)!,widersprechen Sie!},
  perfect = {haben + \textbf{widersprochen}},
  future = {werden + \textbf{widersprechen}},
  pluperfect = {hatten + \textbf{widersprochen}},
  prepositions = {& -- & --\\},
  separable = {& -- & --\\}
}

% Migrated from verbs/b1/w.tex:9
\VerbTable{
  name = {wiederholen},
  english = {to repeat},
  type = {regular},
  auxiliary = {haben},
  style = {default},
  present = {wiederhole,wiederholst,wiederholt,wiederholen,wiederholt,wiederholen},
  preterite = {wiederholte,wiederholtest,wiederholte,wiederholten,wiederholtet,wiederholten},
  subjunctive = {},
  imperative = {wiederhol (du)!,wiederholen wir!,wiederholt (ihr)!,wiederholen Sie!},
  perfect = {haben + \textbf{wiederholt}},
  future = {werden + \textbf{wiederholen}},
  pluperfect = {hatten + \textbf{wiederholt}},
  prepositions = {---},
  separable = {---}
}

% Migrated from verbs/b1/w.tex:10
\VerbTable{
  name = {wiegen},
  english = {to weigh},
  type = {strong / irregular},
  auxiliary = {haben},
  style = {default},
  present = {wiege,wiegst,wiegt,wiegen,wiegt,wiegen},
  preterite = {wog,wogst,wog,wogen,wogt,wogen},
  subjunctive = {},
  imperative = {wieg (du)!,wiegen wir!,wiegt (ihr)!,wiegen Sie!},
  perfect = {haben + \textbf{gewogen}},
  future = {werden + \textbf{wiegen}},
  pluperfect = {hatten + \textbf{gewogen}},
  prepositions = {---},
  separable = {---}
}

% Migrated from verbs/b1/w.tex:11
\VerbTable{
  name = {winken},
  english = {to wave},
  type = {regular},
  auxiliary = {haben},
  style = {default},
  present = {winke,winkst,winkt,winken,winkt,winken},
  preterite = {winkte,winktest,winkte,winkten,winktet,winkten},
  subjunctive = {},
  imperative = {wink (du)!,winken wir!,winkt (ihr)!,winken Sie!},
  perfect = {haben + \textbf{gewinkt}},
  future = {werden + \textbf{winken}},
  pluperfect = {hatten + \textbf{gewinkt}},
  prepositions = {---},
  separable = {---}
}

% Migrated from verbs/w.tex:8
\VerbTable{
  name = {wirken},
  english = {to appear},
  type = {regular},
  auxiliary = {haben},
  style = {default},
  present = {wirke,wirkst,wirkt,wirken,wirkt,wirken},
  preterite = {wirkte,wirktest,wirkte,wirkten,wirktet,wirkten},
  subjunctive = {},
  imperative = {wirk(e) (du)!,wirken wir!,wirkt (ihr)!,wirken Sie!},
  perfect = {haben + \textbf{gewirkt}},
  future = {werden + \textbf{wirken}},
  pluperfect = {hatten + \textbf{gewirkt}},
  prepositions = {
    & \textbf{wirken auf} + Akk & (to influence)\\
    & \textbf{wirken gegen} + Akk & (to work against)\\
    &\textbf{wirken als} + Nom & (to function as)\\
    &\textbf{wirken wie} + Nom & (to seem like)\\
    &\textbf{wirken durch} + Akk & (to have an effect through)\\
    &\textbf{wirken mit} + Dat & (to contribute with)
  },
  separable = {& \textbf{mit$|$wirken} & (to participate to)\\}
}

% Migrated from verbs/w.tex:9
\VerbTable{
  name = {wissen},
  english = {to know},
  type = {irregular},
  auxiliary = {haben},
  style = {default},
  present = {weiß,weißt,weiß,wissen,wisst,wissen},
  preterite = {wusste,wusstest,wusste,wussten,wusstet,wussten},
  subjunctive = {},
  imperative = {wisse (du)!,wissen wir!,wisst (ihr)!,wissen Sie!},
  perfect = {haben + \textbf{gewusst}},
  future = {werden + \textbf{wissen}},
  pluperfect = {hatten + \textbf{gewusst}},
  prepositions = {---},
  separable = {---}
}

% Migrated from verbs/w.tex:10
\VerbTable{
  name = {wohnen},
  english = {to live},
  type = {regular},
  auxiliary = {haben},
  style = {default},
  present = {wohne,wohnst,wohnt,wohnen,wohnt,wohnen},
  preterite = {wohnte,wohntest,wohnte,wohnten,wohntet,wohnten},
  subjunctive = {},
  imperative = {wohn(e) (du)!,wohnen wir!,wohnt (ihr)!,wohnen Sie!},
  perfect = {haben + \textbf{gewohnt}},
  future = {werden + \textbf{wohnen}},
  pluperfect = {hatten + \textbf{gewohnt}},
  prepositions = {---},
  separable = {---}
}

% Migrated from verbs/b1/w.tex:12
\VerbTable{
  name = {wünschen},
  english = {to wish},
  type = {regular},
  auxiliary = {haben},
  style = {default},
  present = {wünsche,wünschst,wünscht,wünschen,wünscht,wünschen},
  preterite = {wünschte,wünschtest,wünschte,wünschten,wünschtet,wünschten},
  subjunctive = {},
  imperative = {wünsch (du)!,wünschen wir!,wünscht (ihr)!,wünschen Sie!},
  perfect = {haben + \textbf{gewünscht}},
  future = {werden + \textbf{wünschen}},
  pluperfect = {hatten + \textbf{gewünscht}},
  prepositions = {---},
  separable = {---}
}

\section{Z}

% Migrated from verbs/b1/z.tex:1
\VerbTable{
  name = {zählen},
  english = {to count},
  type = {regular},
  auxiliary = {haben},
  style = {default},
  present = {zähle,zählst,zählt,zählen,zählt,zählen},
  preterite = {zählte,zähltest,zählte,zählten,zähltet,zählten},
  subjunctive = {},
  imperative = {zähl (du)!,zählen wir!,zählt (ihr)!,zählen Sie!},
  perfect = {haben + \textbf{gezählt}},
  future = {werden + \textbf{zählen}},
  pluperfect = {hatten + \textbf{gezählt}},
  prepositions = {---},
  separable = {---}
}

% Migrated from verbs/z.tex:1
\VerbTable{
  name = {zahlen},
  english = {to pay},
  type = {regular},
  auxiliary = {haben},
  style = {default},
  present = {zahle,zahlst,zahlt,zahlen,zahlt,zahlen},
  preterite = {zahlte,zahltest,zahlte,zahlten,zahltet,zahlten},
  subjunctive = {},
  imperative = {zahle (du)!,zahlen wir!,zahlt (ihr)!,zahlen Sie!},
  perfect = {haben + \textbf{gezahlt}},
  future = {werden + \textbf{zahlen}},
  pluperfect = {hatten + \textbf{gezahlt}},
  prepositions = {
    & \textbf{zahlen für} + Akk & (to pay for)\\
    & \textbf{zahlen an} + Akk & (to pay to)\\
  },
  separable = {---}
}

% Migrated from verbs/b1/z.tex:2
\VerbTable{
  name = {zeichnen},
  english = {to draw},
  type = {regular},
  auxiliary = {haben},
  style = {default},
  present = {zeichne,zeichnest,zeichnet,zeichnen,zeichnet,zeichnen},
  preterite = {zeichnete,zeichnetest,zeichnete,zeichneten,zeichnetet,zeichneten},
  subjunctive = {},
  imperative = {zeichne (du)!,zeichnen wir!,zeichnet (ihr)!,zeichnen Sie!},
  perfect = {haben + \textbf{gezeichnet}},
  future = {werden + \textbf{zeichnen}},
  pluperfect = {hatten + \textbf{gezeichnet}},
  prepositions = {---},
  separable = {---}
}

% Migrated from verbs/z.tex:2
\VerbTable{
  name = {zeigen},
  english = {to show / point out},
  type = {regular},
  auxiliary = {haben},
  style = {acc},
  present = {zeige,zeigst,zeigt,zeigen,zeigt,zeigen},
  preterite = {zeigte,zeigtest,zeigte,zeigten,zeigtet,zeigten},
  subjunctive = {},
  imperative = {zeige (du)!,zeigen wir!,zeigt (ihr)!,zeigen Sie!},
  perfect = {haben + \textbf{gezeigt}},
  future = {werden + \textbf{zeigen}},
  pluperfect = {hatten + \textbf{gezeigt}},
  prepositions = {
    & \textbf{zeigen auf} + Akk & (to point to / indicate)\\
    & \textbf{zeigen für} + Akk & (to show for / demonstrate)\\
    & \textbf{zeigen mit} + Dat & (to show with)\\
    & \textbf{zeigen von} + Dat & (to show of / about)\\
  },
  separable = {
    & \textbf{vor$|$zeigen} & (to demonstrate / exhibit)\\
    & \textbf{auf$|$zeigen} & (to point out / indicate)\\
    & \textbf{nach$|$zeigen} & (to show afterward / demonstrate later)\\
  }
}

% Migrated from verbs/b1/z.tex:3
\VerbTable{
  name = {zelten},
  english = {to camp},
  type = {regular},
  auxiliary = {haben},
  style = {default},
  present = {zelte,zeltest,zeltet,zelten,zeltet,zelten},
  preterite = {zeltete,zeltetest,zeltete,zelteten,zeltetet,zelteten},
  subjunctive = {},
  imperative = {zelte (du)!,zelten wir!,zeltet (ihr)!,zelten Sie!},
  perfect = {haben + \textbf{gezeltet}},
  future = {werden + \textbf{zelten}},
  pluperfect = {hatten + \textbf{gezeltet}},
  prepositions = {---},
  separable = {---}
}

% Migrated from verbs/z.tex:3
\VerbTable{
  name = {zerstören},
  english = {to destroy},
  type = {regular, inseparable},
  auxiliary = {haben},
  style = {default},
  present = {zerstöre,zerstörst,zerstört,zerstören,zerstört,zerstören},
  preterite = {zerstörte,zerstörtest,zerstörte,zerstörten,zerstörtet,zerstörten},
  subjunctive = {},
  imperative = {zerstöre (du)!,zerstören wir!,zerstört (ihr)!,zerstören Sie!},
  perfect = {haben + \textbf{zerstört}},
  future = {werden + \textbf{zerstören}},
  pluperfect = {hatten + \textbf{zerstört}},
  prepositions = {
    & \textbf{zerstören durch} + Akk & (to destroy through / by)\\
    & \textbf{zerstören mit} + Dat & (to destroy with)\\
  },
  separable = {---}
}

% Migrated from verbs/z.tex:4
\VerbTable{
  name = {ziehen},
  english = {to pull},
  type = {irregular},
  auxiliary = {haben},
  style = {default},
  present = {ziehe,ziehst,zieht,ziehen,zieht,ziehen},
  preterite = {zog,zogst,zog,zogen,zogt,zogen},
  subjunctive = {},
  imperative = {zieh(e) (du)!,ziehen wir!,zieht (ihr)!,ziehen Sie!},
  perfect = {haben + \textbf{gezogen}},
  future = {werden + \textbf{ziehen}},
  pluperfect = {hatten + \textbf{gezogen}},
  prepositions = {
    & \textbf{ziehen an} + Dat & (to pull on / at)\\
    & \textbf{ziehen nach} + Dat & (to move to / after)\\
    & \textbf{ziehen in} + Akk & (to move into)\\
    & \textbf{ziehen aus} + Dat & (to move out of)\\
    & \textbf{ziehen für} + Akk & (to pull for / draw for)\\
  },
  separable = {
    & \textbf{ein$|$ziehen} & (to move in / enlist / draw in)\\
    & \textbf{aus$|$ziehen} & (to move out / undress / extract)\\
    & \textbf{um$|$ziehen} & (to move / relocate / change clothes)\\
    & \textbf{zurück$|$ziehen} & (to pull back / withdraw)\\
    & \textbf{ab$|$ziehen} & (to subtract / withdraw / remove)\\
    & \textbf{auf$|$ziehen} & (to wind up / draw up / open)\\
    & \textbf{durch$|$ziehen} & (to pull through / carry out / complete)
  }
}

% Migrated from verbs/b1/z.tex:4
\VerbTable{
  name = {zubereiten},
  english = {to prepare},
  type = {separable, regular},
  auxiliary = {haben},
  style = {default},
  present = {bereite zu,bereitest zu,bereitet zu,bereiten zu,bereitet zu,bereiten zu},
  preterite = {bereitete zu,bereitetest zu,bereitete zu,bereiteten zu,bereitetet zu,bereiteten zu},
  subjunctive = {},
  imperative = {bereite (du) zu!,bereiten wir zu!,bereitet (ihr) zu!,bereiten Sie zu!},
  perfect = {haben + \textbf{zubereitet}},
  future = {werden + \textbf{zubereiten}},
  pluperfect = {hatten + \textbf{zubereitet}},
  prepositions = {---},
  separable = {---}
}

% Migrated from verbs/b1/z.tex:5
\VerbTable{
  name = {zugehen},
  english = {to approach / close},
  type = {separable, strong / irregular},
  auxiliary = {sein},
  style = {default},
  present = {gehe zu,gehst zu,geht zu,gehen zu,geht zu,gehen zu},
  preterite = {ging zu,gingst zu,ging zu,gingen zu,gingt zu,gingen zu},
  subjunctive = {},
  imperative = {geh (du) zu!,gehen wir zu!,geht (ihr) zu!,gehen Sie zu!},
  perfect = {sein + \textbf{zugegangen}},
  future = {werden + \textbf{zugehen}},
  pluperfect = {waren + \textbf{zugegangen}},
  prepositions = {---},
  separable = {---}
}

% Migrated from verbs/b1/z.tex:6
\VerbTable{
  name = {zuhören},
  english = {to listen to},
  type = {separable, regular},
  auxiliary = {haben},
  style = {default},
  present = {höre zu,hörst zu,hört zu,hören zu,hört zu,hören zu},
  preterite = {hörte zu,hörtest zu,hörte zu,hörten zu,hörtet zu,hörten zu},
  subjunctive = {},
  imperative = {hör (du) zu!,hören wir zu!,hört (ihr) zu!,hören Sie zu!},
  perfect = {haben + \textbf{zugehört}},
  future = {werden + \textbf{zuhören}},
  pluperfect = {hatten + \textbf{zugehört}},
  prepositions = {---},
  separable = {---}
}

% Migrated from verbs/b1/z.tex:7
\VerbTable{
  name = {zumachen},
  english = {to close},
  type = {separable, regular},
  auxiliary = {haben},
  style = {default},
  present = {mache zu,machst zu,macht zu,machen zu,macht zu,machen zu},
  preterite = {machte zu,machtest zu,machte zu,machten zu,machtet zu,machten zu},
  subjunctive = {},
  imperative = {mach (du) zu!,machen wir zu!,macht (ihr) zu!,machen Sie zu!},
  perfect = {haben + \textbf{zugemacht}},
  future = {werden + \textbf{zumachen}},
  pluperfect = {hatten + \textbf{zugemacht}},
  prepositions = {---},
  separable = {---}
}

% Migrated from verbs/b1/z.tex:8
\VerbTable{
  name = {zunehmen},
  english = {to increase},
  type = {separable, strong / irregular},
  auxiliary = {haben},
  style = {default},
  present = {nehme zu,nimmst zu,nimmt zu,nehmen zu,nehmt zu,nehmen zu},
  preterite = {nahm zu,nahmst zu,nahm zu,nahmen zu,nahmt zu,nahmen zu},
  subjunctive = {},
  imperative = {nimm (du) zu!,nehmen wir zu!,nehmt (ihr) zu!,nehmen Sie zu!},
  perfect = {haben + \textbf{zugenommen}},
  future = {werden + \textbf{zunehmen}},
  pluperfect = {hatten + \textbf{zugenommen}},
  prepositions = {---},
  separable = {---}
}

% Migrated from verbs/b1/z.tex:9
\VerbTable{
  name = {zurechtkommen},
  english = {to cope},
  type = {separable, strong / irregular},
  auxiliary = {sein},
  style = {default},
  present = {komme zurecht,kommst zurecht,kommt zurecht,kommen zurecht,kommt zurecht,kommen zurecht},
  preterite = {kam zurecht,kamst zurecht,kam zurecht,kamen zurecht,kamt zurecht,kamen zurecht},
  subjunctive = {},
  imperative = {komm (du) zurecht!,kommen wir zurecht!,kommt (ihr) zurecht!,kommen Sie zurecht!},
  perfect = {sein + \textbf{zurechtgekommen}},
  future = {werden + \textbf{zurechtkommen}},
  pluperfect = {waren + \textbf{zurechtgekommen}},
  prepositions = {---},
  separable = {---}
}

% Migrated from verbs/b1/z.tex:10
\VerbTable{
  name = {zusagen},
  english = {to accept / promise},
  type = {separable, regular},
  auxiliary = {haben},
  style = {default},
  present = {sage zu,sagst zu,sagt zu,sagen zu,sagt zu,sagen zu},
  preterite = {sagte zu,sagtest zu,sagte zu,sagten zu,sagtet zu,sagten zu},
  subjunctive = {},
  imperative = {sag (du) zu!,sagen wir zu!,sagt (ihr) zu!,sagen Sie zu!},
  perfect = {haben + \textbf{zugesagt}},
  future = {werden + \textbf{zusagen}},
  pluperfect = {hatten + \textbf{zugesagt}},
  prepositions = {---},
  separable = {---}
}

% Migrated from verbs/b1/z.tex:11
\VerbTable{
  name = {zusammenfassen},
  english = {to summarize},
  type = {separable, regular},
  auxiliary = {haben},
  style = {default},
  present = {fasse zusammen,fasst zusammen,fasst zusammen,fassen zusammen,fasst zusammen,fassen zusammen},
  preterite = {fasste zusammen,fasstest zusammen,fasste zusammen,fassten zusammen,fasstet zusammen,fassten zusammen},
  subjunctive = {},
  imperative = {fass (du) zusammen!,fassen wir zusammen!,fasst (ihr) zusammen!,fassen Sie zusammen!},
  perfect = {haben + \textbf{zusammengefasst}},
  future = {werden + \textbf{zusammenfassen}},
  pluperfect = {hatten + \textbf{zusammengefasst}},
  prepositions = {---},
  separable = {---}
}

% Migrated from verbs/b1/z.tex:12
\VerbTable{
  name = {zuschauen},
  english = {to watch},
  type = {separable, regular},
  auxiliary = {haben},
  style = {default},
  present = {schaue zu,schaust zu,schaut zu,schauen zu,schaut zu,schauen zu},
  preterite = {schaute zu,schautest zu,schaute zu,schauten zu,schautet zu,schauten zu},
  subjunctive = {},
  imperative = {schau (du) zu!,schauen wir zu!,schaut (ihr) zu!,schauen Sie zu!},
  perfect = {haben + \textbf{zugeschaut}},
  future = {werden + \textbf{zuschauen}},
  pluperfect = {hatten + \textbf{zugeschaut}},
  prepositions = {---},
  separable = {---}
}

% Migrated from verbs/b1/z.tex:13
\VerbTable{
  name = {zustimmen},
  english = {to agree},
  type = {separable, regular},
  auxiliary = {haben},
  style = {default},
  present = {stimme zu,stimmst zu,stimmt zu,stimmen zu,stimmt zu,stimmen zu},
  preterite = {stimmte zu,stimmtest zu,stimmte zu,stimmten zu,stimmtet zu,stimmten zu},
  subjunctive = {},
  imperative = {stimm (du) zu!,stimmen wir zu!,stimmt (ihr) zu!,stimmen Sie zu!},
  perfect = {haben + \textbf{zugestimmt}},
  future = {werden + \textbf{zustimmen}},
  pluperfect = {hatten + \textbf{zugestimmt}},
  prepositions = {---},
  separable = {---}
}

% Migrated from verbs/b1/z.tex:14
\VerbTable{
  name = {zweifeln},
  english = {to doubt},
  type = {regular},
  auxiliary = {haben},
  style = {default},
  present = {zweifle,zweifelst,zweifelt,zweifeln,zweifelt,zweifeln},
  preterite = {zweifelte,zweifeltest,zweifelte,zweifelten,zweifeltet,zweifelten},
  subjunctive = {},
  imperative = {zweifle (du)!,zweifeln wir!,zweifelt (ihr)!,zweifeln Sie!},
  perfect = {haben + \textbf{gezweifelt}},
  future = {werden + \textbf{zweifeln}},
  pluperfect = {hatten + \textbf{gezweifelt}},
  prepositions = {---},
  separable = {---}
}

% Migrated from verbs/b1/z.tex:15
\VerbTable{
  name = {zwingen},
  english = {to force},
  type = {strong / irregular},
  auxiliary = {haben},
  style = {default},
  present = {zwinge,zwingst,zwingt,zwingen,zwingt,zwingen},
  preterite = {zwang,zwangst,zwang,zwangen,zwangt,zwangen},
  subjunctive = {},
  imperative = {zwing (du)!,zwingen wir!,zwingt (ihr)!,zwingen Sie!},
  perfect = {haben + \textbf{gezwungen}},
  future = {werden + \textbf{zwingen}},
  pluperfect = {hatten + \textbf{gezwungen}},
  prepositions = {---},
  separable = {---}
}


\backmatter
\printindex[german]
\printindex[english]

\end{document}
